\documentclass[11pt]{article}
\usepackage[T1]{fontenc}\usepackage[utf8]{inputenc}
\usepackage{amsmath,amssymb,amsthm,mathtools,lmodern,microtype}
\usepackage[a4paper,margin=25mm,headheight=15pt]{geometry}
\usepackage{booktabs,longtable,array,enumitem,fancyhdr,xurl,hyperref}
\hypersetup{colorlinks=true,linkcolor=black,citecolor=black,urlcolor=blue,
 pdftitle={Rank Ultra-Log-Concavity of Preorder Gamma and Directed Support Polynomials at the Actual Degree},
 pdfauthor={Research manuscripts prepared for Vladimir Reshetnikov (several with OpenAI); merged in ProveIt's batch-73 intake; Section 14 added in batch 77}}
\pagestyle{fancy}\fancyhf{}\fancyhead[L]{\small Preorder gamma and directed support polynomials}\fancyhead[R]{\small Rank-ULC at the actual degree}\fancyfoot[C]{\thepage}
\setlength{\parskip}{3pt}\setlength{\emergencystretch}{2em}
\numberwithin{equation}{section}
\newtheorem{theorem}{Theorem}[section]\newtheorem{lemma}[theorem]{Lemma}\newtheorem{proposition}[theorem]{Proposition}\newtheorem{corollary}[theorem]{Corollary}
\theoremstyle{definition}\newtheorem{definition}[theorem]{Definition}\newtheorem{remark}[theorem]{Remark}
\setcounter{tocdepth}{2}
% Macros of the sources, collected (Section~\ref{pgr:sec:notions}).
\newcommand{\N}{\mathbb Z_{\ge0}}          % sources 16, 20, 22, 23, 25
\newcommand{\R}{\mathbb R}                 % sources 16, 20, 22, 25
\newcommand{\Q}{\mathbb Q}                 % source 25 (the rationals; never a polynomial)
\newcommand{\Ga}{\Gamma}                   % every source
\newcommand{\HH}{\mathbb H}                % sources 15, 16, 20, 22: upper half-plane
\newcommand{\supp}{\operatorname{supp}}    % sources 16, 20, 25
\newcommand{\Disc}{\operatorname{Disc}}    % source 25
\newcommand{\Comp}{\operatorname{Comp}}    % source 25
\newcommand{\wt}{\operatorname{wt}}        % sources 11, 12, 15
\newcommand{\coefge}{\mathrel{\succeq_{\rm coeff}}} % source 12
\newcommand{\MAP}{\operatorname{MAP}}      % source 15
\newcommand{\ee}{\mathrm e}                % sources 06, 10, 14, 34: elementary symmetric e_j
\newcommand{\Qp}{Q^{+}}                    % source 23
\newcommand{\Bp}{B^{+}}                    % source 23
\newcommand{\Tp}{T^{+}}                    % source 23
\newcommand{\binm}[2]{\binom{#1}{#2}}      % source 23
\newcommand{\file}[1]{\nolinkurl{#1}}     % source 23 (there \texttt{\detokenize{#1}}; breakable here)
\newcommand{\nuu}{\nu}                     % source 22
\newcommand{\del}[2]{#1\setminus #2}       % source 22
% Editorial marks of this merge.
\newcommand{\mnote}[1]{[\textit{Merge note:} #1]}
\newcommand{\prp}{\emph{preorder-root-polytopes}}
\newcommand{\mrn}{\emph{matching-rank-normalization}}
\makeatletter
% A reference to a label of a Part that may be written in a later step of
% the intake: prints \ref{#1} once the label exists, the text #2 before.
\newcommand{\pgrfwd}[2]{\@ifundefined{r@#1}{#2}{\ref{#1}}}
\makeatother
\begin{document}
\hypersetup{pageanchor=false}
\title{Rank Ultra-Log-Concavity of Preorder Gamma and\\ Directed Support Polynomials at the Actual Degree\\[4pt]
\large Degree four, independent role activities, small covers\\ and universal-sink clouds}
\author{Research manuscripts prepared for Vladimir Reshetnikov\\
(PDF metadata of sources 11, 12, 15, 16, 20, 22 and 25: ``prepared with OpenAI'')\\[2pt]
\small merged in ProveIt's batch-73 intake; Section~14 added in batch~77}
\date{1 October 2026}
\maketitle
\thispagestyle{empty}
\hypersetup{pageanchor=true}

\begin{abstract}
For a finite loopless directed relation $D$, for instance the off-diagonal
relation of a finite preorder, let $\gamma_k$ count the ordered pairs $(S,T)$
of disjoint $k$-subsets of the ground set that admit a bijection
$S\to T$ along arcs, each pair counted once however many matchings witness
it, and possibly weighted by independent tail and head activities. For a
preorder, $\Ga_D=\sum_k\gamma_kt^k$ is the gamma polynomial of
\prp. That report asks (its Research question~96, \texttt{mr:q:gamma})
whether $\Ga_D$ is ultra-log-concave after normalization by its
\emph{actual} degree, and names degree three as the first open target. This
report merges eleven manuscripts and a supplement (twelve sources) of % ed. (2026-10-02): was "eleven manuscripts"; source 23 is the supplement
ProveIt's batch~73, and one manuscript of batch~77 (source~35), that % batch 77 (2026-10-02): source 35 added
answer the question in three directions.
Part~I (sources 25, 16 and~23) proves it for every finite preorder of actual
degree at most four with unit activities: degree three by a
Gallai--Edmonds reduction and exact integer-population certificates, with
the minimum order fifteen of a nonreal-rooted cubic; degree four by bounded
cores, 76 four-attachment templates and 50 global binomial-square
identities, with a smaller second route for the largest certificate.
Part~II (sources 20, 12, 11, 15 and~22) treats independent tail and head
activities: real stability of the signed monomer polynomial under every two
by two role cover, coefficientwise Rayleigh inequalities for a three-vertex
shore, rank-ULC under every physical vertex cover of size three and on
bipartite physical graphs, a pendant-head stability criterion, and rank-ULC
of every independently weighted preorder of actual degree at most three.
Part~III (sources 14, 10, 34, 06 and~35) treats an arbitrary directed core with a
cloud of universal sinks: a sharp last-coefficient bound for every core
size, and rank-ULC at every actual degree for four-vertex and for % batch 77 (2026-10-02): was "and for five-vertex cores at actual degrees five and four"
five-vertex cores; for five-vertex cores the cubic constant three holds
on the face $\gamma_4=0$ and is sharp there.
The degree-three and degree-four theorems are computer-assisted; the
structural theorems have ordinary proofs. Nothing here is formalized in
Lean or Rocq, and nothing is refereed.
\end{abstract}

\tableofcontents

\section{Guide to this report}\label{pgr:sec:guide}

\subsection{Sources, status and plan}\label{pgr:sec:scope}
This report merges research manuscripts delivered to ProveIt's
incoming-report intake on 1 October 2026 (batch~73, arrival commit
\texttt{f8c3a392a}, placement commit \texttt{8f5c53106}). Sources are
numbered by their batch-73 manuscript numbers; a package delivered in
several archives takes the lowest number: source~16 is the four archives
16--19 (\texttt{preorder-degree-four}, parts 4 to 1), and source~06 is the
two archives 06 and~07 (\texttt{universal-sink-five-core-boundary}, parts 2
and~1; its manuscript is in archive~07). Sixteen archives are behind the % ed. (2026-10-02): was "Eighteen archives are behind the eleven printed sources"
twelve printed sources (eleven manuscripts and source~23's supplement). Two
further archives, making eighteen, are earlier versions of
printed sources and are not sources: 31 (vertex activities, superseded by
source~22) and 13 (a weaker four-core inequality, superseded by source~10).
Appendix~\ref{pgr:sec:provenance} records every archive.

Source~35 arrived later, in ProveIt's batch~77 (arrival commit % batch 77 (2026-10-02)
\texttt{096ee7b87}, 2~October 2026; placement commit \texttt{089d2b825}).
It is batch-77 manuscript~71, the archive
\nolinkurl{universal-sink-five-core-cubic-package}, delivered as six ZIP parts
in one archive. It takes the next number after this report's highest
source number,~34, and is printed as Section~\ref{pgr:s:fq:sec}, the last
section of Part~III. With it the report has thirteen sources.

The report has three Parts.
\begin{itemize}
\item Part~I, unit activities on finite preorders through actual degree
  four: source~25 (degree three), source~16 (degree four, the base of the
  merge) and source~23 (a second route for the largest degree-four
  certificate).
\item Part~II, independent role activities and small covers: sources 20,
  12, 11, 15 and~22.
\item Part~III, universal-sink clouds over an arbitrary directed core:
  sources 14, 10, 34, 06 and~35.
\end{itemize}

Every source is printed with its statements, proofs, examples, remarks,
limitations, verification notes and research questions; each source's
section opens with its abstract as delivered. A result proved by several
sources is printed once, and the other proofs are kept as marked second
routes or, where the proof is the same, replaced by a pointer recorded at
the place. A result that \prp{} or \mrn{} already proves is cited by its
label there and is not presented as new. Editorial additions are marked
``[\textit{Merge note:} \dots]''.

\emph{Status.} The sources are AI-assisted research manuscripts and notes:
the title pages of sources 11, 12, 15, 16, 20, 22 and~25 read ``Research
manuscript prepared for Vladimir Reshetnikov'' and their PDF metadata
``Research manuscript prepared with OpenAI''; sources 06, 10, 14, 34 and~35 read
``Research note''; source~23 has no author line. Several main theorems are
computer-assisted: their proofs include finite exact certificate
computations shipped with this report (Section~\ref{pgr:sec:prov:files}).
The report is unrefereed, and nothing in it is formalized in a proof
assistant. No source claims literature-wide priority.

\subsection{What the neighbouring reports already prove}\label{pgr:sec:prior}
Two reports of the collection precede this one. Their results are cited by
label and are not reprinted. Line numbers are not used: the labels are
stable.
\begin{itemize}
\item \prp{} (in \texttt{enumerative-combinatorics/}). Theorem
  \texttt{gam:thm:main} expresses the $h$-polynomial of a finite preorder
  through the support numbers of this report, and Theorem
  \texttt{gam:thm:deficiency} says $\deg\Ga_\tau=\nu(\Comp(\tau))$.
  Theorem \texttt{lor:thm:weightedULC} proves two-shore ultra-log-concavity
  normalized by the smaller shore, and Corollary \texttt{lor:cor:gammaULC}
  all height-two gamma sequences. Theorems \texttt{ade:mat:thm:stable-bridge}
  and \texttt{ade:mat:thm:vertex-stable} treat stability of signed matching
  supports of bipartite graphs. In Part~XII, Theorem \texttt{mr:thm:first}
  (with Lemma \texttt{mr:lem:MS}) is the universal first inequality for
  directed relations with nonnegative activities, with a merge note
  crediting its vanishing-activity form; Theorem \texttt{mr:thm:rankthree}
  and Corollary \texttt{lc:cor:rankthree} settle bipartite matching rank
  three; Corollary \texttt{mr:cor:preorderfirst} gives the first inequality
  and hence every degree at most two for preorder gamma polynomials; and
  Questions \texttt{mr:q:gamma} (Research question~96) and
  \texttt{mr:q:certificates} are the questions this report answers or bears
  on.
\item \mrn{} (beside this report). Corollary \texttt{mrn:w:rf:cor:rankfive}:
  every weighted bipartite graph of matching number at most five is
  ultra-log-concave of order its matching number. Theorem
  \texttt{mrn:w:sb:thm:main}: a bipartite graph with a vertex cover of two
  vertices on each shore has a real-stable support polynomial. Equation
  \texttt{mrn:s:tv:eq:rayleigh}: the scalar rank-three Rayleigh inequality
  (Wagner). Its Part~V prints batch-73 manuscript~26 (source~77 there):
  Theorem \texttt{mrn:v:rf4:thm:main}, normalization by the actual degree
  through rank four, and Theorem \texttt{mrn:v:rf4:thm:mixed}, real roots
  under a two by two cover.
\end{itemize}
Sources~16 and~25 cite \prp{} at its commit \texttt{d5e863bba}; the labels
they cite are unchanged at the placement commit.

\subsection{The common object}\label{pgr:sec:object}
The following definitions are common to all thirteen sources; each source % ed. (2026-10-02): was "eleven"; batch 77: "twelve"
restates them in its own words, and those restatements are kept in their
sections because they fix each source's precise scope. Let $D$ be a finite
loopless directed relation on a set $V$. For a preorder, $D$ is its
off-diagonal relation: an arc $i\to j$ means $i\ne j$ and $i$ precedes $j$,
and equivalent elements give opposite arcs.

\begin{definition}[Feasible supports]\label{pgr:def:support}
A feasible support of size $k$ is an ordered pair $(S,T)$ of disjoint
subsets of $V$, each of cardinality $k$, for which there exists a bijection
$\pi:S\to T$ with $s\to\pi(s)$ for every $s\in S$. For nonnegative tail
activities $u_i$ and head activities $v_i$ put
\[
 \gamma_k(D)=\sum_{(S,T)\text{ feasible},\ |S|=k}u^Sv^T,\qquad
 u^Sv^T=\prod_{i\in S}u_i\prod_{j\in T}v_j,\qquad
 \Ga_D(t)=\sum_{k\ge0}\gamma_k(D)t^k .
\]
The empty support contributes $\gamma_0=1$. A support contributes once,
regardless of the number of witnessing bijections. With unit activities
$\gamma_k$ is the number of feasible supports of size~$k$.
\end{definition}
The same-ground-set disjointness is essential. Passing to independent tail
and head copies without subsequently excluding their simultaneous use
changes this polynomial. Counting individual matchings also changes it.
A vertex may have different activities in its two roles; with zero
activities the polynomial degree can decrease.

\begin{definition}[Rank-ULC]\label{pgr:def:ulc}
A polynomial $g(t)=\sum_{k=0}^dg_kt^k$ is \emph{rank-ultra-log-concave}, or
rank-ULC, if its coefficients are nonnegative with no internal zeros and
\[
 \left(\frac{g_k}{\binom dk}\right)^2\ge
 \frac{g_{k-1}}{\binom d{k-1}}\frac{g_{k+1}}{\binom d{k+1}}
 \quad(1\le k<d),
\]
where $d$ is the actual degree.
\end{definition}
Normalization by the ground-set size, by a cover size, or by a presumed
degree after some activities vanish is a different assertion, and
Section~\ref{pgr:sec:notions} lists where a source uses one.

\begin{lemma}[Degree; sources 16, 20, 25]\label{pgr:lem:degree}
For unit activities, the actual degree of $\Ga_D$ equals the matching number
of the underlying undirected graph. Its coefficients have no internal
zeros.
\end{lemma}
\begin{proof}
A witnessing bijection yields an undirected matching. Conversely orient each
edge of an undirected matching in a direction present in $D$. Its disjoint
endpoint sets form a feasible support. Submatchings establish positivity at
every smaller degree.
\end{proof}
\mnote{for preorders this is \prp, Theorem \texttt{gam:thm:deficiency}.
Sources 16, 20 and~25 state and prove it as their Lemma ``degree'' with the
proof above (source~25 for preorders, in its own words); source~22 adds a
filtering argument for zero role activities, printed in Part~II
(Lemma~\pgrfwd{pgr:r:ir:lem:degree}{of source~22}). The
statement is not new.}

\begin{lemma}[Motzkin--Straus bound; source 25]\label{pgr:lem:ms}
If a simple graph $H$ has clique number $d\ge1$ and nonnegative vertex
weights $w_v$, then
\[
 \sum_{uv\in E(H)}w_uw_v\le\frac{d-1}{2d}\left(\sum_vw_v\right)^2.
\]
\end{lemma}
\begin{proof}
Normalize the total weight to one, the zero-total case being immediate. If
nonadjacent vertices $u,v$ both have positive weight, keep their sum fixed
and move all of it to the one with the larger weighted neighbor sum. The
objective does not decrease, since it is affine in this division, and its
positive support shrinks. Repetition leaves a clique of size $q\le d$. On
that clique the objective is $(1-\sum_vw_v^2)/2\le(1-1/q)/2\le(d-1)/(2d)$.
This is the classical compression argument of Motzkin and Straus
\cite{MS}.
\end{proof}

\begin{theorem}[Universal first inequality]\label{pgr:thm:first}
Let $D$ be a finite loopless directed relation with nonnegative independent
tail and head activities, and let $d\ge2$ be the actual degree of $\Ga_D$.
Then
\begin{equation}\label{pgr:eq:first}
 \gamma_1^2\ge\frac{2d}{d-1}\,\gamma_2 .
\end{equation}
Transitivity is not needed.
\end{theorem}
\begin{proof}
Discard every arc of zero weight $u_iv_j$; this changes no coefficient.
Make a graph whose vertices are the remaining directed arcs and whose edges
join arcs with four distinct physical endpoints. A clique is a
positive-weight matching, and every positive-weight support has such a
witness, so the clique number is exactly~$d$. Give arc $i\to j$ weight
$u_iv_j$; the total weight is $\gamma_1$. Every feasible support of size two
is witnessed by at least one edge of this graph, and every witness has the
prescribed support weight; a fixed pair of arcs determines only one ordered
support. The weighted edge sum is consequently at least $\gamma_2$, and
Lemma~\ref{pgr:lem:ms} bounds it by $(d-1)\gamma_1^2/(2d)$.
\end{proof}
\mnote{this is \prp, Theorem \texttt{mr:thm:first}, with the same proof;
the form at the surviving degree for vanishing activities is credited in
the merge note after that theorem. It is re-proved, by the same route and
without novelty claim, as source~16's Theorem ``Universal first
inequality'', source~25's Proposition ``first'' (unit activities, with the
lemma above), source~22's Proposition ``first'', source~11's Lemma
``first'' and the first comparisons of sources 10, 34, 06 and~35; sources % batch 77 (2026-10-02): source 35 added twice
10, 11, 34, 06 and~35 prove Lemma~\ref{pgr:lem:ms} by the same mass-shifting
argument, and source~12 gives its triangle-free case for actual degree two.
Those texts are replaced by pointers to this theorem. Source~25 adds, and
we keep, that counting matchings instead of supports would change the
second cubic gap, and that the comparison in this proof cannot be used
backwards to prove it.}

In degrees two, three and four the inequality reads
$\gamma_1^2\ge4\gamma_2$, $\gamma_1^2\ge3\gamma_2$ and
$3\gamma_1^2\ge8\gamma_2$. In particular every $\Ga_D$ of actual degree at
most two is rank-ULC and real-rooted; for preorders this is \prp,
Corollary \texttt{mr:cor:preorderfirst}.

\subsection{Notions and notation}\label{pgr:sec:notions}
The thirteen sources share one object but not one vocabulary. The following % ed. (2026-10-02): was "eleven"; batch 77: "twelve"
conventions hold throughout; every source section uses its own letters
inside its own scope, and the table below says where a letter means
different things.

\begin{center}\small
\begin{longtable}{p{28mm}p{60mm}p{56mm}}
\toprule
Notion & Meaning in this report & Not to be confused with\\\midrule
\endhead
$\Ga_D$, $\gamma_k$ & ordered \emph{disjoint} endpoint supports on the
  \emph{same} ground set, each counted once
  (Definition~\ref{pgr:def:support}) & the matching polynomial of $D$, and
  the support polynomial of the two-copy bipartite enlargement; the latter
  is the $h$-type setting of \prp{} Parts~XI--XIII and of \mrn{}\\
Preorder gamma & for a preorder $\tau$, $\Ga_\tau=\Ga_D$ with $D$ its
  off-diagonal relation (\prp, \texttt{gam:thm:main}) & the $h$-polynomial
  $h_\tau=\sum_k\gamma_kt^k(1+t)^{n-2k}$\\
Unit activities & all $u_i=v_i=1$ (Part~I) & ``unweighted'' in source~16
  means unit activities; its ``weighted quotient'' means positive integer
  class sizes\\
Vertex activities & one activity $x_i=u_i=v_i$ per vertex (archive~31,
  recorded as a corollary in Part~II; source~23) & independent role
  activities\\
Independent role activities & $u_i$ (tail) and $v_i$ (head), independent
  and nonnegative (sources 06, 10, 11, 12, 14, 15, 20, 22, 34) & a zero tail
  activity does not delete a head role (sources 06 and~22)\\
Actual degree & $\deg\Ga_D$ after zero activities are imposed & the
  ground-set size, a cover size, or a degree presumed before activities
  vanish\\
Order-$r$ ULC & log-concavity of $\gamma_k/\binom rk$ for a fixed $r\ge d$
  (source~16's exterior theorem, source~15's shore order) & rank-ULC; the
  two agree only when $r=d$\\
Role cover & sets $P$ (tails), $Q$ (heads) with every arc $i\to j$ having
  $i\in P$ or $j\in Q$; $P$ and $Q$ may overlap (sources 11, 15, 20, 22) &
  a physical vertex cover of the underlying graph (source~11), and the
  bipartite cover of \mrn{} with two vertices on each shore (the disjoint
  bipartite case)\\
Bipartite physical graph & the underlying undirected graph is bipartite;
  arcs may run both ways across (source~15) & a bipartite \emph{role}
  graph (tail and head copies), which every relation has\\
Template, attachment & Gallai--Edmonds normal form: $a=|A|$ attachments,
  exterior types with integer populations (sources 16, 23, 25) & cores
  with exterior types in general; a universal-sink cloud is the case of
  one exterior type\\
Rayleigh & coefficientwise (source~12, in formal activities); scalar on
  the nonnegative orthant (Wagner's rank-three theorem, sources 16, 22,
  23); all-real (the stability criterion, sources 20 and~22) & these are
  three different statements and are never interchanged\\
Real stable & nonvanishing when all variables lie in
  $\HH=\{\operatorname{Im}z>0\}$ & the half-plane property (positive real
  parts), equivalent for homogeneous polynomials up to rotation
  (source~22)\\\bottomrule
\end{longtable}
\end{center}

\paragraph{Letters with several meanings.}
\begin{itemize}
\item $P,Q$: tail and head sets of a role cover in sources 15, 16, 20 and~22,
  and in source~11's Corollary; $Q$ is also a polynomial in sources 16,
  22 and~23 ($Q(t)=1+Lt+Bt^2+Tt^3$, $Q=\Ga_{H-h}$), and $P$ the
  polynomial $\Ga_H-Q$ in source~22. The report \prp{} is never called $P$
  in the text.
\item $A,B$: in Part~III (sources 06, 10, 34) the moment parameters
  $A=\sqrt{p_2}$, $B=p_1-A$ of the sink activities; in source~22 $A_i$,
  $B_{ij}$, $C_0$ are three-core head sums; in sources 16 and~23 $A,B,C,D$
  and $T$ are coefficients of $Q$ and $R$; in source~16's role-cover
  material $B_P,L_P$ are basis and element sums.
\item $K$: the core $V\setminus I$ of the Gallai--Edmonds normal form in
  sources 16, 22 and~25; a number of two-element supports in source~23;
  $K(r,m)$, the sharp last-gap constant of source~14; the kernel
  $p_0q_0+p_1q_1$ in source~20's appendix.
\item $D$: the directed relation, and also the second cubic gap
  $\gamma_2^2-3\gamma_1\gamma_3$ in source~25, the set of deficient
  vertices in the Gallai--Edmonds decomposition (sources 16, 22, 25), and
  a coefficient of $R$ in source~16.
\item $e_j$ and $E_j$ (Part~III): elementary symmetric functions of the core
  tail activities and of the sink head activities. $\ee_j$ is the
  elementary symmetric polynomial itself.
\item $\rho$: the actual degree in the notation of sources 06, 10 and~34;
  $d$ in the others.
\end{itemize}
No symbol of a source has been renamed, except some letters of
source~35, which Section~\ref{pgr:s:fq:sec} renames and lists (added
2~October 2026, batch~77). Where a source's normalization
differs from the actual degree (source~16's Theorem on the exterior
construction and source~15's shore order), the source says so at the
place, and that sentence is kept.

\paragraph{Labels and file names.} Labels of this report carry the prefix
\texttt{pgr:}, with \texttt{pgr:u:}, \texttt{pgr:r:} and \texttt{pgr:s:}
for Parts I, II and~III and a two-letter infix per source (README). Each
source refers to files of its delivered package by their delivered paths;
the report README maps them to the shipped files in \texttt{code/} and
\texttt{data/}.

\part{Unit activities: finite preorders through actual degree four}\label{pgr:part:one}

This Part answers Research question \texttt{mr:q:gamma} of \prp{} for
unit activities through actual degree four. Section~\ref{pgr:u:d3:sec}
(source~25) proves rank-ULC for every finite preorder of actual degree at
most three, determines the least order of a preorder with a nonreal-rooted
cubic gamma polynomial, and gives a general bounded-core normal form.
Section~\ref{pgr:u:d4:sec} (source~16, the base of this report) proves
rank-ULC for every finite preorder of actual degree four. Its two by two
role-cover material is source~20's article and is printed once, in
Part~II. Section~\ref{pgr:u:t26:sec} (source~23) replaces the largest
certificate of source~16 by a smaller certificate-assisted argument.
Together:

\begin{theorem}[Sources 25 and~16]\label{pgr:u:thm:summary}
Every finite preorder whose unweighted gamma polynomial has actual degree
at most four is rank-ULC at that degree.
\end{theorem}
The proof is Theorem~\ref{pgr:u:d3:thm:main} for degrees at most three and
Theorem~\ref{pgr:u:d4:thm:middle} for degree four (which quotes the former).
Both are computer-assisted. Neither extends to activities other than one:
Part~II proves the independently weighted statement through degree three,
and the weighted statement in degree four, and every statement in degree
five or more, are open. Real-rootedness fails already in degree three
(Section~\ref{pgr:u:d3:sec:realroots}).

\section{Degree at most three (source 25)}\label{pgr:u:d3:sec}

Source~25 is \emph{Rank ultra log concavity for preorder gamma polynomials
of degree at most three} (14 pages as delivered; archive
\texttt{preorder-gamma-degree-three-package}). Its code, certificates and
receipts are shipped with the prefix \texttt{25-degree-three-}.

\paragraph{Abstract of source 25, as delivered.}

We prove that the gamma polynomial of every finite preorder of actual degree at most three is ultra-log-concave after normalization by the binomial coefficients of that degree. The essential new case is the second cubic Newton inequality. A Gallai--Edmonds decomposition reduces arbitrary ground-set size to four cases: components of order at most seven; cores of order at most seven with one pendant population; cores of order at most five with two attachment vertices; and seventeen three-vertex-cover templates. Exact finite enumeration and rational polynomial certificates settle these cases for all nonnegative integer populations. The certificate package includes independent coverage and support-count checks, rather than only numerical tests of small instances. The distinction from real-rootedness is essential: an explicit fifteen-element height-two poset has gamma polynomial $1+24z+162z^2+208z^3$ with negative discriminant. A separate exhaustive use of the same structural reduction proves that fifteen is the minimum order of such a nonreal-rooted cubic: every degree-three preorder on at most fourteen vertices is real-rooted, and at order fifteen the obstruction is unique up to isomorphism and order reversal. We also record a real-rooted infinite ordinal-sum family and describe the scope and reproducibility of the computation.

\subsection{Statement and mathematical setting}

A finite preorder $\tau$ is a reflexive transitive relation on a finite set $V$. Its off-diagonal directed relation has an arc $u\to v$ precisely when $u\ne v$ and $u\mathrel\tau v$. Antisymmetry is not assumed: distinct equivalent vertices give arcs in both directions. Write $G=\Comp(\tau)$ for the simple undirected graph with edge $uv$ whenever at least one of $u\to v$ and $v\to u$ occurs.

\begin{definition}[Disjoint directed supports]
A support of size $k$ is an ordered pair $(S,T)$ of disjoint subsets of $V$, with $|S|=|T|=k$, for which a bijection $\pi:S\longrightarrow T$ satisfies $s\to\pi(s)$ for every $s\in S$. Let $\gamma_k(\tau)$ be the number of these ordered pairs and put
\[
 \Ga_\tau(z)=\sum_{k\ge0}\gamma_k(\tau)z^k,\qquad \gamma_0(\tau)=1.
\]
A support is counted once even when several bijections realize it.
\end{definition}

This is the gamma polynomial identified in Part IV of the supplied root-polytope report \cite{PRP}, in particular its theorem labeled \texttt{gam:thm:main}. The present paper proves its statements directly for this explicitly defined polynomial. Thus the combinatorial result is independent of how that identification is implemented in the larger report. The convention differs both from counting ordinary matchings and from counting arbitrary row--column supports in two copies of the ground set: here the two endpoint sets must be disjoint in the original set.


\mnote{source~25's Lemma ``degree'' (for a preorder $\tau$:
$d=\deg\Ga_\tau=\nu(G)$, and all $\gamma_k$ with $0\le k\le d$ are
positive) is Lemma~\ref{pgr:lem:degree}, printed once in
Section~\ref{pgr:sec:object} with the same proof; for preorders it is
\prp, Theorem \texttt{gam:thm:deficiency}.}


\begin{definition}[Normalization by actual degree]
A polynomial $g(z)=\sum_{k=0}^d g_kz^k$ with nonnegative coefficients and no internal zeros is \emph{rank-ultra-log-concave}, or rank-ULC, if
\begin{equation}\label{pgr:u:d3:eq:ulc}
 \left(\frac{g_k}{\binom dk}\right)^2\ge
 \frac{g_{k-1}}{\binom d{k-1}}\frac{g_{k+1}}{\binom d{k+1}}
 \quad(1\le k<d).
\end{equation}
The parameter is the actual polynomial degree, not $|V|$ or $\lfloor|V|/2\rfloor$.
\end{definition}

\begin{theorem}[Main theorem]\label{pgr:u:d3:thm:main}
Every finite preorder $\tau$ with $\deg\Ga_\tau\le3$ has a rank-ULC gamma polynomial. In actual degree three this says
\begin{equation}\label{pgr:u:d3:eq:cubicgaps}
 \gamma_1^2\ge3\gamma_2,\qquad
 \gamma_2^2\ge3\gamma_1\gamma_3.
\end{equation}
In actual degree two the required inequality is the stronger
$\gamma_1^2\ge4\gamma_2$.
\end{theorem}

The proof has a classical structural input, elementary support arguments, and finite exact certificates. It is a computational proof with independently checked generators and identities; it is not a proof-assistant formalization. In particular, the conclusion quantifies over arbitrarily large finite preorders. No upper bound on the population sizes is used to prove Theorem~\ref{pgr:u:d3:thm:main}.

The original report \cite[Part XII, question \texttt{mr:q:gamma}]{PRP} asks whether every preorder gamma polynomial is log-concave or rank-ULC and identifies degree three beyond the bipartite case as the first unresolved target. The theorem closes that degree-three target. It does not settle degree four or higher, a weighted version, or real-rootedness. A global literature-priority claim is not made.

\mnote{this target is now closed further. Degree four with unit activities
is source~16's Theorem~\ref{pgr:u:d4:thm:middle}. The weighted version in
degree three, with independent tail and head activities, is source~22's
Theorem~\pgrfwd{pgr:r:ir:thm:main}{(source~22, Part~II)}; its case
$u_i=v_i=1$ is Theorem~\ref{pgr:u:d3:thm:main}, and the proof printed
below, by integer-population certificates, and source~22's proof, by a
ledger of role certificates, are independent routes. Real-rootedness fails
(Section~\ref{pgr:u:d3:sec:realroots}).}



\subsection{The universal first inequality}

The first inequality is already established in the report \cite[\texttt{mr:thm:first}]{PRP}. Source~25 includes its short proof to make the normalization and logical dependencies explicit; that proof is printed once, in Section~\ref{pgr:sec:object}.


\mnote{source~25's Lemma ``ms'' (the Motzkin--Straus bound, with its
compression proof) and Proposition ``first'' (for any finite off-diagonal
directed relation with $d=\nu(G)\ge2$, $\gamma_1^2\ge\frac{2d}{d-1}\gamma_2$;
transitivity is not needed) are Lemma~\ref{pgr:lem:ms} and the unit-weight
case of Theorem~\ref{pgr:thm:first}, printed once in
Section~\ref{pgr:sec:object}. Source~25's proof of the proposition is the
proof given there: different supports have disjoint sets of realizing
pairs of arcs, so $\gamma_2$ is at most the number of edges of the
arc-disjointness graph.}


It follows immediately that all degree-two cases are rank-ULC and real-rooted, and that the first degree-three inequality always holds. Degrees zero and one have no internal inequalities. The remaining problem is precisely
\begin{equation}\label{pgr:u:d3:eq:gapD}
 D=\gamma_2^2-3\gamma_1\gamma_3\ge0
 \quad\hbox{when }d=3.
\end{equation}
Counting matchings instead of supports would change $D$, and the comparison in the proof of Proposition~\ref{pgr:thm:first} cannot be used backwards to prove this second inequality.

\subsection{A canonical bounded core from Gallai and Edmonds}

For any finite simple graph $G$, let $D_G$ be the vertices omitted by at least one maximum matching, let $A=N(D_G)\setminus D_G$, and let $C=V(G)\setminus(A\cup D_G)$. The Gallai--Edmonds theorem states that the components $D_i$ of $G[D_G]$ are factor-critical, $G[C]$ has a perfect matching, and every maximum matching matches the vertices of $A$ into distinct $D_i$, with internal near-perfect matchings in all $D_i$ and a perfect matching in $C$. We use the theorem in the form proved by West \cite[Theorem 5]{West}. In particular,
\begin{equation}\label{pgr:u:d3:eq:budget}
 \nu(G)=|A|+\frac{|C|}{2}+\sum_i\frac{|D_i|-1}{2}.
\end{equation}

Assume from now on that $\nu(G)=3$. Let $I$ be the union of the singleton components of $G[D_G]$ and put $K=V(G)\setminus I$. We call $K$ the core. The set $I$ is independent and $N(I)\subseteq A$. Write $a=|A|$, $|C|=2c$, and $|D_i|=2d_i+1$ for each nonsingleton component, where $d_i\ge1$. Equation~\eqref{pgr:u:d3:eq:budget} gives
\begin{equation}\label{pgr:u:d3:eq:corebound}
 a+c+\sum_i d_i=3,\qquad
 |K|=a+2c+\sum_i(2d_i+1)\le9-2a.
\end{equation}
Here $2d_i+1\le3d_i$ and $2c\le3c$. Hence $a\in\{0,1,2,3\}$.

\subsubsection{Why the preorder boundary has no equivalence ambiguity}

An arbitrary vertex cover can cut through a preorder equivalence class. That would permit bidirectional exterior incidences and invalidate a fixed-sign exterior classification. The canonical choice above avoids this issue.

\begin{lemma}\label{pgr:u:d3:lem:exterior}
Every vertex of $I$ is a singleton preorder equivalence class. At each attachment vertex $v\in A$, all arcs incident with nonisolated vertices of $I$ have the same direction.
\end{lemma}
\begin{proof}
Equivalent vertices $u,v$ have identical relations to every third vertex by transitivity. In $G$ they are adjacent true twins, so their transposition is an automorphism. The sets $D_G,A,C$ are automorphism-invariant by their definitions. If an exterior vertex $u\in I$ were equivalent to $v\ne u$, then $v$ would also lie in $D_G$ and be adjacent to $u$, contradicting that $u$ is a singleton component of $G[D_G]$.

There are therefore no bidirectional arcs crossing from $I$ to $K$. If distinct $x,y\in I$ satisfied $x\to v\to y$, transitivity would force $x\to y$, contradicting independence. This excludes opposite directions at each fixed attachment.
\end{proof}

After discarding isolated exterior vertices, each exterior vertex is determined by a nonempty neighborhood $T\subseteq A$, together with the direction fixed at each member of $T$. There are at most $2^a-1$ possible types. Let $n_T\in\N$ be their populations.

\begin{lemma}[One-vertex compatibility is sufficient]\label{pgr:u:d3:lem:compatibility}
Fix a core preorder and a direction for each attachment. Keep precisely the nonempty neighborhood types for which adjoining one vertex gives a preorder. Then arbitrary nonnegative populations of all these types give a preorder.
\end{lemma}
\begin{proof}
A transitivity failure has a witness $u\to v\to w$ with $u\not\to w$. A witness involving at most one exterior vertex is already excluded by the core and single-type checks. A witness involving two distinct exterior vertices would have an attachment as its middle vertex and require both exterior directions there, which is forbidden. There are no exterior-to-exterior arcs, so these are all possibilities.
\end{proof}

Thus enlarging the family by allowing all nonnegative integer populations is safe. The canonical decomposition can impose additional neighborhood-surplus conditions, but there is no need to impose them for a positivity proof on the larger family.

\begin{center}
\small\begin{tabular}{cll}
\toprule
$a$ & Structural consequence & Verification needed\\
\midrule
0 & Every nontrivial component has at most 7 vertices & Finite real-rootedness lemma\\
1 & $|K|\le7$, one pendant-twin population & 2050 pairs of core polynomials\\
2 & $|K|\le5$, at most three exterior types & 1084 templates, 448 polynomial families\\
3 & $K=A$, $|K|=3$, at most seven types & 17 templates\\
\bottomrule
\end{tabular}
\end{center}

\subsubsection{A reusable normal form at every fixed degree}

\begin{theorem}[General bounded-core normal form]\label{pgr:u:d3:thm:normalform}
Let $\tau$ be a finite preorder with comparability matching number $r$. The same canonical construction gives an attachment set $A$, $a=|A|\le r$, and an independent exterior $I$ such that
\[
 |K|\le3r-2a,\qquad \nu(G[K\setminus A])=r-a.
\]
Every exterior point is a singleton equivalence class, every attachment has one fixed exterior direction, and there are at most $2^a-1$ nonisolated exterior types. The gamma coefficients have the finite expansion \eqref{pgr:u:d3:eq:binomialkernel}, of total population degree at most $a$.

Conversely, fix a core of size at most $3r-2a$, a marked set $A$ of size $a$, fixed directions, and the compatible exterior types. If $\nu(G[K\setminus A])\le r-a$, every population blow-up has matching number at most $r$. Thus all degree-$r$ preorders are represented by finitely many bounded-core polynomial families with unbounded nonnegative integer populations.
\end{theorem}
\begin{proof}
Replace $3$ by $r$ in \eqref{pgr:u:d3:eq:budget}. Writing $r=a+c+\sum d_i$ gives the same bound $|K|\le a+3(c+\sum d_i)=3r-2a$. The graph on $K\setminus A$ is the disjoint union of $G[C]$ and the nonsingleton factor-critical components, with matching number $c+\sum d_i=r-a$. Lemmas~\ref{pgr:u:d3:lem:exterior} and \ref{pgr:u:d3:lem:compatibility} do not depend on $r$. In any matching, each selected exterior vertex requires its own attachment, proving the quota bound and \eqref{pgr:u:d3:eq:binomialkernel}. For the converse, at most $a$ matching edges touch $A$ and every other edge lies in $K\setminus A$, so the matching size is at most $a+\nu(G[K\setminus A])\le r$. For fixed $r$, there are only finitely many core relations, marked sets, and direction choices.
\end{proof}

The restriction on $\nu(G[K\setminus A])$ matters in the converse: the core-size bound alone does not imply matching rank at most $r$. The normal form reduces any fixed-rank coefficient inequality to finitely many population-polynomial families; it does not assert that arbitrary integer-polynomial positivity is decidable or that higher-degree ULC holds. At specializations where the actual degree drops, the normalization must be changed accordingly.

\subsection{Finite enumeration and exact positivity methods}\label{pgr:u:d3:sec:methods}

This section specifies what the computer checks and why the checks have the needed quantifiers. The data and source programs in the accompanying archive are part of the certificate.

\subsubsection{Covering all finite preorder cores}

Every finite preorder is uniquely a poset of equivalence classes with positive block sizes. Choose a linear extension of the quotient poset and label its $q$ elements $0,\ldots,q-1$ in that order. Every strict comparison then points forward. We use two complete quotient enumerations:
\begin{enumerate}[itemsep=2pt]
\item Enumerate all subsets of the $\binom q2$ forward pairs and keep the transitive relations.
\item Recursively add a last maximal element. Its predecessors must be an order ideal of the previous poset, and each such ideal gives one extension.
\end{enumerate}
For every quotient, enumerate all positive ordered compositions of each allowed total core size into $q$ block sizes. Expand each block to mutually equivalent vertices and each strict comparison to all corresponding directed arcs. Every preorder is represented up to relabeling. Repetitions are allowed; the recorded counts are not counts of isomorphism classes or of labeled preorders.

The naturally labeled quotient-poset counts for $q=1,\ldots,7$ are
\[
 1,\ 2,\ 7,\ 40,\ 357,\ 4824,\ 96428.
\]
The expanded-core counts for total size $n=1,\ldots,7$ are
\begin{equation}\label{pgr:u:d3:eq:corecounts}
 1,\ 3,\ 12,\ 68,\ 568,\ 7090,\ 131645.
\end{equation}
For a fixed $n$, the count equals $\sum_q Q_q\binom{n-1}{q-1}$, where $Q_q$ is the corresponding quotient count.

\subsubsection{Counting supports without counting matching multiplicities}

For each ordered pair of disjoint equally sized endpoint sets, support feasibility is a Boolean perfect-matching test in the directed incidence relation $S\to T$. The package cross-checks three implementations: recursive matching search; Hall's condition for every subset of $S$; and bottom-up construction of directed matchings with endpoint pairs stored in a set. The last method explicitly discards multiple realizations of the same support.

For a fixed core and exterior types, let $q_T$ be the number of selected exterior vertices of type $T$. For fixed labeled representatives of those selected vertices, enumerate their source/target roles and the unused/source/target roles of core vertices; retain a role assignment exactly when a realizing matching exists. If this gives $c_{k,\boldsymbol q}$ feasible supports, then
\begin{equation}\label{pgr:u:d3:eq:binomialkernel}
 \gamma_k(\boldsymbol n)=
 \sum_{\boldsymbol q} c_{k,\boldsymbol q}
 \prod_T\binom{n_T}{q_T}.
\end{equation}
A support selects its exterior subsets uniquely, so this is an exact identity. Every used exterior vertex needs a different attachment in a realizing matching. Thus $\sum_Tq_T\le a$, making the sum finite regardless of the populations. In particular, $\gamma_k$ has total population degree at most $a$, and $\gamma_1$ has degree at most one.

Multiplication in the binomial basis can be done with integer arithmetic using
\begin{equation}\label{pgr:u:d3:eq:binprod}
 \binom na\binom nb=
 \sum_{j=0}^{\min(a,b)}
 \frac{(a+b-j)!}{j!(a-j)!(b-j)!}\binom n{a+b-j}.
\end{equation}
This counts pairs of subsets by the size $j$ of their intersection. Taking products of this identity over the types gives exact binomial-basis expansions for $D$.

\subsubsection{Finite certificates for an infinite population domain}

A polynomial with nonnegative coefficients in the multivariate binomial basis is nonnegative on $\N^r$. For other polynomials we partition $\N^r$ according to the set $F$ of positive coordinates and substitute
\begin{equation}\label{pgr:u:d3:eq:faces}
 n_i=0\ (i\notin F),\qquad n_i=1+x_i\ (i\in F),\qquad x_i\ge0.
\end{equation}
Every nonnegative integer population vector lies on exactly one such face. A certificate is an exact identity expressing the substituted polynomial as a nonnegative rational combination of nonnegative monomials and terms $x^{\boldsymbol m}Z(x)^2$ with $Z\in\Q[x]$.

Some square terms are stored compactly as three monomials $ux^\alpha+vx^\beta+wx^\eta$, where $u,v>0$, $w<0$, $\alpha+\beta=2\eta$, and $w^2=4uv$. Factoring $x^{\min(\alpha,\beta)}$ leaves a square: the remaining endpoint exponents are even. All exponent vectors are nonnegative integers. General square terms include the rational polynomial $Z$ as additional data. The verifier expands every term, checks every sign and exponent, and compares coefficients exactly.

Numerical linear programming and eigenvector-based column generation helped discover some certificates. They are not used to verify positivity. Verification contains no floating-point tolerance, approximate roots, or finite-population extrapolation. The face substitution proves an integer-population statement; it does not assert positivity for coordinates strictly between zero and one.

\subsection{No attachments and one attachment}

\subsubsection{The finite real-rootedness lemma}

\begin{lemma}[Exact finite lemma]\label{pgr:u:d3:lem:small}
For every preorder on at most seven vertices, $\Ga_\tau$ has only real negative roots, with multiplicities allowed; the constant polynomial is included.
\end{lemma}
\begin{proof}[Computational proof]
The complete enumeration of Section~\ref{pgr:u:d3:sec:methods} gives the core counts \eqref{pgr:u:d3:eq:corecounts}. There are no supports larger than three. A quadratic $1+bz+cz^2$ is real-rooted exactly when $b^2-4c\ge0$. A cubic $1+bz+cz^2+ez^3$, $e>0$, is real-rooted exactly when
\begin{equation}\label{pgr:u:d3:eq:disc}
 \Delta=b^2c^2-4c^3-4b^3e-27e^2+18bce\ge0.
\end{equation}
Every enumerated coefficient vector passes the applicable exact integer test. Positivity of the coefficients excludes nonnegative roots. The generator and independent Hall checker are supplied in \texttt{code/small-check/}; representatives of all distinct gamma vectors are included. Direct enumeration of all reflexive relations through four vertices gives the known labeled-preorder counts $1,4,29,355$ and exactly the same coefficient sets. These checks supplement the coverage proof, rather than replace it.
\end{proof}

The literature already reports computer verification through seven vertices \cite[Section 5]{AC}; Lemma~\ref{pgr:u:d3:lem:small} is independently reproduced here, not claimed as a new size range. The supplied final enumerator also verifies every degree-at-most-three preorder on eight vertices. It skips degree-four cases and is not evidence for a claim about their gamma polynomials.

If $a=0$, there are no edges from $D_G$ to $C$ or between different $D_i$. The total order of $C$ is at most six, and each nonsingleton $D_i$ has at most seven vertices by the matching budget. Thus every nontrivial connected component has at most seven vertices. Support polynomials multiply over graph components: restricting a realizing matching gives equally sized tail/head sets in each component, and taking unions reverses this restriction. Lemma~\ref{pgr:u:d3:lem:small} therefore makes the whole polynomial real-rooted. Newton's inequalities imply rank-ULC with its total actual degree. This proves the $a=0$ branch without a separate ULC convolution argument.

\subsubsection{One attachment and the support identity}

Suppose $a=1$, with attachment $v$. The core has at most seven vertices. After deleting isolated exterior vertices, every exterior vertex is a pendant neighbor of $v$, and all have the same direction. Reverse all arcs if needed to orient them $v\to x$. For $m>0$, transitivity implies that no other core vertex has an arc into $v$, since $u\to v\to x$ would force an additional neighbor of $x$. Thus $v$ is a singleton minimal element of the core preorder.

If the exterior population is $m$, a support either avoids the exterior or uses exactly one exterior vertex. In the latter case the matching must use the pendant edge, and deletion of its endpoints gives a support in $K-v$. Conversely, every such support extends in $m$ ways. Consequently
\begin{equation}\label{pgr:u:d3:eq:pendant}
 \Ga_m(z)=\Ga_K(z)+mz\Ga_{K-v}(z).
\end{equation}
If $m>0$, a matching in $K-v$ extends by a pendant edge, so $\nu(K-v)\le2$. The case $m=0$ is already Lemma~\ref{pgr:u:d3:lem:small}.

Write $\Ga_K=1+p_1z+p_2z^2+p_3z^3$ and $\Ga_{K-v}=1+q_1z+q_2z^2$. The two cubic gaps in \eqref{pgr:u:d3:eq:pendant} are quadratics in $m$ with coefficient triples $(A,B,C)$:
\begin{align}
 D_1(m)&: &&(1,\ 2p_1-3q_1,\ p_1^2-3p_2),\label{pgr:u:d3:eq:pendant1}\\
 D_2(m)&: &&(q_1^2-3q_2,\ 2p_2q_1-3p_1q_2-3p_3,\ p_2^2-3p_1p_3).
 \label{pgr:u:d3:eq:pendant2}
\end{align}
Every certificate is either $A,B,C\ge0$, or $A>0$ and $4AC-B^2\ge0$, using
\begin{equation}\label{pgr:u:d3:eq:quadraticcert}
 4A(Am^2+Bm+C)=(2Am+B)^2+(4AC-B^2).
\end{equation}
These prove nonnegativity for every real $m\ge0$, a strengthening of what the preorder application requires.

The complete marked-core enumeration has $198469$ valid marked expansions and $2050$ distinct tuples $(p_1,p_2,p_3,q_1,q_2)$. By total core size the marked counts are
\[
 1,\ 3,\ 16,\ 115,\ 1160,\ 16786,\ 180388.
\]
All $2050$ first gaps and all $2050$ second gaps pass: respectively $1114$ and $1697$ use nonnegative coefficients, and the remaining $936$ and $353$ use \eqref{pgr:u:d3:eq:quadraticcert}. Where $p_3=q_2=0$, the actual degree can be two for all $m$; all $136$ such tuples also certify $(p_1+m)^2-4(p_2+mq_1)\ge0$. At $m=0$ with $p_3=0$, the corresponding quadratic normalization is checked separately.

The independent C++ checker generates quotient posets by ideal extension and constructs supports by endpoint-mask dynamic programming. It reproduces every one of the $2050$ tuple multiplicities, not merely one representative or the total count. A separate Python checker validates representatives and all $4236$ stored quadratic certificates. This closes $a=1$.

\subsection{Two attachments}

Let $a=2$, so $|K|\le5$. Label the attachments $v_0,v_1$. Exterior neighborhoods are the nonempty subsets of this pair, encoded as masks $1,2,3$. For every core preorder, pair of attachments, and two signs, the generator keeps exactly the types allowed by Lemma~\ref{pgr:u:d3:lem:compatibility}. It then identifies templates under permutations of the core respecting the attachment pair, interchange of the two attachments, and simultaneous reversal of all arcs.

Every resulting blow-up has matching number at most three, even when it is not a canonical Gallai--Edmonds instance. Indeed, a matching using $r$ exterior edges has $r\le2$ and size at most
\[
 r+\left\lfloor\frac{5-r}{2}\right\rfloor\le3.
\]
There are at most three population variables. By \eqref{pgr:u:d3:eq:binomialkernel}, every gamma coefficient has population degree at most two and $D$ has degree at most four.

\begin{proposition}[Exact two-attachment certificate]\label{pgr:u:d3:prop:a2}
For every core preorder of order at most five, every pair of attachment vertices, every fixed direction choice, and all nonnegative integer populations of the compatible exterior types, $D\ge0$.
\end{proposition}
\begin{proof}[Computational proof]
The complete generator produces $1084$ structural templates, with $5,23,141,915$ templates for core orders $2,3,4,5$, respectively. They yield $448$ distinct full gamma-polynomial families. Grouping is used only when all coefficient polynomials agree in their ordered population variables. Of these, $284$ have a coefficientwise nonnegative binomial expansion of $D$. Each of the other $164$ is supplied with all its zero/positive face certificates, totaling $994$ faces. Each face identity has the exact nonnegative form described in Section~\ref{pgr:u:d3:sec:methods}. Thus every nonnegative integer population vector is covered.

An independent audit enumerates all labeled preorders through core order five, obtaining $4,29,355,6942$ at orders $2,3,4,5$, respectively. It reproduces the complete set of $1084$ canonical templates, reconstructs support polynomials using directed matchings deduplicated by endpoint sets, recomputes $D$, and checks every rational face identity. The audit does not call the generator's helper routines. An additional independent Hall implementation checks $19084$ labeled blow-ups on the grid $\{0,1,2,3\}$ and recovers the binomial coefficients by finite differences; the known degree bound makes this polynomial reconstruction exact. This is an additional support-count check, not the proof of positivity on the unbounded domain.
\end{proof}

The files \texttt{templates.json}, \texttt{gap2\_certificates.json}, the exact verifier, and both independent checking routes are in \texttt{code/two-attachment/}. A template with no allowable exterior type contributes only its small core and is covered by Lemma~\ref{pgr:u:d3:lem:small}; omitting it from the nonempty-type export loses no case. If a specialization has actual degree at most two, Proposition~\ref{pgr:thm:first} supplies the correct degree-two normalization. Hence Proposition~\ref{pgr:u:d3:prop:a2} closes the $a=2$ branch.

\subsection{Three attachments and seventeen templates}\label{pgr:u:d3:sec:a3}

For $a=3$, equation~\eqref{pgr:u:d3:eq:budget} forces $C=\varnothing$ and every $D_i$ to be a singleton. Thus $K=A$ consists of exactly three vertices. The canonical exterior has singleton equivalence classes by Lemma~\ref{pgr:u:d3:lem:exterior}. Every attachment has an exterior neighbor because a maximum matching matches all of $A$ into the singleton $D_i$.

There are $29$ labeled three-element preorders. Test all eight active sign vectors and all seven nonempty neighborhood masks. Requiring every attachment to have an allowed neighbor gives $130$ labeled templates and $17$ classes under core permutations and order reversal. The full list is below. A minus sign denotes core-to-exterior arcs, a plus sign exterior-to-core arcs. Mask $t$ encodes the subset of $\{0,1,2\}$ whose bits are one. The listed core arcs exclude loops and include transitive arcs.

\begin{center}\small
\begin{tabular}{rlll}
\toprule
Class & Core arcs & Signs & Allowed masks\\\midrule
0 & $\varnothing$ & $---$ & $1,2,3,4,5,6,7$\\
1 & $\varnothing$ & $--+$ & $1,2,3,4$\\
2 & $01$ & $---$ & $1,3,4,5,7$\\
3 & $01$ & $--+$ & $1,3,4$\\
4 & $01$ & $-+-$ & $1,2,3,4,5$\\
5 & $01,02$ & $---$ & $1,3,5,7$\\
6 & $01,02$ & $--+$ & $1,3,4,5$\\
7 & $01,02$ & $-++$ & $1,2,3,4,5,6,7$\\
8 & $01,02$ & $+++$ & $2,4,6,7$\\
9 & $01,10$ & $---$ & $3,4,7$\\
10 & $01,10$ & $--+$ & $3,4$\\
11 & $01,02,12$ & $---$ & $1,3,7$\\
12 & $01,02,12$ & $--+$ & $1,3,4,5,7$\\
13 & $01,02,10,12$ & $---$ & $3,7$\\
14 & $01,02,10,12$ & $--+$ & $3,4,7$\\
15 & $01,02,10,12$ & $+++$ & $4,7$\\
16 & $01,02,10,12,20,21$ & $---$ & $7$\\\bottomrule
\end{tabular}
\end{center}

\subsubsection{The bipartite template}

Class 0 is a bipartite graph with a source shore of size three and a target shore consisting of the exterior. Its directed disjoint supports are precisely the ordinary bipartite matching supports. The ambient-binomial normalization theorem of R\"ohrle and Ulirsch \cite[Theorem A and Example 2.4]{RU} applies: for a relation on disjoint shores of sizes $m,n$, the support numbers normalized by $\binom{\min(m,n)}k$ are log-concave. If the actual degree is three here, the smaller shore has size three, so this is exactly $D\ge0$. Lower actual degrees are already settled by Proposition~\ref{pgr:thm:first}. \mnote{class~0 is also a case of \prp, Theorem
\texttt{lor:thm:weightedULC} (two-shore ultra-log-concavity normalized by
the smaller shore), proved there with the same transversal-matroid
device.}

For clarity, the relation support numbers fit the published theorem by assigning an independent indeterminate to each edge of the incidence matrix. A minor is nonzero precisely when its row and column sets support a matching: different matching terms have different monomials and cannot cancel. These minors define a representable bimatroid. No theorem about normalizing an arbitrary bimatroid by its actual rank is being invoked. This published ambient-normalization theorem is the only non-elementary polynomial-inequality input for the seventeen-template step.

\subsubsection{Compact product and extension certificates}

In classes $1,3,10$, the polynomial factors as
\[
 \Ga(z)=(1+Az+Bz^2)(1+Cz),\qquad C=n_4.
\]
Consequently
\begin{equation}\label{pgr:u:d3:eq:productcert}
 D=(B-AC/2)^2+\tfrac34 C^2(A^2-4B).
\end{equation}
The remaining expressions are
\begin{align*}
\text{class 1: }\quad
&A=n_1+n_2+2n_3,\quad
B=n_1n_2+(n_1+n_2)n_3+\binom{n_3}{2},\\
&A^2-4B=(n_1-n_2)^2+2n_3^2+2n_3;\\[2pt]
\text{class 3: }\quad
&A=n_1+2n_3+1,\quad B=n_1n_3+\binom{n_3}{2},\\
&A^2-4B=n_1^2+2n_1+2n_3^2+6n_3+1;\\[2pt]
\text{class 10: }\quad
&A=2n_3+2,\quad B=\binom{n_3}{2},\quad
A^2-4B=2n_3^2+10n_3+4.
\end{align*}
These are nonnegative on the required domain. For class 9, subtract the class-10 gap. All coefficients of the difference are nonnegative except possibly those in the pure $n_7$ contribution. That whole univariate contribution is exactly
\begin{equation}\label{pgr:u:d3:eq:class9}
 4\binom{n_7}{1}+41\binom{n_7}{2}
 +57\binom{n_7}{3}+18\binom{n_7}{4},
\end{equation}
which is nonnegative for integer $n_7\ge0$. The coefficient assertion and identity are verified exactly in \texttt{verify\_cover3\_compact.py}. Classes $11,13,15,16$ have nonnegative multivariate-binomial expansions of $D$ directly.

\subsubsection{The complete face certificates}

Classes $5,6,8,12,14$ have a certificate for every zero/positive face using monomials and monomial-weighted binomial squares. Class 2 has such certificates on $31$ of its $32$ faces. The remaining face has only $n_3,n_4$ positive; its gap is exactly that of class 3 at $n_1=0$, so \eqref{pgr:u:d3:eq:productcert} applies.

Classes 4 and 7 need some general rational square polynomials as well. Class 4 has $23$ binomial-square face certificates and $9$ general-square certificates, covering all $32$ faces. Class 7 has $73+55=128$ certificates, covering all its faces. The hardest saved class-7 identity has $210$ nonzero square/monomial weights; the checker expands its $330$ coefficient equations over the rationals.

The independent script \texttt{audit\_cover3.py} separately enumerates the $29$ core preorders, reproduces the $130$ signed templates and all $17$ canonical keys, and reconstructs every coefficient in \eqref{pgr:u:d3:eq:binomialkernel} by Hall's condition on ordered endpoint masks. It recomputes every gap from the gamma coefficients and verifies the face partitions, missing-face reductions, exponent domains, and duplicate guards. The three exact identity verifiers then check all saved rational identities. Thus the certificates are linked to the intended support counts, not merely to an unchecked polynomial export. This proves $D\ge0$ throughout the $a=3$ branch.

\subsection{Completion of the main theorem}

\begin{proof}[Proof of Theorem~\ref{pgr:u:d3:thm:main}]
Lemma~\ref{pgr:lem:degree} identifies the degree and excludes internal zeros. Degrees zero and one are immediate; degree two follows from Proposition~\ref{pgr:thm:first}. In degree three the same proposition proves the first inequality. Apply the Gallai--Edmonds decomposition to the comparability graph. Equation~\eqref{pgr:u:d3:eq:corebound} gives exactly four possibilities. The $a=0$ branch follows from the finite real-rootedness lemma and component multiplication; the $a=1$ branch from \eqref{pgr:u:d3:eq:pendant} and its complete quadratic certificates; the $a=2$ branch from Proposition~\ref{pgr:u:d3:prop:a2}; and the $a=3$ branch from Section~\ref{pgr:u:d3:sec:a3}. These prove the second inequality in every case. All population sizes are unrestricted nonnegative integers throughout these four arguments.
\end{proof}

The reduction is essential to the quantifier. A finite scan of all preorders up to a fixed number of vertices would not prove the theorem. Here the graph structure bounds the core, and algebraic certificates handle the unbounded exterior populations. The finite base computation for $a=0$ is sufficient because each component is provably small, even though isolated components may occur in arbitrarily large number.

\subsection{Real roots and an explicit obstruction}\label{pgr:u:d3:sec:realroots}

Rank-ULC is strictly weaker than real-rootedness even in actual degree three. Take three minimal elements $u_0,u_1,u_2$ and three disjoint four-element sets $X,Y,Z$ of maximal elements. Let the strict relations be exactly
\[
 u_0<X\cup Y,\qquad u_1<X\cup Z,\qquad u_2<Y\cup Z.
\]
This is a height-two poset on fifteen vertices. Its three minimal vertices form a vertex cover, and a matching saturating them exists, so the actual degree is three.

There are $24$ arcs. For each choice of two minimal vertices, all pairs of maximal vertices are feasible except pairs contained in either of the two four-element types that see only one chosen minimal vertex. Thus
\[
 \gamma_2=3\left(\binom{12}{2}-2\binom42\right)=162.
\]
A triple of maximal vertices fails to match all three minima exactly when all three lie in the same four-element type. Hall's condition has no other obstruction because every maximal vertex sees two minima. Hence
\[
 \gamma_3=\binom{12}{3}-3\binom43=208.
\]
Therefore
\begin{equation}\label{pgr:u:d3:eq:counterexample}
 \Ga(z)=1+24z+162z^2+208z^3,\qquad \Delta=-2592<0.
\end{equation}
Its two cubic Newton gaps are $90$ and $11268$, both positive, consistent with Theorem~\ref{pgr:u:d3:thm:main}. A direct endpoint-support enumerator independently reproduces \eqref{pgr:u:d3:eq:counterexample}; no population formula is used in that check.

\mnote{the same example, with the same numbers, is computed in source~22
(Section ``A cubic with nonreal zeros'', where it is noted that it lies in
the pure three-tail family of source~22's three-core theorem and needs no
role asymmetry), in source~16 (``The preceding research \dots{} exhibited
the nonreal-rooted cubic'') and in archive~31. It is printed once, here.}


\subsubsection{The exact minimum order in degree three}

\begin{theorem}[Minimum order and uniqueness]\label{pgr:u:d3:thm:minimum}
Every finite preorder on at most fourteen vertices whose gamma polynomial has actual degree three is real-rooted. Among fifteen-vertex preorders of actual degree three with a nonreal gamma zero, there are exactly two ordinary isomorphism classes, represented by the height-two poset in \eqref{pgr:u:d3:eq:counterexample} and its dual. Equivalently, there is one class modulo isomorphism and global order reversal.
\end{theorem}
\begin{proof}[Computer-assisted proof with global coverage]
Apply the canonical four-branch reduction. The $a=0$ branch is real-rooted at every allowed order by Lemma~\ref{pgr:u:d3:lem:small} and component multiplication. For $a=1,2,3$, enumerate every population specialization whose effective total order is at most fifteen and whose cubic coefficient is positive. Test the exact integer discriminant \eqref{pgr:u:d3:eq:disc}. The complete counts are
\begin{center}
\begin{tabular}{lrrr}
\toprule
Branch & Evaluations & Actual cubics & Negative discriminants\\\midrule
One attachment & 19430 & 17785 & 0\\
Two attachments & 84191 & 58800 & 0\\
Three attachments & 128726 & 121510 & 1\\\bottomrule
\end{tabular}
\end{center}
The sole negative case has total order fifteen, class 0 of Section~\ref{pgr:u:d3:sec:a3}, and populations $(n_1,\ldots,n_7)=(0,0,4,0,4,4,0)$. It is exactly \eqref{pgr:u:d3:eq:counterexample}. An independently written scanner regenerates every per-size count and the complete negative-case list using arbitrary-precision integers.

There are two important coverage points. In the two-attachment scan we retain all $1084$ structural templates and their actual core sizes, not only the $448$ polynomial families. In the one-attachment scan we group by $(p_1,p_2,p_3,q_1,q_2)$ but retain a \emph{minimum-order} realization of every tuple. The complete marked-core generator explicitly minimizes this size over all occurrences; all $2050$ stored minima are verified by a fresh enumeration. Replacing a core by this minimum-size representative preserves \eqref{pgr:u:d3:eq:pendant} and can only reduce total order. Hence any example through a bound appears at some scanned order through that bound. Isolated vertices may be deleted or added without changing gamma. The displayed counts are therefore representation counts, not a per-order census of polynomials or preorder isomorphism classes.

For uniqueness, any nonreal cubic on at most fifteen vertices must lie in the canonical $a=3$ branch, after deleting isolates. Up to a permutation of the three core vertices and possibly global duality, its core, signs, and populations occur in the complete seventeen-template scan. The sole negative specialization has effective order fifteen, so no isolates were deleted. Its populations determine the directed relation up to permutations within types. Thus the preorder is the displayed poset or its dual. They are not isomorphic: one has three minimal and twelve maximal elements, and its dual has twelve minimal and three maximal elements.
\end{proof}


\mnote{this theorem supplies the minimality statement that source~22
explicitly does not make (``No minimal-size claim is made''). It does not
conflict with the nonreal gamma polynomials of \prp: the eight-vertex
example $\Theta_3$ there has gamma degree four, and the only degree-three
member of that family is real-rooted.}

This theorem concerns actual degree three only. It makes no claim about higher-degree preorders on at most fourteen vertices. Its bounded discriminant scan is logically separate from the all-population second-gap certificates used for Theorem~\ref{pgr:u:d3:thm:main}. The scripts and receipts are in \texttt{code/minimum/}.


\subsection{An infinite family with three negative roots}\label{pgr:u:d3:sec:ordinal}

A complementary positive result holds for a large antichain whose complement has three vertices. Let $\mathcal A_m$ be an antichain of size $m$, and let $\oplus$ denote ordinal sum, with all elements of an earlier summand strictly below all elements of a later one.

\begin{theorem}\label{pgr:u:d3:thm:ordinal}
If $m\ge3$ and $P,Q$ are arbitrary finite preorders with $|P|+|Q|=3$, allowing either to be empty, then $\Ga_{P\oplus\mathcal A_m\oplus Q}$ has degree three and three distinct negative real roots. Every ordinal sum of antichains of actual gamma degree three also has three distinct negative real roots.
\end{theorem}
\begin{proof}
Every edge meets the three-element complement of $\mathcal A_m$, and its vertices can be matched to three distinct antichain vertices, so the degree is three. Reversal reduces the possible splits to two forms.

For $\mathcal A_m\oplus R$, $|R|=3$, let $e$ be the number of off-diagonal arcs in $R$ and let $t$ be the number of its vertices with at least one outgoing off-diagonal arc. Then
\begin{equation}\label{pgr:u:d3:eq:outer}
 (\gamma_1,\gamma_2,\gamma_3)=
 \left(3m+e,\ 3\binom m2+tm,\ \binom m3\right).
\end{equation}
The feasible $(e,t)$ are $(0,0),(1,1),(2,1),(2,2),(3,2),(4,2),(4,3),(6,3)$. For $1\oplus\mathcal A_m\oplus R_2$, $|R_2|=2$, let $j\in\{0,1,2\}$ be its number of off-diagonal arcs. Then
\begin{equation}\label{pgr:u:d3:eq:middle}
 (\gamma_1,\gamma_2,\gamma_3)=
 \left(3m+2+j,\ 5\binom m2+(j+1)m,\ 3\binom m3\right).
\end{equation}
To justify \eqref{pgr:u:d3:eq:outer}, a size-two support with two antichain points chooses two of the three core heads, while one with a single antichain point chooses a core tail with an outgoing arc. A size-three support uses all three core heads. For \eqref{pgr:u:d3:eq:middle}, two chosen antichain points admit five size-two role patterns: one with both as tails, and four with one tail and one head. One antichain point gives $j+1$ patterns, and three give three choices for the head matched from the lower singleton. These are Boolean support counts.

For all eleven forms, put $x=m-3$. The coefficients of $12\Delta/(x+3)$ are listed in Appendix~\ref{pgr:u:d3:app:disc}; every one is positive. Thus $\Delta>0$ for $m\ge3$. The cubic has three distinct real roots and, by positive coefficients, they are negative.

For an ordinal sum of antichains with total size $n$ and maximum layer size $M$, its comparability graph is complete multipartite and has matching number
$\min(\lfloor n/2\rfloor,n-M)$. If $M\ge n/2$, match all outside vertices into the largest layer. If $M<n/2$, pairing vertices from two largest remaining parts constructs a matching missing at most one vertex; the invariant that no part exceeds the ceiling of half the remaining total is preserved. Degree three consequently means either $n\in\{6,7\}$ and $M\le n-3$, or $n\ge8$ and $M=n-3$. Whenever $M=n-3$, the preceding result applies. The remaining compositions have $n=6$ with all layers at most two, or $n=7$ with all layers at most three: respectively $13$ and $44$ cases. Their exact cubic discriminants are positive in the supplied table \texttt{layered\_certificates.json}. An independent layer-quota recurrence and direct support enumeration verify all $57$ cases.
\end{proof}

Two specializations are already in the literature. The two-antichain case $(e,t)=(0,0)$ is covered by Athanasiadis and Chapoton \cite[Proposition 9.2]{AC}. The case $(e,t)=(6,3)$, a three-element equivalence block above singleton leaves, is covered by Athanasiadis, Xiao, and Yan \cite[Remark 2.3]{AXY}. Neither specialization is claimed as new. The eleven-template statement is included as a proved complementary family; a complete literature-priority audit for its precise scope has not been performed.

\subsection{Reproducibility and scope}

The archive contains the editable LaTeX source, this PDF, exact source code, certificate data, checking receipts, and a checksum manifest. It contains no required compiled binaries or numerical-optimizer state. The top-level \texttt{README.md} gives the dependency versions and commands. The primary checks need Python 3 with SymPy and a C++17 compiler. No optimization package is required to verify the saved certificates.

The command \texttt{python3 verify\_all.py} runs the exact checking routes and writes fresh logs under \texttt{receipts/}. It checks the finite-core lemma; the full marked-core enumeration and all one-population quadratic identities; independent two-attachment coverage and all its rational certificates; the seventeen three-cover templates and their exact identities; the ordinal-sum family; the bounded scans proving Theorem~\ref{pgr:u:d3:thm:minimum}; and the explicit nonreal-root example. The finite generators can also be rerun separately using the commands in the README. The main theorem depends on the published Gallai--Edmonds theorem and the class-0 ambient-normalization theorem as clearly stated in the text, not on a numerical solver.


\mnote{file names in this section are those of source~25's delivered
package; the files are shipped as \texttt{code/25-degree-three-\dots} and
\texttt{data/25-degree-three-\dots}, and the README maps the names and
gives rerun instructions. At intake every step of \texttt{verify\_all.py}
reproduced its recorded outputs (on a copy, compared modulo line endings
and timing fields) except the order-eight census of
\texttt{small-check/exhaust}, whose optimized build crashed after finishing
its enumeration and was not rerun; the C++ sources also need
\texttt{-include cassert} with GCC~16. Of the research directions below,
the first is answered by source~16 (Section~\ref{pgr:u:d4:sec}) and the
fourth, in degree three, by source~22
(\pgrfwd{pgr:r:ir:sec}{Part~II}); the second, third and fifth remain
open.}

The following directions remain meaningful beyond this result.
\begin{enumerate}[itemsep=3pt]
\item \textbf{Actual degree four.} The same matching-budget argument gives bounded canonical cores, but more exterior types and higher-degree inequalities. A direct certificate approach must control that growth rather than infer higher-rank behavior from the cubic case.
\item \textbf{A conceptual second-gap proof.} The rational certificates may reveal common decompositions that replace the finite case analysis. At present the general cubic second inequality genuinely uses the supplied exact certificate computation.
\item \textbf{Equality in the second inequality.} The first-gap equality theory in the original report is not automatically a classification for the second gap. Extracting and reconciling the zero sets of the certificates is a separate problem.
\item \textbf{Weighted variants.} The present theorem uses unweighted supports. Replacing populations by arbitrary real activities changes the domain and the coefficient model; the integer-face certificates do not by themselves establish such a theorem.
\item \textbf{Formal verification.} The small support kernels and rational identities are natural candidates for a proof-assistant checker. Exact arithmetic and independent implementations are valuable checks, but are not a formalized proof.
\end{enumerate}


\subsection{Explicit discriminant certificates for the ordinal family}\label{pgr:u:d3:app:disc}

For \eqref{pgr:u:d3:eq:outer} and \eqref{pgr:u:d3:eq:middle}, the table lists the coefficients in ascending powers of $x$ of $12\Delta/(x+3)$, with $x=m-3$. Thus each row gives a standalone integer positivity certificate. The common factor $x+3$ is positive for $x\ge0$.

\begin{center}\small
\begin{tabular}{llrrrrrr}
\toprule
Form & Parameters & $x^0$ & $x^1$ & $x^2$ & $x^3$ & $x^4$ & $x^5$\\\midrule
Outer & $(0,0)$ & 8640 & 13248 & 7884 & 2286 & 324 & 18\\
Outer & $(1,1)$ & 22484 & 30336 & 15577 & 3789 & 432 & 18\\
Outer & $(2,1)$ & 30148 & 37836 & 17888 & 4014 & 432 & 18\\
Middle & $j=0$ & 41412 & 46188 & 18456 & 3378 & 288 & 6\\
Outer & $(2,2)$ & 45376 & 55296 & 25472 & 5490 & 540 & 18\\
Outer & $(3,2)$ & 60804 & 69480 & 29613 & 5877 & 540 & 18\\
Middle & $j=1$ & 76356 & 78552 & 28221 & 4473 & 324 & 6\\
Outer & $(4,2)$ & 76880 & 83328 & 33268 & 6174 & 540 & 18\\
Middle & $j=2$ & 129156 & 123780 & 40644 & 5718 & 360 & 6\\
Outer & $(4,3)$ & 107300 & 112980 & 44092 & 7938 & 648 & 18\\
Outer & $(6,3)$ & 163620 & 158076 & 55080 & 8766 & 648 & 18\\\bottomrule
\end{tabular}
\end{center}

The generator \texttt{core\_templates.py} independently enumerates all $29$ labeled three-element preorders and all $82$ compatible choices of the core vertices below the antichain, recovering these eleven polynomial forms. It symbolically counts roles with multiplicity $\binom ma\binom{m-a}b$ for $a$ antichain tails and $b$ antichain heads. The layer checker processes antichain layers from bottom to top: if $b$ unmatched earlier tails are available and $k$ tails have been selected so far, choosing $a$ new tails and $c$ new heads is feasible exactly when $c\le b$. It has multiplicity $\binom{n_i}a\binom{n_i-a}c$ and updates $(b,k)$ to $(b+a-c,k+a)$. Finishing with $b=0$ gives the support counts. This recurrence is compared with literal disjoint-support enumeration for all $96$ compositions of six and seven.


\section{Degree four (source 16, the base)}\label{pgr:u:d4:sec}

Source~16 is \emph{Degree four preorder support polynomials} (17 pages as
delivered), the four-part package \texttt{preorder-degree-four} (archives
16--19); it is the base of this report. Its files are shipped with the
prefix \texttt{16-degree-four-}, except those that are byte copies of
files of sources 20, 22 and~23, which are shipped once under those
sources' prefixes. Its abstract describes, among other things, the role-cover
theorem of its Section~3, which is source~20's article; see
Section~\ref{pgr:u:d4:sec:rolecover}.

\paragraph{Abstract of source 16, as delivered.}

We study the polynomial counting disjoint directed matching supports of a finite preorder, with each feasible pair of endpoint sets counted once. We prove that the unweighted support polynomial of every finite preorder is rank-ultra-log-concave whenever its actual degree is at most four. The proof reduces arbitrary vertex counts to finitely many bounded cores with unbounded integer populations, and certifies the resulting inequalities by exact rational identities. Independently, we prove a structural theorem for every loopless directed relation whose arcs are covered by at most two selected tail vertices and two selected head vertices, with overlap allowed: its signed physical monomer polynomial is real stable, even with arbitrary nonnegative independent tail and head activities. Its support polynomial therefore has only negative real roots. The proof couples two rank-two transversal matroids through a stability-preserving operator and then merges the two copies of each physical vertex without allowing double use. For an unbalanced one-source, three-sink core, an ordinary proof of the last Newton inequality combines Lorentzian polarization with the Rayleigh property of rank-three matroids. A general exterior-only Lorentzian construction, an exact Gallai--Edmonds reduction, and independently checked rational certificates explain both the mathematical scope and the computational component. The final four-attachment ledger covers all 76 canonical templates; its last certificate consists of 2,273 rational square orbits and has been independently replayed. The result is computer-assisted, with ordinary structural proofs and finite algebraic verification distinguished throughout.

\subsection{The support polynomial and the questions}
A finite preorder is a reflexive transitive relation on a finite set $V$. Its off-diagonal relation is regarded as a directed graph $D$: an arc $i\to j$ means $i\ne j$ and $i$ precedes $j$. Equivalent elements give opposite arcs. Several structural results below apply to directed relations without transitivity.

\begin{definition}
A feasible support of size $k$ is an ordered pair $(S,T)$ of disjoint subsets of $V$, each of cardinality $k$, for which there exists a bijection $\pi:S\to T$ with $s\to\pi(s)$ for every $s\in S$. Write
\[
 \gamma_k(D)=\#\{(S,T): (S,T)\text{ is feasible of size }k\},\qquad
 \Ga_D(t)=\sum_{k\ge0}\gamma_k(D)t^k.
\]
The empty support contributes $\gamma_0=1$. A support contributes once, regardless of the number of witnessing bijections.
\end{definition}
The same-ground-set disjointness is essential. Passing to independent tail and head copies without subsequently excluding their simultaneous use changes this polynomial. Counting individual matchings also changes it.

For nonnegative tail activities $u_i$ and head activities $v_i$, replace the contribution of $(S,T)$ by
\[
 u^S v^T=\prod_{i\in S}u_i\prod_{j\in T}v_j.
\]
We use the same notation when the weights are specified. A vertex may have different activities in its two roles. With zero activities, the polynomial degree can decrease.

Preorder polytopes and their lattice-point enumerators were introduced and studied by Chapoton and Athanasiadis \cite{AC}. The supplied ProveIt report identifies the unweighted support polynomial with its preorder gamma polynomial \cite{PRP}. Our proofs concern the explicit support definition, so their validity does not require a particular formalization of that identification.


\mnote{source~16's Lemma ``degree'' is Lemma~\ref{pgr:lem:degree}, printed
once in Section~\ref{pgr:sec:object} with source~16's own proof.}


\begin{definition}
A polynomial $g(t)=\sum_{k=0}^d g_k t^k$ is \emph{rank-ultra-log-concave}, or rank-ULC, if its coefficients are nonnegative with no internal zeros and
\[
 \left(\frac{g_k}{\binom dk}\right)^2\ge
 \frac{g_{k-1}}{\binom d{k-1}}\frac{g_{k+1}}{\binom d{k+1}}
 \quad(1\le k<d),
\]
where $d$ is the actual degree.
\end{definition}
In degree four the inequalities are
\begin{equation}\label{pgr:u:d4:eq:gaps}
 3\gamma_1^2\ge8\gamma_2,\qquad
 4\gamma_2^2\ge9\gamma_1\gamma_3,\qquad
 3\gamma_3^2\ge8\gamma_2\gamma_4.
\end{equation}
Normalization by the ground-set size, by a larger cover size, or by a presumed degree after some weights vanish is a different assertion.


\mnote{source~16's Theorem ``Universal first inequality'' (if the
underlying graph has matching number at most $r\ge2$, then
$\gamma_1^2\ge\frac{2r}{r-1}\gamma_2$, including nonnegative role
activities) is Theorem~\ref{pgr:thm:first}, printed once in
Section~\ref{pgr:sec:object} with source~16's Motzkin--Straus proof; it is
\prp, Theorem \texttt{mr:thm:first}.}

The preceding research established rank-ULC for every finite preorder of actual degree at most three by a separately auditable finite reduction (source~25, Theorem~\ref{pgr:u:d3:thm:main}). It also exhibited the nonreal-rooted cubic
$1+24t+162t^2+208t^3$, whose discriminant is $-2592$ (Section~\ref{pgr:u:d3:sec:realroots}). Thus real-rootedness cannot be assumed in the general preorder problem.

\subsection{An exterior-only Lorentzian construction}
We use standard facts about Lorentzian polynomials: matroid basis polynomials are Lorentzian; products, nonnegative linear substitutions, coordinate upper truncations, and polarization preserve this class \cite{BH,RSW}. For a homogeneous bivariate polynomial $\sum a_k z^{r-k}t^k$, the Lorentzian condition gives log-concavity of $a_k/\binom rk$ with interval support.

\begin{theorem}\label{pgr:u:d4:thm:exterior}
Let $P,Q,I$ be disjoint finite sets, and let the only arcs be $P\to I$ and $I\to Q$. Put $r=|P|+|Q|$. With independent nonnegative exterior role activities, the polynomial
\[
 H(z,t)=\sum_{k=0}^r F_k z^{r-k}t^k
\]
is Lorentzian, where $F_k$ counts disjoint supports of size $k$. In particular $F_k/\binom rk$ is log-concave. This is actual-degree ULC when $\deg F=r$.
\end{theorem}
\begin{proof}
For the $P$ side form a transversal matroid on the physical exterior vertices and private dummy elements $d_p$, one for each $p\in P$. Its presentation has right side $P$; an exterior vertex has its given neighbors, and $d_p$ is adjacent only to $p$. The dummies form a basis, so the rank is exactly $|P|$.

For $P'\subseteq P$ and $R\subseteq I$, the set
$R\cup\{d_p:p\in P\setminus P'\}$ is a basis precisely when $|R|=|P'|$ and $R$ can be matched to $P'$. This encodes an endpoint support once, rather than counting its matchings. Construct the analogous $Q$ matroid with distinct dummies and the same physical variables $x_i$.

Their squarefree basis polynomials $B_P$ and $B_Q$ are Lorentzian. Multiply them, then discard every monomial in which some $x_i$ has exponent greater than one. Coordinate upper truncation preserves Lorentzianity \cite[Proposition 3.3]{RSW}. It removes exactly the choices using a physical exterior vertex in both roles. The all-dummy monomial survives.

Finally set every dummy variable equal to $z$ and every physical variable equal to $t$. A surviving pair of bases specifies the support
$S=P'\cup L$, $T=Q'\cup R$, where $L,R\subseteq I$ are disjoint. It contributes $z^{r-|L|-|R|}t^{|L|+|R|}$. This is exactly $H$. Exterior activities are introduced by scaling $x_i$ separately in the two factors before multiplication. Every operation retains homogeneous degree $r$.
\end{proof}
\begin{remark}
If $H$ has a common factor $z$, dividing that factor out is not a generally valid Lorentzian closure operation. The theorem does not silently strengthen its order-$r$ normalization to a smaller actual degree. It also does not assert real-rootedness.
\end{remark}
For $r=4$, differentiation of $H$ twice gives the three constants in \eqref{pgr:u:d4:eq:gaps}. For example $\partial_t^2H=2F_2z^2+6F_3zt+12F_4t^2$, whose discriminant yields $3F_3^2\ge8F_2F_4$.


\subsection{The two by two role cover}\label{pgr:u:d4:sec:rolecover}
Source~16's Section~3, ``Real-rootedness for a complete balanced core''
(its Theorems ``balanced'', ``twoedge'', ``missing'' and ``rolecover'', its
Lemmas ``quadratic'' and ``asano'', its Corollary ``freeproduct'', the
corollary on the gamma expansion, the interpolation remark, and its
Appendix~A of seven Rayleigh identities for two disjoint core edges), is
source~20's article \emph{Real stable monomers under a two by two role
cover}: 176 of the 218 sentences of source~20 recur verbatim in source~16,
with the same labels. It is printed once, from source~20, in Part~II
(Section~\pgrfwd{pgr:r:rc:sec}{on source~20}; the main theorem is
Theorem~\pgrfwd{pgr:r:rc:thm:rolecover}{``Two by two role cover''}). Source~16's
text differs in four places, all printed there as marked second routes:
\begin{enumerate}[leftmargin=*]
\item where source~20 invokes the stability-preserver theorem of Borcea and
  Br\"and\'en, source~16 proves each operator elementarily: an elementary
  Asano contraction lemma, a polarized auxiliary quadratic for the
  rank-two coupling, and the copy merge by diagonalization and
  differentiation;
\item the two-edge core covers canonical template~6, and the core with one
  missing edge covers canonical template~21, of the degree-four
  classification below;
\item the gamma-transfer corollary is stated with $N=|V|$; source~20 allows
  any $N\ge2d$;
\item the appendix of source~20 has a typographical slip (``labels0,1'')
  that source~16 does not.
\end{enumerate}
The degree-four ledger below discharges some templates by ordinary
theorems, among them the two-edge core and the core with one missing edge
(templates 6 and~21); the package records seven ordinary approvals, all
for the last gap, of templates 1, 4, 6, 9, 10, 21 and~34.

\subsection{The unbalanced one-source three-sink family}
Let $p$ be one source, let $Q$ consist of three sinks, and let $I$ be independent. Allow arbitrary arcs $p\to I$ and $I\to Q$, and any subset $p\to J$ of core arcs, where $J\subseteq Q$. No other arcs are present. All activities in this section are one.

\begin{theorem}\label{pgr:u:d4:thm:star}
Every relation in this family, and every reversal of such a relation, satisfies
$3\gamma_3^2\ge8\gamma_2\gamma_4$.
\end{theorem}
Let $Q(t)=1+Lt+Bt^2+Tt^3$ be the exterior-to-three-sink polynomial. Let $H$ be the exterior neighbors of $p$, $m=|H|$, and set
\[
 R(t)=\sum_{x\in H}Q_{-x}(t)=m+At+Ct^2+Dt^3.
\]
Let $Q^+(t)$ be $Q$ with one new exterior tail of neighborhood $J$. Then
\begin{equation}\label{pgr:u:d4:eq:star}
 \Ga_D(t)=Q^+(t)+tR(t).
\end{equation}
Indeed, supports with no exterior head belong to $Q^+$; otherwise their unique exterior head $x$ is matched from $p$, leaving exactly a $Q_{-x}$ support. This excludes physical double use and introduces no matching multiplicity.


\mnote{source~23 (Section~\ref{pgr:u:t26:sec}) treats this family with
the internal sink-side arc $q_2\to q_0$ added (canonical template~26),
where Lemma~\ref{pgr:u:d4:lem:average} is not available, and replaces it
by a deletion-mixture lemma.}

\begin{lemma}\label{pgr:u:d4:lem:average}
The homogeneous degree-three polynomial $mz^3+Az^2t+Czt^2+Dt^3$ is Lorentzian. In particular $C^2\ge3AD$.
\end{lemma}
\begin{proof}
Use the rank-three transversal matroid with private dummies. Identify dummy variables with $z$, the $m$ variables indexed by $H$ with $s$, and the other exterior variables with $t$. The resulting polynomial is Lorentzian. Polarize $s$ into $m$ slots using the degree bound $m$ \cite[Proposition 3.1]{BH}. Set one slot equal to zero and the other slots equal to $t$. A term using $k$ elements of $H$ is multiplied by
$\binom{m-1}{k}/\binom mk=(m-k)/m$.
Multiplication by $m$ therefore gives the deletion sum. Vertices of $H$ with no sink neighbor are matroid loops; they still count as unused polarization slots. If $m=0$ the assertion is immediate. No homogenizing monomial is divided out.
\end{proof}

\begin{lemma}\label{pgr:u:d4:lem:ratio}
For a three-center bipartite support polynomial, adjoining an exterior vertex cannot decrease the ratio of its cubic to quadratic coefficient, wherever that ratio is defined. Without division the assertion is
$T_{\rm new}B-TB_{\rm new}\ge0$.
\end{lemma}
\begin{proof}
Every rank-three matroid is Rayleigh \cite[Theorem 1.1]{WagnerR3}. Adjoin the new element $e$ to the transversal matroid with three private dummy elements $d_i$. Let $f$ be its basis polynomial. At existing exterior variables equal to one and all dummy variables and $e$ equal to zero, we have
\[
 f=T,\quad \partial_e f=\Delta T,\quad
 \sum_i\partial_{d_i}f=B,\quad
 \sum_i\partial_e\partial_{d_i}f=\Delta B.
\]
Summing the Rayleigh inequalities
$(\partial_e f)(\partial_{d_i}f)-f\partial_e\partial_{d_i}f\ge0$
gives $\Delta T B-T\Delta B\ge0$. Boundary evaluations follow by continuity. An added loop changes neither coefficient.
\end{proof}
\begin{proof}[Proof of Theorem~\ref{pgr:u:d4:thm:star}]
If $D=0$, the desired inequality is immediate. Otherwise $B,T,B^+,T^+,C$ are positive. For each deleted exterior vertex $x$, Lemma~\ref{pgr:u:d4:lem:ratio} applied when $x$ is restored and when the new core tail is added gives
$T_xB\le B_xT$ and $TB^+\le BT^+$. Thus $T_xB^+\le B_xT^+$, without dividing by a possibly zero $B_x$. Summing gives $DB^+\le CT^+$. Together with Lemma~\ref{pgr:u:d4:lem:average} and \eqref{pgr:u:d4:eq:star},
\begin{align*}
3\gamma_3^2-8\gamma_2\gamma_4
 &=3(C+T^+)^2-8(A+B^+)D\\
 &\ge \frac{C^2}{3}-2CT^++3(T^+)^2
 =\frac{(C-3T^+)^2}{3}\ge0.
\end{align*}
\end{proof}
The real-rootedness conclusion of the balanced theorem does not extend to the full unbalanced family. For example the checked polynomial
$1+40t+510t^2+2380t^3+2704t^4$ has discriminant $-4112937216$ and two real roots. It is realized with four exterior vertices of each of the three two-sink neighborhoods, thirteen source-only exterior vertices, and all three core arcs. Thus there are twenty-five exterior vertices and twenty-nine vertices in total. The package contains the exact population specification and a direct support verifier.


\subsection{A finite normal form for bounded matching rank}
Let $G$ be the comparability graph of a preorder of matching number $r$. In the canonical Gallai--Edmonds decomposition, $D$ is the set of vertices missed by some maximum matching, $A=N(D)\setminus D$, and $C=V\setminus(A\cup D)$. The components $D_i$ of $G[D]$ are factor-critical, $G[C]$ has a perfect matching, and
\begin{equation}\label{pgr:u:d4:eq:ge}
 r=|A|+\frac{|C|}{2}+\sum_i\frac{|D_i|-1}{2}.
\end{equation}
These standard facts are proved, for example, by West \cite{West}.

\begin{proposition}\label{pgr:u:d4:prop:normal}
Let $I$ be the union of the singleton components of $G[D]$, and put $K=V\setminus I$, $a=|A|$. Then $I$ is independent, all its neighbors lie in $A$, and
\[
 |K|\le3r-2a.
\]
Moreover, $\nu(G[K\setminus A])=r-a$. Every nonisolated exterior vertex is a singleton preorder equivalence class. Each attachment has one strict orientation toward all its exterior neighbors. Hence at most $2^a-1$ nonempty exterior types occur, with arbitrary integer populations.
\end{proposition}

\mnote{this proposition is source~25's Theorem~\ref{pgr:u:d3:thm:normalform}
(its first half), with the same proof: the matching budget
\eqref{pgr:u:d4:eq:ge} gives $|K|\le3r-2a$ and $\nu(G[K\setminus A])=r-a$,
and automorphism invariance of the canonical sets together with
transitivity gives singleton exterior classes and one orientation per
attachment (Lemma~\ref{pgr:u:d3:lem:exterior} there). Source~16's
four-paragraph proof is replaced by this pointer.}

For $r=4$ the core bounds are $10,8,6,4$ for $a=1,2,3,4$. When $a=0$, each nontrivial component has at most nine vertices. Polynomials multiply over comparability components because every support splits uniquely. Rank-ULC is preserved by products: homogenize each factor at its actual degree, use the bivariate Lorentzian characterization and product closure, then specialize. This uses no change of normalization.

\subsection{Exact support kernels and coverage}
For a fixed core and a list of legal exterior types, let $n_t$ be the population of type $t$. If $q_t$ selected exterior labels are prescribed, let $c_{k,q}$ count feasible supports in the induced relation on the core and precisely those prescribed labels, with every prescribed label required to occur. Then
\begin{equation}\label{pgr:u:d4:eq:kernel}
 \gamma_k(\mathbf n)=\sum_{|q|\le\min(k,a)}c_{k,q}\prod_t\binom{n_t}{q_t}.
\end{equation}
Every selected exterior vertex needs a different core partner, so at most $a$ can occur. The binomial factors choose their physical labels. Feasibility depends only on the type quota and the selected core endpoint sets; unselected vertices cannot help match a fixed support. Thus \eqref{pgr:u:d4:eq:kernel} is an exact all-population formula, not interpolation from bounded numerical tests.

Independent kernel checks enumerate disjoint tail and head sets and apply all Hall inequalities to the induced directed bipartite relation. The producer uses a different Boolean matching routine. Both count one feasible support even when it has several witnesses.

Every preorder core is obtained from a quotient poset and positive sizes of its equivalence classes. Posets can be generated by adjoining a maximal element with every possible predecessor ideal. Canonicalization merges only complete adjacency matrices related by an explicitly tried permutation; global reversal exchanges tail and head supports and is also allowed. The four-core audit independently reconstructs the domain from upper-triangular quotient relations, block compositions, and all relabelings.

Exterior legality is checked with two identical copies of a type and with the union of all legal types. Any failed transitivity relation involves at most two exterior vertices, so these tests cover arbitrary multiplicities. For structural pruning, let $H$ be a component of the graph formed with one clone of every legal type. Its largest matching rank over all its populations is
\[
 |A_H|+\nu\bigl((K\cap H)\setminus A_H\bigr),
\]
where $A_H$ is its set of active attachments. The upper bound follows by removing edges incident with $A_H$. Enough clones attain it by matching the attachments to distinct exterior vertices before matching the residual core. This is an all-population argument, not a finite sampling rule.

A connected bipartite preorder component with at least three vertices has no equivalence edge and no directed two-edge path on distinct vertices, since either would force a triangle. It therefore has genuine source and sink shores. Its support polynomial is the bipartite shore-support polynomial. The previously established weighted bipartite rank-at-most-four theorem applies \mnote{source~16 cites here a weighted bipartite
rank-four package that it embeds as a dependency; that package is
batch-73 manuscript~26, which this intake prints in \mrn{} as its
source~77 (Theorem \texttt{mrn:v:rf4:thm:main}). The stronger
\mrn, Corollary \texttt{mrn:w:rf:cor:rankfive}, already covers every
weighted bipartite graph of matching number at most five.}; an isolated two-element equivalence component contributes $1+2t$.

\subsection{The bounded branches and the degree four theorem}
All activities in the computational classification are one. The term \emph{weighted quotient} below refers only to positive integer equivalence-class sizes, not to arbitrary role activities. The following finite counts describe exact certificate domains, not random tests.
\begin{center}\small
\begin{tabular}{p{24mm}p{36mm}p{82mm}}
\toprule
Branch & Core domain & Exact reduction\\\midrule
$a=0$ & Components of order $\le9$ & Finite degree-four check and lower-degree theorem\\
$a=1$ & Marked core of order $\le10$ & $412622$ coefficient pairs and $1237866$ quadratic checks\\
$a=2$ & Core of order $\le8$ & $89863$ distinct gamma polynomials\\
$a=3$ & Core of order $\le6$ & $6213$ templates and $3437$ distinct gamma polynomials\\
$a=4$ & Four attachment vertices & $76$ canonical nonbipartite templates\\\bottomrule
\end{tabular}
\end{center}
The complete ten-element quotient scan checks $5049656$ weighted quotient cases, $686481$ of actual degree four. Smaller preorders are included by padding with isolated vertices. All actual-degree-four cases satisfy \eqref{pgr:u:d4:eq:gaps}.

For $a=1$, transitivity makes the attachment a singleton source or sink. Its $m$ pendant clones give the exact identity
\[
 \Ga_{K_m}(t)=\Ga_K(t)+mt\Ga_{K-v}(t).
\]
Each gap is a quadratic $A m^2+B m+C$. The certificate verifies $A,C\ge0$ and either $B\ge0$ or $B^2\le4AC$, proving nonnegativity for all real $m\ge0$. Independent support enumeration reproduces all $412622$ coefficient pairs and the complete quadratic ledger.

For $a=2$, the core condition $\nu(K\setminus A)\le2$ bounds the degree by four at every population. The three population variables have coefficient degree at most two. Exact enumeration produces $89863$ distinct gamma polynomials and $179726$ remaining-gap instances. Independent positivity replay covers $29215$ normalized targets and all required face aliases.

For $a=3$, $90042$ legal marked specifications reduce structurally to $6213$ templates and $3437$ distinct gamma polynomials. Independent Hall recount verifies $188715$ quota vectors and $943575$ coefficients. All $204388$ population faces are covered by exact rational identities, including $53722$ normalized nontrivial targets.

For $a=4$, independent enumeration gives $355$ labeled four-element preorders, $4882$ legal signed specifications, and $76$ canonical retained templates. It checks $32516$ quota vectors, $162580$ coefficients, and the complete $258153$-term gap data. These counts establish coverage and algebra before any positivity claim is made.

\begin{theorem}[Rank ultra log concavity through degree four]\label{pgr:u:d4:thm:middle}
For every finite preorder whose unweighted support polynomial $\Ga_D$ has actual degree $d\le4$, the sequence
\[
 \left(\gamma_k/\binom dk\right)_{k=0}^d
\]
is log-concave. In particular, if $d=4$, all three inequalities in \eqref{pgr:u:d4:eq:gaps} hold.
\end{theorem}
\begin{proof}
Degrees at most three follow from the separately verified lower-degree theorem (source~25, Theorem~\ref{pgr:u:d3:thm:main}). Suppose $d=4$. The first inequality follows from Theorem~\ref{pgr:thm:first}. The normal form is exhaustive. The $a\le3$ sectors satisfy all three gaps by the preceding exact certificates. For the $a=4$ sector, the complete independently replayed ledger covers both remaining gaps for all $76$ templates and every nonnegative integer population, using the structural proofs above and the exact certificate schema below. Structural pruning and product closure handle discarded specifications and disconnected preorders. Every product is normalized by its own actual degree before homogeneous Lorentzian product closure is applied. Thus no vertex-count bound or unproved change of normalization enters the conclusion.
\end{proof}
Theorem~\ref{pgr:u:d4:thm:middle} is computer-assisted. Its proof requires exact replay of a finite, explicitly supplied collection of identities; floating-point optimization is not a logical dependency. The final template, numbered $26$, has core arcs $2\to0$, $3\to0$, $3\to1$, $3\to2$, with vertex $3$ the unique outward attachment. Its certificate contains $2273$ nonnegative rational square orbits. Independent expansion verifies the full unscaled last-gap polynomial, thereby covering its $1785$ remaining population faces. The complete ledger has no outstanding template or face obligations.

\mnote{template~26 has a second, smaller route for its last gap: source~23,
Section~\ref{pgr:u:t26:sec}. The original route's certificate
(\texttt{global\_integer\_sextic\_shifted\_gmp\_26.json}, 29.4~MB, the 2273
square orbits above) is not shipped with this report; its provenance
record and replay log are (\texttt{data/16-degree-four-provenance-\dots}),
and the README says how to retrieve the certificate from the arrival
commit. It was replayed exactly at intake.}


\subsubsection{A hybrid last-gap certificate}
For a four-core template, delete all internal core arcs and let $F_k$ be the exterior-only coefficients. Each size-four support uses all four core vertices on exterior edges, so $\gamma_4=F_4$. Write $B=\gamma_2-F_2$ and $C=\gamma_3-F_3$; these count disjoint support classes using core edges and are nonnegative. Then
\begin{equation}\label{pgr:u:d4:eq:hybrid}
 3\gamma_3^2-8\gamma_2\gamma_4
 =(3F_3^2-8F_2F_4)+E,\qquad
 E=6F_3C+3C^2-8BF_4.
\end{equation}
The first summand is nonnegative by Theorem~\ref{pgr:u:d4:thm:exterior}, at cover order four even if the exterior-only degree is smaller. For templates 9,36,37,50 the package proves $E\ge0$ by exact nonnegative binomial or binomial-square identities. Its degree is at most five, smaller than the full sextic gap. This is a case-specific reduction: $E$ can be negative for other templates, and those cases use full-gap certificates instead.

\subsection{Integer population certificates}
Write $B_\alpha(\mathbf n)=\prod_i\binom{n_i}{\alpha_i}$. A global certificate has the form
\begin{equation}\label{pgr:u:d4:eq:cert}
 P(\mathbf n)=\sum_j\frac{w_j}{|G|}\sum_{g\in G}B_{m_j}(g\mathbf n)
       \left(\sum_\alpha c_{j,\alpha}B_\alpha(g\mathbf n)\right)^2
       +\sum_\alpha r_\alpha B_\alpha(\mathbf n),
\end{equation}
where $w_j,r_\alpha\ge0$ are rational and $G$ is a specified coordinate-permutation group. Each summand is nonnegative on $\N^d$. These are identities of rational polynomials, but positivity is asserted on integer populations: for instance $\binom{x}{2}$ is negative for $0<x<1$.

A standard-library checker expands every symmetry action and every product using
\[
 \binom na\binom nb=
 \sum_{h=0}^{\min(a,b)}\frac{(a+b-h)!}{h!(a-h)!(b-h)!}
                         \binom n{a+b-h}.
\]
This counts pairs of subsets by their intersection size. Tensor products give the multivariate multiplication rule. Comparing every coefficient, including omitted zeros, proves the identity; checking all rational signs proves the inequality. Solver status and numerical tolerances are irrelevant to this final step.

For some smaller domains, split each coordinate into $n_i=0$ or $n_i=1+y_i$. These finitely many faces cover all integer populations. Each resulting polynomial is certified by nonnegative monomials and nonnegative monomial multiples of exact rational squares. Unused coordinates, permutations, and positive scaling factors are explicitly checked in the face ledger.

\subsubsection{Shifted binomial discovery}
The successful higher-dimensional search writes a square root accompanying $B_m$ in the basis $B_q(\mathbf n-m)$. When an integer population is outside $\mathbf n\ge m$, the multiplier vanishes. Within that orthant, shifted coordinates reveal boundary zeros. The final output removes every shift using
\[
 \binom{n-m}{q}=\sum_{a=0}^q\binom{-m}{q-a}\binom na,
\]
and is verified in the unshifted schema \eqref{pgr:u:d4:eq:cert}.

If a target vanishes on the entire lower integer box $m+s$, $0\le s\le q$, any square with multiplier $B_m$ must vanish there. Triangular Newton interpolation then eliminates its coefficients indexed by $a\le q$. An isolated zero imposes only its evaluation equation, not the deletion of a single basis coefficient. These restrictions aid discovery; only the final full identity is trusted.

\subsubsection{Reproducing the finite verification}

\mnote{source~16's package is not shipped in its delivered layout: its
files are renamed into \texttt{code/} and \texttt{data/} with the prefix
\texttt{16-degree-four-}, the files larger than 2~MB (the 93.7~MB
three-attachment certificates among them) stay in the arrival commit, and
the two embedded dependency archives are exact copies of source~25's
package and of batch-73 manuscript~26. The delivered \texttt{verify.py}
first checks every one of its 260 manifest entries, so it runs only on the
reconstructed package (README, ``Reconstructing what is not shipped''). On
Windows its path substitution writes backslashes into Python source; run
it on Linux or WSL, or in a copy that uses \texttt{as\_posix()}.}

The accompanying directory \texttt{reproducibility/} contains the exact data, independent checkers, approved coverage inputs, and a portable driver. From that directory run
\begin{quote}
\verb|python verify.py --mode fast|
\end{quote}
This replays the fixed rational identities, gamma-to-gap reconstructions, face aliases, and complete template coverage. Its four-attachment manifest has $50$ global identities: $24$ middle-gap identities, $23$ whole last-gap identities, and $3$ identities for a core correction added to a Lorentzian baseline. They contain $15408$ square orbits in total. Together with the ordinary proofs and coefficientwise-positive cases, the complete coverage check discharges all $39367$ recorded outstanding face tasks and all $76$ canonical templates.

For independent exhaustive kernel regeneration as well, run
\begin{quote}
\verb|python verify.py --mode full|
\end{quote}
The full mode adds the bounded-core support enumerations and Hall-based coefficient checks; its ten-element enumeration is the expensive step. Python~3.11 or later suffices for fixed certificate replay; full regeneration also requires GNU \texttt{g++} with C++17 support (the enumerators use GNU headers and 128-bit integers). Neither mode needs a numerical optimization package. The final package requires no partial-result override: a successful final coverage assertion demands all $76$ templates and zero remaining tasks. A separate \verb|--mode hashes| checks integrity only and does not establish the mathematics.

The previous degree-three and bipartite rank-four packages are included as explicit dependencies. The driver verifies their archive hashes; their own documented commands replay their earlier computational components. These are not silently re-enumerated by the new wrapper. Ordinary proof sources are supplied and pinned, but checking a source hash is likewise not a mathematical proof. The wrapper runs fresh copies of the independent checkers in a temporary directory, removes historical positivity receipts, and derives its final coverage union from the newly replayed identities and the specified ordinary theorem approvals. The supplied README describes the small path-portability adaptations and where fresh logs and results are written.

\subsection{Verification boundary and remaining research}
The current article separates ordinary structural proofs from finite computer-assisted results. The role-cover theorem includes arbitrary independent role activities; the outward-star last-gap theorem is presently stated for unit activities. The exterior-only theorem has cover-order normalization, which is actual-degree normalization only at full cover rank. None is silently substituted for the general preorder theorem.

The completed degree-four ledger is included in the package. All 76 canonical four-attachment templates have both remaining gaps covered, and the lower-attachment branches and finite support kernels are independently checked. These finite identities, together with the ordinary reductions, establish Theorem~\ref{pgr:u:d4:thm:middle} for arbitrary finite preorders of actual degree at most four. No Lean formalization or global priority claim is made here.

The most useful further questions are:
\begin{enumerate}[leftmargin=*]
\item Replace the largest integer-population certificates, particularly the final one-source template with an internal sink-side edge, by ordinary structural proofs. Can its weighted deletion average be proved rank-ultra-log-concave directly?
\item Determine whether arbitrary independent role activities preserve rank-ULC for every degree-three preorder; a separate program proves the one-activity-per-vertex version and all unbounded branches for independent roles, but the remaining finite independent-role cases require their own coverage proof.
\item Determine whether actual-degree ULC holds for every two-by-three role cover. An unrestricted real-rootedness extension is impossible: the rank-three nonreal example, with two independent edges adjoined, already gives a nonreal degree-five product in that cover class.
\item Characterize when the exterior top-coefficient ratio remains monotone in ranks above three. General transversal matroids are not all Rayleigh, so the rank-three proof has a genuine boundary.
\item Classify equality in the degree-four inequalities and determine whether equality forces a factorization into lower-rank components.
\item Formalize the finite support kernel and the rational certificate checker in a proof assistant. This would strengthen the trust boundary without requiring formalization of the numerical discovery algorithm.
\item Find practical degree-five extensions of the bounded-core method, and identify the first true failures under actual-degree normalization rather than cover-size normalization.
\end{enumerate}

\mnote{status of these questions in this report. (1)~Source~23
(Section~\ref{pgr:u:t26:sec}) gives a smaller proof of the last gap of the
final one-source template, by a deletion-mixture lemma, but it is still
certificate-assisted (ten quartic identities) and does not prove the
weighted deletion average rank-ULC directly. (2)~Answered: source~22
(Part~II) proves rank-ULC for every degree-three preorder with arbitrary
independent role activities. (3)~Open in general; settled for role covers
with $P\subseteq Q$ by source~11's corollary, for independent physical
cores with saturation by source~15, and by real stability for two by two
role covers (source~20); sources 11, 12 and~15 state the general
two-tail/three-head case as open. (4)--(6) remain open. (7)~Part~III proves
actual-degree five for five-vertex cores with a universal-sink cloud
(source~34); degree five for preorders in general is open.
Added 2~October 2026, batch~77: with sources 06 and~35, Part~III proves
rank-ULC at every actual degree for five-vertex cores with a
universal-sink cloud (Corollary~\ref{pgr:s:fq:cor:all}).}



\section{A second route for the last gap of template 26 (source 23)}\label{pgr:u:t26:sec}

Source~23 is \emph{A smaller proof of the last Newton gap for template 26.
A certificate assisted supplement} (6 pages as delivered; archive
\texttt{template26-structural-supplement}; no author line). Its files are
shipped with the prefix \texttt{23-template26-}.

\mnote{source~23 proves again, by a different and smaller route, the one
entry of source~16's ledger with the largest certificate: canonical
template~26, last gap (Section~\ref{pgr:u:d4:sec}). Its family and the
decomposition $\Ga=\Qp+tR$ are those of source~16's
Theorem~\ref{pgr:u:d4:thm:star} with the internal arc $q_2\to q_0$ added;
its Rayleigh step and its final square are source~16's
Lemma~\ref{pgr:u:d4:lem:ratio} and the last display of the proof of
Theorem~\ref{pgr:u:d4:thm:star}, with an internal-edge correction. It
answers source~16's first research question in part. Its dependency
``the independently approved vertex-weighted theorem at actual degree at
most three'' was delivered as archive~31, an earlier version of source~22;
the statement is the case $u_i=v_i$ of source~22's theorem,
Corollary~\pgrfwd{pgr:r:ir:cor:vertex}{of source~22 in Part~II}. Both
routes stay: source~16's certificate for template~26 is a direct
binomial-square identity; source~23's route needs ten quartic identities,
that theorem and Wagner's rank-three Rayleigh theorem.}

The text below is source~23's, from its opening paragraph.

This supplement proves the last order-four Newton inequality for the integer-population family called template 26. The argument uses a support decomposition, a deletion-mixture lemma, and monotonicity of a ratio of two top coefficients. Its finite algebraic input consists of ten exact quartic identities. The proof also depends on the independently approved vertex-weighted theorem at actual degree at most three \mnote{that theorem is the case $u_i=v_i$ of source~22's
Theorem~\pgrfwd{pgr:r:ir:thm:main}{(Part~II)}, recorded as
Corollary~\pgrfwd{pgr:r:ir:cor:vertex}{of source~22}.} and Wagner's rank-three Rayleigh theorem. Thus it is a smaller certificate-assisted proof, not a wholly ordinary proof.

\subsection{Supports and the precise family}
Let $P$ be a finite preorder on a physical ground set $V$. A size-$k$ support is an \emph{ordered} pair $(U,W)$ of subsets of that same set such that
\[
U\cap W=\varnothing,\qquad |U|=|W|=k,
\]
and some bijection $U\to W$ uses off-diagonal arcs of $P$. Count each feasible ordered pair once, regardless of the number of matching witnesses. Distinct ordered pairs with the same union remain distinct. Set
\[
\Gamma_P(t)=\sum_{k\ge0}\gamma_k(P)t^k,\qquad \gamma_0(P)=1.
\]
All activities in the template-26 conclusion are one. Diagonal preorder arcs do not contribute to supports because $U$ and $W$ are disjoint.

There are four core vertices $q_0,q_1,q_2,p$. Their off-diagonal arcs are
\[
q_2\longrightarrow q_0,\qquad p\longrightarrow q_i\quad(0\le i\le2).
\]
Exterior vertices are independent. An exterior vertex $x$ may point to any subset $N(x)\subseteq\{q_0,q_1,q_2\}$ satisfying
\[
q_2\in N(x)\ \Longrightarrow\ q_0\in N(x).
\]
Independently, the arc $p\to x$ may be present. In bit order $q_0,q_1,q_2,p$, the first three bits describe $x\to q_i$, while the fourth describes $p\to x$. The nonisolated types are exactly
\begin{equation}\label{pgr:u:t26:eq:types}
1,2,3,5,7,8,9,10,11,13,15.
\end{equation}
Each type has an arbitrary nonnegative integer population. Isolates may be omitted. This is the canonical template with core rows $[0,0,1,7]$ and orientation mask $8$.

\begin{theorem}\label{pgr:u:t26:thm:main}
For every such population, writing $\Gamma(t)=\sum_k\gamma_k t^k$, one has
\begin{equation}\label{pgr:u:t26:eq:target}
3\gamma_3^2\ge8\gamma_2\gamma_4.
\end{equation}
At actual degree four, this is the last Newton inequality for $\gamma_k/\binom4k$. If the degree drops below four, the displayed inequality is automatic.
\end{theorem}
There can be at most three core heads and at most one exterior head, since only $p$ points to exterior vertices. Thus the actual degree never exceeds four.

\subsection{The support decomposition}
Delete $p$ and call the induced preorder $Q$. Its nonisolated exterior types are $1,2,3,5,7$; let their populations be $m=(a,b,c,d,e)\in\N^5$. A former type-$8$ vertex is now isolated. Write
\[
Q(t)=1+A t+B t^2+T t^3.
\]
Let $H$ be the exterior vertices with $p\to x$, including the possible $Q$-isolates. For each $x\in H$, let $Q_{-x}$ be the support polynomial after deleting that \emph{physical vertex}, and put
\[
R(t)=\sum_{x\in H}Q_{-x}(t)=\sum_{j=0}^{3}R_jt^j.
\]
Let $\Qp$ be $Q$ with one new exterior tail adjacent to all three core/head vertices. This is still a preorder, of support degree at most three. Write $\Bp=[t^2]\Qp$ and $\Tp=[t^3]\Qp$.

\begin{proposition}\label{pgr:u:t26:prop:decomp}
The full support polynomial satisfies
\begin{equation}\label{pgr:u:t26:eq:decomp}
\Gamma(t)=\Qp(t)+tR(t).
\end{equation}
Consequently,
\begin{equation}\label{pgr:u:t26:eq:coeffs}
\gamma_2=\Bp+R_1,\qquad\gamma_3=\Tp+R_2,\qquad\gamma_4=R_3.
\end{equation}
\end{proposition}
\begin{proof}
If the head set contains no exterior vertex, the support is exactly a support of $\Qp$, with its new universal tail identified with $p$. If it contains an exterior vertex $x$, then $x$ is unique: only $p$ can match to exterior heads. The pair $p\to x$ is forced in every matching witness. Removing $p$ and $x$ gives precisely a support of $Q_{-x}$. Conversely, adjoining $p\to x$ to such a support recovers a full support. Distinct $x$ give distinct head sets. These are bijections of supports, so no matching multiplicity is introduced.
\end{proof}

For reference, the three coefficients of $Q$ are the following polynomials:
\begin{align}
A={}&1+a+b+2c+2d+3e,\label{pgr:u:t26:eq:A}\\
B={}&ab+ac+ad+2ae+bc+2bd+2be\notag\\
 &+3cd+3ce+3de+\binom c2+\binom d2+3\binom e2+b+c+e,\label{pgr:u:t26:eq:B}\\
T={}&b\left((d+e)(a+c)+\binom{d+e}{2}\right)\notag\\
 &+a\left(cd+ce+de+\binom e2\right)
 +\binom{c+d+e}{3}-\binom c3-\binom d3.\label{pgr:u:t26:eq:T}
\end{align}
Here binomial expressions are ordinary binomial polynomials, evaluated at integer populations. The accompanying independent Hall-support checker reconstructs these formulas directly: all $56$ exterior quota vectors of total at most three are enumerated, and all $42$ nonzero coefficient entries agree. The quota bound follows because every selected exterior vertex needs a distinct core head. The decomposition \eqref{pgr:u:t26:eq:decomp} is also independently replayed from direct ordered disjoint support counts, rather than assumed by the algebra checker.

\subsection{Reducing deletion mixtures to pairs}
For $i\in\{a,b,c,d,e\}$ with $m_i>0$, write $Q_i=Q(m-e_i)$, where $e_i$ is the corresponding coordinate unit vector, and write $Q_*=Q(m)$. We require the following statement.

\begin{lemma}[Deletion mixtures]\label{pgr:u:t26:lem:mixtures}
Every nonnegative real linear combination $S$ of $Q_*$ and the available $Q_i$ satisfies
\begin{equation}\label{pgr:u:t26:eq:mixture}
S_2^2\ge3S_1S_3.
\end{equation}
The populations $m$ are nonnegative integers. In particular, \eqref{pgr:u:t26:eq:mixture} applies to $R$ from Section 2: same-type deletions coincide, and deletion of a $Q$-isolate contributes $Q_*$.
\end{lemma}

\paragraph{The two-point principle.}
Consider any finite list of triples $(A_i,B_i,T_i)$ with $A_i>0$. For a nonzero nonnegative mixture, normalize by $A=\sum_iw_iA_i$ and let $\lambda_i=w_iA_i/A$. Then
\[
\left(\frac BA,\frac TA\right)
=\sum_i\lambda_i\left(\frac{B_i}{A_i},\frac{T_i}{A_i}\right),
\qquad \lambda_i\ge0,\quad\sum_i\lambda_i=1.
\]
At each fixed horizontal coordinate, a highest point of this finite convex hull is a combination of at most two listed points. To see this, take a maximizing combination with smallest positive support. If it has at least three positive coefficients, a nonzero perturbation $\delta$ satisfies $\sum\delta_i=\sum\delta_i(B_i/A_i)=0$. Both sufficiently small signs are feasible. Maximality forces $\sum\delta_i(T_i/A_i)=0$; extend one sign until a coefficient vanishes, contradicting minimality. This also covers repeated horizontal coordinates.

It follows that every mixture satisfies $B^2\ge3AT$ if every two-point mixture does: all upper boundary points then lie below $y=x^2/3$, as does the entire hull. This does \emph{not} assume that the subgraph of this convex parabola is convex. The zero mixture is immediate. In our family every $A_i\ge1$, because the internal arc remains.

\paragraph{Undeleted and deleted pairs.}
Give one marked physical exterior vertex $x$ of type $i$ activity $s\in[0,1]$, keeping every other vertex activity equal to one. A support has weight equal to the product of the activities of its physical vertices. Therefore its polynomial is exactly
\begin{equation}\label{pgr:u:t26:eq:activity}
Q^{(s)}(t)=sQ_*(t)+(1-s)Q_i(t).
\end{equation}
Indeed, a support either uses $x$ once or avoids it. The independently approved vertex-weighted actual-degree-at-most-three theorem \mnote{that theorem is the case $u_i=v_i$ of source~22's
Theorem~\pgrfwd{pgr:r:ir:thm:main}{(Part~II)}, recorded as
Corollary~\pgrfwd{pgr:r:ir:cor:vertex}{of source~22}.} gives $B(s)^2\ge3A(s)T(s)$ when the surviving degree is three. At lower degree $T(s)=0$, so the same displayed inequality is automatic, including $s=0$. This is a change in one vertex's activity, not a substitution of a fractional population into \eqref{pgr:u:t26:eq:A}--\eqref{pgr:u:t26:eq:T}.

For $u,v\ge0$ with $u+v>0$, choose $s=u/(u+v)$ and multiply \eqref{pgr:u:t26:eq:activity} by $u+v$. Homogeneity gives the inequality for $uQ_*+vQ_i$. The same theorem supplies every endpoint gap $B_i^2-3A_iT_i\ge0$. It remains to check the ten distinct-deletion pairs.

\subsection{The ten exact quartic identities}
Write $(A_i,B_i,T_i)$ for the coefficients of $Q_i$ and define
\[
C_{ij}(m)=2B_iB_j-3(A_iT_j+A_jT_i)\qquad(i<j).
\]
Since both deletion types are available, write $m=x+e_i+e_j$ with $x\in\N^5$. Thus the two coefficient triples come from $Q(x+e_j)$ and $Q(x+e_i)$.

\begin{lemma}[Computer-checked quartic identities]\label{pgr:u:t26:lem:quartics}
For each of the ten unordered pairs below, the supplied exact data give an identity
\begin{equation}\label{pgr:u:t26:eq:certificate}
C_{ij}(x+e_i+e_j)
=\sum_{\ell}q_{\ell}\binom{x}{\alpha_{\ell}}f_{\ell}(x)^2
 +\sum_{\beta}r_{\beta}\binom{x}{\beta},
\end{equation}
where $q_{\ell},r_{\beta}>0$ are rational,
$\binom{x}{\alpha}=\prod_{k=1}^{5}\binom{x_k}{\alpha_k}$,
all exponent vectors belong to $\N^5$, and each $f_{\ell}$ is an explicitly listed rational polynomial in the same binomial basis. Consequently $C_{ij}(m)\ge0$ on its integer population domain.
\end{lemma}
\begin{center}
\begin{tabular}{@{}lrr@{}}\toprule
Deleted population coordinates & Square terms & Remainder terms\\\midrule
$a,b$ & 15 & 84\\
$a,c$ & 21 & 81\\
$a,d$ & 18 & 83\\
$a,e$ & 22 & 83\\
$b,c$ & 24 & 81\\
$b,d$ & 19 & 88\\
$b,e$ & 21 & 85\\
$c,d$ & 18 & 89\\
$c,e$ & 19 & 86\\
$d,e$ & 19 & 86\\\midrule
Total & 196 & 846\\\bottomrule
\end{tabular}
\end{center}

\noindent\emph{Certificate format and verification.}
The data file \file{distinct_deletion_cross_certificates.json} records, for each pair, its deletion indices ($0,\ldots,4$ correspond to $a,\ldots,e$), the population shift, and the two sums in \eqref{pgr:u:t26:eq:certificate}. A square record contains a rational weight, a multiplier exponent vector, and a list of exponent--coefficient pairs specifying $f_{\ell}$. A remainder record is an exponent--coefficient pair. All ten identities are checked coefficientwise with exact rational arithmetic, together with every positivity and integer-exponent condition. Each summand is nonnegative for $x\in\N^5$. Integer scope matters: ordinary binomial polynomials need not be nonnegative at arbitrary nonnegative real arguments.

The independent checker reconstructs the targets from Hall support feasibility and imports no producer code. It also derives binomial multiplication constants by counting pairs of subsets with prescribed union. Thus the finite input is an exact identity certificate, rather than a sample of population values or a floating-point optimization result.

\begin{proof}[Completion of Lemma \ref{pgr:u:t26:lem:mixtures}]
Write $G(S)=S_2^2-3S_1S_3$, so $G_i=G(Q_i)\ge0$. For $u,v\ge0$,
\[
G(uQ_i+vQ_j)=u^2G_i+uvC_{ij}+v^2G_j\ge0.
\]
Lemma \ref{pgr:u:t26:lem:quartics} settles all distinct-deletion pairs. Section 3 settles the undeleted/deleted pairs; identical deletion types need no extra check. The two-point principle proves \eqref{pgr:u:t26:eq:mixture}, hence $R_2^2\ge3R_1R_3$.
\end{proof}

\subsection{The top ratio with the internal edge present}
Remove $q_2\to q_0$ from $Q$ and call the exterior-only support polynomial $F$. Set $B_0=F_2$, $T_0=F_3$, and let $L$ be the number of exterior vertices adjacent to $q_1$. At support level,
\begin{equation}\label{pgr:u:t26:eq:internal}
B=B_0+L,\qquad T=T_0.
\end{equation}
A size-two support using the internal edge must match its sole exterior tail to $q_1$, because $q_2$ is a tail and $q_0$ a head. A size-three support cannot use the internal edge: it would require three core heads while $q_2$ is already a tail.

Adjoin one exterior vertex $x$ of any allowed outgoing neighborhood, possibly empty. Set
\[
\beta=\Delta B_0,\qquad \tau=\Delta T_0,\qquad \ell=\Delta L\in\{0,1\}.
\]
We prove the division-free top-ratio inequality
\begin{equation}\label{pgr:u:t26:eq:ratio}
T_{\rm new}B-TB_{\rm new}
=\tau(B_0+L)-T_0(\beta+\ell)\ge0.
\end{equation}

\paragraph{Exterior-only contribution.}
Form the transversal matroid whose elements are the old exterior vertices, $x$, and three private dummies $d_0,d_1,d_2$, each $d_i$ adjacent only to $q_i$. Its independent sets are those matchable injectively into the three heads; the dummies ensure rank exactly three. Let $f$ be its multiaffine basis polynomial, counting each feasible basis once. Wagner's theorem \cite{WagnerR3} gives, for distinct elements $e,h$,
\[
f_e f_h-f f_{eh}\ge0
\]
for positive weights and hence, by continuity, for nonnegative weights. Subscripts here mean partial derivatives.

Evaluate with all old exterior weights one and the weights of $x,d_0,d_1,d_2$ zero. Then
\[
f=T_0,\qquad f_x=\tau,\qquad
\sum_{i=0}^{2}f_{d_i}=B_0,\qquad
\sum_{i=0}^{2}f_{x d_i}=\beta.
\]
For example, a basis using $d_i$ corresponds to an exterior two-tail support on the complementary pair of heads; that omitted head is unique. This preserves support counting even when a tail pair admits several choices of head pair. Summing the three Rayleigh inequalities for $x,d_i$ gives
\begin{equation}\label{pgr:u:t26:eq:rayleigh}
\tau B_0-T_0\beta\ge0.
\end{equation}

\paragraph{Contribution of the internal edge.}
Let $K$ count exterior two-element supports on the fixed head pair $\{q_0,q_2\}$. Then
\begin{equation}\label{pgr:u:t26:eq:LK}
T_0\le LK.
\end{equation}
For every feasible exterior triple choose one matching witness, and record the vertex assigned to $q_1$ and the remaining pair. These data determine the triple, giving an injection into the pairs counted by $LK$. If $\ell=1$, adjoining $x$ to each pair counted by $K$ creates a distinct feasible triple, so $\tau\ge K$ and $\tau L-T_0\ell\ge LK-T_0\ge0$. If $\ell=0$, this expression is simply $\tau L\ge0$. Adding this to \eqref{pgr:u:t26:eq:rayleigh} proves \eqref{pgr:u:t26:eq:ratio}. An isolated new vertex gives equality. No positivity or division is needed here.

\subsection{Compatibility and the final square}
If $R_3=0$, then \eqref{pgr:u:t26:eq:coeffs} gives $\gamma_4=0$, so \eqref{pgr:u:t26:eq:target} is immediate. Assume $R_3>0$. Some $Q_{-x}$ has a positive cubic coefficient. Restoring a vertex preserves its supports; restricting a matching witness gives a quadratic support. Hence $T>0$ and $B>0$.

For $x\in H$, write $B_x=[t^2]Q_{-x}$ and $T_x=[t^3]Q_{-x}$. Apply \eqref{pgr:u:t26:eq:ratio} first when restoring $x$, then when adjoining the universal exterior vertex that defines $\Qp$:
\[
T_xB\le B_xT,\qquad T\Bp\le B\Tp.
\]
Multiply by nonnegative factors to obtain
\[
T_xB\Bp\le B_xT\Bp\le B_xB\Tp.
\]
Cancel only the strictly positive factor $B$. Summing over $x\in H$ yields
\begin{equation}\label{pgr:u:t26:eq:compatibility}
R_3\Bp\le R_2\Tp.
\end{equation}
The argument covers $B_x=0$ or $T_x=0$ without division by either. An isolate deletion gives equality in the first comparison.

By Lemma \ref{pgr:u:t26:lem:mixtures}, $R_1R_3\le R_2^2/3$. Combining this with \eqref{pgr:u:t26:eq:coeffs} and \eqref{pgr:u:t26:eq:compatibility} proves
\begin{align*}
3\gamma_3^2-8\gamma_2\gamma_4
 &=3(R_2+\Tp)^2-8(R_1+\Bp)R_3\\
 &\ge 3(R_2+\Tp)^2-\frac83 R_2^2-8R_2\Tp\\
 &=\frac13(R_2-3\Tp)^2\ge0.
\end{align*}
This completes the proof of Theorem \ref{pgr:u:t26:thm:main}.

\subsubsection*{Scope and reproducibility}
This supplement proves the last gap. Together with the separately approved middle-gap proof and general first-gap bound, it gives the smaller proof for this family. The original degree-four theorem and frozen certificate archive are unchanged. There is no claim of a general weighted degree-four theorem, fractional-population theorem, real-rootedness, or novelty.

\mnote{the middle-gap proof meant here is source~16's certificate for
template~26, gap two (shipped as \texttt{data/16-degree-four-a4-gram\_za\_26.json}),
and the first-gap bound is Theorem~\ref{pgr:thm:first}. The embedded
\file{weighted-degree-three-package.zip} is an exact copy of archive~31,
superseded by source~22, and is not shipped; because source~23's
\texttt{verify.py} checks it against its checksum list, the delivered
wrapper runs only on the reconstructed package (README). The same Windows
path problem as in source~16's wrapper applies.}


The tradeoff is explicit: ten quartic certificates replace a much larger direct last-gap certificate, using the already approved vertex-weighted degree-three theorem. The supplement includes the certificates, independent exact checkers, approval receipts, and \file{weighted-degree-three-package.zip}. From the extracted supplement root, run \verb|python verify.py| for the two local independent replays and the pinned dependency archive. Run \verb|python verify.py --verify-dependency| to extract and replay the full portable verifier of that dependency as well.
The local replays check the exact quartic certificates and the support decomposition coefficientwise. The full mathematical assembly has independent approval, including the convex-hull argument, degree-drop endpoints, internal-edge correction, zero cases, and final square.


\part{Independent role activities and small covers}\label{pgr:part:two}

This Part allows independent nonnegative tail and head activities $u_i,v_i$
(Definition~\ref{pgr:def:support}) on arbitrary loopless directed
relations, transitive or not. Five sources contribute.
\begin{itemize}[leftmargin=*]
\item Section~\ref{pgr:r:rc:sec} (source~20): if every arc has its tail in
  a set $P$ or its head in a set $Q$ with $|P|,|Q|\le2$ (overlap allowed),
  the signed monomer polynomial is real stable, so $\Ga_D$ is
  real-rooted.
\item Section~\ref{pgr:r:tc:sec} (source~12): coefficientwise Rayleigh
  inequalities in the core monomers when the physical graph is bipartite
  with a three-vertex shore (computer-assisted boundary inequality).
\item Section~\ref{pgr:r:pc:sec} (source~11): rank-ULC whenever the
  underlying graph has a vertex cover of at most three vertices, internal
  core arcs allowed; it uses source~12.
\item Section~\ref{pgr:r:ps:sec} (source~15): a Lorentzian theorem for
  bipartite physical graphs, and real stability for pendant uncovered heads
  over a stable transversal seed.
\item Section~\ref{pgr:r:ir:sec} (source~22): rank-ULC for every finite
  preorder of actual degree at most three under independent role
  activities, by a Gallai--Edmonds reduction, ordinary structural theorems,
  a complete ledger of role certificates and half-plane-property side
  matroids.
\end{itemize}
Several statements of source~22 are special cases of statements of sources
11, 15 and~20, and source~25's theorem (Part~I) is the unit case of
source~22's; these are recorded where they occur. In the bipartite case,
with every vertex in one role, these results specialize to theorems of
\prp{} and \mrn: two-shore ultra-log-concavity (\prp,
\texttt{lor:thm:weightedULC}), weighted matching rank three and five
(\prp, \texttt{lc:cor:rankthree}; \mrn, \texttt{mrn:w:rf:cor:rankfive}),
scalar rank-three Rayleigh (\mrn, \texttt{mrn:s:tv:eq:rayleigh}) and
stability under a two by two cover (\mrn, \texttt{mrn:w:sb:thm:main}). The
directed, overlapping and coefficientwise forms are new.

\section{Two by two role covers (source 20)}\label{pgr:r:rc:sec}

Source~20 is \emph{Real stable monomers under a two by two role cover}
(9 pages as delivered; archive \texttt{role-cover-package}). Its files are
shipped with the prefix \texttt{20-role-cover-}. Source~16's Section~3 and
Appendix~A are this article, with the differences printed below as
``Second route (source~16)''; exact copies of the whole package were also
embedded in source~22 and archive~31.

\paragraph{Abstract of source 20, as delivered.}

We count ordered disjoint tail and head supports of directed matchings, with every feasible support counted once. For any finite loopless directed relation whose bipartite role graph has a vertex cover comprising at most two tail copies and at most two head copies, we prove real stability of the signed physical monomer polynomial. The selected physical sets may overlap, and arbitrary nonnegative independent tail and head activities are allowed. Consequently the support polynomial has only negative real roots and satisfies Newton's inequalities at its actual surviving degree. The proof treats every subset of a two-source, two-sink core, using rank-two stable quadratics, stability-preserving operators, and explicit all-real Rayleigh square identities, before merging role copies to enforce physical disjointness. We also obtain a rank-two matroid free-product corollary and an elementary root transfer to gamma expansions. These are ordinary proofs; the accompanying exact algebra and finite Hall-support checks are supplementary. No proof-assistant formalization or literature-wide priority claim is made.

\subsection{Support polynomials and normalization}
\begin{definition}
A feasible support of size $k$ is an ordered pair $(S,T)$ of disjoint subsets of $V$, each of cardinality $k$, for which there exists a bijection $\pi:S\to T$ with $s\to\pi(s)$ for every $s\in S$. Write
\[
 \gamma_k(D)=\#\{(S,T): (S,T)\text{ is feasible of size }k\},\qquad
 \Ga_D(t)=\sum_{k\ge0}\gamma_k(D)t^k.
\]
The empty support contributes $\gamma_0=1$. A support contributes once, regardless of the number of witnessing bijections.
\end{definition}
The same-ground-set disjointness is essential. Passing to independent tail and head copies without subsequently excluding their simultaneous use changes this polynomial. Counting individual matchings also changes it.

For nonnegative tail activities $u_i$ and head activities $v_i$, replace the contribution of $(S,T)$ by
\[
 u^S v^T=\prod_{i\in S}u_i\prod_{j\in T}v_j.
\]
We use the same notation when the weights are specified. A vertex may have different activities in its two roles. With zero activities, the polynomial degree can decrease.


\mnote{source~20's Lemma ``degree'' is Lemma~\ref{pgr:lem:degree}, printed
once in Section~\ref{pgr:sec:object} with this proof.}


\begin{definition}
A polynomial $g(t)=\sum_{k=0}^d g_k t^k$ is \emph{rank-ultra-log-concave}, or rank-ULC, if its coefficients are nonnegative with no internal zeros and
\[
 \left(\frac{g_k}{\binom dk}\right)^2\ge
 \frac{g_{k-1}}{\binom d{k-1}}\frac{g_{k+1}}{\binom d{k+1}}
 \quad(1\le k<d),
\]
where $d$ is the actual degree.
\end{definition}

Our main theorem concerns directed relations satisfying the two by two role-cover condition in Theorem~\ref{pgr:r:rc:thm:rolecover}; transitivity is not required. Throughout, the monomer variable of a vertex records that the vertex is unused. Real-rootedness is always interpreted at the degree remaining after zero activities are imposed.
\subsection{Real-rootedness for a complete balanced core}
Let $P=\{p_1,p_2\}$, $Q=\{q_1,q_2\}$, and let $I$ be an independent exterior set. Include all four arcs $P\to Q$. Each $x\in I$ may have arbitrary incoming neighbors in $P$ and outgoing neighbors in $Q$, including an empty neighbor set on either side. There are no other arcs. This directed family includes every nonnegative integer population of all fifteen nonempty exterior types.

A real polynomial is \emph{real stable} if it is nonzero whenever every variable lies in the open upper half-plane $\HH$. We use the finite-degree stability-preserver theorem of Borcea and Br\"and\'en \cite{BB}; Wagner's survey gives the same plus-sign algebraic-symbol convention \cite[Theorem 5.2]{WagnerSurvey}.

\begin{theorem}\label{pgr:r:rc:thm:balanced}
For the complete balanced core described above, with arbitrary nonnegative independent tail and head activities, the signed multivariate monomer polynomial
\[
 \mu_D(\mathbf z)=\sum_{(S,T)\text{ feasible}}(-1)^{|S|}u^S v^T
                   \prod_{i\notin S\cup T}z_i
\]
is real stable. Consequently every root of $\Ga_D$ is negative real, and $\Ga_D$ is rank-ULC at its actual degree. No transitivity assumption is required.
\end{theorem}

\subsubsection{A stable coupling of two rank-two matroids}
A loopless rank-two matroid has parallel classes with variable sums $r_1,\ldots,r_j$. Its basis polynomial and element sum are
$B=e_2(r_1,\ldots,r_j)$ and $L=\sum_i r_i$.

\begin{lemma}\label{pgr:r:rc:lem:quadratic}
The polynomial $f(s)=s^2+\sqrt2 Ls+B$ is real stable.
\end{lemma}
\begin{proof}
The identity
\[
 f(s)=(s+L/\sqrt2)^2-\frac12\sum_i r_i^2
\]
shows that its quadratic form in $s,r_1,\ldots,r_j$ has exactly one positive direction. It has nonnegative coefficients and is positive on the positive orthant. For an alleged zero at $a+ib$ with $b$ strictly positive, its imaginary part says that $a$ and $b$ are orthogonal for the form. Since $f(b)>0$, the orthogonal complement of $b$ is nonpositive. Thus $f(a)\le0$, contradicting the real-part equation $f(a)=f(b)$. Replacing each $r_i$ by its parallel-class sum preserves stability.
\end{proof}
Apply the lemma to two matroids on disjoint variable sets. On auxiliary monomials of separate degrees at most two define the coefficientwise linear operator
\[
 \mathcal T(1)=1,\quad \mathcal T(st)=-\tfrac12,\quad
 \mathcal T(s^2t^2)=1,
\]
and let it vanish on the other six auxiliary monomials. Its symbol is
\[
 \mathcal T[(s+a)^2(t+b)^2]=(ab-1)^2.
\]
This symbol is stable: $ab=1$ is impossible for $a,b\in\HH$, since $1/a$ lies in the lower half-plane. The full symbol, including all untouched variables, acquires only stable factors of the form $(z+w)^\kappa$. Therefore
\begin{equation}\label{pgr:r:rc:eq:coupling}
 B_PB_Q-L_PL_Q+1
\end{equation}
is real stable. The constant term shows that it is not identically zero.

\paragraph{Second route (source 16): elementary operators.} Source~16
does not invoke the stability-preserver theorem. It writes: ``We use the
standard stability closures under products, diagonalization,
differentiation, nonnegative linear substitutions, and real specialization
(with zero outputs allowed) \cite{WagnerSurvey}. The special contraction
operators needed here have elementary realizations; they also fit the
general stability-preserver framework \cite{BB}.'' It then proves the
coupling as follows.

\begin{lemma}[Elementary Asano contraction; source 16]\label{pgr:u:d4:lem:asano}
For a stable polynomial affine in each of $a,b$, write
$f=abA+aB+bC+D$. The coefficientwise operation $f\mapsto D-A$ preserves
stability or gives zero. So does $f\mapsto f-\partial_a\partial_bf$.
\end{lemma}
\begin{proof}
For $t\in\HH$, both $t$ and $-1/t$ lie in $\HH$. Therefore
\[
 t f(t,-1/t)=Bt^2+(D-A)t-C
\]
is stable in $t$ and all untouched variables. Differentiate in $t$ and
specialize $t=0$ to obtain $D-A$. The standard boundary-specialization
closure permits the zero polynomial. For the second assertion, first
replace $a,b$ by $a+x,b+y$, which preserves stability, and contract the
fresh variables $x,y$. Their constant coefficient is $f(a,b)$ and their
product coefficient is $\partial_a\partial_bf$.
\end{proof}

For each rank-two side use the polarized auxiliary quadratic
\[
 F_P(s_1,s_2)=s_1s_2+\frac{L_P}{\sqrt2}(s_1+s_2)+B_P.
\]
It is stable by the same one-positive-direction argument as
Lemma~\ref{pgr:r:rc:lem:quadratic}: it equals the original auxiliary
quadratic at $(s_1+s_2)/2$, minus $(s_1-s_2)^2/4$. The extra independent
direction is negative, and every coefficient of the expanded polynomial is
nonnegative. Construct $F_Q(t_1,t_2)$ similarly.

Apply Lemma~\ref{pgr:u:d4:lem:asano} to the stable product $F_PF_Q$, first
in $(s_1,t_1)$ and then in $(s_2,t_2)$. With $a=L_P/\sqrt2$ and
$b=L_Q/\sqrt2$, the result is $B_PB_Q-2ab+1$. Consequently
\eqref{pgr:r:rc:eq:coupling} is real stable. Its constant coefficient one
rules out a zero output at either contraction stage. This proves the
coupling without invoking the general classification of
stability-preserving linear operators.


\begin{corollary}\label{pgr:r:rc:cor:freeproduct}
If $M$ and $N$ are loopless rank-two matroids on disjoint ground sets, then the basis polynomial of the free product $M\mathbin{\square}N^*$ is real stable.
\end{corollary}
\begin{proof}
Write the ground sets as $S,T$, with variables $\mathbf x,\mathbf y$. By the basis characterization of the free product \cite[Proposition 1 and the following paragraph]{CS}, its bases are $I\cup(T\setminus K)$, where $I$ is independent in $M$, $K$ is independent in $N$, and $|I|=|K|\le2$. Indeed, $T\setminus K$ spans $N^*$ exactly when $K$ is independent in $N$, and the total basis size then forces the cardinalities to agree. Consequently its basis polynomial is
\[
 \left(\prod_{j\in T}y_j\right)
 \left[1+L_M(\mathbf x)L_N(\mathbf y^{-1})
          +B_M(\mathbf x)B_N(\mathbf y^{-1})\right].
\]
This is \eqref{pgr:r:rc:eq:coupling} after substituting $-1/y_j$ for each $N$-variable and multiplying by $\prod_{j\in T}y_j$. Multiaffinity cancels the denominators. Since $-1/y_j\in\HH$ for $y_j\in\HH$, the resulting polynomial is real stable.
\end{proof}

\subsubsection{The two-copy identity}
Temporarily split each exterior vertex into a head copy, adjacent from $P$, and a tail copy, adjacent to $Q$. Keep isolated copies. For each two-center side use its rank-two transversal matroid on the nonisolated copies and two private core dummies. Isolated copies contribute only their factors to $Z$ and are omitted from both $B$ and $L$. The private-neighborhood copies and the corresponding dummy form a parallel class; each common-neighborhood copy forms its own class.

Give a core dummy its monomer variable $z_p$ or $z_q$, and give an active exterior copy the value $-1/z_x$. These substitutions preserve the upper half-plane. Let $Z$ be the product of all exterior-copy monomer variables. The unweighted signed monomer polynomial of the split relation is
\begin{equation}\label{pgr:r:rc:eq:copy}
 \mu_{\rm copy}(\mathbf z)=Z(B_PB_Q-L_PL_Q+1).
\end{equation}
The possible denominators cancel because each exterior-copy variable is multiaffine before substitution.

Here is the support-count explanation. The number of core arcs in a witnessing matching is determined by its endpoint support. Zero core arcs gives $ZB_PB_Q$. One core arc gives $-ZL_PL_Q$: on each side, the remaining core monomer can be either dummy, while an exterior singleton used together with both core vertices is counted once, not by its neighbor degree. Two core arcs use all four core vertices and no exterior vertex, contributing $Z$ once. Although the complete core admits two perfect matchings, it has only one such endpoint support. Completeness of $P\to Q$ makes the separate side conditions sufficient.

Formula \eqref{pgr:r:rc:eq:copy} and \eqref{pgr:r:rc:eq:coupling} prove split monomer stability. Positive activities $w_i$ on split vertices are introduced by
\[
 \mu_w(\mathbf z)=\left(\prod_i w_i\right)
                  \mu_{\rm copy}\bigl((z_i/w_i)_i\bigr).
\]
Assign the original tail activity to a tail copy, the head activity to a head copy, $u_p$ to a source core vertex, and $v_q$ to a sink core vertex.

\subsubsection{Merging the physical vertices}
For the two monomer variables $a,b$ of one physical exterior vertex apply
\begin{equation}\label{pgr:r:rc:eq:merge}
 ab\longmapsto z,\qquad a\longmapsto1,\qquad b\longmapsto1,
 \qquad1\longmapsto0.
\end{equation}
Both unused copies become one unused physical vertex. Exactly one used copy becomes a used vertex. Two used copies are discarded. Contributions for the two different one-used roles add, as the support definition requires. The algebraic symbol is
$\mathcal M[(a+r)(b+s)]=z+r+s$, which is stable. Including untouched variables only adds stable factors. Thus every merge preserves stability. After all merges we obtain precisely $\mu_D$.

\paragraph{Second route (source 16): the merge by differentiation.} There
is also an elementary realization of this operator. For $f=abA+aB+bC+D$,
with coefficients independent of $a,b$, diagonalization gives
$f(t,t)=t^2A+t(B+C)+D$. Differentiate in $t$ and substitute $t=z/2$ to
obtain $zA+B+C$, exactly \eqref{pgr:r:rc:eq:merge}. Diagonalization,
differentiation, and positive scaling preserve real stability or give
zero. Thus this particular merge needs only these standard closure
properties, and the empty-support monomial excludes zero. (Source~15's
Lemma~\pgrfwd{pgr:r:ps:lem:merge}{``Core merge''} gives the same elementary
realization; source~22 uses the directional-derivative form.)


Zero activities follow by continuity from positive ones. The empty-support monomial has coefficient one throughout, so the limit is nonzero. Finally, with $N$ physical vertices,
\[
 \mu_D(s,\ldots,s)=\sum_k(-1)^k\gamma_k s^{N-2k}
                  =s^N\Ga_D(-s^{-2}).
\]
The left side has only real zeros. A nonreal or positive zero of $\Ga_D$ would produce a nonreal solution of $s^2=-1/\rho$ and hence a nonreal zero of this monomer polynomial. Since $\Ga_D(0)=1$, all its roots are negative. Newton's inequalities finish the proof of Theorem~\ref{pgr:r:rc:thm:balanced}.

\subsubsection{A two-edge core extension}
The same physical-copy construction also works when the core has only the two arcs $p_2\to q_1,p_2\to q_2$.
\begin{theorem}\label{pgr:r:rc:thm:twoedge}
For this two-edge core, arbitrary exterior neighborhoods and arbitrary nonnegative independent role activities again give a real-stable signed monomer polynomial. The support polynomial is therefore real-rooted and rank-ULC at its actual degree.
\end{theorem}
\begin{proof}
With the same rank-two side matroids, put $A_P=\partial_{d_{p_2}}B_P$. This is the dummy $d_{p_1}$ plus the sum of all exterior head-copy variables adjacent to $p_1$, including common-neighbor copies once. The split kernel is
\[
 H=B_PB_Q-A_PL_Q.
\]
At most one core edge can occur. If it does, its tail is forced to be $p_2$; the remaining $P$-side support is exactly $A_P$, and the remaining $Q$-side support is exactly $L_Q$. This proves the identity at support level.

For the parallel-class sums of the $Q$ matroid, $B_Q+sL_Q=e_2(r_1,\ldots,r_j,s)$ is stable. Fix its original arguments in $\HH$. Both $B_Q$ and $L_Q$ are nonzero, and its root $-B_Q/L_Q$ cannot be strictly upper. Hence $-L_Q/B_Q$ has nonnegative imaginary part. Since $B_P$ is affine in $d_{p_2}$,
\[
 H=B_Q\,B_P\bigl(d_{p_2}-L_Q/B_Q\bigr),
\]
with all other arguments unchanged. The shifted dummy remains strictly upper, so this expression is nonzero. The reciprocal substitutions, activity scalings, physical-copy merges, and monomer-to-support root transfer from Theorem~\ref{pgr:r:rc:thm:balanced} now apply unchanged. Zero activities follow by the same nonzero limit argument.
\end{proof}
The proof does not assume that adding a polynomial with positive coefficients preserves stability.
\mnote{source~16 adds: ``This includes canonical template~6 of the
degree-four classification'' (Section~\ref{pgr:u:d4:sec}).}

\subsubsection{A core with one missing edge}
\begin{theorem}\label{pgr:r:rc:thm:missing}
Keep the three core arcs $p_1\to q_1,p_1\to q_2,p_2\to q_1$, and omit $p_2\to q_2$. With arbitrary exterior neighborhoods and nonnegative independent role activities, the signed monomer polynomial is real stable and the support polynomial is real-rooted at its actual degree.
\end{theorem}
The split kernel is
\begin{equation}\label{pgr:r:rc:eq:missingkernel}
 H=B_PB_Q-L_PL_Q+1+r_Pr_Q,
\end{equation}
where $r_P$ is the parallel-class sum containing dummy $p_1$ and the head copies private to $p_1$, and $r_Q$ is defined similarly using $q_1$. A one-core-edge support is lost from the complete core exactly when its residual source and sink are forced to be $p_2,q_2$. These forced choices contribute $r_Pr_Q$; their former signed contribution was negative, so the correction is positive. Common-neighbor copies do not force that residual choice. The two-core-edge endpoint support survives through the off-diagonal matching, and contributes one.

We prove the stability of \eqref{pgr:r:rc:eq:missingkernel} for arbitrary loopless rank-two matroids with a distinguished class on each side. Write their independent class variables as $p_0,\ldots,p_a$ and $q_0,\ldots,q_b$, with $r=p_0$, $s=q_0$, and set
\[
 F=e_2(\mathbf p),\quad G=e_2(\mathbf q),\quad
 L=\sum p_i,\quad M=\sum q_j,\quad H=FG-LM+1+rs.
\]
For a real multiaffine polynomial the all-real Rayleigh criterion says that stability is equivalent to nonnegativity of every
$\Delta_{xy}=(\partial_xH)(\partial_yH)-H\partial_x\partial_yH$
at every real assignment \cite[Theorem 5.6]{BrandenHPP}; see also \cite[Theorem 3.1]{WagnerSurvey}. Testing only positive assignments would not suffice. Here five identities cover every pair type, up to exchange of sides and permutation of nondistinguished classes.

\paragraph{One distinguished variable on the same side.}
For $x=p_0,y=p_i$ with $i>0$, put
$A=\sum_{k\ne0,i}p_k$, $U=\sum_{k\ne0,i}p_k^2$,
$B=\sum_{l>0}q_l$, $V=\sum_{l>0}q_l^2$. Then
\[
 2\Delta_{xy}=(GA-B)^2+G^2U+V.
\]
\paragraph{Two nondistinguished variables on the same side.}
For $x=p_i,y=p_j$ with distinct $i,j>0$, put
$A=\sum_{k\ne i,j}p_k$, $U=\sum_{k\ne0,i,j}p_k^2$,
$V=\sum_{l>0}q_l^2$. Then
\[
 2\Delta_{xy}=(GA-M)^2+(Gr-s)^2+G^2U+V.
\]
\paragraph{Both distinguished variables on opposite sides.}
For $x=p_0,y=q_0$, put
$A=\sum_{k>0}p_k$, $U=\sum_{k>0}p_k^2$,
$B=\sum_{l>0}q_l$, $V=\sum_{l>0}q_l^2$. Then
\[
 2\Delta_{xy}=A^2V+B^2U.
\]
\paragraph{One distinguished variable on opposite sides.}
For $x=p_0,y=q_j$ with $j>0$, put
$A=\sum_{k>0}p_k$, $U=\sum_{k>0}p_k^2$,
$B=\sum_{l\ne j}q_l$, $V=\sum_{l\ne0,j}q_l^2$. Here $B$ includes $s$. Then
\[
 4\Delta_{xy}=[A(B+s)-2]^2+A^2V+U[(B-s)^2+V].
\]
\paragraph{Neither variable distinguished on opposite sides.}
For $x=p_i,y=q_j$ with $i,j>0$, put
$A=\sum_{k\ne i}p_k$, $U=\sum_{k\ne0,i}p_k^2$,
$B=\sum_{l\ne j}q_l$, $V=\sum_{l\ne0,j}q_l^2$. Thus $A,B$ include $r,s$. Then
\[
 4\Delta_{xy}=(AB-rs-2)^2+(As-rB)^2
             +U(B^2+s^2)+V(A^2+r^2)+UV.
\]
Each identity follows by expansion and
$2e_2(\mathbf z)=(\sum z_i)^2-\sum z_i^2$.
Every expression is nonnegative for all real variables: $U,V$ are sums of squares, so their displayed products are sums of squares as well. Empty sums are zero. The constant term one makes $H$ nonzero. The Rayleigh criterion proves stability, and substitution of parallel-class sums proves it in the individual matroid variables. The reciprocal, activity, and merge steps of Theorem~\ref{pgr:r:rc:thm:balanced} complete the proof of Theorem~\ref{pgr:r:rc:thm:missing}.
\mnote{source~16 adds: ``This covers canonical template~21.''}

\begin{remark}[Interpolation]
The kernels $H_\theta=FG-LM+1+\theta rs$ are stable for $0\le\theta\le1$. In the independent class variables, every Rayleigh difference of the monomial $rs$ is zero. Thus $\Delta_{xy}(H_\theta)$ is the convex combination of the two endpoint differences. The same linear monomer and copy-merge transformations show that the corresponding interpolated support polynomials remain real-rooted. This is a statement about this particular interpolation, not a general convexity claim for stable polynomials.
\end{remark}

\subsubsection{Arbitrary cross-core patterns and role covers}
For two disjoint core edges $p_1\to q_1,p_2\to q_2$, the split kernel is
\[
 H=B_PB_Q-L_PL_Q+1+r_{P1}r_{Q2}+r_{P2}r_{Q1},
\]
where each $r_{Pi},r_{Qj}$ is the private-plus-dummy parallel-class sum. Each omitted cross edge removes its uniquely forced residual one-core-edge supports. Those two lost sets are disjoint, and the diagonal perfect support survives. Section~\ref{pgr:r:rc:app:disjoint} proves this kernel stable by seven all-real Rayleigh identities. The weighted and physical-merge conclusions follow as before.

The empty core gives the stable product $B_PB_Q$. A single core edge gives
\[
 (1-\partial_{d_{p_i}}\partial_{d_{q_j}})(B_PB_Q).
\]
On polynomials affine in these two dummy variables, the operator's symbol is
$(a+x)(b+y)-1$, which is stable because two upper-half-plane numbers cannot have product one. These cases, the two-edge star and its reversal, the disjoint pair, the one-missing-edge case, and the complete core exhaust all sixteen subsets of the four possible core arcs. Thus every such subset has a stable split monomer polynomial.

\paragraph{Second route (source 16).} Lemma~\ref{pgr:u:d4:lem:asano},
applied after translation by fresh variables, proves that this operator
preserves stability or gives zero. The all-dummy monomial excludes zero.


\begin{theorem}[Two by two role cover]\label{pgr:r:rc:thm:rolecover}
Let $D$ be any finite loopless directed relation on $V$. Suppose there are sets $P,Q\subseteq V$, each of size at most two, such that every arc $i\to j$ has $i\in P$ or $j\in Q$. The two sets may overlap. Then for arbitrary nonnegative independent tail and head activities, the signed physical monomer polynomial is real stable. Its support polynomial has only negative real roots and is rank-ULC at its actual degree.
\end{theorem}
\begin{proof}
Replace every physical vertex by distinct tail and head copies. The selected tail copies of $P$ and head copies of $Q$ form a bipartite vertex cover. They are disjoint as copies even when their original physical sets overlap. All other copies are exterior; no edge joins two exterior copies. The core is an arbitrary subset of the four possible cross edges, so the preceding classification proves split monomer stability. If a cover side is smaller than two, pad it with isolated auxiliary vertices. Removing their monomer factors preserves nonvanishing in the upper half-plane.

Give every tail copy its tail activity and every head copy its head activity. Merge the two copies of every physical vertex by \eqref{pgr:r:rc:eq:merge}, including core/core and core/exterior pairs. An original disjoint support lifts uniquely to its tail and head copies. A split support survives precisely when it does not use both copies of any physical vertex. Matching feasibility and its product weight are unchanged. Distinct original role assignments can have the same unused-vertex monomial and their coefficients correctly add. This is an exact support bijection, not a count of matching witnesses. Each merge preserves stability, and the empty-support monomial excludes the zero polynomial. The root transfer and zero-activity limit finish the proof.
\end{proof}
The criterion permits arbitrary arcs inside each of two disjoint two-vertex core sides, as well as arbitrary $P\to Q$ core arcs and $P\to I\to Q$ exterior attachments. It does not permit an arbitrary backward $Q\to P$ arc. It also shows why an actual degree below four causes no normalization problem: real-rootedness applies directly to the surviving polynomial.

\mnote{for a bipartite graph oriented from one shore to the other, every
vertex has one role, a cover with two vertices on each shore is a role
cover with $P\cap Q=\varnothing$, vertex activities are role activities,
and $\Ga_D$ is the matching-support polynomial. That case is \mrn,
Theorem \texttt{mrn:w:sb:thm:main}, proved there by a different route
(vertex sums, expected determinants and Grace--Walsh--Szeg\H{o}
polarization); batch-73 manuscript~26, printed in \mrn{} as its source~77
(Theorem \texttt{mrn:v:rf4:thm:mixed}), proves the same undirected
statement with the weaker conclusion of real-rootedness by a third route. Theorem~\ref{pgr:r:rc:thm:rolecover} generalizes it to
directed relations, overlapping $P$ and $Q$, arcs inside $P$ and inside
$Q$, and independent role activities. Source~22 uses it in its role-cover
corollary and in the three-attachment branch of its global proof
(Section~\pgrfwd{pgr:r:ir:sec}{on source~22}); source~16's degree-four ledger
uses its core cases for templates 6 and~21. The free-product corollary has
no counterpart in \prp{} or \mrn.}


\begin{corollary}[Transfer to the gamma expansion]\label{pgr:r:rc:cor:transfer}
Under Theorem~\ref{pgr:r:rc:thm:rolecover}, let $d=\deg\Ga_D$ and choose any integer $N\ge2d$. Put
\[
 h_D(t)=\sum_k\gamma_k t^k(1+t)^{N-2k}.
\]
Then $h_D$ has only negative real roots. The choice $N=|V|$ is permitted. The same conclusion applies to any polynomial independently identified with this gamma expansion; no geometric convention is assumed here.
\end{corollary}
\begin{proof}
Write $\Ga_D(z)=\prod_{i=1}^d(1+\lambda_i z)$ with $\lambda_i>0$. Since $2d\le N$,
\[
 h_D(t)=(1+t)^{N-2d}\prod_{i=1}^d\bigl[t^2+(2+\lambda_i)t+1\bigr].
\]
Each quadratic has discriminant $\lambda_i(\lambda_i+4)>0$, negative root sum, and positive root product, so both roots are negative. The remaining roots are $-1$. This is the standard elementary root transfer, not a separate novelty claim.
\end{proof}

\mnote{source~16 states this corollary with $N=|V|$ only and ends: ``In
particular this applies to a preorder $h$-polynomial when its gamma
expansion is the support expansion identified in the source report''
(\prp, Theorem \texttt{gam:thm:main}).}


\subsection{Proof and verification boundary}
The arguments above are ordinary mathematical proofs. The all-real Rayleigh identities hold for any number of parallel classes; their validity is not inferred from a bounded computation. The finite-degree stability-preserver theorem, the all-real multiaffine Rayleigh criterion, and elementary rank-two matroid facts are the external mathematical dependencies.

The supplementary archive contains the approved proof notes and scoped review receipts, two independent exact symbolic checkers for the twelve aggregate-variable identities, and four independent Hall-based support checkers. Each Hall checker compares weighted multivariate monomer coefficients for all $4845$ exterior-type multisets through four vertices, including isolated copies and zero activities: $14535$ activity-weighted identities per core pattern. These finite checks test the combinatorial formulas; they are not a substitute for the arbitrary-population proofs. The archive includes no numerical optimization or binary executable as a logical dependency. Its portable README describes how to rebuild the article and rerun the exact checks.

\mnote{source~20's files are shipped with the prefix \texttt{20-role-cover-}
(proof notes at the top of the report directory, checkers in
\texttt{code/}, receipts and logs in \texttt{data/}). At intake its
\texttt{verify.py} passed on a copy of the delivered package in 73~s; with
GCC~16 its four C++ checkers need \texttt{-include cassert}, because they
call \texttt{assert} without including \texttt{<cassert>}.}


\subsection{Appendix of source 20: Rayleigh identities for two disjoint core edges}\label{pgr:r:rc:app:disjoint}

\mnote{the typographical slip ``labels0,1'' of source~20 is corrected as in source~16.}

Use independent parallel-class variables $p_0,p_1,p_2,\ldots$ and $q_0,q_1,q_2,\ldots$. Label the first two $Q$ classes so that the two correction terms are $p_0q_0+p_1q_1$. Put
\[
 F=e_2(\mathbf p),\quad G=e_2(\mathbf q),\quad
 L=\sum p_i,\quad M=\sum q_j,\quad K=p_0q_0+p_1q_1,
 \quad H=FG-LM+1+K.
\]
The following seven cases cover every pair under side interchange, simultaneous exchange of the partnered labels $0,1$, and permutations of the remaining classes. Every displayed square sum is nonnegative on all real assignments. The identities follow from the elementary symmetric identity used above, with no bound on the number of classes.

\paragraph{Same side with both distinguished.}
For $x=p_0,y=p_1$, take $A=\sum_{k\ge2}p_k$, $U=\sum_{k\ge2}p_k^2$, $B=\sum_{l\ge2}q_l$, $V=\sum_{l\ge2}q_l^2$. Then
\[
 2\Delta_{xy}=(GA-B)^2+G^2U+V.
\]
\paragraph{Same side with one distinguished.}
For $x=p_0,y=p_i$, $i\ge2$, take $A=\sum_{k\ne0,i}p_k$, $U=\sum_{k\ge2,k\ne i}p_k^2$, $V=\sum_{l\ge2}q_l^2$. Then
\[
 2\Delta_{xy}=[GA-(M-q_0)]^2+(Gp_1-q_1)^2+G^2U+V.
\]
\paragraph{Same side with neither distinguished.}
For $x=p_i,y=p_j$, $i,j\ge2$ and $i\ne j$, take $A=\sum_{k\ne i,j}p_k$, $U=\sum_{k\ge2,k\ne i,j}p_k^2$, $V=\sum_{l\ge2}q_l^2$. Then
\[
 2\Delta_{xy}=(GA-M)^2+(Gp_0-q_0)^2+(Gp_1-q_1)^2+G^2U+V.
\]
\paragraph{Opposite sides with partnered distinguished variables.}
For $x=p_0,y=q_0$, take $A=\sum_{k\ne0}p_k$, $U=\sum_{k\ge2}p_k^2$, $B=\sum_{l\ne0}q_l$, $V=\sum_{l\ge2}q_l^2$. Then
\[
 2\Delta_{xy}=(Aq_1-Bp_1)^2+A^2V+B^2U.
\]
\paragraph{Opposite sides with unpartnered distinguished variables.}
For $x=p_0,y=q_1$, let $u=p_1$, $v=q_0$, $A=\sum_{k\ne0}p_k$, $U=\sum_{k\ge2}p_k^2$, $B=\sum_{l\ne1}q_l$, $V=\sum_{l\ge2}q_l^2$. Then
\[
 4\Delta_{xy}=(AB+uB+vA-uv-2)^2+U(B-v)^2+V(A-u)^2+UV.
\]
\paragraph{Opposite sides with one distinguished variable.}
For $x=p_0,y=q_j$, $j\ge2$, let $r=p_1$, $s=q_1$, $v=q_0$, $A=\sum_{k\ne0}p_k$, $U=\sum_{k\ge2}p_k^2$, $B=\sum_{l\ne j}q_l$, $V=\sum_{l\ge2,l\ne j}q_l^2$. Then
\begin{align*}
4\Delta_{xy}={}&[A(B+v)-rs-2]^2+[As-r(B-v)]^2\\
 &+V(A^2+r^2)+U[(B-v)^2+s^2]+UV.
\end{align*}
\paragraph{Opposite sides with neither distinguished.}
For $x=p_i,y=q_j$, $i,j\ge2$, take $A=\sum_{k\ne i}p_k$, $U=\sum_{k\ge2,k\ne i}p_k^2$, $B=\sum_{l\ne j}q_l$, $V=\sum_{l\ge2,l\ne j}q_l^2$, $R_2=p_0^2+p_1^2$, $S_2=q_0^2+q_1^2$. Then
\begin{align*}
4\Delta_{xy}={}&(AB-K-2)^2+(Aq_0-Bp_0)^2+(Aq_1-Bp_1)^2\\
 &+(p_0q_1-p_1q_0)^2+U(B^2+S_2)+V(A^2+R_2)+UV.
\end{align*}
The all-real Rayleigh criterion now proves stability. Parallel-class sum substitution transfers it to individual variables. These explicit identities, rather than a finite numerical test, establish the arbitrary-class statement.


\section{Coefficientwise Rayleigh inequalities for a three-vertex shore (source 12)}\label{pgr:r:tc:sec}

Source~12 is \emph{Coefficientwise Rayleigh inequalities for a three vertex
shore} (6 pages as delivered; archive \texttt{three-core-rayleigh-result});
source~11 embedded an exact copy of it as its dependency. Its files are
shipped with the prefix \texttt{12-three-core-rayleigh-}.

\paragraph{Abstract of source 12, as delivered.}

For a directed relation whose physical graph is bipartite with a three-vertex shore, we prove coefficientwise Rayleigh inequalities for the polynomial of feasible ordered endpoint supports. Opposite arcs and independent tail/head activities are allowed, and each feasible support is counted once. The boundary inequality reduces to 48 role-monomial types and 17,376 finite directed graphs; two independent exact checkers verify every coefficient. A private-head argument then proves coefficientwise positivity of the full Rayleigh differences in the three core monomers. Consequently, use of any two core vertices is negatively correlated in the endpoint-support measure under all nonnegative role activities and positive core monomer activities. We give a four-core example, already admitting a two-tail/three-head role cover, where coefficientwise positivity fails despite an exact positive two-square formula. The boundary theorem is computer-assisted; the corollaries and obstruction are ordinary arguments. No full-stability, internal-core extension, proof-assistant formalization, or priority claim is made.

\subsection{Supports and the main inequality}
Let $D$ be a finite loopless directed relation on the physical vertex set
$V=C\sqcup I$, where $C=\{0,1,2\}$, and assume every arc crosses this partition.
Both orientations of an edge may be present. Work in the polynomial ring in
independent commuting tail and head activities $u_v,v_v$ for $v\in V$.

A \emph{feasible support} of size $k$ is an ordered pair $(S,T)$ of disjoint
physical vertex sets of size $k$ admitting a directed perfect matching from
$S$ to $T$. Its weight is $\wt(S,T)=u^Sv^T$. It contributes once, regardless
of how many bijections witness feasibility. Since all arcs cross the physical
bipartition, every size-$k$ support uses exactly $k$ vertices in each shore.

Define $a_i$ as the sum of weights of size-one supports with used core set
$\{i\}$, define $b_{ij}=b_{ji}$ similarly for size-two supports with used core
set $\{i,j\}$, and let $c$ be the size-three coefficient. The core activities
are included in these polynomials. The polynomial recording unused core
vertices is
\begin{equation}\label{pgr:r:tc:eq:F}
 F(\mathbf z)=z_0z_1z_2+\sum_{i\in C}a_i z_jz_k
                     +\sum_{i<j}b_{ij}z_k+c,
\end{equation}
where in each summand the omitted indices complete $C$. Exterior monomer
activities are one; exterior role activities remain independent variables.

Write $p\coefge0$ when every coefficient of $p$ in the original activity
variables is a nonnegative integer.
\begin{theorem}[Boundary inequality]\label{pgr:r:tc:thm:boundary}
For every distinct $i,j,k\in C$,
\begin{equation}\label{pgr:r:tc:eq:boundary}
 b_{ij}b_{ik}-a_ic\coefge0.
\end{equation}
This holds for every finite exterior set $I$ and every directed relation
satisfying the physical bipartition assumption.
\end{theorem}
Section~\ref{pgr:r:tc:sec:finite} gives the complete localization argument and exact
finite certificate for this computer-assisted theorem. Section~\ref{pgr:r:tc:sec:full}
derives the full core-variable Rayleigh property without another enumeration.
Rank-three matroid Rayleigh theory \cite{WagnerR3} provides relevant context,
but the proof here works directly with directed endpoint supports and their
independent role activities; no identification with an ordinary unweighted
matroid basis measure is assumed.


\mnote{in the scalar case with one orientation (all arcs from the core
$C$ to the exterior), the boundary inequality is source~22's inequality
$B_{ij}B_{ik}\ge A_iC_0$ (Section~\pgrfwd{pgr:r:ir:sec:rayleigh}{on
source~22's three-core theorem}), which source~22 derives from Wagner's
rank-three Rayleigh theorem; for undirected graphs its scalar form is
\mrn, Equation \texttt{mrn:s:tv:eq:rayleigh}. The coefficientwise form,
with both orientations and independent role activities, is new.}

\subsection{Complete coefficient localization}\label{pgr:r:tc:sec:finite}
By relabeling the core, it suffices to prove
$G=b_{01}b_{02}-a_0c\coefge0$.

\begin{lemma}\label{pgr:r:tc:lem:localization}
Every potentially negative coefficient of $G$ is determined by an induced
relation on the three core vertices and either three or four exterior
physical vertices. Their exterior physical multiplicities are respectively
$(2,1,1)$ or $(1,1,1,1)$.
\end{lemma}
\begin{proof}
A negative coefficient must occur in $a_0c$. The $a_0$ factor uses one exterior
vertex and the $c$ factor uses three distinct exterior vertices. Their union
therefore has exactly one of the two stated patterns. The physical core
multiplicities are $(2,1,1)$ in either product, and total exterior degree is
four. All other monomials have coefficient zero in $a_0c$, so their gap
coefficients are nonnegative automatically.

For a fixed monomial, no support contributing to its coefficient can use a
vertex absent from that monomial. Nor can it use an arc whose tail or head
role variable is absent. Removing such vertices and arcs preserves the
coefficient exactly. The remaining finite graph is determined solely by the
roles and physical labels in the monomial.
\end{proof}

Use role bit $1$ for a tail and $0$ for a head. At doubled core vertex $0$,
let $h\in\{0,1,2\}$ be the exponent of $u_0$; its head exponent is $2-h$.
Let $\epsilon_1,\epsilon_2\in\{0,1\}$ be the roles at vertices $1,2$.
The allowed ordered allocations of the two roles at a doubled vertex are
\begin{equation}\label{pgr:r:tc:eq:allocations}
 \mathcal A(0)=\{(0,0)\},\quad
 \mathcal A(1)=\{(0,1),(1,0)\},\quad
 \mathcal A(2)=\{(1,1)\}.
\end{equation}
Thus identical-role occurrences have one allocation, not two. The total tail
degree in both products is four: the factor sizes are $2+2$ or $1+3$.

\paragraph{Four distinct exterior vertices.}
Each exterior vertex has one fixed role, and exactly
$4-h-\epsilon_1-\epsilon_2$ of them are tails. Up to exterior relabeling,
put these tails first and the heads afterward. The twelve choices of
$(h,\epsilon_1,\epsilon_2)$ give exactly twelve role-monomial cases.

\paragraph{Three exterior vertices.}
Label the repeated exterior vertex $0'$ and the other two $1',2'$. Let
$\ell\in\{0,1,2\}$ be its tail exponent, and let
$\delta_1,\delta_2\in\{0,1\}$ be the other exterior roles. Retain precisely
the tuples satisfying
\begin{equation}\label{pgr:r:tc:eq:balance}
 h+\epsilon_1+\epsilon_2+\ell+\delta_1+\delta_2=4.
\end{equation}
There are thirty-six such tuples, including the core choices.

For each role-monomial case, list every cross-partition directed arc whose
tail and head roles occur in the monomial, and enumerate all subsets of that
list. If both endpoint role types occur, both directions of the same physical
edge are separate arc bits. The exact sizes are
\begin{center}
\begin{tabular}{lrrr}
\toprule
Exterior multiplicities & Role types & Most arc bits & Graphs\\
\midrule
$1,1,1,1$ & 12 & 12 & 11\,392\\
$2,1,1$   & 36 & 10 & 5\,984\\
\midrule
Total     & 48 & & 17\,376\\
\bottomrule
\end{tabular}
\end{center}
These are complete finite case counts, not a sample of arbitrary larger
graphs. Lemma~\ref{pgr:r:tc:lem:localization} explains why they cover every potentially
negative coefficient for every exterior size.

\subsubsection*{Computing one coefficient exactly}
For fixed endpoint sets and fixed roles, the producer checker lists the at
most six bijections between the selected core and exterior vertices. The
feasibility predicate is true if \emph{at least one} bijection uses present,
role-compatible directed arcs. The number of successful bijections is never
used as a coefficient.

With four exterior vertices, the coefficient in $b_{01}b_{02}$ is obtained
by choosing the two exterior vertices assigned to the first factor (six
choices), assigning the complement to the second, and summing over the
distinct role allocations \eqref{pgr:r:tc:eq:allocations} at core $0$. A choice
contributes one exactly when both endpoint supports are feasible. For $a_0c$,
choose the singleton exterior vertex (four choices), give the remaining
three to $c$, and use the same core-role allocation rule.

With three exterior vertices, both positive factors must contain $0'$.
Choose which of $1',2'$ accompanies it in $b_{01}$ (two choices) and put the
other in $b_{02}$. Sum over the distinct role allocations at core $0$ and
exterior $0'$. In the negative product, the exterior singleton of $a_0$ must
be $0'$, while $c$ uses all three exterior vertices. Again sum over precisely
those two allocation sets. Each choice specifies an ordered pair of endpoint
supports uniquely, so subtracting these two counts is the desired coefficient.

All $17\,376$ cases give nonnegative coefficients; the recorded values in
these potentially negative cases are $0,1,2,3$. The complete per-case
histograms and executable checker are in the package.

\subsubsection*{Independent reconstruction}
A second checker does not import the producer code or use its allocation
routine. It constructs actual directed physical matchings by adjoining
vertex-disjoint arcs, deduplicates their ordered endpoint sets, and separates
them according to their used core subsets. It encodes the formal role
monomials in base three: each exponent in a product of two supports is at most
two, so this representation has no carries. It then extracts the positive
and negative product coefficients directly.

This independent enumeration checks $228\,264$ support instances across the
same $48$ role types and $17\,376$ graphs. Every coefficient histogram agrees
with the producer's output. Together with the complete localization argument,
these finite checks prove Theorem~\ref{pgr:r:tc:thm:boundary}. They are the finite part
of the proof, not merely floating-point evidence. No random search is used.

\subsection{Full Rayleigh differences and negative correlation}\label{pgr:r:tc:sec:full}
For a multiaffine polynomial $F$, define
\[
 \Delta_{jk}(F)=(\partial_{z_j}F)(\partial_{z_k}F)
                   -F\partial_{z_j}\partial_{z_k}F.
\]
\begin{theorem}\label{pgr:r:tc:thm:full}
For $F$ in \eqref{pgr:r:tc:eq:F}, every $\Delta_{jk}(F)$ is coefficientwise
nonnegative in the core monomers and original independent role activities.
Consequently, after every nonnegative role-activity specialization, $F$ is
Rayleigh as a polynomial of its three core monomers at nonnegative arguments.
\end{theorem}
\begin{proof}
For distinct $i,j,k$, direct cancellation gives
\begin{align}\label{pgr:r:tc:eq:delta}
 \Delta_{jk}(F)={}&(a_ja_k-b_{jk})z_i^2\notag\\
 &+(a_jb_{ik}+a_kb_{ij}-c-a_ib_{jk})z_i
 +(b_{ij}b_{ik}-a_ic).
\end{align}
The constant coefficient is nonnegative by Theorem~\ref{pgr:r:tc:thm:boundary}.
Also $a_ja_k\coefge b_{jk}$: every size-two support using $j,k$ has at
least one pair of single-edge witnesses counted by the product, with the
same role weight; repeated exterior vertices and additional witnesses only
add terms to that product.

To treat the middle coefficient, add a fresh exterior physical vertex $x$
and only the arc $i\to x$, and put $L=u_iv_x$. Deleting that forced edge is
a bijection of endpoint supports; any residual matching-witness multiplicity
remains ignored. The enlarged relation has coefficients
\[
 a_i'=a_i+L,\qquad b_{ij}'=b_{ij}+La_j,\qquad
 b_{ik}'=b_{ik}+La_k,\qquad c'=c+Lb_{jk}.
\]
It still satisfies the bipartition hypotheses of Theorem~\ref{pgr:r:tc:thm:boundary}.
Its boundary gap is
\begin{align}\label{pgr:r:tc:eq:private}
 b_{ij}'b_{ik}'-a_i'c'={}&b_{ij}b_{ik}-a_ic\notag\\
 &+L(a_jb_{ik}+a_kb_{ij}-c-a_ib_{jk})
 +L^2(a_ja_k-b_{jk}).
\end{align}
The fresh variable $v_x$ is absent from the original coefficients. Its
coefficient in \eqref{pgr:r:tc:eq:private} is $u_i$ times the middle coefficient of
\eqref{pgr:r:tc:eq:delta}. The enlarged boundary gap is coefficientwise nonnegative,
and multiplication by the monomial $u_i$ shifts exponents injectively.
The middle coefficient is therefore nonnegative as well.
\end{proof}

The statement concerns nonnegative monomer arguments. It is not an all-real
Rayleigh criterion, and no real-stability conclusion follows here. It also
does not assert Rayleigh inequalities in the exterior or activity variables.

\begin{corollary}\label{pgr:r:tc:cor:covariance}
Fix nonnegative role activities and strictly positive core monomer activities.
Give each feasible ordered support probability proportional to
\[
 u^Sv^T\prod_{i\in C\setminus(S\cup T)}z_i.
\]
Then use of any two distinct core vertices is negatively correlated. The
same is true of their unused indicators.
\end{corollary}
\begin{proof}
The normalizing constant is $F(\mathbf z)>0$, since the empty support
contributes $z_0z_1z_2>0$. If $U_j$ is the indicator that $j$ is unused,
differentiation gives
\[
 \operatorname{Cov}(U_j,U_k)
 =\frac{z_jz_k(F\partial_{z_j}\partial_{z_k}F
                    -\partial_{z_j}F\partial_{z_k}F)}{F^2}
 =-\frac{z_jz_k\Delta_{jk}(F)}{F^2}\le0.
\]
Replacing each indicator by one minus itself leaves their pairwise covariance
unchanged. This proves the assertion for used indicators too.
\end{proof}
This probability measure is on feasible endpoint supports, not on individual
matching witnesses. Zero role activities cause no problem: the coefficientwise
inequalities specialize directly, and the positive empty-support monomial
still normalizes the measure.

\subsection{A scalar cubic consequence}
Write
\[
\Ga_D(t)=1+At+Bt^2+ct^3,\qquad
 A=a_0+a_1+a_2,\quad B=b_{01}+b_{02}+b_{12}.
\]
Rank-ultra-log-concavity means log-concavity of
$\gamma_k/\binom dk$ at the actual surviving degree $d$. Positive-weight
submatchings ensure that the coefficient sequence has no internal zeros.
At nonnegative activities,
\begin{equation}\label{pgr:r:tc:eq:three-squares}
 B^2-3Ac=
 \frac12\bigl[(b_{01}-b_{02})^2+(b_{01}-b_{12})^2
                         +(b_{02}-b_{12})^2\bigr]
 +3\sum_i(b_{ij}b_{ik}-a_ic)\ge0.
\end{equation}
Moreover, $b_{ij}\le a_ia_j$ gives
$A^2\ge3\sum_{i<j}a_ia_j\ge3B$.

\mnote{identity~\eqref{pgr:r:tc:eq:three-squares} is printed once, here. Source~11 uses
it in its equality proposition, and source~22's identity for its
three-core theorem is this identity with $b_{ij}$, $a_i$ and $c$ replaced by
$u_iu_jb_{ij}$, $u_ia_i$ and $u_0u_1u_2C_0$.}
Thus the cubic Newton inequalities hold. Equation~\eqref{pgr:r:tc:eq:three-squares}
is a scalar nonnegativity identity; its squared differences must not be
mistaken for coefficientwise nonnegative polynomials.


\mnote{source~12 continues with a paragraph proving the stronger
normalization $A^2\ge4B$ after a drop to actual degree two, by
Motzkin--Straus compression on the triangle-free graph of positive-weight
arcs; this is the case $d=2$ of Theorem~\ref{pgr:thm:first} and is
replaced by that pointer. Degrees zero and one have no inequalities to
check.}


Consequently $\Ga_D$ is rank-ultra-log-concave at its actual surviving degree,
including zero activities. This scalar conclusion is weaker than the main
coefficientwise and negative-correlation results.

\subsection{An exact obstruction at four cores}\label{pgr:r:tc:sec:four}
The coefficientwise assertion cannot be extended without qualification to a
four-vertex shore, even inside the two-tail/three-head role-cover class.

Take tails $0,1,2,3$ and five distinct heads $a,b,c,d,e$, and only the arcs
\begin{equation}\label{pgr:r:tc:eq:fourgraph}
 N(0)=\{c,d,e\},\quad N(1)=\{a,b,e\},\quad
 N(2)=\{b,d\},\quad N(3)=\{e\}.
\end{equation}
The tail roles $\{0,1\}$ and head roles $\{b,d,e\}$ cover every arc. A
matching $0\to c,1\to a,2\to d,3\to e$ shows degree four.
Set tail activities to one and use the head labels as their activities.
Let $B_J$ count head supports matchable to the selected tail set $J$, once
per feasible head set. Boolean matching tests give
\begin{align*}
 B_{012}&=abc+abd+abe+acd+ade+bcd+bce+bde+cde,\\
 B_{013}&=ace+ade+bce+bde,\\
 B_{01}&=ac+ad+ae+bc+bd+be+ce+de,\\
 B_{0123}&=abce+abde+acde+bcde.
\end{align*}
Their boundary gap is
\begin{align}\label{pgr:r:tc:eq:fourgap}
 B_{012}B_{013}-B_{01}B_{0123}
 &=e^2(a^2d^2-abcd+abd^2+b^2c^2+b^2cd+b^2d^2)\notag\\
 &=e^2\left[\left(ad+\frac{b(d-c)}2\right)^2
                      +\frac34b^2(c+d)^2\right].
\end{align}
The coefficient of $abcde^2$ is $-1$. Restoring tail activities multiplies
the entire gap by $u_0^2u_1^2u_2u_3$, so this remains a negative coefficient
in the original role-activity ring.

Nevertheless, the second line of \eqref{pgr:r:tc:eq:fourgap} is nonnegative at every
real head-activity assignment. This example refutes the coefficientwise
extension only. It is not a counterexample to scalar Rayleigh nonnegativity,
real-rootedness, or rank-ULC. No minimality assertion is made.

\subsection{Proof status and reproducibility}
The boundary theorem consists of a complete all-size localization proof and
an exhaustive finite integer check. The full core-variable Rayleigh theorem,
negative-correlation corollary, scalar consequence, and four-core obstruction
are ordinary arguments once that boundary theorem is established. No
proof-assistant verification or literature-wide novelty claim is made.

The package includes the producer's finite checker and all $48$ histograms,
the independently written direct-support checker and its receipts, exact
four-core verifiers, mathematical audits, and editable source. The replay
driver first checks file hashes, then runs copied checkers in temporary
directories and compares deterministic receipt fields without altering the
frozen package. All computations use Python's standard library and exact
integer or rational arithmetic.

The physical core must remain independent. No result here permits internal
core arcs. The general rank-ULC problem for a cover of at most two tail roles
and three head roles, with arbitrary independent activities and physical
overlap, remains open. Section~\ref{pgr:r:tc:sec:four} gives a specific reason that
the strongest coefficientwise approach cannot simply be extrapolated to a
larger shore.

\mnote{in this report the two-tail/three-head problem is settled in
part: by real stability when $|Q|\le2$ (source~20,
Theorem~\ref{pgr:r:rc:thm:rolecover}), for $P\subseteq Q$ by source~11's
Corollary~\pgrfwd{pgr:r:pc:cor:rolecover}{``rolecover''}, and for an
independent physical core at its cover order by source~15's
Corollary~\pgrfwd{pgr:r:ps:cor:overlap}{``overlap''}. The general case,
with internal core arcs and a physical union of four or five vertices,
remains open.}



\section{A physical vertex cover of size three (source 11)}\label{pgr:r:pc:sec}

Source~11 is \emph{Rank ultra log concavity under a three vertex physical
cover} (5 pages as delivered; archive \texttt{physical-cover-three-result}).
Its files are shipped with the prefix \texttt{11-physical-cover-}; the copy
of source~12 that it embeds is not shipped, because source~12's own files
are.

\paragraph{Abstract of source 11, as delivered.}

We prove actual-degree ultra-log-concavity of the endpoint-support polynomial
for every finite loopless directed relation with a physical vertex cover of
size at most three. Internal core arcs are arbitrary, transitivity is
unnecessary, and independent nonnegative tail/head activities are allowed.
The proof combines a verified bipartite boundary inequality, support-level
witness bounds, and an exact two-square decomposition. Each support is
counted once, with zero activities normalized at the surviving degree.
A separate bipartite-only section characterizes strict core negative
correlation and equality in the upper cubic Newton inequality. The ordinary
arguments depend on a replayable finite boundary certificate included in the
package. No real-rootedness, proof-assistant formalization, or priority claim
is asserted.

\subsection{Statement and normalization}
Let $D$ be a finite loopless directed relation on physical vertex set $V$.
Opposite arcs are allowed. Give every vertex tail activity $u_i\ge0$ and
head activity $v_i\ge0$.
A feasible support of size $k$ is an ordered pair $(S,T)$ of disjoint subsets
of size $k$ admitting a directed perfect matching from $S$ to $T$. It has
weight $\wt(S,T)=u^Sv^T$ and is counted once, regardless of how many matching
witnesses realize it. Set
\[
 \gamma_k=\sum_{\substack{(S,T)\text{ feasible}\\|S|=k}}u^Sv^T,
 \qquad \Ga_D(t)=\sum_{k\ge0}\gamma_k t^k,
 \qquad\gamma_0=1.
\]
The actual degree $d$ is the matching number of the underlying graph of the
positive-weight relation $D_+$, which retains precisely the arcs
$i\to j$ with $u_iv_j>0$. Taking submatchings gives no internal zeros.
Rank-ultra-log-concavity, or rank-ULC, means
\begin{equation}\label{pgr:r:pc:eq:rankulc}
 k(d-k)\gamma_k^2\ge(k+1)(d-k+1)\gamma_{k-1}\gamma_{k+1}
 \qquad(1\le k<d),
\end{equation}
at this actual surviving degree.

\begin{theorem}\label{pgr:r:pc:thm:main}
If the underlying undirected graph of $D$ has a vertex cover of size at most
three, then $\Ga_D$ is rank-ULC for every nonnegative choice of independent
tail and head activities. The internal core arcs are unrestricted and
transitivity is not required.
\end{theorem}


\mnote{Theorem~\ref{pgr:r:pc:thm:main} contains source~22's three-core
theorem (Theorem~\pgrfwd{pgr:r:ir:thm:three}{``three'' of source~22}: three
core tails and no arcs into the core, so the core is a physical vertex
cover) and the rank-ULC half of source~22's role-matching corollary (a role
cover of size three has at most three physical vertices); source~22's
ordinary proof of the former is kept there as a second route. For bipartite
graphs with every vertex in one role it is the case $\nu\le3$ of \prp,
Corollary \texttt{lc:cor:rankthree}, itself contained in \mrn, Corollary
\texttt{mrn:w:rf:cor:rankfive}.}

\begin{corollary}\label{pgr:r:pc:cor:rolecover}
Suppose every arc $i\to j$ satisfies $i\in P$ or $j\in Q$, with
$|P|\le2$, $|Q|\le3$, and $P\subseteq Q$. Then $\Ga_D$ is rank-ULC at its
actual surviving degree, with arbitrary independent activities.
\end{corollary}
\begin{proof}
Every underlying edge meets the physical set $Q$, so
Theorem~\ref{pgr:r:pc:thm:main} applies.
\end{proof}
This closes the two-tail/three-head role-cover case in which both selected
tail vertices also belong to the selected head set. Covers with a physical
union of four or five vertices are not settled by this theorem.


\mnote{source~11's Lemma ``first'' (Universal first inequality: for any
directed relation of actual degree $d\ge2$,
$\gamma_1^2\ge\frac{2d}{d-1}\gamma_2$), proved by moving weight between
nonadjacent positive-weight arcs of the compatibility graph, is
Theorem~\ref{pgr:thm:first} with Lemma~\ref{pgr:lem:ms}; it is replaced by
that pointer.}

The lemma supplies the first cubic inequality and the stronger degree-two
inequality after a degree drop. Only the last cubic inequality remains.

\subsection{The exterior boundary input}
Assume now that $C=\{0,1,2\}$ is a physical vertex cover and put
$I=V\setminus C$. The exterior $I$ is independent. Delete all internal core
arcs, obtaining the bipartite physical relation $D_0$ on $C\sqcup I$.
Define, with all original role activities included,
\begin{itemize}[nosep]
\item $a_i$: the weight sum of size-one supports of $D_0$ using core $\{i\}$;
\item $b_{ij}$: the weight sum of size-two supports using core $\{i,j\}$;
\item $c$: the weight sum of size-three supports, using all three cores.
\end{itemize}
Set $A=\sum_i a_i$ and $B=\sum_{i<j}b_{ij}$.

The following is the explicit dependency of the present proof:
\begin{lemma}[Verified three-core boundary theorem]\label{pgr:r:pc:lem:boundary}
For every distinct $i,j,k\in C$,
\begin{equation}\label{pgr:r:pc:eq:boundary}
 b_{ij}b_{ik}-a_ic\succeq_{\rm coeff}0
\end{equation}
in the original independent role-activity variables.
\end{lemma}
This theorem is proved in the accompanying manuscript. \mnote{that manuscript is source~12, printed in
Section~\ref{pgr:r:tc:sec}; the statement above is its
Theorem~\ref{pgr:r:tc:thm:boundary}. Source~12's files are shipped once,
with its own prefix.}
Its complete all-size localization reduces every potentially negative
coefficient to three or four exterior physical vertices. There are $48$
role-monomial types and $17\,376$ relevant directed graphs. Two independently
written exact checkers reproduce every coefficient histogram. That proof,
its audits and the checkers are included and replayed in this package; the
dependency is not replaced by a random sample. We use only its valid
nonnegative specialization below.

Summing \eqref{pgr:r:pc:eq:boundary} gives
\begin{equation}\label{pgr:r:pc:eq:Ac}
 Ac\le b_{01}b_{02}+b_{01}b_{12}+b_{02}b_{12}.
\end{equation}
We also need the elementary coefficientwise inequality
\begin{equation}\label{pgr:r:pc:eq:remove}
 c\le a_kb_{ij}.
\end{equation}
For a size-three support, choose a witnessing matching. Its edge at core $k$
is a singleton support, and deleting it leaves a support counted by $b_{ij}$.
The product has the original weight. Additional decompositions and intersecting
products only add nonnegative contributions.

\subsection{Restoring every internal core arc}
For an unordered core pair $\{i,j\}$, with remaining core $k$, write
\[
 h_{ij}=\mathbf1_{i\to j\in D}u_iv_j+
         \mathbf1_{j\to i\in D}u_jv_i,\qquad
 E=h_{01}+h_{02}+h_{12},\qquad w_{ij}=h_{ij}a_k.
\]
Let $H$ be the weight sum of size-two supports using all three physical core
vertices. These supports use one internal core arc and one exterior arc.
The other size-two supports use precisely two core vertices and belong to
$D_0$. A size-three matching cannot contain an internal core edge, since
that would leave only one core vertex for its other two disjoint edges.
Thus, exactly at support level,
\begin{equation}\label{pgr:r:pc:eq:decomposition}
 \gamma_1=A+E,\qquad \gamma_2=B+H,\qquad\gamma_3=c.
\end{equation}

\begin{lemma}[Internal witness bounds]\label{pgr:r:pc:lem:witnesses}
At all nonnegative activities,
\begin{equation}\label{pgr:r:pc:eq:witnessbounds}
 0\le w_{ij}\le H,\qquad w_{01}+w_{02}+w_{12}\le2H.
\end{equation}
Both inequalities hold coefficientwise as support-counting statements.
\end{lemma}
\begin{proof}
A product in $h_{ij}a_k$ chooses an internal arc on $\{i,j\}$ and an exterior
arc at $k$. Their physical endpoints are disjoint. For this fixed unordered
core pair, every resulting ordered endpoint support is counted exactly once:
opposite orientations of the internal arc yield different role assignments,
and the exterior edge is then forced. Hence $w_{ij}\le H$.

The sum of the three $w$'s counts the internal-edge witnesses of each
$H$-support. Such a support has three core vertices and one exterior vertex.
If two core vertices are tails and one is a head, the internal edge must run
from one of those two tails to that head. There are at most two choices.
The case of one core tail and two core heads is the reverse case. Each
witness has the same prescribed support weight, so the total is at most
$2H$. In particular, fork-induced duplicate witnesses are retained in this
bound rather than silently counted as distinct supports in $H$.
\end{proof}
Multiplying \eqref{pgr:r:pc:eq:remove} by $h_{ij}$ and summing yields
\begin{equation}\label{pgr:r:pc:eq:Ec}
 Ec\le\sum_{i<j} w_{ij}b_{ij}.
\end{equation}

At the fixed activity specialization, order the three $b$-values as
$x\ge y\ge z\ge0$, and let $w_x,w_y,w_z$ be their corresponding $w$-values.
The witness bounds imply
\begin{equation}\label{pgr:r:pc:eq:weightedbound}
 xw_x+yw_y+zw_z\le H(x+y),
\end{equation}
because its slack is exactly
\begin{equation}\label{pgr:r:pc:eq:slack}
 (x-y)(H-w_x)+y(2H-w_x-w_y-w_z)+(y-z)w_z\ge0.
\end{equation}
Combining \eqref{pgr:r:pc:eq:Ac}, \eqref{pgr:r:pc:eq:Ec} and \eqref{pgr:r:pc:eq:weightedbound} gives
\[
 \gamma_1\gamma_3\le xy+xz+yz+H(x+y),\qquad
 \gamma_2=x+y+z+H.
\]
Consequently
\begin{align}\label{pgr:r:pc:eq:twosquares}
 \gamma_2^2-3\gamma_1\gamma_3
 &\ge(x+y+z+H)^2-3\bigl[xy+xz+yz+H(x+y)\bigr]\notag\\
 &=\left(z+H-\frac{x+y}{2}\right)^2+\frac34(x-y)^2\ge0.
\end{align}
Together with Lemma~\ref{pgr:thm:first}, this proves Theorem~\ref{pgr:r:pc:thm:main} at
actual degree three. Degrees at most two were already handled by that lemma,
so zero-weight degree drops cause no normalization gap.

For a completely explicit accounting of the slack, the proof gives the exact
identity
\begin{align}\label{pgr:r:pc:eq:exact}
 \gamma_2^2-3\gamma_1\gamma_3={}&
 \left(z+H-\frac{x+y}{2}\right)^2+\frac34(x-y)^2\notag\\
 &+3\sum_i(b_{ij}b_{ik}-a_ic)
   +3\sum_{i<j}h_{ij}(a_kb_{ij}-c)\notag\\
 &+3\bigl[(x-y)(H-w_x)+y(2H-\textstyle\sum w)
                                  +(y-z)w_z\bigr].
\end{align}
Every term is nonnegative. The ordering of $x,y,z$ is numerical, after
activity specialization; \eqref{pgr:r:pc:eq:exact} is not a coefficientwise positivity
claim for the full Newton gap.

\subsection{Bipartite only strictness results}\label{pgr:r:pc:sec:strictness}
The remainder of this section assumes that \emph{there are no internal core
arcs}. It must not be applied to the unrestricted internal-core theorem.
Let $D_+$ denote the positive-weight relation, and give the three core
monomers strictly positive activities $z_0,z_1,z_2$. The core polynomial is
\[
 F=z_0z_1z_2+\sum_i a_i z_jz_k+\sum_{i<j}b_{ij}z_k+c.
\]
The boundary theorem and a private-head insertion prove (source~12, Theorem~\ref{pgr:r:tc:thm:full})
that all three coefficients of
\begin{align}\label{pgr:r:pc:eq:rayleigh}
 \Delta_{jk}(F)={}&(a_ja_k-b_{jk})z_i^2\notag\\
 &+(a_jb_{ik}+a_kb_{ij}-c-a_ib_{jk})z_i+(b_{ij}b_{ik}-a_ic)
\end{align}
are coefficientwise nonnegative. Here
$\Delta_{jk}=(\partial_jF)(\partial_kF)-F\partial_j\partial_kF$.

\begin{proposition}\label{pgr:r:pc:prop:strict}
Give each endpoint support probability proportional to
\[
 u^Sv^T\prod_{i\in C\setminus(S\cup T)}z_i.
\]
Then the covariance of the used
indicators of two distinct core vertices is strictly negative exactly when
they are in the same connected component of the underlying graph of $D_+$.
Otherwise it is zero. The unused indicators have the same covariance.
\end{proposition}
\begin{proof}
Supports and weights factor across different physical components, so those
components are independent. If $j,k$ are connected, a simple alternating
path between them has length two or four, since there are only three core
vertices.

For a path $j-x-k$, choose a positive available arc on each physical edge.
Their product is a positive monomial of $a_ja_k$ with exterior vertex $x$
used twice. It cannot occur in $b_{jk}$, whose supports require two distinct
exterior vertices. Thus the quadratic coefficient of \eqref{pgr:r:pc:eq:rayleigh}
is strictly positive at these activities.

For a path $j-x-i-y-k$, the two supports formed by edges
$\{j-x,i-y\}$ and $\{i-x,k-y\}$ contribute a positive monomial to
$b_{ij}b_{ik}$ on only the two exterior vertices $x,y$. It cannot occur in
$a_ic$, because $c$ requires three distinct exterior vertices. The constant
coefficient is therefore strictly positive. Each support separately is
physically disjoint; role reuse across the product is allowed.

In either case, coefficientwise nonnegativity prevents cancellation and
$\Delta_{jk}(F)>0$. The empty support makes $F>0$, and the unused indicators
$U_j,U_k$ satisfy
\[
 \operatorname{Cov}(U_j,U_k)=-\frac{z_jz_k\Delta_{jk}(F)}{F^2}<0.
\]
Replacing both indicators by one minus themselves preserves covariance.
\end{proof}

\begin{proposition}\label{pgr:r:pc:prop:equality}
Assume actual degree three, so $c>0$, and still no internal core arcs. Then
$B^2=3Ac$ if and only if, after deleting zero-weight arcs and isolates, the
physical graph is a union of three separate stars centered at the core
vertices, with equal positive total directed arc activities. In that case
$\Ga_D(t)=(1+at)^3$ for their common activity total $a$.
\end{proposition}
\begin{proof}
The exact identity~\eqref{pgr:r:tc:eq:three-squares} of source~12, which source~11 prints again,
uses the three unordered core pairs as $e,f$. Equality forces the three
$b$-values to agree and all boundary gaps to vanish.

Suppose an exterior vertex $x$ has at least two core neighbors, say $j,k$.
If the remaining core $i$ has a neighbor $y\ne x$, the supports
$\{i-y,j-x\}$ and $\{i-y,k-x\}$ give a positive monomial of the boundary
gap for $i$ on only two exterior vertices. It cannot occur in $a_ic$.
If $i$ has no neighbor other than $x$, then $x$ is adjacent to all three
cores. A saturating matching exists because $c>0$; it uses $i-x$ and must
match $j$ to some $y\ne x$. Use $j$ as the doubled core in the same
two-support construction. In either case, one boundary gap is strictly
positive, contradicting equality.

Thus no exterior vertex is shared by different cores. The graph is the
stated union of stars and supports factor:
$b_{ij}=a_ia_j$, $c=a_0a_1a_2$, and
$\Ga_D=\prod_i(1+a_it)$. Each $a_i>0$, so equality of the three $b$-values
forces $a_0=a_1=a_2$. The converse is immediate. Opposite arcs simply
contribute their two distinct singleton-support weights to a star's total.
\end{proof}

\subsection{Proof status and independent replay}
The only computational dependency is the three-core bipartite boundary
theorem. The internal-core assembly and bipartite-only strictness results
are ordinary arguments from that input. No real-rootedness or signed-monomer
stability is asserted for arbitrary internal core arcs.

The independent audit checks all $4\,096$ one-exterior graphs for witness
bounds, $1\,600$ three-exterior cases for \eqref{pgr:r:pc:eq:remove}, and $10\,080$
exact activity evaluations for the assembly and actual-degree inequalities.
The producer checks $10\,496$ weighted cases, including all $64$ core
patterns, $6\,478$ zero-activity cases and $504$ larger activity-induced
degree drops, and verifies \eqref{pgr:r:pc:eq:exact}. These checks corroborate the
ordinary all-size proof; they impose no exterior-population bound.

The archive includes the full boundary manuscript, source, audits and both
finite checkers, plus the new independent and producer assembly checks.
The replay driver verifies hashes and deterministic receipts using temporary
copies. All code uses Python's standard library and exact arithmetic.

The general two-tail/three-head problem with four or five physical core
vertices remains open. No proof-assistant formalization or priority claim
is made.


\section{Bipartite physical graphs and pendant heads (source 15)}\label{pgr:r:ps:sec}

Source~15 is \emph{Lorentzian and real stable directed support polynomials}
(6 pages as delivered; archive \texttt{bipartite-and-pendant-support-result}).
Its files are shipped with the prefix \texttt{15-pendant-support-}.

\paragraph{Abstract of source 15, as delivered.}

We count ordered disjoint endpoint supports of directed matchings once per feasible support, with independent nonnegative tail and head activities. For any directed relation whose physical graph is bipartite, we prove a multiaffine Lorentzian theorem at the physical shore order. It gives actual-degree ultra-log-concavity when a positive-weight matching saturates a shore, and accounts for overlap among selected core roles. We also prove that a stable selected-head transversal seed remains stable after adding pendant heads at arbitrarily many tail hubs and merging physical copies. This yields negative real roots and actual-degree ultra-log-concavity even on zero-weight degree drops. Singleton-or-universal neighborhoods give an explicit stable seed in arbitrary rank. The proofs are ordinary arguments; exact independent checks are supplementary. The general two-tail/three-head question remains open. No proof-assistant formalization or priority claim is made.

\subsection{Supports and the normalization issue}
Let $D$ be a finite loopless directed relation on $V$. Opposite arcs are allowed; transitivity is not assumed. Assign tail activities $u_i\ge0$ and head activities $v_i\ge0$.
\begin{definition}
A \emph{feasible support} of size $k$ is an ordered pair $(S,T)$ of disjoint subsets of $V$, with $|S|=|T|=k$, admitting a directed perfect matching from $S$ to $T$. Its weight is
\[
 \wt(S,T)=u^S v^T=\prod_{i\in S}u_i\prod_{j\in T}v_j.
\]
Every feasible pair contributes once, regardless of how many matching witnesses it has. Define
\[
 \gamma_k=\sum_{\substack{(S,T)\text{ feasible}\\|S|=k}}\wt(S,T),
 \qquad \Ga_D(t)=\sum_{k\ge0}\gamma_k t^k.
\]
The empty support has weight one, so $\gamma_0=1$.
\end{definition}
The two activities of the same physical vertex can differ. Replacing physical vertices by independent role copies without later excluding their simultaneous use changes the polynomial.

Let $D_+$ retain exactly the arcs $i\to j$ for which $u_iv_j>0$. The actual degree $d$ of $\Ga_D$ is the matching number of the underlying undirected graph of $D_+$. Indeed, every positive-weight support supplies such a matching; conversely, orient each edge of such a matching in an available direction. Taking submatchings also shows that $\gamma_k>0$ for $0\le k\le d$.

\begin{definition}
For an integer $r\ge d$, the coefficients satisfy \emph{order-$r$ ultra-log-concavity} if
\begin{equation}\label{pgr:r:ps:eq:ulc}
 k(r-k)\gamma_k^2\ \ge\ (k+1)(r-k+1)\gamma_{k-1}\gamma_{k+1}
 \qquad(1\le k<r),
\end{equation}
with coefficients beyond degree $d$ set to zero. \emph{Rank-ULC} means \eqref{pgr:r:ps:eq:ulc} with $r=d$, the actual surviving degree.
\end{definition}
These distinctions matter. The binary polynomial
\[
 s^3+3s^2t+3st^2=s(s^2+3st+3t^2)
\]
is Lorentzian, but its degree-two quotient is not: $3^2<4\cdot1\cdot3$. Thus division by a common monomial cannot be used to upgrade an order bound to actual-degree ULC.

Write $\HH=\{z\in\mathbb C:\operatorname{Im}z>0\}$. A real polynomial is \emph{real stable} if it is nonzero whenever every variable lies in $\HH$.

We use the classical Lorentzian theory of Br\"and\'en and Huh \cite{BH}. In particular, matroid basis polynomials are Lorentzian; products, nonnegative linear substitutions and multiaffine-part extraction preserve this class. For a nonzero binary homogeneous polynomial $\sum_{k=0}^r a_k s^{r-k}t^k$, Lorentzianity is equivalent to the log-concavity of $a_k/\binom rk$, with interval support \cite[Example 2.26]{BH}. This paragraph invokes that theory, rather than reproving it.

\subsection{A polynomial for bipartite physical graphs}
\begin{theorem}\label{pgr:r:ps:thm:bipartite}
Suppose $V=C\sqcup I$ is a bipartition of the underlying undirected graph of $D$. Every arc crosses the partition, but may have either orientation and may coexist with its reverse. Then
\begin{equation}\label{pgr:r:ps:eq:physicalF}
 F_C(\mathbf z,\mathbf x)=
 \sum_{(S,T)\text{ feasible}}\wt(S,T)
 \prod_{c\in C\setminus(S\cup T)}z_c
 \prod_{i\in I\cap(S\cup T)}x_i
\end{equation}
is a nonzero multiaffine Lorentzian polynomial, homogeneous of degree $|C|$.
\end{theorem}

\begin{lemma}[Core merge]\label{pgr:r:ps:lem:merge}
On polynomials affine in $a,b$, the coefficientwise operation
\begin{equation}\label{pgr:r:ps:eq:merge}
 ab\longmapsto z,\qquad a\longmapsto1,\qquad b\longmapsto1,
 \qquad1\longmapsto0
\end{equation}
preserves Lorentzianity, allowing zero output. It lowers homogeneous degree by one. It also preserves real stability, allowing zero output.
\end{lemma}
\begin{proof}
The Lorentzian assertion is precisely \cite[Lemma 3.3]{BH}, after renaming its surviving variable. Alternatively, the operator is homogeneous of degree $-1$, preserves coefficient nonnegativity, and has algebraic symbol
\[
 T[(a+r)(b+s)]=z+r+s.
\]
This symbol is homogeneous and real stable, hence Lorentzian; the full symbol includes only additional factors $(w_j+y_j)^{\kappa_j}$ for untouched variables. The finite-degree symbol theorem \cite[Theorem 3.2]{BH} applies.

For an elementary stability proof, write $f=Aab+Ba+Cb+D$. Then the image $Az+B+C$ equals
\[
 \left.\frac{\partial}{\partial t}f(t,t)\right|_{t=z/2}.
\]
Diagonalization, differentiation and positive rescaling preserve real stability or give zero; see \cite{BB,WagnerSurvey}. Every surviving monomial in \eqref{pgr:r:ps:eq:merge} loses exactly one degree.
\end{proof}


\mnote{this is the copy merge of source~20, \eqref{pgr:r:rc:eq:merge},
with source~16's elementary realization (Section~\ref{pgr:r:rc:sec}).
The same operation glues two undirected graphs at a vertex in \prp,
Lemma \texttt{ade:mat:lem:vertex-sum} (support-exact vertex sum).}

\begin{proof}[Proof of Theorem~\ref{pgr:r:ps:thm:bipartite}]
For the outgoing side, form a transversal matroid $M_+$ with matching positions indexed by $C$, and ground set consisting of an exterior head copy for each $i\in I$ and a private dummy $a_c$ for each $c\in C$. The exterior copy of $i$ is adjacent to $c$ exactly when $c\to i$ is an arc; $a_c$ is adjacent only to $c$. The dummies form a basis, so the rank is exactly $|C|$.

A set consisting of exterior copies $R\subseteq I$ and dummies $C\setminus U$, for $U\subseteq C$, is a basis precisely when $|U|=|R|$ and $U$ can be matched to $R$. Its basis monomial records the endpoint sets once, not their matching witnesses. Similarly form $M_-$ from arcs $i\to c$, using exterior tail copies and private dummies $b_c$. Its bases encode $(L,W)$ with $L\subseteq I$, $W\subseteq C$, and a matching $L\to W$.

The squarefree basis polynomials $B_+,B_-$ are Lorentzian \cite[Theorem 3.10]{BH}. First assume positive activities. In $B_+$ substitute $a_c/u_c$ for $a_c$ and $v_ix_i$ for the exterior variable, and multiply by $\prod_{c\in C}u_c$. A basis $(U,R)$ now has coefficient $u^Uv^R$. In $B_-$ substitute $b_c/v_c$ and $u_iy_i$, and multiply by $\prod_{c\in C}v_c$. Its coefficient for $(L,W)$ is $u^Lv^W$.

Multiply the two polynomials, identify $x_i=y_i$, and take the multiaffine part. Squared exterior variables correspond exactly to vertices selected in both exterior roles, so this discards precisely $R\cap L\ne\varnothing$. All other variables were already multiaffine. The result remains Lorentzian by \cite[Theorem 2.10, Corollaries 2.32 and 3.5]{BH}, and has degree $2|C|$.

For each $c\in C$, merge $a_c,b_c$ by Lemma~\ref{pgr:r:ps:lem:merge}. If both dummies occur, $c$ is unused in both roles and becomes one physical monomer $z_c$. If exactly one dummy occurs, $c$ is used in exactly one role and the dummy disappears. Neither dummy occurs exactly when $c$ is used in both roles, and that term is discarded. The degree falls by one at each of the $|C|$ merges.

A surviving pair of bases determines
\[
 S=U\cup L,\qquad T=R\cup W,
 \quad U\cap W=\varnothing,\quad R\cap L=\varnothing.
\]
Conversely, intersecting a feasible ordered support with $C$ and $I$ determines these four sets uniquely. Every witnessing matching splits into the two orientations across the bipartition, which establish the two basis conditions. This is a bijection of supports and surviving pairs of bases. Distinct role assignments may have the same physical-variable monomial, but their weights should add, because they are distinct ordered supports.

A support of size $k$ uses exactly $k$ vertices in each shore. It therefore contributes $x$-degree $k$ and $z$-degree $|C|-k$, as required in \eqref{pgr:r:ps:eq:physicalF}. The all-dummy term survives as $\prod_{c\in C}z_c$, with coefficient one. Finally, nonnegative activities follow by limits in the closed Lorentzian cone; the same coefficient-one term prevents zero output.
\end{proof}


\mnote{with every arc oriented from $C$ to $I$ this is the two-shore
setting of \prp, Theorem \texttt{lor:thm:weightedULC}, normalized by the
shore order. Corollary~\ref{pgr:r:ps:cor:shore} is an order-$|C|$
statement; it gives rank-ULC only under the saturation hypothesis stated
there.}

\begin{corollary}\label{pgr:r:ps:cor:shore}
The coefficients of $\Ga_D$ satisfy order-$|C|$ ULC, and likewise order-$|I|$ ULC. If the positive-weight relation admits a matching saturating either shore, they satisfy actual-degree rank-ULC.
\end{corollary}
\begin{proof}
Set every $z_c=s$ and every $x_i=t$ in \eqref{pgr:r:ps:eq:physicalF}. The result is
$\sum_k\gamma_ks^{|C|-k}t^k$, which remains Lorentzian. Apply the binary characterization, and interchange the shores for the second statement.
\end{proof}

For a disconnected graph, each component's bipartition may be reversed independently. Placing the smaller shore of each nontrivial component in $C$, and isolated vertices in $I$, gives the order
\begin{equation}\label{pgr:r:ps:eq:componentorder}
 r=\sum_{K\text{ nontrivial component}}\min\{|C_K|,|I_K|\}.
\end{equation}
The theorem gives actual-degree ULC if every chosen smaller shore is saturated in $D_+$. The same observation can be applied after zero-weight arcs are removed.

\begin{corollary}[Overlapping independent role covers]\label{pgr:r:ps:cor:overlap}
Suppose every arc $i\to j$ satisfies $i\in P$ or $j\in Q$, where $P,Q\subseteq V$ may overlap. If the physical core $C=P\cup Q$ is independent in the underlying graph, then $\Ga_D$ satisfies order-$|P\cup Q|$ ULC. A positive-weight matching saturating $C$ makes this actual-degree rank-ULC.
\end{corollary}
\begin{proof}
There are no exterior-to-exterior arcs by the role-cover condition and no core-to-core arcs by independence. Thus $C$ and $V\setminus C$ are a physical bipartition. Apply Corollary~\ref{pgr:r:ps:cor:shore}.
\end{proof}
For a cover of two tail roles and three head roles, this improves the cover order from five to four or three when the selected physical sets intersect in one or two vertices. It does not allow internal physical core arcs.

\subsection{Pendant heads attached to a stable transversal seed}
\begin{theorem}\label{pgr:r:ps:thm:pendant}
Fix $Q\subseteq V$, and suppose every vertex $j\notin Q$ has at most one incoming neighbor in $D$. Form the transversal matroid $M_Q$ with positions $Q$, all tail copies $i^+$ as ground elements adjacent to $N_D^+(i)\cap Q$, and a private dummy $d_q$ adjacent only to each $q\in Q$. If its ordinary basis polynomial $B_{M_Q}$ is real stable, then the signed physical monomer polynomial
\begin{equation}\label{pgr:r:ps:eq:monomer}
 \mu_D(\mathbf z)=\sum_{(S,T)\text{ feasible}}
 (-1)^{|S|}\wt(S,T)\prod_{i\notin S\cup T}z_i
\end{equation}
is real stable for every choice of nonnegative independent activities. Consequently $\Ga_D$ has only negative real zeros and is rank-ULC at its actual surviving degree.
\end{theorem}

\mnote{source~22's one-tail side criterion
(Theorem~\pgrfwd{pgr:r:ir:thm:oneside}{``oneside'' of source~22}) is the
case of this theorem with a single tail hub $p$: there every head outside
$Q$ has at most the one in-neighbour $p$, and the matroid is the same.
Source~22's proof (seed identity, pendant insertion, physical merge) is
the proof below and is replaced there by a pointer. The seed construction
is also used by batch-73 manuscript~08, printed in \mrn.}

There is no bound on $|Q|$, on the number of distinct tail hubs, or on physical overlap. If a role cover $P^+\cup Q^-$ is specified, the condition says that every uncovered head has degree at most one; the unique parents may be different elements of $P$. Internal physical core arcs are permitted.

\begin{proof}
Split every physical vertex into tail and head copies. First retain all tail copies and only the head copies in $Q$, with the specified arcs into $Q$. Write $a_i$ for tail-copy monomers and $b_q$ for these head-copy monomers. For positive activities, the split seed has signed monomer polynomial
\begin{equation}\label{pgr:r:ps:eq:seed}
 \mu_0(\mathbf a,\mathbf b)=
 \left(\prod_{i\in V}a_i\right)\left(\prod_{q\in Q}v_q\right)
 B_{M_Q}\left(x_i=-\frac{u_i}{a_i},\ d_q=\frac{b_q}{v_q}\right).
\end{equation}
A basis consists of selected tails $S$ and dummies for the unused heads $Q\setminus T$; its basis condition says that the selected tail copies $S$ can be matched to the head copies $T$ in the split role graph. No physical disjointness is imposed yet. Its transformed monomial is
\[
 (-1)^{|S|}u^Sv^T\prod_{i\notin S}a_i\prod_{q\in Q\setminus T}b_q.
\]
Thus \eqref{pgr:r:ps:eq:seed} counts supports once. All denominators cancel by multiaffinity. For upper-half-plane monomer arguments, both $-u_i/a_i$ and $b_q/v_q$ lie in the upper half-plane. Stability of the basis polynomial therefore gives stability of $\mu_0$.

Now insert each missing head copy $j^-$, $j\notin Q$. An isolated head contributes the factor $b_j$. Otherwise let $i^+$ be its unique parent and put $w=u_iv_j$. The exact support recurrence is
\begin{equation}\label{pgr:r:ps:eq:pendantrec}
 \mu_{\rm new}=b_j\mu_{\rm old}-w\partial_{a_i}\mu_{\rm old}
 =b_j\mu_{\rm old}\left(a_i-\frac{w}{b_j}\right),
\end{equation}
with other arguments unchanged. Indeed, a support using the new head must contain the forced edge $i^+\to j^-$; removing its endpoints leaves exactly a support in which $i^+$ is unused. This proves the recurrence without matching multiplicity. Its second equality uses affinity in $a_i$.

For $b_j\in\HH$ and $w\ge0$, the number $-w/b_j$ has nonnegative imaginary part, so the shifted $a_i$ stays strictly upper. Equation~\eqref{pgr:r:ps:eq:pendantrec} proves stability after each insertion, at any sequence of tail hubs. Multiaffinity prevents a previously used hub from being used again.

Once all head copies are present, merge each physical pair by \eqref{pgr:r:ps:eq:merge}. Exactly the supports with disjoint physical endpoint sets survive, giving \eqref{pgr:r:ps:eq:monomer}. The coefficient-one all-unused monomial survives every stage. Nonnegative activities follow by limits and this monomial excludes an identically zero limit.

With $n=|V|$, diagonal specialization yields
\begin{equation}\label{pgr:r:ps:eq:roottransfer}
 \mu_D(s,\ldots,s)=s^n\Ga_D(-s^{-2}).
\end{equation}
The left side has only real zeros. A zero $\rho$ of $\Ga_D$ that is not negative real would give a nonreal solution of $s^2=-1/\rho$, contradicting \eqref{pgr:r:ps:eq:roottransfer}. Since $\Ga_D(0)=1$, all zeros are strictly negative. Newton's inequalities apply at the actual remaining degree, including after weights vanish.
\end{proof}

\subsection{An explicit stable seed in arbitrary rank}
\begin{corollary}\label{pgr:r:ps:cor:elementary}
Under the pendant-head assumption of Theorem~\ref{pgr:r:ps:thm:pendant}, suppose every nonempty set $N_D^+(i)\cap Q$ is a singleton or all of $Q$. Then its conclusions hold, without any additional matroid hypothesis.
\end{corollary}
\begin{proof}
Write $r=|Q|\ge2$. For $q\in Q$, let
\[
 R_q=d_q+\sum_{\{i:\,N_D^+(i)\cap Q=\{q\}\}}x_i,
\]
and let $g_1,\ldots,g_m$ be the separate variables for tails whose neighborhood is all of $Q$. Empty-neighborhood elements are loops. The basis polynomial is
\begin{equation}\label{pgr:r:ps:eq:elementary}
 B_{M_Q}=e_r\bigl((R_q)_{q\in Q},g_1,\ldots,g_m\bigr).
\end{equation}
To see this, singleton elements with the same neighbor, together with their private dummy, form one parallel class. An $r$-element set can be matched precisely when it uses at most one element from each such class. Its remaining universal elements can then be assigned to the unused positions. This is exactly \eqref{pgr:r:ps:eq:elementary}; distinct universal elements must not be merged into a single parallel class.

The elementary symmetric polynomial is stable: differentiate the stable product $\prod_j(s+y_j)$ sufficiently many times in $s$, then specialize $s=0$. Replacing variables by nonnegative class sums preserves stability. For $r=1$ the basis polynomial is linear, with no separate universal class; for $r=0$ the seed is edgeless and the pendant construction starts from the product of tail monomers.
\end{proof}
If a physical tail belongs to $Q$, a universal neighborhood would include a loop and is therefore disallowed by the original loopless hypothesis. There is no conflict with the corollary: the condition is checked on the actual relation.

\paragraph{An overlapping non-bipartite example.}
Let $V=\{0,\ldots,8\}$, $P=\{0,1\}$, and $Q=\{1,2,3\}$. Include only
\[
 0\to1,2,3,7;\qquad 1\to0,2,8;\qquad
 i\to1,2,3\quad(i=4,5,6).
\]
The physical graph contains the triangle $0,1,2$, so Theorem~\ref{pgr:r:ps:thm:bipartite} does not apply. The selected role cover overlaps at vertex $1$. Every uncovered head has at most one parent, and every $Q$-neighborhood is empty, singleton or universal. Corollary~\ref{pgr:r:ps:cor:elementary} therefore gives real-rootedness for all independent nonnegative activities. At unit activities,
\begin{equation}\label{pgr:r:ps:eq:example}
 \Ga_D(t)=1+16t+47t^2+31t^3+4t^4.
\end{equation}
There are $99$ feasible supports, but $163$ directed matching witnesses. The independent checkers recover \eqref{pgr:r:ps:eq:example} by deduplicating endpoint sets.

\subsection{Reproducibility and the remaining question}
The proofs above are unbounded ordinary mathematical arguments. The package includes two independently written audit programs, their exact-check receipts, and the producer's separate checkers. They corroborate support identities, boundary cases and examples; finite tests are not the premise of either theorem.

For the bipartite construction, the independent audit checks $5\,120$ graph--weight instances, including all $4\,096$ directed subrelations of $K_{2,3}$. It compares direct physical support enumeration with Hall-based transversal constructions, tests symmetric basis exchange, and performs $13\,152$ exact rational Hessian-inertia checks. The separate producer check covers $12\,800$ graph--weight--shore tests, including $4\,345$ zero-activity cases.

For the pendant criterion, the independent audit checks $1\,884$ weighted cases, including all $1\,024$ eligible loopless four-vertex relations with $|Q|=2$, and random seed ranks up to five. These include $1\,734$ zero-activity cases and $880$ activity-induced degree drops. It checks exact operator/support identities, the elementary-symmetric Hall identity, negative real roots at the surviving degree, and $70\,752$ exact Rayleigh evaluations. The producer separately checks $480$ random weighted cases and \eqref{pgr:r:ps:eq:example} by rational Sturm sequences. Rayleigh evaluations at finitely many points are supplementary checks, not stability certificates.

The two structural conclusions have different normalization guarantees. The bipartite theorem gives the shore order, or \eqref{pgr:r:ps:eq:componentorder}; saturation is required to identify that order with the actual degree. The stable-seed theorem instead gives real-rootedness and hence actual-degree ULC directly, even on degree-dropping zero-weight faces.

\paragraph{Open question.}
Does every finite directed relation whose bipartite role graph has a vertex cover comprising at most two tail copies and three head copies have a rank-ULC support polynomial for all nonnegative independent activities, with physical overlap allowed? The present results cover independent physical cores with suitable saturation, and stable transversal seeds with pendant uncovered heads. They do not resolve the general internal-core case or all deficient-rank cases. Real-rootedness is a stronger property and is not asserted for the entire class.

\mnote{the case of two tail and two head roles is settled in the
strongest sense by source~20 (Theorem~\ref{pgr:r:rc:thm:rolecover}), which
source~15 does not cite; the case $P\subseteq Q$ of the question by
source~11 (Corollary~\ref{pgr:r:pc:cor:rolecover}). The question as stated,
with $|Q|=3$, internal core arcs and a physical union of four or five
vertices, remains open.}


For comparison, the bimatroid ULC theorem of R\"ohrle and Ulirsch \cite[Theorem A]{RU} is normalized by the smaller ground-side size, not by the actual bimatroid rank. Likewise, the current revision of Alexandersson and Jal \cite[Corollary 4.1]{AJ} uses the number of board rows. Neither is used here as an unproved actual-degree shortcut. All stability assumptions in Theorem~\ref{pgr:r:ps:thm:pendant} are explicit; arbitrary transversal matroids are not assumed to have the half-plane property.


\section{Independent role activities through actual degree three (source 22)}\label{pgr:r:ir:sec}

Source~22 is \emph{Independently weighted preorder support polynomials of
actual degree at most three} (21 pages as delivered; archive
\texttt{independent-role-degree-three-package}). Its files are shipped with
the prefix \texttt{22-independent-role-}; its nested audit package is
shipped as the delivered zip, the only carrier of its role certificates.
Archive~31, \emph{Vertex-weighted preorder support polynomials of actual
degree at most three}, is an earlier version of this manuscript, with one
activity per vertex; its theorem is Corollary~\ref{pgr:r:ir:cor:vertex}
below, and it is not printed separately.

\paragraph{Abstract of source 22, as delivered.}

For a finite preorder, we count ordered disjoint tail and head supports admitting a directed matching, counting each support once. Assign arbitrary nonnegative activities independently to the tail and head role of every vertex. We prove ultra-log-concavity with binomial normalization at the actual surviving degree whenever that degree is at most three, with no bound on the number of vertices. The proof combines the Gallai--Edmonds decomposition with ordinary structural inequalities and a complete exact verification through seven vertices. Its main ingredients are Motzkin--Straus, an all-rank articulation theorem, rank-three matroid Rayleigh inequalities, two-attachment compression and moment cones, and rational nonnegative identities. A half-plane-property side-matroid criterion resolves the final finite cases. One ten-element side matroid is certified by twenty half-plane-property minors and an exact rank-32 rational Gram identity for a Rayleigh difference, using Wagner--Wei's criterion. Published certified matroid data are explicit dependencies. Portable exact checkers and all local verification inputs are included. General real-rootedness fails already in degree three, as an explicit example shows. This is a computer-assisted proof, not a proof-assistant formalization or a literature-wide priority claim.

\subsection{Definitions and the main theorem}\label{pgr:r:ir:sec:main}
Let $R$ be a preorder on a finite set $V$: it is reflexive and transitive, but need not be antisymmetric. Its directed arcs are the \emph{off-diagonal} pairs $i\to j$ with $i\ne j$ and $iRj$. Reflexive loops never participate in a matching.

\begin{definition}
A feasible support of size $k$ is an ordered pair $(S,T)$ of subsets of $V$ such that $S\cap T=\varnothing$, $|S|=|T|=k$, and some bijection $\pi:S\to T$ satisfies $i\to\pi(i)$ for every $i\in S$. For independent nonnegative tail and head activities $\mathbf u=(u_i)_{i\in V}$ and $\mathbf v=(v_i)_{i\in V}$, put
\begin{equation}\label{pgr:r:ir:eq:roledef}
 \gamma_k=\sum_{(S,T)\text{ feasible},\ |S|=k}u^Sv^T,
 \qquad \Ga_R(z;\mathbf u,\mathbf v)=\sum_{k\ge0}\gamma_k z^k,
\end{equation}
where $u^Sv^T=\prod_{i\in S}u_i\prod_{j\in T}v_j$. The empty support contributes $\gamma_0=1$.
\end{definition}
A support is counted once, irrespective of the number of witnessing bijections. However, \emph{different ordered pairs with the same physical union are different summands}. Thus the object is neither an unordered endpoint-set enumerator nor a matching-witness enumerator. Physical disjointness $S\cap T=\varnothing$ is essential.

The single-vertex-activity model is the specialization $u_i=v_i=x_i$. It weights a support by $\prod_{i\in S\cup T}x_i$. The theorem below includes this specialization but does not require it.

\begin{definition}
For a nonnegative sequence $(g_0,\ldots,g_d)$ with $g_0,g_d>0$ and no internal zeros, \emph{rank-ultra-log-concavity}, abbreviated rank-ULC, means
\begin{equation}\label{pgr:r:ir:eq:ulc}
 \left(\frac{g_k}{\binom dk}\right)^2\ge
 \frac{g_{k-1}}{\binom d{k-1}}\frac{g_{k+1}}{\binom d{k+1}}
 \qquad(1\le k<d).
\end{equation}
Here $d$ is the actual polynomial degree, not a degree obtained by adding zero coefficients.
\end{definition}

\begin{theorem}[Independent-role degree-three theorem]\label{pgr:r:ir:thm:main}
For every finite preorder $R$ and every assignment $u_i,v_i\ge0$, if the actual surviving degree $d=\deg\Ga_R(z;\mathbf u,\mathbf v)$ is at most three, then $\Ga_R(z;\mathbf u,\mathbf v)$ is rank-ULC. In particular,
\begin{alignat}{2}
 d=2:&\qquad \gamma_1^2\ge4\gamma_2,\label{pgr:r:ir:eq:d2}\\
 d=3:&\qquad \gamma_1^2\ge3\gamma_2,
 &\qquad \gamma_2^2\ge3\gamma_1\gamma_3.\label{pgr:r:ir:eq:d3}
\end{alignat}
The conditions in degrees zero and one are vacuous.
\end{theorem}

\subsubsection{What the theorem does and does not say}
The number of vertices is unrestricted. In particular, this is not merely an enumeration of preorders on seven vertices. The reduction to bounded cores, with arbitrary exterior populations and heterogeneous activities, is part of the proof.

Several structural ingredients below hold even for directed relations that are not transitive. Their scopes are stated separately. The global theorem uses transitivity in its bounded-core reduction. It does not imply real-rootedness; see Section~\ref{pgr:r:ir:sec:nonreal}. In particular, stability of the auxiliary matroids used for certain finite targets is a sufficient local tool, not an assertion about every preorder support polynomial.

The accompanying package separates ordinary proofs, finite polynomial identities, exhaustive finite-domain coverage, and audit receipts. An audit receipt identifies exactly what was checked; it is not a substitute for the mathematical reduction. Numerical optimization was used in discovery, but no numerical solver is needed to verify any certificate.

\subsection{Degree, zero activities, and the universal first gap}\label{pgr:r:ir:sec:first}
Let $G_R$ be the simple undirected graph whose edge $ij$ is present when at least one of $i\to j$ and $j\to i$ is present. Write $\nu(G_R)$ for its matching number.

\begin{lemma}\label{pgr:r:ir:lem:degree}
If all relevant activities are positive, then $\deg\Ga_R=\nu(G_R)$ and the coefficients up to the actual degree are positive. With nonnegative independent role activities, there is a preorder with positive activities and the same support polynomial, whose physical matching number equals the original surviving degree.
\end{lemma}
\begin{proof}
A witnessing bijection gives an undirected matching. Conversely, orient each edge of an undirected matching along a direction available in $R$. Its disjoint tail and head sets form a feasible support with positive weight. Submatchings give positivity at every lower degree.

For the second statement let $U=\{i:u_i>0\}$ and $W=\{j:v_j>0\}$. Retain exactly the original off-diagonal arcs $i\to j$ with $i\in U$, $j\in W$, and then add all reflexive loops. This relation is transitive. Indeed, a surviving nontrivial path $i\to j\to k$ has $i\in U$, $k\in W$ and the original transitive relation $iRk$; its shortcut survives when $i\ne k$, and otherwise the required loop was added. Steps containing loops are harmless.

A positive-weight original support uses only retained arcs. Conversely every support of the retained relation uses only previously positive role activities. Set a previously zero activity to $1$ in its now-unused role. No new support is created, so the polynomial is unchanged and all activities are positive. The first statement now identifies its matching number with the original surviving degree.
\end{proof}
This filtering argument is stronger than vertex deletion: a vertex may still be usable in one role when its other role has zero activity. It proves the actual-degree scope even if the original, unfiltered comparability graph has matching number larger than three. In the vertex specialization ordinary deletion of zero-activity vertices is enough.

\mnote{the first statement of this lemma is Lemma~\ref{pgr:lem:degree}
(\prp, Theorem \texttt{gam:thm:deficiency}, for preorders). The second,
the transitivity-preserving filter for zero role activities, is source~22's
own; source~22's earlier version (archive~31) mentioned it only as a
remark.}


\mnote{source~22's Proposition ``Universal first inequality'' (for any
finite loopless directed relation with positive independent role
activities and actual degree $d\ge2$, $\gamma_1^2\ge\frac{2d}{d-1}\gamma_2$,
and the same at the surviving degree for nonnegative activities) is
Theorem~\ref{pgr:thm:first}; its proof by Motzkin--Straus on the graph of
arcs with four distinct endpoints is the proof printed there.}

For degree two the quadratic $1+\gamma_1 z+\gamma_2z^2$ thus has negative real roots. Degree-one polynomials do as well. This observation will handle disconnected graphs without needing a general convolution theorem for ULC sequences.

\begin{lemma}[Products over components]\label{pgr:r:ir:lem:products}
If $R$ has underlying connected components $R_1,\ldots,R_m$, then $\Ga_R=\prod_i\Ga_{R_i}$. At total degree at most three, rank-ULC follows if the unique degree-three component, when present, is rank-ULC; otherwise the product is real-rooted.
\end{lemma}
\begin{proof}
Every feasible support splits uniquely among components, and witnesses can be combined. If a component has degree three, all others have degree zero. Otherwise every component has degree at most two and is real-rooted by Proposition~\ref{pgr:thm:first}. Newton's inequalities apply to their product at its actual degree.
\end{proof}

\subsection{An ordinary all-rank articulation argument}\label{pgr:r:ir:sec:articulation}
This structural result does not require transitivity and is useful beyond the main theorem.
\begin{theorem}\label{pgr:r:ir:thm:articulation}
Let $D$ be a loopless directed relation and $h$ a vertex such that every component of the underlying graph of $D-h$ has matching number at most one. For arbitrary nonnegative independent role activities, $\Ga_D$ has only negative real zeros. In particular it is rank-ULC at its actual degree.
\end{theorem}
\begin{proof}
Let the components after deleting $h$ be $C_1,\ldots,C_m$. Write
\[
 Q_i(z)=\Ga_{C_i}(z)=1+a_i z,
 \qquad
 P_i(z)=\Ga_{C_i\cup\{h\}}(z)-Q_i(z)=c_i z+d_i z^2.
\]
All coefficients are nonnegative, and
\begin{equation}\label{pgr:r:ir:eq:artbound}
 d_i\le a_i c_i.
\end{equation}
Indeed, a size-two support counted by $d_i$ has a witness consisting of an arc incident to $h$ and a disjoint arc inside $C_i$. The product $a_i c_i$ contains that support's weight at least once. Extra witnesses and intersecting arc pairs only add nonnegative terms. This is a coefficientwise inequality in the role activities.

At support level, the exact identity is
\begin{equation}\label{pgr:r:ir:eq:artdecomp}
 \Ga_D=\prod_i Q_i+\sum_i P_i\prod_{j\ne i}Q_j.
\end{equation}
A support omitting $h$ contributes to the product. If it uses $h$, precisely one component contains an odd number of its selected endpoints: the component paired with $h$. Every other component has an even number. Thus that component is determined by the support itself, independent of any choice of witness; there is no duplicate counting in the sum.

For $a_i>0$,
\begin{equation}\label{pgr:r:ir:eq:pickart}
 \frac{P_i(z)}{Q_i(z)}=\frac{d_i}{a_i}z+\left(c_i-\frac{d_i}{a_i}\right)\frac{z}{1+a_i z}.
\end{equation}
Both coefficients are nonnegative by \eqref{pgr:r:ir:eq:artbound}. If $a_i=0$, then $d_i=0$ and the quotient is $c_i z$. For $z\in\HH=\{z:\operatorname{Im}z>0\}$,
\[
 \operatorname{Im}\frac{z}{1+a z}=\frac{\operatorname{Im}z}{|1+a z|^2}>0\qquad(a\ge0).
\]
Hence $1+\sum_iP_i/Q_i$ is either the constant $1$ or has positive imaginary part on $\HH$. It has no zero there, and neither does any $Q_i$. Equation~\eqref{pgr:r:ir:eq:artdecomp} excludes upper-half-plane zeros of $\Ga_D$; conjugation excludes lower-half-plane zeros. Its nonnegative coefficients and constant term $1$ exclude nonnegative real roots. The argument includes zero activities directly.
\end{proof}
The odd-endpoint argument in \eqref{pgr:r:ir:eq:artdecomp} is important: simply adding local matching polynomials would generally count witnesses instead of supports.

\subsection{Three active tails: a rank-three Rayleigh proof}\label{pgr:r:ir:sec:rayleigh}
\begin{theorem}\label{pgr:r:ir:thm:three}
Let a loopless directed relation have a three-vertex set $C=\{0,1,2\}$ and an independent exterior $X$. Allow arbitrary internal arcs in $C$ and arbitrary arcs from $C$ to $X$, but no arcs from $X$ to $C$. Then its support polynomial is rank-ULC for arbitrary nonnegative independent role activities. The order-dual assertion also holds.
\end{theorem}
Only the second cubic inequality needs proof, by Proposition~\ref{pgr:thm:first}. Transitivity is not assumed in this theorem.

\mnote{this theorem is a special case of source~11's
Theorem~\ref{pgr:r:pc:thm:main}: the three core vertices form a physical
vertex cover, since there are no exterior-to-exterior arcs. Source~11's
proof rests on source~12's computer-assisted boundary inequality; the proof
below is a second route for this special case, ordinary apart from
Wagner's rank-three Rayleigh theorem. Its inequality $B_{ij}B_{ik}\ge
A_iC_0$ is the scalar one-orientation case of source~12's
Theorem~\ref{pgr:r:tc:thm:boundary}.}


\subsubsection{Exterior head sets and private dummies}
Let $N_i\subseteq X$ be the exterior heads available to tail $i$. Define
\begin{align*}
 A_i&=\sum_{x\in N_i}v_x,\\
 B_{ij}&=\sum_{\substack{\{x,y\}\subseteq X\text{ matchable from }\{i,j\}}}v_xv_y,\\
 C_0&=\sum_{\substack{\{x,y,z\}\subseteq X\text{ matchable from }\{0,1,2\}}}v_xv_yv_z.
\end{align*}
Each head set is counted once. For distinct $i,j,k$ we use
\begin{equation}\label{pgr:r:ir:eq:exteriorineq}
 B_{ij}B_{ik}\ge A_i C_0,
 \qquad A_k B_{ij}\ge C_0.
\end{equation}
For the first inequality, adjoin private dummy heads $d_i$, each adjacent only to tail $i$. The resulting transversal matroid on heads and dummies has rank exactly three. Its basis polynomial, after substituting the exterior activities, is
\begin{equation}\label{pgr:r:ir:eq:dummy}
 M(\mathbf d)=d_0d_1d_2+\sum_i A_i d_jd_k+\sum_{i<j}B_{ij}d_k+C_0.
\end{equation}
Wagner's theorem states that every rank-three matroid is Rayleigh \cite[Theorem 1.1]{WagnerR3}. Thus
$\partial_jM\,\partial_kM-M\partial_{jk}M\ge0$ on the nonnegative orthant. Setting the three dummy variables to zero, by continuity, gives $B_{ij}B_{ik}\ge A_i C_0$.

The second inequality is coefficientwise. For each feasible three-head set choose a witness and its head paired with $k$. The remaining pair is feasible for $\{i,j\}$ and gives the same monomial in $A_k B_{ij}$. Different three-head sets give different squarefree monomials; repeated-head monomials or extra witnesses cause no difficulty.

\subsubsection{Restoring every internal arc}
Let $\epsilon_{ij}=1$ when $i\to j$ is an internal arc, and $0$ otherwise. For distinct $i,j,k$, let
\[
 H_j=\sum_{x\in (N_k\text{ if }\epsilon_{ij}=1)\,\cup\,(N_i\text{ if }\epsilon_{kj}=1)}v_x.
\]
An absent set in this notation means the empty set. This union, rather than a sum with possible duplication, counts exterior heads which can be used alongside core head $j$. Set
\[
 a_i=A_i+\epsilon_{ij}v_j+\epsilon_{ik}v_k,
 \qquad b_{ij}=B_{ij}+v_kH_k.
\]
Every tail lies in $C$, so direct support counting gives
\begin{equation}\label{pgr:r:ir:eq:threecoeff}
 \gamma_1=\sum_i u_i a_i,
 \qquad \gamma_2=\sum_{i<j}u_i u_j b_{ij},
 \qquad \gamma_3=u_0u_1u_2C_0.
\end{equation}
In particular a size-three support uses every core vertex as a tail and every head in the exterior.

The Rayleigh-type inequalities survive the restoration:
\begin{align}\label{pgr:r:ir:eq:restore}
 b_{ij}b_{ik}-a_iC_0={}&(B_{ij}B_{ik}-A_iC_0)\notag\\
 &+v_j(B_{ij}H_j-\epsilon_{ij}C_0)
 +v_k(B_{ik}H_k-\epsilon_{ik}C_0)+v_jv_kH_jH_k\ge0.
\end{align}
If $\epsilon_{ij}=0$ its bracket is nonnegative. If it is $1$, then $H_j\ge A_k$, and the second inequality of \eqref{pgr:r:ir:eq:exteriorineq} proves the same conclusion. The other bracket is symmetric; all remaining terms are nonnegative.

Put $U=u_0u_1u_2$ and $x_{ij}=u_i u_jb_{ij}$. Expansion yields the exact identity
\[
 \gamma_2^2-3\gamma_1\gamma_3
 =\frac12\bigl[(x_{01}-x_{02})^2+(x_{01}-x_{12})^2+(x_{02}-x_{12})^2\bigr]
 +3U\sum_i u_i(b_{ij}b_{ik}-a_iC_0).
\]
\mnote{this is the three-squares identity \eqref{pgr:r:tc:eq:three-squares}
of source~12 with $b_{ij}$, $a_i$, $c$ replaced by $x_{ij}$, $u_ia_i$,
$UC_0$; source~22's label for it is redirected there.}
Equations~\eqref{pgr:r:ir:eq:restore}--\eqref{pgr:r:tc:eq:three-squares} prove Theorem~\ref{pgr:r:ir:thm:three}. If the degree drops below three, the stronger surviving-degree inequality is supplied by Proposition~\ref{pgr:thm:first}, not by padding the cubic formula.

\subsection{Role covers and a structural corollary}\label{pgr:r:ir:sec:rolecover}
The \emph{bipartite role graph} has a tail copy and a head copy of every physical vertex, with an edge $i_{\rm tail}j_{\rm head}$ for every arc $i\to j$. Its matching number can exceed the physical matching number: a role matching can use the two copies of a physical vertex simultaneously.

We use the following independently verified companion theorem, whose full ordinary proof and portable algebraic checker are included as a separate dependency (source~20's package; \mnote{source~22 restates here, as its Theorem ``Two-tail/two-head cover'', source~20's Theorem~\ref{pgr:r:rc:thm:rolecover} (Section~\ref{pgr:r:rc:sec}); the statement is not repeated, and source~22's references to it point there. Its embedded copy of source~20's package is not shipped.}).
Here real stability means nonvanishing whenever all variables lie in $\HH$. The companion proof first splits roles, treats every subset of the four possible core arcs, and then merges the copies of each physical vertex. In monomer variables the local merge sends $ab\mapsto z$, $a,b\mapsto1$, $1\mapsto0$. Its stable symbol is $z+r+s$. Thus simultaneous use of both copies is removed rather than silently allowed. Positive activities enter by scaling before merging. Zero activities follow by continuity; the empty-support monomial prevents a zero-polynomial limit. No assertion that an arbitrary transversal matroid is stable is used.

To pass from monomer stability to support real-rootedness, diagonalize:
$\mu(t,\ldots,t)=t^n\Ga(-t^{-2})$. The diagonal is real-rooted. A nonzero nonreal or nonnegative root of $\Ga$ would yield a forbidden nonreal root in this identity. The constant coefficient is $1$, so zero is not a root of $\Ga$. This also handles surviving-degree losses.

\begin{corollary}\label{pgr:r:ir:cor:rolerank}
Any finite loopless directed relation whose bipartite role graph has matching number at most three has a rank-ULC support polynomial under arbitrary nonnegative independent role activities.
\end{corollary}
\begin{proof}
By K\H onig's matching-cover theorem, choose a role-vertex cover of size at most three. If both shores occur, each contributes at most two vertices and Theorem~\ref{pgr:r:rc:thm:rolecover} applies, even if their physical sets overlap. If just one shore occurs, there are at most three active tails or at most three active heads. Theorem~\ref{pgr:r:ir:thm:three} or its dual applies. Missing active vertices can be supplied as isolated, zero-activity placeholders. Lower surviving degrees are covered by Proposition~\ref{pgr:thm:first}.
\end{proof}

\mnote{the rank-ULC conclusion of this corollary is contained in source~11's
Theorem~\ref{pgr:r:pc:thm:main}, since the physical images of a role cover
of size three cover every underlying edge. When both shores of the role
cover occur, the proof gives real-rootedness through source~20, which is
more. For bipartite relations with one role per vertex the corollary is
\prp, Theorem \texttt{mr:thm:rankthree} and Corollary
\texttt{lc:cor:rankthree}.}

This corollary does not by itself prove Theorem~\ref{pgr:r:ir:thm:main}: the bipartite role matching number may exceed three even when the physical matching number is three.

\subsection{The finite theorem through seven vertices}\label{pgr:r:ir:sec:finite}
\begin{theorem}[Finite independent-role theorem]\label{pgr:r:ir:thm:finite}
Every preorder on at most seven vertices has a rank-ULC support polynomial under arbitrary nonnegative independent role activities, at its actual surviving degree.
\end{theorem}
Since these orders have physical matching number at most three, only the second cubic gap requires new work. We describe both the exact finite domain and the scope of its certificates.

\subsubsection{A complete enumeration rather than a sample}
The equivalence relation $i\sim j$ defined by $iRj$ and $jRi$ partitions any preorder into equivalence classes. The quotient is a poset. Choose a linear extension and label its classes $1,\ldots,m$. Every quotient arc then goes forward. Conversely, a transitive forward relation on $m$ labels and a positive composition $n=n_1+\cdots+n_m$ determine a preorder on $n$ vertices by expanding each class into $n_i$ equivalent vertices.

The verifier enumerates every forward relation by finite bit masks, tests transitivity, and enumerates every positive composition. This is redundant, but completeness is immediate from the quotient construction. Isomorphisms and global order duality are then checked explicitly; weights are transported with their vertices, and duality exchanges the tail and head activity families.

For the connected, nonbipartite, nontotal, physical-rank-three seven-vertex domain there are
\begin{center}
\begin{tabular}{lr}
Quotient/composition configurations&$131\,645$\\
Relevant configurations&$75\,160$\\
Isomorphism/duality representatives&$1\,686$
\end{tabular}
\end{center}
The entire catalog and all maps are retained. The independent replay regenerates the domain; it does not infer completeness from the number of files in a certificate directory.

\subsubsection{Branches outside the residual catalog}
Degree at most two is universal. Disconnected cases are handled by Lemma~\ref{pgr:r:ir:lem:products}. A connected bipartite preorder with more than two vertices is a height-two poset: a directed path on three distinct vertices forces a triangle, and a nontrivial equivalence pair with a third neighbor also forces a triangle. Its physical and role matching numbers therefore agree. Corollary~\ref{pgr:r:ir:cor:rolerank} applies.

A total preorder has complete underlying graph, so positive-activity degree three means $n=6$ or $7$. The $32+64=96$ positive compositions reduce, by the extremal-block identities below and duality, to $24$ role-weighted targets. Exact rational identities verify every target. For the nontotal six-vertex residual domain, the same quotient procedure gives $7\,090$ configurations, $3\,059$ relevant configurations, and $223$ representatives. Their independent-role identities are all included.

\subsubsection{Exact certificates and independent support reconstruction}\label{pgr:r:ir:sec:certificates}
Let $w$ denote all independent role variables. A typical certificate for the second cubic gap is
\begin{equation}\label{pgr:r:ir:eq:sos}
 p(w):=\gamma_2(w)^2-3\gamma_1(w)\gamma_3(w)
 =\sum_\ell q_\ell w^{\alpha_\ell}
       \left(\sum_\beta c_{\ell\beta}w^\beta\right)^2
  +\sum_\eta r_\eta w^\eta,
\end{equation}
where $q_\ell>0$, $c_{\ell\beta}\in\mathbb Q$, $r_\eta\ge0$, and all exponents are nonnegative integers. The remainder is reconstructed by exact subtraction and checked coefficientwise, rather than accepted from a numerical residual.

The independent checkers reconstruct the feasible supports by two complementary methods. The main hybrid replay and the two-attachment auditors enumerate physical directed matchings and deduplicate their ordered support pairs. The separate six-vertex and total-preorder routines use Hall's criterion
\[
 |N(A)\cap T|\ge |A|\qquad\text{for every }A\subseteq S
\]
for disjoint equal-sized sets $(S,T)$, with explicit bijection cross-checks. In either method, a feasible pair contributes its monomial once. In the vertex specialization, different pairs that specialize to the same monomial have their coefficients \emph{added}. This is where losing ordered-support multiplicity would invalidate the computation. Matching enumeration is used only to construct the set of supports, never to retain the number of matching witnesses.

All arithmetic in the replay uses integers and rational numbers. Floating-point optimization proposed square directions during discovery; those optimization programs and their binaries are not part of the proof package or needed by the verifier.

\subsubsection{Two structural reductions used in the selected ledger}
The finite ledger uses several ordinary theorems already proved above, as well as the two-attachment theorem of Section~\ref{pgr:r:ir:sec:a2}. Two other reductions are worth stating explicitly.

\paragraph{A marked five-vertex deletion lemma.}
Let $H$ be a preorder on five vertices with factor-critical underlying graph, and mark $h$. Write
\[
 Q=\Ga_{H-h}=1+az+bz^2,\qquad
 P=\Ga_H-Q=cz+dz^2.
\]
The finite role-weighted identity is
\begin{equation}\label{pgr:r:ir:eq:marked}
 acd-bc^2-d^2\ge0.
\end{equation}
There are $96$ marked isomorphism/duality types; $48$ margins are coefficientwise nonnegative and $48$ have rational square certificates. The collection contains $258$ square terms and $6\,971$ positive remainder monomials. The independent enumeration also checks the factor-critical condition, rather than presupposing it from a label.

For positive activities $b,c,d>0$ and $Q$ has negative real roots by Proposition~\ref{pgr:thm:first}. Equation~\eqref{pgr:r:ir:eq:marked} says $Q(-c/d)\le0$, so the root of $c+dz$ lies between the roots of $Q$. Thus $(c+dz)/Q$ has nonnegative residues at its negative poles, and $P/Q=z(c+dz)/Q$ maps $\HH$ into its closure. Repeated roots are handled by cancellation and limits.

Consequently the decomposition \eqref{pgr:r:ir:eq:artdecomp} proves real-rootedness whenever every block $C_i\cup\{h\}$ is a factor-critical preorder on three or five vertices. This is an all-rank gluing statement. In a seven-vertex factor-critical graph with a cut vertex $h$, the components after deleting $h$ have even sizes $2+4$ or $2+2+2$. Each block is factor-critical: for $v\in C_i$, parity forces $h$ to be matched into $C_i\setminus\{v\}$ in a perfect matching of $G-v$. Restricting that matching to $(C_i\setminus\{v\})\cup\{h\}$ proves the required deletion property; deleting $h$ itself leaves the perfectly matched $C_i$. Hence every such seven-vertex cut-vertex case is covered.

\paragraph{Refining small extremal equivalence classes.}
A minimal equivalence class has identical outgoing neighborhoods outside the class and no incoming outside arcs. For a class of size two, replacing equivalence by the chain $1<2$, keeping tails $u_1,u_2$, and setting
\begin{equation}\label{pgr:r:ir:eq:block2}
 v'_2=v_2+u_2v_1/u_1
\end{equation}
preserves the support polynomial. The first head activity in the chain is unused.

For a class of size three, sort the tail activities as $x\le y\le z$, carrying their head activities $v_1,v_2,v_3$ with them. Put
\[
 a=\sum_{i\ne j}u_i v_j,\qquad
 c=yzv_1+xzv_2+xyv_3,\qquad
 \Delta=z(x+y)-xy.
\]
Replace the class by a strict chain with the same tails and effective heads
\begin{equation}\label{pgr:r:ir:eq:block3}
 v'_2=\frac{(x+y)c-xy a}{x\Delta},\qquad
 v'_3=\frac{za-c}{\Delta}.
\end{equation}
These are nonnegative since
\begin{align*}
 (x+y)c-xy a&=y^2(z-x)v_1+x^2(z-y)v_2,\\
 za-c&=z^2(v_1+v_2)+\Delta v_3,
\end{align*}
and $\Delta>0$ for positive tails. They satisfy
\[
 xv'_2+(x+y)v'_3=a,
 \qquad xzv'_2+xyv'_3=c.
\]
These are precisely the local weights with one selected head and respectively one or two selected tails. Pure-tail states are unchanged. No other local head state is possible in a three-element minimal class. The surplus tails have identical outside neighborhoods, so equality of these local states preserves every continuation. Maximal classes are handled by duality. Disjoint chosen classes can be refined successively; each is refined once, and other classes' positive denominator activities are unchanged.

This transfer holds for any minimal or maximal class of size two or three, not just a boundary class of a total preorder. The full seven-vertex refinement manifest has $648$ checked relational maps, each strictly reducing the number of equivalent pairs. Dependencies are therefore acyclic.

\begin{remark}[Activity and dependency safeguards]\label{pgr:r:ir:rem:scope}
The transformations \eqref{pgr:r:ir:eq:block2}--\eqref{pgr:r:ir:eq:block3} use strictly positive original denominator activities. They may create zero effective heads, which are permitted by the closed-orthant target theorems. Original zero activities are handled by Lemma~\ref{pgr:r:ir:lem:degree}, or by continuity of the original gap; no finite sorted-chain replacement at arbitrary zero denominators is asserted. The transformations generally destroy $u=v$. Every refinement here terminates in a proof valid for fully independent role activities. No vertex-only target is used. Strict decrease of equivalent pairs makes the dependency graph acyclic, and the verifier requires a proved terminal target.
\end{remark}

\subsubsection{The complete selected independent-role ledger}
The final manifest partitions the $1\,686$ residual targets as follows. Labels identify the proof selected, rather than every applicable proof.
\begin{center}
\begin{tabular}{lr}
\toprule Selected justification&Number\\\midrule
Articulation real-rootedness&168\\
Two-tail/two-head role-cover real-rootedness&478\\
Three-core role-ULC&28\\
Factor-critical cut-vertex real-rootedness&83\\
Direct rational role certificate&812\\
Two-attachment role-ULC&43\\
HPP side matroid by small-matroid classification&27\\
HPP side matroid by published nine-element list&6\\
HPP side matroid by the final exact Rayleigh identity&1\\
Fully role-justified extremal refinement&40\\\midrule
Total&1\,686\\\bottomrule
\end{tabular}
\end{center}
The HPP branches are proved in Section~\ref{pgr:r:ir:sec:oneside} below. All refinements end in unconditional role-valid proofs; there are no remaining conditional sources or independent obligations.

Four selected direct certificates, with catalog IDs $722,736,756,758$, use a positive multiplier. They certify \eqref{pgr:r:ir:eq:sos} for $H(w)p(w)$ rather than $p(w)$, where
\[
 H=(u_4+u_5)(v_0+v_1)\quad\text{for }722,736,756,
 \qquad H=(u_5+u_6)(v_0+v_1)\quad\text{for }758.
\]
The exact multiplier is reconstructed and checked. It is strictly positive when all role variables are positive, so division proves $p\ge0$ there; the original polynomial gap extends to the closed orthant by continuity. This is a declared positive-orthant argument and is distinct from the all-real Rayleigh identity in Section~\ref{pgr:r:ir:sec:last}.

For reference, the six-vertex role certificates contain $586$ square terms and $31\,094$ positive remainders. The $24$ total-preorder role targets contain $793$ squares and $50\,660$ positive remainders. The full local catalog, exact identities and all selected terminal witnesses are included; counts describe proof data rather than sampled assignments. With the HPP arguments below, this establishes Theorem~\ref{pgr:r:ir:thm:finite}.

\subsection{Half-plane-property side matroids and the final cases}\label{pgr:r:ir:sec:oneside}
The final finite targets have a useful structural presentation. It leads to a general real-rootedness criterion, valid in every degree and without transitivity. The criterion is conditional on a matroid half-plane property; we verify that condition exactly in every application.

\subsubsection{The one-tail side criterion}
\begin{theorem}\label{pgr:r:ir:thm:oneside}
Let $D$ be a finite loopless directed relation. Suppose a physical vertex $p$ and a set $Q$ satisfy
\[
 i\to j\quad\Longrightarrow\quad i=p\ \text{or}\ j\in Q.
\]
The vertex $p$ may belong to $Q$. Form a transversal matroid $M$ with centers $Q$, one ground element for every tail copy $i$ with neighborhood $N(i)\cap Q$, and one private dummy element $d_j$ with neighborhood $\{j\}$ for each $j\in Q$. If the basis-generating polynomial $F_M$ is real stable, then the signed physical monomer polynomial of $D$ is real stable for arbitrary nonnegative independent role activities. Hence $\Ga_D$ has only negative real zeros, with rank-ULC at its actual degree.
\end{theorem}

The private dummies are an independent basis, so $M$ has rank $|Q|$. A basis is counted once, not once for each matching that witnesses it. For homogeneous real polynomials, the half-plane property (nonvanishing when all variables have positive real parts) is equivalent to real stability by a simultaneous rotation of variables. We use the abbreviation HPP.


\begin{proof}[Proof (pointer)]
Every head copy outside $Q$ has no possible predecessor except $p$, so the
hypothesis of source~15's Theorem~\ref{pgr:r:ps:thm:pendant} holds with
the single hub $p$, and the matroid $M$ is the matroid $M_Q$ there. That
theorem gives the conclusion.
\end{proof}
\mnote{source~22's own proof is source~15's proof in this special case:
the split seed identity $\mu_B=(\prod_ia_i)(\prod_{j\in Q}v_j)
F_M(d_j=b_j/v_j,\ t_i=-u_i/a_i)$ (source~15's
\eqref{pgr:r:ps:eq:seed}), the forced pendant insertion
$f\mapsto b_xf-u_pv_x\partial_{a_p}f$ (\eqref{pgr:r:ps:eq:pendantrec}),
the physical merge, here realized as the directional derivative
$T_if=(\partial_{a_i}+\partial_{b_i})f|_{a_i=b_i=z_i/2}$, and the diagonal
root transfer. It is replaced by this pointer. Source~22 notes that
successive pendant additions cannot reuse $p$, and that the evolving
empty-support term keeps coefficient one through all merges.}

Reversing every arc and exchanging $u$ with $v$ gives the dual one-covered-head criterion.

\subsubsection{Simplification and the published HPP input}
An empty-neighborhood tail element is a loop. A singleton-neighborhood element with neighbors $\{j\}$ is parallel to the dummy $d_j$. These are the only parallel classes: two nonempty columns are a dependent pair exactly when their neighborhood union has size one. Repeated multi-neighbor columns must therefore \emph{not} be merged.

The side simplification consists of one private representative for each center and one separate element for each tail whose neighborhood has size at least two. HPP lifts back from the simplification by replacing a private representative's variable with the sum over its true parallel class; sums of upper-half-plane variables stay in that half-plane. Loops cause no difficulty.

We use two precisely bounded published inputs from Kummer--Sert \cite{KS,KSdata}:
\begin{enumerate}[leftmargin=*,itemsep=3pt]
\item The complete exceptions on at most seven elements are the five rank-three matroids $F_7,F_7^{-},F_7^{--},F_7^{-3},M(K_4)+e$ and their rank-four duals. Every new eight-element excluded minor is rank-four sparse paving. Thus an eight-element matroid whose one-element deletions and contractions are HPP is HPP if its rank is not four or it is not sparse paving \cite[Section 5, Theorem 5.2]{KS}.
\item The published file \texttt{n9r4Hpp.txt} contains $4\,125$ certified-positive rank-four nine-element matroids. Section 6 of \cite{KS} explains their rational positive-semidefinite Gram certifications. The separate lists of candidates and undetected cases are not used.
\end{enumerate}
We do not assume every small transversal matroid has HPP. In particular, the rank-four dual of $F_7^{-3}$ is transversal and is an exception. Our higher-rank side matroids are checked against the full exception/minor tests, including duality.

The small-matroid side screen supplies 27 selected catalog targets. It checks 58 candidate presentations, including 47 eight-element sides and 752 one-element minors. Each witness stores the working relation, direction reversal, $p,Q$, every tail-neighborhood mask, retained simplified columns, and the complete basis set. The independent audit reconstructs these objects and checks exact matroid isomorphisms.

Six further targets use the published nine-element positive list. Its basis fingerprints use increasing colex order on four-element subsets, equivalently increasing subset bit masks; an asterisk marks a basis. Rank-five sides are dualized to rank four before comparison. The selected applications are
\begin{center}
\begin{tabular}{lr}
\toprule Catalog targets&One-based published line\\\midrule
1030, 1474, 1480&1975\\
1057, 1486&1536\\
1495&1751\\\bottomrule
\end{tabular}
\end{center}
The data retain explicit ground-element permutations, not just invariant matches. All 126 basis indicators are compared after the permutation. Both possible side presentations were independently recounted, giving twelve comparisons of complete side basis systems.

The pinned official file has SHA256
\begin{center}\footnotesize\texttt{a3c7a6eaebe02ef50998710282542be2751b218a6713853cdd2eecadbe1f1f86}.\end{center}

\mnote{the published file \texttt{n9r4Hpp.txt} of Kummer and Sert is
distributed under CC~BY~4.0 (Zenodo record 6108027) and is shipped
unmodified, four times, inside source~22's nested audit zip; the report
README carries the attribution and licence notice.}

Its MD5 matches the publisher's \texttt{c73d80267112a960df5485f159ab56e0}. The exact data file is included for offline membership replay. This is an application of the published certified-positive list, not a local rerun of all of its original Gram certificates.

\subsubsection{The final ten-element side matroid}\label{pgr:r:ir:sec:last}
The final independent target, catalog ID 1565, has seven off-diagonal row masks
\begin{equation}\label{pgr:r:ir:eq:lastrows}
 (0,1,3,7,39,23,63).
\end{equation}
Choose $p=6$ and $Q=\{0,1,2,4,5\}$ in this order. Its side simplification is the rank-five transversal matroid $M$ on ten elements whose center-neighborhood masks are
\begin{equation}\label{pgr:r:ir:eq:lastmasks}
 (1,2,4,8,16,3,7,23,15,31).
\end{equation}
We index these elements $0,\ldots,9$. The first five private columns form a basis, and exact matching enumeration gives $198$ bases in total.

All twenty one-element deletions and contractions have HPP. Loops, coloops, genuine parallel classes and series classes are reduced with exact basis checks. HPP is restored by the standard loop/coloop and parallel operations, with series reduction obtained by duality. After reduction, the dependencies group as follows; ``small'' means the previously stated at-most-eight criterion, and a number is a one-based line of the positive nine-element file, with dualization when the reduced rank is five.
\begin{center}
\begin{tabular}{lcc}
\toprule Elements&Deletion&Contraction\\\midrule
$0,1,5$&3368&small\\
$2,6$&1751&2020\\
$3,4,7,8$&small&1536\\
$9$&small&2385\\\bottomrule
\end{tabular}
\end{center}
Eight minors use the small criterion and twelve use the published list. The independent replay verifies $10\,240$ minor rank-function identities, every reduction, and all explicit published-basis comparisons. These checks establish the proper-minor hypothesis only; one further all-real inequality is needed.

Let $F$ be the basis polynomial of $M$, and write $F_{ab}$ for the coefficient of $x_0^a x_9^b$, $a,b\in\{0,1\}$. Its Rayleigh difference is
\begin{equation}\label{pgr:r:ir:eq:lastrayleigh}
 q=\Delta_{09}F=F_{10}F_{01}-F_{00}F_{11}.
\end{equation}
It is independent of $x_0,x_9$. In the remaining element order $1,\ldots,8$, it is homogeneous of degree eight and has 397 nonzero monomials. Let $\mathcal B$ be the sorted set of the 37 squarefree degree-four exponent vectors $a$ for which the coefficient of $x^{2a}$ in $q$ is nonzero. Define
\begin{equation}\label{pgr:r:ir:eq:gramformula}
 N_e=\#\{(a,b)\in\mathcal B^2:a+b=e\},
 \qquad G_{ab}=\frac{[x^{a+b}]q}{N_{a+b}}.
\end{equation}
The exact support check gives $\operatorname{supp}(q)\subseteq\mathcal B+\mathcal B$. With $h=(x^a)_{a\in\mathcal B}$, the identity $q=h^TGh$ therefore follows by collecting coefficients. The matrix entries are only
\[
 0,\quad \frac18,\quad\frac16,\quad\frac14,\quad\frac13,\quad\frac12,\quad1.
\]

\begin{lemma}[Exact final Gram certificate]\label{pgr:r:ir:lem:lastgram}
The rational matrix in \eqref{pgr:r:ir:eq:gramformula} is positive semidefinite of rank 32. Equivalently, the supplied exact vectors $c_1,\ldots,c_{32}\in\mathbb Q^{37}$ and positive rational weights $w_1,\ldots,w_{32}$ satisfy
\begin{equation}\label{pgr:r:ir:eq:lastsos}
 G=\sum_{\ell=1}^{32}w_\ell c_\ell c_\ell^T,
 \qquad q=\sum_{\ell=1}^{32}w_\ell(c_\ell^T h)^2.
\end{equation}
In particular $q$ is nonnegative on all of $\mathbb R^8$.
\end{lemma}
\begin{proof}[Exact certificate verification]
The supplied complementary standard-library checker reconstructs all 198 bases by explicit assignment to the five centers, forms the four coefficient polynomials $F_{ab}$, and expands \eqref{pgr:r:ir:eq:lastrayleigh}. It independently reconstructs $\mathcal B$, every multiplicity $N_e$ and the uniform Gram formula. Exact expansion of both identities in \eqref{pgr:r:ir:eq:lastsos} has zero remainder. Finally it performs rational Schur elimination, obtaining 32 strictly positive pivots and a zero residual matrix, independently confirming the rank and positive semidefiniteness. All vectors and weights are included in the machine-readable certificate. The independent audit separately reconstructs the full ten-variable derivatives, verifies both the Gram expansion and the positive-square factorization, and confirms rank 32 by exact row elimination.

This is an all-real sum of ordinary polynomial squares. There is no uncontrolled-sign monomial prefactor and no merely coefficientwise-positive remainder. Thus it supplies precisely the all-real domain required below. Numerical projection found the matrix, but no floating-point output or tolerance is used in its exact verification.
\end{proof}

Wagner--Wei's criterion \cite[Theorem 3(c)]{WW} says that a multiaffine polynomial with positive coefficients is stable if every one-variable deletion and contraction is strongly Rayleigh and one Rayleigh difference is nonnegative on all real values of the remaining variables. For homogeneous matroid basis polynomials, HPP, real stability and the strong Rayleigh property are equivalent. The twenty proper minors and Lemma~\ref{pgr:r:ir:lem:lastgram} satisfy these hypotheses. Hence $M$ has HPP, and Theorem~\ref{pgr:r:ir:thm:oneside} proves real-rootedness for target 1565 under every nonnegative independent role assignment.

The existing approved extremal-class refinement maps the last conditional source, ID 1611, to this target. Effective zero activities are permitted by Theorem~\ref{pgr:r:ir:thm:oneside}. This closes the final two entries in the selected finite ledger.

\subsection{Two attachments and arbitrary exterior populations}\label{pgr:r:ir:sec:a2}
The following theorem supplies the two-attachment infinite-family input to the global reduction.
\begin{theorem}\label{pgr:r:ir:thm:a2}
Let a preorder of physical matching number at most three have a retained core $K$ of at most five vertices, with two marked attachments $i,j$, and an independent exterior whose neighbors lie in $\{i,j\}$. Suppose each attachment has a fixed direction toward all of its exterior neighbors and no exterior vertex belongs to a nontrivial equivalence class. In the canonical Gallai--Edmonds core family, the support polynomial is rank-ULC under arbitrary nonnegative independent role activities and arbitrary heterogeneous exterior populations.
\end{theorem}
The canonical-core qualification fixes the finite catalog used in the proof. We do not replace it by an unchecked assertion about arbitrary marked five-vertex directed cores. The catalog has $1\,084$ structural templates, including harmless overcoverage. Its generation and coverage are independently checked in the included two-attachment dependency. Degree at most two and the first cubic gap are settled by Proposition~\ref{pgr:thm:first}; below we prove the second cubic gap.

\subsubsection{Opposite signs: exact two-vertex compression}
Assume exterior arcs at $i$ point outward and those at $j$ inward. The exterior types are head-only neighbors of $i$, tail-only neighbors of $j$, and shared vertices $i\to s\to j$. Put
\[
 H=\sum_{s:\,i\to s}v_s,
 \qquad T=\sum_{s:\,s\to j}u_s,
 \qquad W=\sum_{s:\,i\to s\to j}u_s v_s.
\]
The exact support identity is
\begin{align}\label{pgr:r:ir:eq:mixedcompress}
 \Ga_R={}&\Ga_K+T v_jz\Ga_{K-j}+H u_i z\Ga_{K-i}\notag\\
 &+(HT-W)u_i v_j z^2\Ga_{K-\{i,j\}}.
\end{align}
A selected exterior tail must be matched to $j$ and a selected exterior head to $i$. Two selected exterior vertices must have distinct physical identities; $HT-W$ removes exactly the forbidden shared-vertex collisions. Removing their forced matched endpoints leaves a core support. This proves \eqref{pgr:r:ir:eq:mixedcompress} without counting witnessing matchings.

We have $0\le W\le HT$. If the shared type is legal, so are both private types: after deleting either its incoming or its outgoing incidence, no two-step transitivity requirement through that exterior vertex remains, and core closure of the retained neighborhood is inherited. Legality is also checked directly in the finite catalog. For $T>0$, use one shared vertex with
\[
 u_s=T,\quad v_s=W/T,
\]
and one private head at $i$ of activity $H-W/T$. This preserves $H,T,W$ and hence the whole polynomial. If $T=0$, then $W=0$ and only the head sum remains. If the shared type is not legal, sum the two private clouds separately. The replacement uses at most two exterior vertices, so has at most seven vertices.

Crucially, this compression need not preserve $u=v$. Its finite targets are proved under the stronger role weighting, independently of the global assembly. The finite certificate targets are replayed directly, so this use creates no circular dependence on the two-attachment branch of the selected ledger. Among the $494$ mixed-sign templates, the checked compressed branches are
\begin{center}
\begin{tabular}{lr}
At most six vertices&366\\
Physical matching rank at most two&19\\
Disconnected&49\\
Connected bipartite&3\\
Residual seven-vertex references&57
\end{tabular}
\end{center}
The last $57$ references use $56$ distinct direct role certificates, with $1\,135$ square terms and $6\,445$ positive remainders. Explicit relabeling/duality witnesses and all compression identities are included. The three bipartite cases can use Corollary~\ref{pgr:r:ir:cor:rolerank}; no higher-rank bipartite theorem is required here.

\subsubsection{Equal signs: the exact moment cone}
By duality suppose both attachments point to exterior heads. Let $A,B$ be the private head-activity sums at $i,j$. For the shared cloud define
\[
 U=\sum_s v_s,\qquad E=\sum_{s<t}v_s v_t.
\]
Put
\[
 F_i=zu_i\Ga_{K-i},\quad F_j=zu_j\Ga_{K-j},
 \quad G=z^2u_i u_j\Ga_{K-\{i,j\}}.
\]
Let $F_*$ be the contribution of supports using one newly adjoined common head, whose activity is set to $1$. This is a genuine support kernel: if either attachment can use the same head, a final support is still counted once. In general $F_*\ne F_i+F_j$.

There are at most two exterior vertices in a support, and direct counting gives
\begin{equation}\label{pgr:r:ir:eq:samemoment}
 \Ga_R=\Ga_K+A F_i+B F_j+U F_*+\bigl[AB+(A+B)U+E\bigr]G.
\end{equation}
Two heads from one private cloud cannot be used. The remaining pairs of types contribute $AB,AU,BU,E$, each at support multiplicity one. With two exterior heads selected, the two attachments are used, leaving exactly the core support in $K-\{i,j\}$.

The moment range satisfies
\[
 U\ge0,\qquad 0\le E\le U^2/2,
\]
because $2E=U^2-\sum_s v_s^2$. The whole closed cone is parameterized by
\begin{equation}\label{pgr:r:ir:eq:cone}
 U=\xi+\eta,\qquad E=2\xi\eta,
 \qquad \xi,\eta\ge0.
\end{equation}
Indeed the inverse values are $(U\pm\sqrt{U^2-2E})/2$. This is a positivity reduction on a cone, not an assertion that two ordinary heads realize every shared cloud. In the formal support generator, two common heads of weights $\xi,\eta$ are used and a support selecting both is assigned coefficient $2$. That formal convention produces exactly \eqref{pgr:r:ir:eq:samemoment} after \eqref{pgr:r:ir:eq:cone}; it is not confused with an ordinary two-clone preorder.

If the common type is not legal, at most one private type is legal: a head's incoming neighborhood must be predecessor-closed in the core, and the union of two such legal neighborhoods is again predecessor-closed. Therefore the private-only cases compress to at most six vertices.

There are $590$ same-sign catalog templates: $333$ private-only reductions and $257$ shared-cone cases. Every shared-cone substituted cubic gap has an exact certificate of the form \eqref{pgr:r:ir:eq:sos}, in the core role activities and $A,B,\xi,\eta$. Together these contain $7\,866$ squares and $32\,729$ positive remainder monomials. The verifier independently constructs $F_*$ from deduplicated ordered supports, checks the symbolic moment identity, tests heterogeneous and zero-face substitutions, and checks every rational identity. A separate local Hall-based reconstruction provides an additional check. No unfinished chamber split or numerical multiplier is needed for these $257$ certificates.

\subsubsection{Degree and zero-activity safeguards}
Both formulas preserve the original weighted polynomial, not merely its first few coefficients. Effective activities can be zero; all target polynomial inequalities hold on the closed orthant. Original degree at most two is handled before applying a cubic certificate, with its factor $4$. At original degree three, equality of support polynomials preserves that degree. Thus the argument never deduces degree-two rank-ULC from a degree-three zero-padding convention. This completes Theorem~\ref{pgr:r:ir:thm:a2}.

\subsection{The Gallai--Edmonds reduction and the global proof}\label{pgr:r:ir:sec:ge}
We now pass from bounded cores to arbitrarily many vertices. The graph-theoretic input is the classical Gallai--Edmonds structure theorem; an accessible primary proof is West \cite[Theorem 5]{West}.

For a graph $G$, let $D$ be the vertices omitted by some maximum matching, let $A=N(D)\setminus D$, and let $C=V(G)\setminus(A\cup D)$. Each component of $G[D]$ is factor-critical, $G[C]$ has a perfect matching, and every maximum matching pairs every vertex of $A$ into a distinct component of $D$, with near-perfect matchings inside those components. Hence
\begin{equation}\label{pgr:r:ir:eq:gebudget}
 \nu(G)=|A|+\frac{|C|}{2}+\sum_{D_t}\frac{|D_t|-1}{2}.
\end{equation}

\subsubsection{The canonical exterior and its size bound}
Assume positive activities and $\nu(G_R)=3$. Let $I$ be the union of the singleton components of $G_R[D]$ and let $K=V\setminus I$. Distinct vertices of $I$ are nonadjacent and all their neighbors lie in $A$. Write $a=|A|$, $|C|=2c$, and $|D_t|=2d_t+1$ for the nonsingleton components, with $d_t\ge1$. Then
\[
 a+c+\sum_t d_t=3,
 \qquad
 |K|=a+2c+\sum_t(2d_t+1)\le a+3\left(c+\sum_t d_t\right)=9-2a.
\]
Thus $a\le3$, with core bounds $7,5,3$ when $a=1,2,3$.

The preorder adds two essential restrictions. Equivalent vertices are adjacent true twins, so their transposition is a graph automorphism. The canonical sets $D,A,C$ are invariant under every automorphism. An exterior vertex in a singleton $D$ component therefore cannot be equivalent to another vertex: such a vertex would also lie in $D$ and be adjacent to it. In particular, exterior-to-core arcs are not bidirectional.

At each attachment, all exterior arcs have one direction. Otherwise $x\to h\to y$ with distinct $x,y\in I$ would make two exterior vertices adjacent by transitivity. The same exterior vertex cannot create the exception because bidirectionality has just been excluded. Each exterior type is thus determined by its nonempty neighborhood subset of $A$ and the fixed signs of those attachments. Isolated exterior points have no effect.

\subsubsection{The four attachment cases}
\paragraph{$a=0$.}
There are no attachments. Every nontrivial component is either a factor-critical $D$ component on at most seven vertices or a component in $C$, whose total order is at most six. Apply Theorem~\ref{pgr:r:ir:thm:finite} to each and Lemma~\ref{pgr:r:ir:lem:products} to their product.

\paragraph{$a=1$.}
All nonisolated exterior vertices are equally oriented pendant twins at the sole attachment $h$. At most one occurs in a support. Sum their relevant activities into $S$. Replacing them by one pendant preserves the support polynomial. Assign the new pendant its relevant role activity $S$ and any positive activity in its unused role. Core activities are unchanged. The full independent-role polynomial is preserved; in the vertex specialization both new role activities may be set to $S$.

If $|K|\le6$, the result has at most seven vertices and Theorem~\ref{pgr:r:ir:thm:finite} applies. If $|K|=7$, equality in the budget bound forces $c=0$ and precisely two nonsingleton factor-critical $D$ components, each with $d_t=1$ and hence three vertices. Deleting $h$ leaves these two rank-one components and the isolated exterior points. Theorem~\ref{pgr:r:ir:thm:articulation} applies directly to every population. If the exterior is empty, the seven-vertex finite theorem is already sufficient.

\paragraph{$a=2$.}
The core has at most five vertices and the canonical exterior has only the three possible nonempty neighborhood types. All hypotheses of Theorem~\ref{pgr:r:ir:thm:a2} hold. Its independent-role conclusion applies directly.

\paragraph{$a=3$.}
Equation~\eqref{pgr:r:ir:eq:gebudget} forces $K=A$, with three attachments and no other core vertices. Every attachment has an exterior neighbor, since $A=N(D)\setminus D$ and now every $D$ component is a singleton. Split $A=P\sqcup Q$, with exterior arcs outward at $P$ and inward at $Q$.

There is no core arc $q\to p$ with $q\in Q$ and $p\in P$. Choose exterior neighbors $x\to q$ and $p\to y$. If $x\ne y$, the path $x\to q\to p\to y$ contradicts independence of $I$. If $x=y$, the directed cycle makes that exterior vertex equivalent to a core vertex, contradicting the singleton-equivalence property above. Every core arc is therefore covered by its tail in $P$ or its head in $Q$: arcs within $P$ are covered on the tail side, arcs within $Q$ on the head side, and $P\to Q$ on both. The exterior arcs are covered too.

The role graph consequently has a cover of size three, and Corollary~\ref{pgr:r:ir:cor:rolerank} proves the result. This gives an ordinary conceptual replacement for following all seventeen three-attachment classes individually. The previously checked catalog of nine same-sign and eight mixed-sign classes remains in the supplementary dependency record.

These cases exhaust the Gallai--Edmonds decomposition. Together with Proposition~\ref{pgr:thm:first} and Lemma~\ref{pgr:r:ir:lem:degree}, they prove Theorem~\ref{pgr:r:ir:thm:main}, including arbitrary vertex count and all zero-activity specializations at surviving degree.

\begin{corollary}[Vertex activities; archive 31]\label{pgr:r:ir:cor:vertex}
Every finite preorder with nonnegative vertex activities $x_i$, that is
$u_i=v_i=x_i$, whose support polynomial has actual degree at most three,
is rank-ULC at that degree.
\end{corollary}
\begin{proof}
The case $u_i=v_i=x_i$ of Theorem~\ref{pgr:r:ir:thm:main}.
\end{proof}
\mnote{this is the main theorem of archive~31, the earlier, vertex-activity
version of this manuscript (same skeleton; eleven proof notes and the
role-cover dependency byte-identical). Archive~31 closed its finite
seven-vertex domain with a selected ledger of 97 direct vertex-only
certificates (100 shipped) after the specialization $u=v$; that ledger is
redundant given the role ledger above and survives only in the arrival
commit (README). Source~23 (Section~\ref{pgr:u:t26:sec}) uses exactly this
corollary. With $x_i=1$ it is source~25's Theorem~\ref{pgr:u:d3:thm:main},
whose proof in Part~I is an independent route.}



\subsection{A cubic with nonreal zeros}\label{pgr:r:ir:sec:nonreal}

Source~22 computes the fifteen-vertex example of
Section~\ref{pgr:u:d3:sec:realroots} (three lower vertices, twelve upper
vertices, four of each pair-neighborhood type; all activities one) with
the same counts, $\Ga(z)=1+24z+162z^2+208z^3$, discriminant $-2592$, gaps
$90$ and $11268$; that computation is printed once, in Part~I.

It is rank-ULC, consistently with the theorem. This example requires no independent-role asymmetry and lies in the pure three-tail family of Section~\ref{pgr:r:ir:sec:rayleigh}. No minimal-size claim is made.

\mnote{source~25's Theorem~\ref{pgr:u:d3:thm:minimum} makes the minimality
claim: fifteen vertices is the least order of a preorder of actual degree
three with a nonreal-rooted gamma polynomial, and this example and its
dual are the only ones at that order.}


\subsection[Verification and dependencies]{Verification package and dependency closure}\label{pgr:r:ir:sec:replay}
The deliverable contains this TeX source and PDF, a portable exact checker entry point, frozen certificate archives, the companion role-cover paper and its checker package, and proof notes recording the precise scopes of structural inputs. All required data are local. The full replay uses Python 3 and a C++17 compiler for supplementary role-cover support checks; it needs no network connection, numerical solver, precompiled binary, or third-party Python library.


\mnote{in this report source~22's package is shipped renamed: its proof
notes at the top of the report directory, its checkers in \texttt{code/},
its nested audit package as
\texttt{data/22-independent-role-checker-role-degree-three-audit-package.zip}.
Its embedded copies of source~20's package and article are not shipped
(source~20's files are). Because \texttt{verify.py} first checks every
listed file, the delivered wrapper runs only on the reconstructed package
(README). At intake its sub-checkers passed on copies: the final Rayleigh
certificate (198 bases, rank~32), the coverage and aggregate phases of the
complete hybrid ledger, the 39 role certificates new relative to
archive~31, the two-attachment and marked-five modules and the side-1565
package; the whole wrapper was not run in one process.}

\subsubsection{How to run the check}
With Python 3 and a C++17 compiler available, from any working directory run
\begin{center}\texttt{python /path/to/package/verify.py}\end{center}
The wrapper verifies pinned dependency archives, extracts them into temporary directories, and runs their exact checkers. It verifies the complete role catalog, its rational certificates and multipliers, every structural witness and refinement terminal, all $223$ six-vertex role certificates and all $24$ total-preorder targets. Both two-attachment certificate modules and the marked-five deletion/gluing module are replayed as well. The HPP modules verify the side-matroid bases, explicit published-list isomorphisms, proper-minor reductions and final all-real Gram identity. The companion role-cover algebra and the nonreal example are checked separately.

The command writes a fresh local replay receipt while preserving the original approval receipts. A separate build script compiles the paper with a standard LaTeX installation; the mathematical checker does not require LaTeX. Python checks use the standard library, and the companion role-cover checker compiles included exact C++ sources into temporary storage.

\subsubsection{What is checked exactly}
The independent checkers reconstruct support polynomials from relations, expand all rational squares, and compare coefficients exactly. They verify template generation, structural witnesses, maps under relabeling and duality, population-compression formulas, moment substitutions, and terminal dependency scope. Finite point evaluations supplement these symbolic checks but do not prove positivity.

Ordinary theorems such as Gallai--Edmonds and Wagner's rank-three Rayleigh theorem are cited mathematical inputs, not re-proved by running a Python script. The support-level structural proofs are given in this paper and the included companion. The certificate and coverage layers are independently audited, with their frozen source hashes recorded. These distinctions make the boundary of the computer-assisted proof explicit.

\subsubsection{Dependency closure and published data}
The selected ledger is complete for independent roles. Each extremal refinement has a proved terminal; positive-multiplier division is confined to the positive orthant and followed by continuity. The final Rayleigh certificate instead has an exact sum of squares valid on all real inputs. These different domains are never interchanged.

The classification results of Kummer--Sert and the certified-positive matroid list are cited mathematical dependencies. The local replay verifies membership in the published list by exact basis isomorphism, with the official file hash pinned. It does not purport to rerun the authors' original Gram certificates for every published entry. The new ten-element case has its own fully local rational Gram and square identity. Ordinary results such as Wagner--Wei's criterion and stability preservation under directional differentiation are proved or cited, not inferred from test points.

Some provenance notes retain draft-status wording from the time their mathematical text was frozen. The approval receipts and complete selected ledger specify the final scope. Numerical discovery code is excluded. The earlier vertex-only release is a separate valid specialization and is not a proof dependency of this stronger theorem.

\subsection{Further questions and mathematical boundaries}\label{pgr:r:ir:sec:questions}
\paragraph{Beyond degree three.}
The independent-role extension in degree three is now proved. Degree four requires two interior coefficient comparisons beyond the universal first inequality and has larger retained cores. An unweighted theorem does not imply its weighted counterpart; separate role-sensitive inequalities or exact certificates are needed. The HPP side-matroid criterion applies in every degree, but it is only a sufficient structural condition, not a general solution.

\paragraph{A conceptual replacement for the finite certificates.}
The Rayleigh identity \eqref{pgr:r:tc:eq:three-squares} and the upper-half-plane articulation proof are short structural explanations. It would be useful to find similarly uniform inequalities for the remaining two-connected factor-critical cores, eliminating some or all of the finite square certificates. The marked five-vertex deletion inequality suggests that interlacing under deletion may be stronger and more useful than ULC alone in certain small cores.

\paragraph{Real-rootedness versus rank-ULC.}
The example in Section~\ref{pgr:r:ir:sec:nonreal} rules out general real-rootedness already in degree three. A useful separate problem is to characterize the preorder or role-cover structures that force real-rootedness for every positive activity assignment. The articulation and two-tail/two-head cover theorems give substantial sufficient families but no general characterization.

\paragraph{Equality and higher degree.}
Classifying equality in the second cubic gap could clarify why weighted positivity is robust in some structures and difficult in others. The square identity identifies several simultaneous equality requirements in the pure three-tail case. Extending arbitrary independent-role-weighted rank-ULC to degree four or higher requires additional inequalities and is not implied by this paper. Even a complete unweighted higher-degree theorem would not justify the weighted extension without a new argument.


\mnote{in this report, unit-activity degree four is settled (Part~I),
real-rootedness is forced by the role-cover and articulation theorems and
by source~15's pendant criterion in their families, and Part~III settles
weighted degree four for four-vertex cores with universal sinks and for
the five-vertex boundary. Weighted degree four in general remains open.}

\paragraph{Attribution.}
The first-gap method is an application of Motzkin--Straus; the bounded-core reduction uses Gallai--Edmonds; the pure three-tail proof applies Wagner's rank-three Rayleigh theorem; the role-cover companion uses classical real-stability preservers; the side-matroid applications use Kummer--Sert's published classification and certified data; and the final ten-element case applies Wagner--Wei's criterion. The new exact artifacts establish the statements and scopes recorded here. We make no claim that a literature-wide novelty search or a proof-assistant formalization has been completed.


\part{Universal-sink clouds over an arbitrary directed core}\label{pgr:part:three}

Every source of this Part works in one setting. A core $C$ of $r$ vertices
carries an arbitrary loopless directed relation $H$, transitive or not; a
finite independent set $W$ of \emph{universal sinks} receives every arc
$C\to W$; there are no other arcs. Core vertices have independent tail and
head activities $u_i,v_i\ge0$, sinks have head activities $w_x\ge0$, and
$e_j$, $E_j$ are the elementary symmetric functions of the core tail and
of the sink head activities. When $H$ is the off-diagonal relation of a
preorder, adjoining the diagonal gives a preorder, so each theorem below
covers an unbounded family of weighted preorders. Reversing every arc and
exchanging the roles gives the universal-source versions.
\begin{itemize}[leftmargin=*]
\item Section~\ref{pgr:s:lg:sec} (source~14): the sharp last-coefficient
  bound $\gamma_{r-1}^2\ge K(r,m)\gamma_{r-2}\gamma_r$ for every
  $r\ge4$, with equality and sharpness; it contains the last inequalities
  of sources 10 and~34.
\item Section~\ref{pgr:s:ad:sec} (source~10): four-vertex cores,
  $\gamma_2^2\ge3\gamma_1\gamma_3$ (computer-assisted) and rank-ULC at
  every actual degree.
\item Section~\ref{pgr:s:fc:sec} (source~34): five-vertex cores, rank-ULC
  at actual degree five, and a counterexample to the constant three.
\item Section~\ref{pgr:s:fb:sec} (source~06): five-vertex cores on the face
  $\gamma_5=0$, rank-ULC at actual degree four.
\item Section~\ref{pgr:s:fq:sec} (source~35, batch~77): five-vertex cores
  on the face $\gamma_4=0$, $\gamma_2^2\ge3\gamma_1\gamma_3$
  (computer-assisted), with the constant three sharp; hence, with sources
  06 and~34, rank-ULC at every actual degree.
\end{itemize}
\mnote{relation to Parts I and~II. A universal-sink cloud is a single
exterior type in the sense of source~16's normal form, so with unit
activities and a preorder core these families are special cases of
Theorem~\ref{pgr:u:d4:thm:middle} (actual degree at most four); with a
preorder core and actual degree at most three they are special cases of
source~22's Theorem~\ref{pgr:r:ir:thm:main}. Neither contains the other in
general: here the core need not be transitive, the activities are
independent, and sources 34 and~06 reach five-vertex cores. For five-vertex
cores actual degree three remains open; the stronger cubic constant three
of source~10 fails there (source~34). The first comparison
$(\rho-1)\gamma_1^2\ge2\rho\gamma_2$ that sources 10, 34 and~06 prove is
Theorem~\ref{pgr:thm:first}. On faces where every vertex used as a tail is
never used as a head, the relation is bipartite and \mrn, Corollary
\texttt{mrn:w:rf:cor:rankfive}, applies. Neither \prp{} nor \mrn{} has
universal-sink material.}
\mnote{Added 2~October 2026, batch~77: source~35
(Section~\ref{pgr:s:fq:sec}) proves actual degree three for five-vertex
cores, so the statement above that it remains open is answered. It proves
the constant three on the face $\gamma_4=0$, and that constant is sharp
there, although it fails at actual degree five (source~34's example has
$\gamma_4>0$). With sources 06 and~34, every five-vertex core with a
universal-sink cloud is rank-ULC at every actual degree
(Corollary~\ref{pgr:s:fq:cor:all}). Source~35's first comparison is again
Theorem~\ref{pgr:thm:first}, and for preorder cores its theorem is a
case of source~22's Theorem~\ref{pgr:r:ir:thm:main}.}

\section{The sharp last gap for every core size (source 14)}\label{pgr:s:lg:sec}

Source~14 is \emph{Sharp final coefficient bounds for directed cores with
universal sinks} (4 pages as delivered; archive
\texttt{universal-sink-last-gap-package}; ``Research note''). Its files are
shipped with the prefix \texttt{14-sink-last-gap-}. It was approved first of
the four-core sequence 14, 13, 10; archive~13 pins it by its SHA-256.

\paragraph{Abstract of source 14, as delivered.}

For a loopless directed core of size $r\ge4$ with independent universal sinks, we prove a sharp inequality for the last three coefficients of its matching-support polynomial, with arbitrary nonnegative independent tail and head activities. The coefficient counts each feasible ordered pair of disjoint endpoint sets once, not its matching witnesses. With $m\ge r$ positive sink-head activities, the sharp constant is $2r^2(m-r+2)/[(r-1)^2(m-r+1)]$. The proof is ordinary mathematics: a support partition, two coefficientwise comparisons, and Newton's inequalities. Exact computational checks are supplied only as corroboration. The result proves a strengthened last inequality in weighted degree four for this family; its middle inequality remains open here.

\subsection{Statement and normalization}
Let $C$ be a set of $r\ge4$ vertices and $H$ any loopless directed relation on $C$. Let $W$ be a disjoint independent set. The full relation $D$ has the arcs of $H$ and all arcs $C\to W$, with no other arcs. In particular, no vertex of $W$ can be a tail. No transitivity assumption is imposed.

Assign tail and head activities $u_i,v_i\ge0$ to each core vertex, and head activities $w_x\ge0$ to the sinks. Sink-tail activities are irrelevant. For a set $S\subseteq C$, write $u_S=\prod_{i\in S}u_i$. Define
\[
 \Ga_D(t)=\sum_{k\ge0}\gamma_k t^k,\qquad
 \gamma_k=\sum_{\substack{S\cap T=\varnothing, |S|=|T|=k\\
                         S\text{ admits a directed perfect matching to }T}}
 u_S\prod_{j\in T\cap C}v_j\prod_{x\in T\cap W}w_x.
\]
The summation is over endpoint pairs, so every feasible pair is counted once. As there are at most $r$ tails, $\deg\Ga_D\le r$.

\begin{theorem}\label{pgr:s:lg:thm:sharp}
Suppose $\gamma_r>0$, and let $m$ be the number of strictly positive sink-head activities. Then $m\ge r$, the actual degree is $r$, and
\begin{equation}\label{pgr:s:lg:eq:sharp}
 \gamma_{r-1}^2\ \ge\ K(r,m)\gamma_{r-2}\gamma_r,
 \qquad K(r,m)=\frac{2r^2(m-r+2)}{(r-1)^2(m-r+1)}.
\end{equation}
For fixed $r,m$, the constant is best possible. Equality holds exactly when all core-tail activities are equal, all positive sink-head activities are equal, and every core head with positive combinatorial indegree has zero head activity.
\end{theorem}

The optimal constant uniformly over all finite sink counts is $2r^2/(r-1)^2$. It is approached as $m\to\infty$ and is strictly smaller than $K(r,m)$ for every finite $m$.

\begin{corollary}\label{pgr:s:lg:cor:four}
If $r=4$ and $\gamma_4>0$, then
\[
 \gamma_3^2\ge \frac{32(m-2)}{9(m-3)}\gamma_2\gamma_4
              >\frac{32}{9}\gamma_2\gamma_4
              >\frac83\gamma_2\gamma_4.
\]
Thus $3\gamma_3^2\ge8\gamma_2\gamma_4$, the last rank-ULC inequality at actual degree four, holds with quantitative slack for this family.
\end{corollary}


\mnote{this corollary is the last degree-four inequality that source~10
(and its earlier version, archive~13) proves again in the weaker form
$\gamma_3^2\ge\frac{32}{9}\gamma_2\gamma_4$ by the same route, with the
finite-$m$ factors of \eqref{pgr:s:lg:eq:L0}--\eqref{pgr:s:lg:eq:L2}
replaced by their limits; source~34's
$\gamma_4^2\ge\frac{25}{8}\gamma_3\gamma_5$ is the case $r=5$ in the same
way. Those proofs are replaced by pointers to this theorem
(Sections~\ref{pgr:s:ad:sec} and~\ref{pgr:s:fc:sec}).}

\subsection{The three top coefficients}
Put $e_j=\ee_j((u_i)_{i\in C})$ and $E_j=\ee_j((w_x)_{x\in W})$, where $\ee_j$ denotes the elementary symmetric polynomial. Set $E_0=1$ and $E_j=0$ beyond the number of sinks. Write
\[
 p=e_r,\qquad q=e_{r-1},\qquad e=e_{r-2}.
\]
Let $N^-(j)$ be the in-neighborhood of $j$ in $H$, and define
\begin{align*}
 d&=\sum_{j:N^-(j)\ne\varnothing}v_j u_{C\setminus\{j\}},\\
 c&=\sum_{j\in C}v_j
       \sum_{\substack{S\subseteq C\setminus\{j\},\ |S|=r-2\\
                       S\cap N^-(j)\ne\varnothing}}u_S,\\
 b&=\sum_{\{j,k\}\text{ admissible}}v_jv_k u_{C\setminus\{j,k\}}.
\end{align*}
An unordered pair $\{j,k\}$ is admissible if distinct vertices $i,\ell\in C\setminus\{j,k\}$ satisfy $i\to j$ and $\ell\to k$. Each admissible pair contributes once, even if several such injections exist.

\begin{lemma}\label{pgr:s:lg:lem:coeff}
The exact support identities are
\begin{align}
 \gamma_r&=pE_r,\label{pgr:s:lg:eq:top}\\
 \gamma_{r-1}&=qE_{r-1}+dE_{r-2},\label{pgr:s:lg:eq:next}\\
 \gamma_{r-2}&=eE_{r-2}+cE_{r-3}+bE_{r-4}.\label{pgr:s:lg:eq:third}
\end{align}
\end{lemma}
\begin{proof}
A size-$k$ support has $k$ core tails and $h$ core heads. Physical disjointness gives $k+h\le r$, and it has $k-h$ sink heads. Thus the three top ranks permit at most zero, one, and two core heads, respectively. For one core head, the selected tails must contain an in-neighbor. For two, distinct selected tails must cover the two heads. Once these heads have been covered, all remaining selected tails can be matched to the selected universal sinks. Feasibility is Boolean; no choice of witness enters the coefficient.
\end{proof}

\begin{lemma}\label{pgr:s:lg:lem:core}
For all nonnegative activities,
\begin{equation}\label{pgr:s:lg:eq:core}
 q^2\ge\frac{2r}{r-1}pe,\qquad qd\ge pc,\qquad d^2\ge2pb.
\end{equation}
The last two comparisons hold coefficientwise.
\end{lemma}
\begin{proof}
The first is Newton's inequality in exactly $r$ variables. For a division-free proof of the second, a summand of $c$ has the form $v_j u_{C\setminus\{i,j\}}$, with $i\ne j$, and implies $N^-(j)\ne\varnothing$. The summands $u_{C\setminus\{i\}}$ in $q$ and $v_j u_{C\setminus\{j\}}$ in $d$ have product
\[
 u_{C\setminus\{i\}}v_j u_{C\setminus\{j\}}
   =p v_j u_{C\setminus\{i,j\}}.
\]
Distinct choices of the omitted vertex $i$ and the head $j$ are distinct product contributions, so these terms cover $pc$.

For every admissible head pair $\{j,k\}$, both heads have positive combinatorial indegree. The two cross terms of $d^2$ give
\[
 2v_jv_k u_{C\setminus\{j\}}u_{C\setminus\{k\}}
   =2p v_jv_k u_{C\setminus\{j,k\}}.
\]
All other terms are nonnegative. Both arguments allow zero activities directly.
\end{proof}

\subsection{Proof of the sharp bound}
Delete sinks of zero head activity. By \eqref{pgr:s:lg:eq:top}, the hypothesis $\gamma_r>0$ implies that all core-tail activities are positive and that $m\ge r$. Write $s=m-r\ge0$.

Newton's inequalities say that $A_j=E_j/\binom mj$ is log-concave. Comparing pairs of indices with the same sum gives
\begin{align}
 E_{r-1}^2&\ge L_0 E_{r-2}E_r,
 &L_0&=\frac{r(s+2)}{(r-1)(s+1)},\label{pgr:s:lg:eq:L0}\\
 E_{r-1}E_{r-2}&\ge L_1 E_{r-3}E_r,
 &L_1&=\frac{r(s+3)}{(r-2)(s+1)},\label{pgr:s:lg:eq:L1}\\
 E_{r-2}^2&\ge L_2 E_{r-4}E_r,
 &L_2&=\frac{r(r-1)(s+4)(s+3)}{(r-2)(r-3)(s+2)(s+1)}.\label{pgr:s:lg:eq:L2}
\end{align}
For example, $A_{r-2}A_{r-1}\ge A_{r-3}A_r$ follows from the decreasing consecutive ratios of a positive log-concave sequence. The other two comparisons follow in the same way. The constants above are the respective binomial ratios.

Let $K=K(r,m)=(2r/(r-1))L_0$. Combining Lemma~\ref{pgr:s:lg:lem:core} with \eqref{pgr:s:lg:eq:L0}--\eqref{pgr:s:lg:eq:L2} yields
\begin{align*}
 q^2E_{r-1}^2&\ge KpeE_{r-2}E_r,\\
 2qdE_{r-1}E_{r-2}&\ge 2L_1pcE_{r-3}E_r,\\
 d^2E_{r-2}^2&\ge 2L_2pbE_{r-4}E_r.
\end{align*}
Both latter constants exceed $K$, because
\[
 \frac{2L_1}{K}=\frac{(r-1)^2}{r(r-2)}\frac{s+3}{s+2}>1,
 \qquad
 \frac{2L_2}{K}=\frac{(r-1)^3}{r(r-2)(r-3)}
                     \frac{(s+4)(s+3)}{(s+2)^2}>1.
\]
The first core factor exceeds one by $(r-1)^2-r(r-2)=1$; for the second,
\[
 (r-1)^3-r(r-2)(r-3)=2r^2-3r-1>0\qquad(r\ge4).
\]
Adding the three termwise inequalities and applying \eqref{pgr:s:lg:eq:next}--\eqref{pgr:s:lg:eq:third} proves \eqref{pgr:s:lg:eq:sharp}.

\paragraph{Equality and sharpness.}
With positive core tails, $c=0$ if and only if $d=0$: every head of positive indegree has an $(r-2)$-element tail set containing an in-neighbor. Since $E_{r-3}>0$ and $2L_1>K$, equality in \eqref{pgr:s:lg:eq:sharp} forces $c=d=0$, and then $b=0$. The two remaining Newton inequalities attain equality precisely when all positive variables in their respective families are equal. Thus the equality conditions in Theorem~\ref{pgr:s:lg:thm:sharp} are necessary and sufficient.

For an edgeless core, equal core-tail activities $u>0$, and $m$ equal sink-head activities $w>0$, one has
\[
 \gamma_k=\binom rk\binom mk(uw)^k.
\]
Its last coefficient ratio is exactly $K(r,m)$. This proves sharpness for each finite $m$. Letting $m\to\infty$ proves that $2r^2/(r-1)^2$ is the optimal uniform constant.

\subsection{Scope and exact checks}
If $\gamma_r=0$, the ordinary comparison
\[
 (r-1)\gamma_{r-1}^2\ge2r\gamma_{r-2}\gamma_r
\]
is immediate, but it does not establish the last Newton inequality at a smaller actual degree. Zero core-head activities are fully permitted in the sharp theorem. Zero sink-head activities are removed before defining the effective count $m$. Reversing all arcs and interchanging tail and head activities gives the corresponding universal-source theorem.

For a preorder core, adjoining universal sinks again gives a preorder, so Corollary~\ref{pgr:s:lg:cor:four} supplies an unbounded weighted degree-four subclass. It does not settle
\[
 4\gamma_2^2\ge9\gamma_1\gamma_3,
\]
nor general independent-role rank-ULC in degree four.
\mnote{the middle inequality is settled for four-vertex cores by source~10
(Theorem~\ref{pgr:s:ad:thm:main}, in the stronger form
$\gamma_2^2\ge3\gamma_1\gamma_3$) and for five-vertex cores at actual
degrees five and four by sources 34 and~06; general independent-role
rank-ULC in degree four remains open.}
 No stability, real-rootedness, formalization, or priority claim is made.

The theorem above has an ordinary proof. The accompanying standard-library Python checkers are corroborating tests, not logical dependencies. The producer checker verifies full supports by Hall's condition in 8,312 exact weighted cases, including all 4,096 loopless four-core relations, and separately tests the general formulas and 27 sharp equality examples. An independent checker imports no producer code: its 13,132 weighted instances cover 4,518 configurations and 5,088,496 feasible endpoint pairs, with 10,001 finite-$m$ sharp-bound checks. Both include zero activities and smaller surviving degrees. The supplied driver replays both checkers in fresh temporary directories and reports newly generated receipts; hashes check integrity only.

Newton's inequalities, including their normalization, are classical. A modern primary reference is Br\"and\'en and Huh~\cite{BH}: Proposition~2.2 and Example~2.26 recover the normalized inequalities from the stable bivariate product $\prod_i(z+u_it)$. The support partition and the two coefficientwise comparisons are proved directly here.


\section{Four-vertex cores at every actual degree (source 10)}\label{pgr:s:ad:sec}

Source~10 is \emph{Rank ultra log concavity for four vertex cores with
universal sinks} (5 pages as delivered; archive
\texttt{universal-sink-all-degrees-package}; ``Research note''). Its files
are shipped with the prefix \texttt{10-sink-all-degrees-}. It revises
archive~13 (see the remark after Proposition~\ref{pgr:s:ad:prop:cert}).

\paragraph{Abstract of source 10, as delivered.}

We prove rank ultra log concavity at every actual surviving degree for the matching-support polynomial of any loopless directed four-vertex core with an arbitrary independent set of universal sinks, under independent nonnegative tail and head activities. No transitivity is required, and zero activities are allowed. Each feasible ordered pair of disjoint endpoint sets is counted once, irrespective of its number of matching witnesses. A sharp two-parameter reduction of the sink moments leads to exact rational certificates for the stronger cubic coefficient inequality $\gamma_2^2\ge3\gamma_1\gamma_3$ in all 218 directed core classes. Together with the actual-degree first-coefficient bound and an ordinary last-coefficient argument, this closes degrees zero through four without padding the normalization. The stronger cubic inequality is computer-assisted; all structural reductions and the other inequalities have ordinary proofs.

\subsection{The result and the support convention}
Let $C=\{0,1,2,3\}$ carry an arbitrary loopless directed relation $H$. Let $W$ be a finite independent set disjoint from $C$. The full relation $D$ has the arcs of $H$ and all arcs $C\to W$, with no other arcs. Give each core vertex independent nonnegative tail and head activities $u_i,v_i$, and give each sink head activity $w_x\ge0$. Sink-tail activities never enter.

For a set $S\subseteq C$, write $u_S=\prod_{i\in S}u_i$. Define
\[
 \Ga_D(t)=\sum_{k=0}^4\gamma_k t^k,\qquad
 \gamma_k=\sum_{\substack{S\cap T=\varnothing, |S|=|T|=k\\
                         S\text{ admits a directed perfect matching to }T}}
 u_S\prod_{j\in T\cap C}v_j\prod_{x\in T\cap W}w_x.
\]
There are only four possible tails, so the degree is at most four. The empty support gives $\gamma_0=1$.

\begin{theorem}\label{pgr:s:ad:thm:main}
For every such $D$ and every nonnegative independent activity assignment,
\begin{equation}\label{pgr:s:ad:eq:middle}
 \gamma_2^2\ge3\gamma_1\gamma_3.
\end{equation}
If $\rho=\deg\Ga_D$, then $0\le\rho\le4$ and the sequence
\[
 \left(\frac{\gamma_k}{\binom{\rho}k}\right)_{k=0}^{\rho}
\]
is log-concave. The normalization uses the actual surviving degree, including after activities vanish. In degree four, the last comparison admits the stronger bound $\gamma_3^2\ge(32/9)\gamma_2\gamma_4$.
\end{theorem}

The cubic inequality uses a finite exact certificate collection supplied with this note. The actual-degree first inequality and the quartic last inequality have ordinary proofs below. Reversing every arc and exchanging the two activity roles gives the universal-source version.

\subsection{Boolean support coefficients}
Write $e_j=\ee_j(u_0,u_1,u_2,u_3)$ and $E_j=\ee_j((w_x)_{x\in W})$, with $E_0=1$ and $E_j=0$ beyond the sink count. Let $N^-(j)$ be the in-neighborhood of $j$ within $H$. Put
\begin{align*}
 a&=\sum_{i\to j\text{ in }H}u_i v_j,\\
 b&=\sum_{\{j,k\}\text{ admissible}}v_jv_k u_{C\setminus\{j,k\}},\\
 c&=\sum_{j\in C}v_j
       \sum_{\substack{S\subseteq C\setminus\{j\},\ |S|=2\\
                       S\cap N^-(j)\ne\varnothing}}u_S,\\
 d&=\sum_{j:N^-(j)\ne\varnothing}v_j u_{C\setminus\{j\}}.
\end{align*}
An unordered pair $\{j,k\}$ is admissible when the two remaining distinct core vertices can cover $j,k$ by directed arcs. Such a pair contributes once to $b$, even if both perfect matchings exist.

\begin{lemma}\label{pgr:s:ad:lem:coeff}
The exact coefficient formulas are
\begin{equation}\label{pgr:s:ad:eq:coeff}
\begin{aligned}
 \gamma_1&=a+e_1E_1,&\qquad \gamma_2&=b+cE_1+e_2E_2,\\
 \gamma_3&=dE_2+e_3E_3,& \gamma_4&=e_4E_4.
\end{aligned}
\end{equation}
\end{lemma}
\begin{proof}
A size-$k$ support has $k$ tails in $C$ and at most $4-k$ core heads. Any injection of distinct selected tails onto the selected core heads extends to the selected sink heads, because every remaining tail points to every sink. Size two therefore permits zero, one, or two core heads, giving its three summands. Size three permits zero or one, and size four permits none. Size one is immediate. All internal tests are Boolean existence tests, not counts of injections.
\end{proof}

\subsection{An exact reduction to two sink parameters}
\begin{lemma}\label{pgr:s:ad:lem:moments}
For every finite nonnegative sink-head vector there are $A,B\ge0$ such that
\begin{equation}\label{pgr:s:ad:eq:envelope}
 E_1=A+B,\qquad E_2=AB+\frac{B^2}{2},\qquad
 E_3\le\frac{AB^2}{2}+\frac{B^3}{6}.
\end{equation}
The displayed triples are sharp after taking the closure over all finite sink counts.
\end{lemma}
\begin{proof}
Set $p_j=\sum_x w_x^j$, $A=\sqrt{p_2}$, and $B=p_1-\sqrt{p_2}$. Since $p_1^2\ge p_2$, both parameters are nonnegative. The elementary identities
\[
 E_1=p_1,\qquad 2E_2=p_1^2-p_2,\qquad
 6E_3=p_1^3-3p_1p_2+2p_3
\]
give the first two equations. Every $w_x\le\sqrt{p_2}=A$, so $p_3\le A p_2=A^3$. Substituting $p_1=A+B$ proves the third. The zero sink vector is included.

Conversely, one sink of activity $A$ and $n$ sinks of activity $B/n$ have $p_1=A+B$, $p_2=A^2+B^2/n$, and $p_3=A^3+B^3/n^2$. Their first three elementary symmetric functions tend to the boundary values in \eqref{pgr:s:ad:eq:envelope} as $n\to\infty$. These are limiting triples, not generally exact finite realizations.
\end{proof}


\mnote{this moment envelope is printed once, here. Source~34 restates it
with the same proof, and source~06 in the scaled form
$E_1=A+B$, $2E_2=2AB+B^2$, $6E_3\le3AB^2+B^3$; both now refer to this
lemma.}

With the core activities and $E_1,E_2$ fixed, the coefficient of $E_3$ in $\gamma_2^2-3\gamma_1\gamma_3$ is
\[
 -3(a+e_1E_1)e_3\le0.
\]
Thus replacing $E_3$ by its upper bound in \eqref{pgr:s:ad:eq:envelope} can only decrease the cubic gap. Define the integer polynomials
\begin{align}
 G_1&=a+e_1(A+B),\notag\\
 G_2&=2b+2c(A+B)+e_2(2AB+B^2),\label{pgr:s:ad:eq:G}\\
 G_3&=3d(2AB+B^2)+e_3(3AB^2+B^3),\notag\\
 Q_H&=G_2^2-2G_1G_3.\label{pgr:s:ad:eq:P}
\end{align}
At the boundary substitution, $G_1=\gamma_1$, $G_2=2\gamma_2$, and $G_3=6\gamma_3$. Hence $Q_H$ is four times the boundary cubic gap. Proving $Q_H\ge0$ for all ten nonnegative variables suffices. By the limiting construction, this positivity problem is also equivalent after closure to the cubic inequality for all finite clouds.

\subsection{Complete exact certificates for the reduced polynomials}
There are $2^{12}=4096$ loopless labeled directed relations on four vertices. Applying all $24$ vertex permutations partitions them into $218$ isomorphism classes. Relabeling a core simultaneously permutes its tail and head activities, so one representative per class suffices. The full directed domain includes all $33$ unlabeled preorder cores; no transitivity is used to discard any case.

Use the variable order
\[
 x=(u_0,u_1,u_2,u_3,v_0,v_1,v_2,v_3,A,B).
\]
For each representative $H$, the supplied certificate is an exact identity
\begin{equation}\label{pgr:s:ad:eq:sos}
 Q_H(x)=\sum_{\ell}q_\ell x^{m_\ell}
             \left(\sum_{a\in\mathcal A_\ell}c_{\ell,a}x^a\right)^2+R_H(x),
\end{equation}
where $q_\ell>0$, all scalar coefficients are rational, all exponent vectors have nonnegative integer coordinates, and every coefficient of $R_H$ is nonnegative. Thus each identity proves nonnegativity on the closed nonnegative orthant.

\begin{proposition}[Finite exact verification]\label{pgr:s:ad:prop:cert}
The certificate collection covers all $218$ permutation classes. Exact expansion of every identity \eqref{pgr:s:ad:eq:sos} gives a coefficientwise nonnegative remainder. There are $3532$ weighted polynomial-square terms and $64298$ positive remainder monomials in total. Of the $218$ certificates, $127$ use only binomial squares and $91$ also use larger square polynomials. Each squared polynomial has between two and twenty-four monomials.
\end{proposition}
\begin{proof}[Verification]
The independent standard-library checker enumerates all $4096$ row-mask relations and their permutation orbits. For every labeled core, it reconstructs the grouped support polynomial directly from ordered disjoint endpoint sets and Boolean matching existence, checking ranks zero through four against Lemma~\ref{pgr:s:ad:lem:coeff}. It then makes the boundary substitution and forms \eqref{pgr:s:ad:eq:P} independently of the producer code.

For each representative it reads every rational term, checks its signs and exponent domains, expands the squares exactly, and subtracts the complete expansion from $Q_H$. Every resulting coefficient, including coefficients at exponents absent from the original target, is checked to be nonnegative. The remainder is reconstructed rather than trusted as an input. The same run checks that every core orbit is represented exactly once. The accompanying checker and certificate files supply the reproducible finite calculation.
\end{proof}


\mnote{archive~13, \emph{Weighted degree four rank ultra log concavity for
universal sink extensions}, is the earlier version of this note: the same
218 core representatives in the same order, the same reduction and
coefficient lemma, and the same last-gap proof, but the weaker middle
inequality $4\gamma_2^2\ge9\gamma_1\gamma_3$ (target $2G_2^2-3G_1G_3$),
certified by 2,830 binomial squares with 75,413 positive remainder
monomials (45 classes coefficientwise nonnegative), and rank-ULC only at
actual degree four. Source~10's inequality implies archive~13's on every
face. Archive~13 also proves the degree-four first gap by grouping tails
($\gamma_2\le e_2(R_0,\ldots,R_3)\le\frac38\gamma_1^2$ with
$R_i=u_i\cdot$(head mass out of~$i$)), an argument valid only for that
constant. Its certificates survive in the arrival commit (README).}

The discovery program used floating-point linear programming to propose the square weights. Those weights were reconstructed as rational numbers, using exact active-system solutions where needed, and accepted only after exact expansion. No numerical optimization result, sampled activity assignment, or numerical tolerance is part of the verification. Combining Proposition~\ref{pgr:s:ad:prop:cert} with Lemma~\ref{pgr:s:ad:lem:moments} proves \eqref{pgr:s:ad:eq:middle} for every finite weighted sink cloud.

\subsection{Ordinary proofs at the actual degree}

\mnote{source~10 first proves the actual-degree inequality
$(\rho-1)\gamma_1^2\ge2\rho\gamma_2$ for every finite loopless directed
relation with nonnegative independent role activities, by Motzkin--Straus
on the graph of positive-weight arcs (removing arcs of zero product
activity but retaining their physical vertices in every remaining role)
with the mass-shifting proof of the bound: this is
Theorem~\ref{pgr:thm:first}, \eqref{pgr:eq:first}. It then proves the last
inequality $\gamma_3^2\ge\frac{32}{9}\gamma_2\gamma_4$ from the three core
comparisons $q^2\ge\frac83pe_2$, $qd\ge pc$, $d^2\ge2pb$ (with $p=e_4$,
$q=e_3$) and the sink inequalities $E_3^2\ge\frac43E_2E_4$,
$E_2E_3\ge2E_1E_4$, $E_2^2\ge6E_4$: this is source~14's proof at $r=4$ with
the finite-$m$ factors dropped, and the inequality follows from source~14's
Corollary~\ref{pgr:s:lg:cor:four} when $\gamma_4>0$ (it is trivial when
$\gamma_4=0$). Both proofs are replaced by these pointers; ``the
last-coefficient estimate'' below is that corollary.}

We now use the correct normalization in each surviving degree:
\begin{itemize}
\item If $\rho=0$ or $\rho=1$, rank ultra log concavity has no nontrivial inequalities.
\item If $\rho=2$, \eqref{pgr:eq:first} gives $\gamma_1^2\ge4\gamma_2$.
\item If $\rho=3$, \eqref{pgr:eq:first} gives $\gamma_1^2\ge3\gamma_2$, and \eqref{pgr:s:ad:eq:middle} gives $\gamma_2^2\ge3\gamma_1\gamma_3$.
\item If $\rho=4$, \eqref{pgr:eq:first} gives $3\gamma_1^2\ge8\gamma_2$; the stronger cubic bound \eqref{pgr:s:ad:eq:middle} implies $4\gamma_2^2\ge9\gamma_1\gamma_3$; and the last-coefficient estimate implies $3\gamma_3^2\ge8\gamma_2\gamma_4$.
\end{itemize}
Submatchings of a positive-weight top support ensure positivity of every earlier coefficient. These degree-specific inequalities prove Theorem~\ref{pgr:s:ad:thm:main}, including all zero-activity faces.

\subsection{Reproduction and scope}
The package includes all $218$ rational certificates, the complete core catalog, the discovery scripts, the independent verifier, and the separate ordinary-reduction and algebraic audits. Run
\begin{center}\texttt{python verify.py}\end{center}
with Python~3.10 or later. Verification uses only the standard library, imports no producer code, and needs neither a numerical solver nor any previous research archive. The driver checks file integrity and runs the independent certificate checker afresh, producing a new receipt. Integrity checking alone is not mathematical verification.

The finite algebraic verification is a logical part of the stronger cubic-inequality proof. In contrast, the moment reduction, the support formulas, and the first and last inequalities are ordinary arguments.

All activities may vanish. The certified cubic inequality holds directly on every nonnegative face, and \eqref{pgr:eq:first} is applied at the resulting actual degree. A vertex with zero tail activity can still be used as a head; it is never deleted merely because one role is inactive. No padded quartic inequality is used to infer a cubic one. No earlier theorem about three active tail roles is needed as a dependency.


\mnote{for a preorder core at actual degree at most three this is also a
case of source~22's Theorem~\ref{pgr:r:ir:thm:main}, and with unit
activities of source~16's Theorem~\ref{pgr:u:d4:thm:middle}; source~10
needs neither. Source~10's proof note \texttt{ACTUAL\_DEGREE\_CLOSURE.md}
likewise says it closes the degree drop without the earlier
role-matching-rank-three theorem (source~22's corollary).}

When $H$ is the off-diagonal relation of a preorder, adjoining the diagonal to $D$ yields a preorder. Thus the theorem supplies an unbounded weighted preorder subclass at every surviving degree. It does not concern non-universal exterior neighborhoods or larger cores, and therefore does not resolve general weighted degree-four preorders. No real-rootedness, stability, formalization, external publication, or novelty/priority claim is made.

The remaining classical input is Newton's inequality for elementary symmetric functions. As a modern primary reference, Proposition~2.2 and Example~2.26 of Br\"and\'en and Huh~\cite{BH} recover the normalized inequalities from stable products of linear factors. The moment estimate itself is proved above directly from power sums.


\section{Five-vertex cores at actual degree five (source 34)}\label{pgr:s:fc:sec}

Source~34 is \emph{Rank five ultra log concavity for five vertex cores with
universal sinks} (6 pages as delivered; archive
\texttt{universal-sink-five-core-package}; ``Research note''). Its files are
shipped with the prefix \texttt{34-five-core-}; its 9,610 certificate files
are shipped in the derived container
\texttt{data/34-five-core-certificate-data.tar.xz}, whose members are
byte-identical to the delivered ones (README).

\paragraph{Abstract of source 34, as delivered.}

We prove rank ultra log concavity at actual degree five for every loopless directed five-vertex core with an arbitrary finite independent set of universal sinks, under independent nonnegative tail and head activities. Each feasible ordered pair of disjoint endpoint sets is counted once, regardless of its matching witnesses. A sharp two-parameter bound on the first three sink moments reduces the first interior coefficient inequality to exact rational certificates for all 9,608 directed core classes. The second interior inequality has an ordinary coefficientwise proof, as do the first and last comparisons. Transitivity is not required. The theorem does not infer the stronger normalizations needed on lower-degree faces. We also give an exact preorder example showing that the stronger cubic constant proved for four-vertex cores fails for five-vertex cores.

\subsection{The theorem and the counting convention}
Let $C=\{0,1,2,3,4\}$ carry an arbitrary loopless directed relation $H$. Let $W$ be a finite independent set disjoint from $C$. The relation $D$ consists of $H$ and every arc $C\to W$, with no other arcs. Core tail activities $u_i$, core head activities $v_i$, and sink head activities $w_x$ are independent and nonnegative. Sink-tail activities never enter. Put
\[
 \Gamma_D(t)=\sum_{k=0}^5\gamma_k t^k,\qquad
 \gamma_k=\sum_{\substack{S\cap T=\varnothing,\ |S|=|T|=k\\
 S\text{ admits a directed perfect matching to }T}}
 \prod_{i\in S}u_i\prod_{j\in T\cap C}v_j\prod_{x\in T\cap W}w_x.
\]
Here every tail lies in $C$, the empty support has weight $\gamma_0=1$, and each ordered endpoint pair $(S,T)$ contributes once.

\begin{theorem}\label{pgr:s:fc:thm:main}
For every such core and every nonnegative activity assignment,
\begin{equation}\label{pgr:s:fc:eq:interior}
 \gamma_2^2\ge2\gamma_1\gamma_3,
 \qquad \gamma_3^2\ge2\gamma_2\gamma_4.
\end{equation}
If the actual surviving degree is five, then
$\bigl(\gamma_k/\binom5k\bigr)_{k=0}^5$ is log-concave. Moreover,
\begin{equation}\label{pgr:s:fc:eq:last}
 \gamma_4^2\ge\frac{25}{8}\gamma_3\gamma_5.
\end{equation}
The first inequality in \eqref{pgr:s:fc:eq:interior} is computer-assisted by an exact finite certificate collection. The second, the first-coefficient bound, and \eqref{pgr:s:fc:eq:last} have ordinary proofs.
\end{theorem}

All five core-tail activities must be positive and at least five sink-head activities must be positive for the actual degree to be five; these conditions are also sufficient. Core-head activities may vanish. Reversing all arcs and interchanging the activity roles gives the corresponding universal-source theorem. When $H$ is the off-diagonal relation of a preorder, adjoining the diagonal to $D$ gives a preorder, so this is an unbounded weighted preorder subclass as well.

\subsection{Exact Boolean coefficients}
Write $u_S=\prod_{i\in S}u_i$, $v_J=\prod_{j\in J}v_j$, and
$e_r=\ee_r((u_i)_{i\in C})$, $E_r=\ee_r((w_x)_{x\in W})$.
Define $F_{kr}$ to be the sum of $u_Sv_J$ over disjoint core sets $S,J$ of sizes $k,r$ for which distinct selected tails can cover all selected heads $J$. This is an existence test, not a count of injections. Put
\[
 a=F_{11},\quad b=F_{22},\quad c=F_{21},\quad
 d=F_{31},\quad f=F_{32},\quad g=F_{41}.
\]
A support of size $k$ has at most $5-k$ core heads. Any injection covering those heads extends to the selected sink heads. Therefore
\begin{equation}\label{pgr:s:fc:eq:coeff}
\begin{aligned}
 \gamma_1&=a+e_1E_1,& \gamma_2&=b+cE_1+e_2E_2,\\
 \gamma_3&=fE_1+dE_2+e_3E_3,&
 \gamma_4&=gE_3+e_4E_4,\qquad \gamma_5=e_5E_5.
\end{aligned}
\end{equation}
In particular $\gamma_5>0$ has exactly the characterization stated above.

For later use write $c=\sum_jv_jc_j$, $d=\sum_jv_jd_j$, and $g=\sum_jv_jg_j$.
If $N_j$ is the nonempty core in-neighborhood of $j$, then
\[
 c_j=\sum_{\substack{S\subset C\setminus j,\ |S|=2\\S\cap N_j\ne\varnothing}}u_S,
 \quad d_j=\sum_{\substack{S\subset C\setminus j,\ |S|=3\\S\cap N_j\ne\varnothing}}u_S,
 \quad g_j=u_{C\setminus j};
\]
all three are zero for an empty neighborhood. Similarly write
$b=\sum_{j<k}v_jv_kb_{jk}$ and $f=\sum_{j<k}v_jv_kf_{jk}$.
The polynomial $f_{jk}$ is either zero or $u_{C\setminus\{j,k\}}$; it is nonzero exactly when $b_{jk}$ is nonzero.

\subsection{The first interior inequality by exact certificates}

For any nonnegative finite sink vector, let $p_r=\sum_xw_x^r$,
$A=\sqrt{p_2}$, and $B=p_1-\sqrt{p_2}$. By Lemma~\ref{pgr:s:ad:lem:moments}
of source~10, $A,B\ge0$ and the moment envelope \eqref{pgr:s:ad:eq:envelope}
holds, sharp after closure over finite sink counts.
\mnote{source~34 restates the envelope with the same proof (each
$w_x\le\sqrt{p_2}$, so $p_3\le p_2^{3/2}$; $6E_3=p_1^3-3p_1p_2+2p_3$ and
$2E_2=p_1^2-p_2$) and the same limiting construction, one sink of activity
$A$ and $n$ sinks of activity $B/n$; it is printed once, in
Section~\ref{pgr:s:ad:sec}.}


With $E_1,E_2$ and the core activities fixed, the coefficient of $E_3$ in
$\gamma_2^2-2\gamma_1\gamma_3$ is $-2(a+e_1E_1)e_3\le0$.
It therefore suffices to prove nonnegativity after substituting the upper bound in \eqref{pgr:s:ad:eq:envelope}. Define
\begin{align*}
 G_1&=a+e_1(A+B),\\
 G_2&=2b+2c(A+B)+e_2(2AB+B^2),\\
 G_3&=6f(A+B)+3d(2AB+B^2)+e_3(3AB^2+B^3),\\
 Q_H&=3G_2^2-4G_1G_3.
\end{align*}
At this substitution $G_1=\gamma_1,G_2=2\gamma_2,G_3=6\gamma_3$, so
\begin{equation}\label{pgr:s:fc:eq:scale}
 Q_H=12(\gamma_2^2-2\gamma_1\gamma_3).
\end{equation}
By the limiting construction, nonnegativity of $Q_H$ is also necessary for this inequality to hold for all finite clouds with the fixed core.

There are $2^{20}=1,048,576$ labeled loopless directed relations on five vertices. Simultaneous relabeling of core tails and heads partitions them into $9,608$ classes, including all $139$ preorder classes. In the variable order
\[
 x=(u_0,u_1,u_2,u_3,u_4,v_0,v_1,v_2,v_3,v_4,A,B),
\]
the supplied certificate for each class is an exact identity
\begin{equation}\label{pgr:s:fc:eq:certificate}
 Q_H(x)=\sum_\ell q_\ell x^{m_\ell}
              (\alpha_\ell x^{r_\ell}+\beta_\ell x^{s_\ell})^2+R_H(x),
\end{equation}
where $q_\ell>0$, all scalar coefficients are rational, all exponents are nonnegative integers, and every coefficient of $R_H$ is nonnegative.

\begin{proposition}[Finite exact verification]\label{pgr:s:fc:prop:cert}
All $9,608$ identities \eqref{pgr:s:fc:eq:certificate} have coefficientwise nonnegative remainders. Their permutation orbits cover all labeled relations exactly once. The collection contains $57,089$ weighted binomial squares; $2,391$ certificates need no square terms.
\end{proposition}
\begin{proof}[Verification]
The independent Python checker reconstructs the Boolean support kernels, forms $Q_H$, expands each rational square, and subtracts it exactly. It checks every remainder coefficient, including at exponents absent from the target. It also checks complete labeled-core coverage under all $120$ vertex permutations. No producer import, numerical optimization tolerance, or sampled activity assignment is used to establish any identity. The certificate data and checker supply the reproducible finite calculation.
\end{proof}


\mnote{dated note, 1~October 2026 (batch-73 intake): the shipped note
\nolinkurl{34-five-core-ORDINARY_PROOFS.md} says at its line~13 ``No claim
that this reduced polynomial is always nonnegative is made by this note''.
That sentence describes the note's own scope before the certificates; the
reduced polynomial $Q_H$ is proved nonnegative by
Proposition~\ref{pgr:s:fc:prop:cert}. The delivered note is kept
unchanged.}

Floating-point linear programming was used only to discover candidates. Exact rational reconstruction and expansion determine acceptance. Proposition~\ref{pgr:s:fc:prop:cert}, \eqref{pgr:s:fc:eq:scale}, and the monotonicity direction above prove the first inequality in \eqref{pgr:s:fc:eq:interior}, including every nonnegative activity face.

\subsection{An ordinary proof of the second interior inequality}
The following comparisons hold coefficientwise in all core activities:
\begin{equation}\label{pgr:s:fc:eq:corelemmas}
 3fd\ge bg,\qquad
 4fe_3+3d^2\ge be_4+4cg,\qquad
 2de_3\ge2e_2g+ce_4.
\end{equation}
Complete ordinary proofs are given in the appendix.
\mnote{the three core comparisons are proved without enumeration in
source~34's appendix; source~06's appendix proves the same comparisons
with the middle one strengthened to $4fe_3+2d^2\ge be_4+4cg$, which
implies the form used here because $d^2\ge0$. The appendix is printed
once, in source~06's stronger form (Section~\ref{pgr:s:fb:app:core}).}
 For the sink moments, Newton's normalized inequalities and direct elementary expansions give
\begin{equation}\label{pgr:s:fc:eq:sinklemmas}
\begin{gathered}
 E_1E_2\ge3E_3,\quad E_2^2\ge\tfrac32E_1E_3,\quad
 E_1E_3\ge4E_4,\\
 E_2E_3\ge2E_1E_4,\qquad E_3^2\ge\tfrac43E_2E_4.
\end{gathered}
\end{equation}
The finite-variable factors in Newton's inequalities only strengthen these bounds. Vanishing-moment cases follow directly. Newton for five tail variables gives $e_3^2\ge2e_2e_4$.

Expand the square using \eqref{pgr:s:fc:eq:coeff}. Its terms obey
\begin{align*}
 2fdE_1E_2&\ge6fdE_3\ge2bgE_3,\\
 2fe_3E_1E_3+d^2E_2^2
 &\ge(2fe_3+\tfrac32d^2)E_1E_3\\
 &\ge(\tfrac12be_4+2cg)E_1E_3
 \ge2be_4E_4+2cgE_1E_3,\\
 2de_3E_2E_3&\ge2e_2gE_2E_3+ce_4E_2E_3
 \ge2e_2gE_2E_3+2ce_4E_1E_4,\\
 e_3^2E_3^2&\ge\tfrac83e_2e_4E_2E_4\ge2e_2e_4E_2E_4.
\end{align*}
Adding these comparisons and discarding $f^2E_1^2\ge0$ gives
\[
 \gamma_3^2\ge2(b+cE_1+e_2E_2)(gE_3+e_4E_4)=2\gamma_2\gamma_4.
\]
This argument has no computational dependency.


\subsection{The first and last comparisons}

\mnote{source~34 proves the first comparison $(\rho-1)\gamma_1^2\ge2\rho\gamma_2$
for every loopless directed relation with independent nonnegative role
activities by the weighted Motzkin--Straus bound with mass shifting,
noting that physical vertices with a zero tail role are never deleted from
their possible head role: this is Theorem~\ref{pgr:thm:first},
\eqref{pgr:eq:first}. It proves the last comparison
$\gamma_4^2\ge\frac{25}{8}\gamma_3\gamma_5$ from $q^2\ge\frac52pe_3$,
$qg\ge pd$, $g^2\ge2pf$ ($p=e_5$, $q=e_4$; for $qg\ge pd$ a feasible
three-tail summand $v_ju_S$ of $d$ is supplied by $u_{C\setminus i}$ from
$q$ times $v_ju_{C\setminus j}$ from $g$, and for $g^2\ge2pf$ the two
ordered cross terms of a feasible head pair) and the sink bounds
$E_4^2\ge\frac54E_3E_5$, $E_3E_4\ge\frac53E_2E_5$,
$E_3^2\ge\frac{10}{3}E_1E_5$: this is source~14's proof at $r=5$ with the
finite-$m$ factors dropped, and \eqref{pgr:s:fc:eq:last} follows from
Theorem~\ref{pgr:s:lg:thm:sharp}, since $K(5,m)>\frac{25}{8}$, when
$\gamma_5>0$ (it is trivial when $\gamma_5=0$). Both proofs are replaced
by these pointers. The package's \texttt{ORDINARY\_PROOFS.md} also cites
``the ordinary general universal-sink last-gap theorem''.}

At actual degree five, \eqref{pgr:eq:first} gives $2\gamma_1^2\ge5\gamma_2$.
Together with \eqref{pgr:s:fc:eq:interior} and \eqref{pgr:s:fc:eq:last}, this gives all four binomial-normalized comparisons. A positive-weight top support has positive-weight submatchings of every smaller size, so there are no internal zeros. This proves Theorem~\ref{pgr:s:fc:thm:main}.

\subsection{Scope and an exact limitation of a stronger claim}
The inequalities proved here hold on closed nonnegative activity faces, but their constants do not automatically prove rank ultra log concavity at smaller actual degrees. Degree four requires $\gamma_2^2\ge\tfrac94\gamma_1\gamma_3$ and $\gamma_3^2\ge\tfrac83\gamma_2\gamma_4$; degree three requires $\gamma_2^2\ge3\gamma_1\gamma_3$. A fifth core vertex with zero tail activity may remain a nonuniversal head, so it cannot simply be deleted to invoke a four-core result.

In fact, the stronger all-cloud cubic bound $\gamma_2^2\ge3\gamma_1\gamma_3$ fails for five-core preorders. Take only the core arc $0\to1$, six universal sinks of head activity one, core tails $(1,1,100,100,100)$ and core heads $(0,240,0,0,0)$. Direct Boolean endpoint-pair enumeration gives
\[
 (\gamma_0,\ldots,\gamma_5)
 =(1,2052,891015,129206000,4830450000,6000000),
\]
and
\[
 \gamma_2^2-3\gamma_1\gamma_3=-1,484,405,775<0.
\]
The relation becomes a preorder upon adjoining the diagonal. Its actual degree is five, and it satisfies the rank-five theorem. This example refutes only the proposed stronger constant three.
\mnote{this does not contradict source~10's Theorem~\ref{pgr:s:ad:thm:main}:
the example has five positive core tail activities, and source~10 proves
the constant three for four-vertex cores only. It was reproduced at intake
by direct enumeration.}
\mnote{Added 2~October 2026, batch~77: nor does it contradict source~35's
Theorem~\ref{pgr:s:fq:thm:main}, which proves the constant three for
five-vertex cores only on the face $\gamma_4=0$; here
$\gamma_4=4{,}830{,}450{,}000>0$. The lower-degree closure that this
section does not claim is now supplied by source~06 at actual degree four
and by source~35 at actual degrees at most three
(Corollary~\ref{pgr:s:fq:cor:all}).}


The package's \texttt{python verify.py} command performs a fresh exact replay. File integrity checking alone is not mathematical verification. Optional discovery programs and numerical-search receipts are not proof dependencies. The ordinary core proofs below, the moment reduction, and the first and last arguments are separate from the finite exact calculation. The result makes no claim about arbitrary weighted preorders outside the universal-sink family, no lower-degree closure claim for five cores, and no real-rootedness or stability claim.


\mnote{source~34's appendix, ``Ordinary coefficientwise core comparisons'',
proves \eqref{pgr:s:fc:eq:corelemmas} without finite enumeration. It is
source~06's appendix (Section~\ref{pgr:s:fb:app:core}) except in the
comparison involving two head roles: source~34 proves
$4fe_3+3d^2\ge be_4+4cg$ where source~06 proves the stronger
$4fe_3+2d^2\ge be_4+4cg$, and correspondingly its coefficient~$6$ of
$d_jd_k$, of $\epsilon_P\epsilon_Q$, of $\epsilon_PB_V$, of $B_UB_V$ and of
$sB_V$ is~$4$ in source~06, and its final $5B_UB_V$ is $3B_UB_V$. Its
other two comparisons and its auxiliary facts on three tails are the same
text. It is printed once, in source~06's form; the strengthening is
corroborated exactly by source~06's \texttt{check\_ordinary\_lemmas.py}.}

\section{Five-vertex cores at the degree-four boundary (source 06)}\label{pgr:s:fb:sec}

Source~06 is \emph{The rank four boundary for five vertex cores with
universal sinks} (7 pages as delivered; archives 07 and~06, the two parts
of \texttt{universal-sink-five-core-boundary-package}; ``Research note'').
Its files are shipped with the prefix \texttt{06-five-core-boundary-}. Its
bounded-sink certificates (two archives of 13.0 and 15.6~MB) and their
manifest are not shipped; they are in the arrival commit (README).

\paragraph{Abstract of source 06, as delivered.}

We prove rank ultra log concavity at actual degree four for every arbitrary loopless directed five-vertex core with a finite independent cloud of universal sinks and independent nonnegative tail and head activities. This closes the degree-four boundary of the previously studied full-rank-five family. Each feasible ordered pair of disjoint endpoint sets is counted once. The first interior comparison is certified exactly on two finite core domains: 3,044 classes with a retained fifth head-only vertex, and 9,608 classes with a bounded four-sink moment cone. The last comparison has ordinary proofs in both boundary branches. Zero tail activities never justify deleting an otherwise active head role. The normalization is the actual degree four; no general actual-degree-three conclusion is inferred.

\subsection{The boundary theorem}
Let $C=\{0,1,2,3,4\}$ carry any loopless directed relation $H$. Let $W$ be a finite independent set disjoint from $C$. The full relation $D$ consists of $H$ and all arcs $C\to W$, with no other arcs. Assign independent nonnegative core tail activities $u_i$, core head activities $v_i$, and sink head activities $w_x$. Define
\[
 \Gamma_D(t)=\sum_{k=0}^5\gamma_k t^k,\qquad
 \gamma_k=\sum_{\substack{S\cap T=\varnothing,\ |S|=|T|=k\\
 S\text{ admits a directed perfect matching to }T}}
 \prod_{i\in S}u_i\prod_{j\in T\cap C}v_j\prod_{x\in T\cap W}w_x.
\]
Every nonempty tail set lies in $C$, $\gamma_0=1$, and a feasible ordered pair $(S,T)$ is counted once, independently of how many matching witnesses it has.

\begin{theorem}\label{pgr:s:fb:thm:main}
On the entire nonnegative activity face $\gamma_5=0$,
\begin{equation}\label{pgr:s:fb:eq:gaps}
 \gamma_2^2\ge\frac94\gamma_1\gamma_3,
 \qquad \gamma_3^2\ge\frac83\gamma_2\gamma_4.
\end{equation}
Consequently, whenever the actual surviving degree is four,
$\bigl(\gamma_k/\binom4k\bigr)_{k=0}^4$ is log-concave.
The first inequality in \eqref{pgr:s:fb:eq:gaps} uses exact rational certificates; the second and the first-coefficient comparison have ordinary proofs.
\end{theorem}

No transitivity is assumed. In particular the theorem applies when $H$ is the off-diagonal relation of a preorder, since adjoining the diagonal to $D$ then gives a preorder. Reversing all arcs and exchanging tail and head roles gives the universal-source version.

\subsection{Boolean coefficients and the two boundary branches}
For disjoint $S,J\subseteq C$, put $u_S=\prod_{i\in S}u_i$, $v_J=\prod_{j\in J}v_j$. Let $F_{kr}$ sum $u_Sv_J$ over $|S|=k$, $|J|=r$ such that a matching from $S$ covers $J$. This is Boolean feasibility, not the number of injections. Write
\[
 a=F_{11},\ b=F_{22},\ c=F_{21},\ d=F_{31},\ f=F_{32},\ g=F_{41},
 \qquad e_r=\ee_r(u),\quad E_r=\ee_r(w).
\]
Every injection covering the selected core heads extends to the selected sink heads. Partitioning supports by the number of core heads therefore gives
\begin{equation}\label{pgr:s:fb:eq:coeff}
\begin{aligned}
 \gamma_1&=a+e_1E_1,& \gamma_2&=b+cE_1+e_2E_2,\\
 \gamma_3&=fE_1+dE_2+e_3E_3,&
 \gamma_4&=gE_3+e_4E_4,\qquad \gamma_5=e_5E_5.
\end{aligned}
\end{equation}
Because the activities are nonnegative, $\gamma_5=0$ means either some core tail is zero or at most four sink heads are positive. These give the two proof domains:
\begin{enumerate}
\item A zero core-tail activity, relabeled $u_4=0$, with an arbitrary finite sink cloud. The vertex $4$ remains available as a head. Arcs out of it have zero product weight and may be ignored; its incoming arcs cannot be ignored.
\item At most four positive sink-head activities, with all five core-tail activities allowed. Zeros may be padded to give a four-entry sink vector.
\end{enumerate}
Their union contains the entire boundary, including overlaps and further zero activities. At actual degree four there are at least four active tails; if all five are active, there must be three or four active sinks. The proofs below do not require positivity of $\gamma_4$.

\subsection{Two moment reductions for the first interior comparison}

Let $p_r=\sum_xw_x^r$, $A=\sqrt{p_2}$, and $B=p_1-\sqrt{p_2}$. Then $A,B\ge0$ and, by Lemma~\ref{pgr:s:ad:lem:moments} of source~10 in source~06's scaled form,
\[
 E_1=A+B,\qquad 2E_2=2AB+B^2,\qquad 6E_3\le3AB^2+B^3.
\]
\mnote{source~06 proves this display directly ($p_1^2\ge p_2$, and each
$w_x\le\sqrt{p_2}$ implies $p_3\le p_2^{3/2}$; substitute into
$2E_2=p_1^2-p_2$ and $6E_3=p_1^3-3p_1p_2+2p_3$), as its equation
``envelope''; references to it point to source~10's lemma.}

The coefficient of $E_3$ in $\gamma_2^2-\tfrac94\gamma_1\gamma_3$ is
$-\tfrac94(a+e_1E_1)e_3\le0$, so replacing $E_3$ by its displayed upper bound can only decrease the gap.

For the zero-tail branch use the variable order
\[
 (u_0,\ldots,u_4,v_0,\ldots,v_4,A,B),\qquad u_4=0,
\]
and define
\begin{align*}
 G_1&=a+e_1(A+B),\\
 G_2&=2b+2c(A+B)+e_2(2AB+B^2),\\
 G_3&=6f(A+B)+3d(2AB+B^2)+e_3(3AB^2+B^3).
\end{align*}
All core polynomials here are specialized only in the tail role $u_4$; $v_4$ remains an independent variable.

For the bounded-sink branch, Cauchy's inequality on four entries gives $p_1^2\le4p_2$, hence $A\ge B$. Write $A=B+C$ with $C\ge0$. Then
\begin{equation}\label{pgr:s:fb:eq:bounded}
 E_1=C+2B,\qquad 2E_2=2BC+3B^2,\qquad 6E_3\le3CB^2+4B^3.
\end{equation}
Use the variable order $(u_0,\ldots,u_4,v_0,\ldots,v_4,C,B)$ and the same definitions of $G_k$, with the three scaled moments replaced by \eqref{pgr:s:fb:eq:bounded}. This enlarged nonnegative cone is a sufficient bound; no claim of sharpness or equivalence for a fixed four-sink count is needed.

In both cases the boundary substitution satisfies $G_k=k!\gamma_k$ for $k=1,2,3$, and
\begin{equation}\label{pgr:s:fb:eq:Q}
 Q=2G_2^2-3G_1G_3
   =8\left(\gamma_2^2-\frac94\gamma_1\gamma_3\right).
\end{equation}
Thus nonnegativity of the corresponding $Q$ proves the required inequality.

\subsection{Finite exact certificates}
The zero-tail domain has $16$ relevant arc bits: tails $0,1,2,3$ may point to any different core vertex, including the head-only vertex $4$. All $2^{16}=65,536$ labeled relations split into $3,044$ classes under permutations of the four active-tail labels fixing $4$. The full five-tail domain has $2^{20}=1,048,576$ relations, split into $9,608$ classes under all five-label permutations.

For each representative in each domain the certificate is an exact identity
\begin{equation}\label{pgr:s:fb:eq:cert}
 Q(x)=\sum_\ell q_\ell x^{m_\ell}
          \left(\sum_{r\in\mathcal R_\ell}c_{\ell,r}x^r\right)^2+R(x),
\end{equation}
where $q_\ell>0$, all scalar coefficients are rational, every exponent is a nonnegative integer, and $R$ is coefficientwise nonnegative. Such an identity proves nonnegativity on the closed nonnegative orthant.

\begin{proposition}[Finite exact verification]\label{pgr:s:fb:prop:cert}
The two supplied collections cover all $3,044$ zero-tail classes and all $9,608$ bounded-sink classes. Every exact identity \eqref{pgr:s:fb:eq:cert} has a coefficientwise nonnegative remainder. The zero-tail collection contains $212,319$ binomial-square terms, and $158$ targets are already coefficientwise nonnegative.
\end{proposition}
\begin{proof}[Verification]
The independent checker reconstructs the endpoint-support kernels by Boolean matching existence, forms the appropriate scaled moment polynomial and \eqref{pgr:s:fb:eq:Q}, then expands and subtracts every rational square. It checks every remainder coefficient, including exponents absent from the target. For the zero-tail domain it also checks that $u_4$ occurs in neither targets nor certificate terms. Independently generated relabeling orbits are disjoint and cover the complete labeled graph domains. No producer code is imported. The accompanying data and standard-library checker supply the reproducible finite calculation.
\end{proof}


\mnote{source~06's proposition states square counts for the zero-tail
collection only. Its README and \texttt{audits/four-sink-receipt.json}
record those of the bounded-sink collection: 1,755,474 polynomial squares
of two to 67 terms, with 57 targets already coefficientwise nonnegative.
At intake all 12,652 certificate members of both collections were
authenticated against the audit manifests and 40 certificates were
replayed; the bounded-sink certificates are not shipped with this report
(README).}

Numerical linear programming served only to propose candidates. Rational reconstruction, exact active-system solutions where needed, and exact expansion determine validity; no numerical tolerance is a proof premise. Proposition~\ref{pgr:s:fb:prop:cert}, the moment bounds, and the correct monotonicity direction prove the first inequality in \eqref{pgr:s:fb:eq:gaps} on the entire boundary.

\subsection{An ordinary last-gap proof with a retained head-only vertex}
Assume the active-tail set is $I=\{0,1,2,3\}$ and the remaining core vertex is $h=4$. Initially take the four tail activities positive. Put $p=\prod_{i\in I}u_i$, $x_i=1/u_i$, and $t=v_h$. If $h$ has no in-neighbor in $I$, set $t=0$, since its head role then contributes no weight. Define
\[
 Z=E_4+tE_3.
\]
For each $i\in I$, let $\alpha_i,\beta_i,\delta_i$ indicate whether $I\setminus\{i\}$ can cover, respectively, the head sets $\{h\},\{i\},\{h,i\}$. Set
\[
 A_i=E_3+(t\alpha_i+v_i\beta_i)E_2+tv_i\delta_iE_1.
\]
For $i\ne j$, let $a,b,c,d,e,f$ indicate whether $S=I\setminus\{i,j\}$ can cover, respectively,
$\{h\}$, $\{i\}$, $\{j\}$, $\{h,i\}$, $\{h,j\}$, $\{i,j\}$, and put
\[
 B_{ij}=E_2+(ta+v_ib+v_jc)E_1+tv_id+tv_je+v_iv_jf.
\]
These local Boolean letters are independent of the earlier global core-polynomial names. Partitioning by omitted active tails gives
\begin{equation}\label{pgr:s:fb:eq:reciprocal}
 \gamma_4=pZ,\qquad \gamma_3=p\sum_i x_iA_i,
 \qquad \gamma_2=p\sum_{i<j}x_ix_jB_{ij}.
\end{equation}

\begin{lemma}\label{pgr:s:fb:lem:pair}
For every pair $i\ne j$, $A_iA_j\ge ZB_{ij}$.
\end{lemma}
\begin{proof}
If $t=0$ because the extra head has no incoming arc, the comparisons below involving $t$ vanish and the remaining ones follow by retaining matching witnesses. Otherwise the extra head has a nonempty in-neighborhood $N\subseteq I$.
The relevant Boolean implications are
\begin{gather*}
 \alpha_i+\alpha_j\ge1+a,\quad \alpha_i\alpha_j\ge a,
 \quad \beta_i\ge b,\quad\beta_j\ge c,\\
 \delta_i\ge d,\quad\beta_i\alpha_j\ge d,
 \quad\delta_i+\beta_i\alpha_j\ge b,
 \quad\delta_i\alpha_j\ge d,\\
 \delta_j\ge e,\quad\beta_j\alpha_i\ge e,
 \quad\delta_j+\beta_j\alpha_i\ge c,
 \quad\delta_j\alpha_i\ge e,\\
 \beta_i\beta_j\ge f,\qquad
 \delta_i\beta_j+\delta_j\beta_i\ge f.
\end{gather*}
Most assertions follow by enlarging a feasible tail set. For $\delta_i+\beta_i\alpha_j\ge b$, choose a source $r\in S$ for $i$. If $\alpha_j=1$, the second term suffices. Otherwise $N=\{j\}$, and $j\to h,r\to i$ prove $\delta_i=1$. The symmetric assertion is identical. For the last implication, if $f=1$ choose distinct $r,s\in S$ with $r\to i,s\to j$. For any $z\in N$, if $z=i$ use $i\to h,s\to j$; if $z=j$ use $j\to h,r\to i$; and if $z\in S$, pair its arc to $h$ with the other matching source. Thus at least one of $\delta_i,\delta_j$ is one.

The coefficients of $A_iA_j-ZB_{ij}$ as a polynomial in $t,v_i,v_j$ are listed below; the rows containing only $v_i$ have symmetric $v_j$ counterparts.
\[
\begin{array}{c|l}
 1&E_3^2-E_2E_4\\
 t&(\alpha_i+\alpha_j-1)E_2E_3-aE_1E_4\\
 v_i&\beta_iE_2E_3-bE_1E_4\\
 t^2&\alpha_i\alpha_jE_2^2-aE_1E_3\\
 tv_i&\delta_iE_1E_3+\beta_i\alpha_jE_2^2-bE_1E_3-dE_4\\
 v_iv_j&\beta_i\beta_jE_2^2-fE_4\\
 t^2v_i&\delta_i\alpha_jE_1E_2-dE_3\\
 tv_iv_j&(\delta_i\beta_j+\delta_j\beta_i)E_1E_2-fE_3\\
 t^2v_iv_j&\delta_i\delta_jE_1^2.
\end{array}
\]
Every coefficient is nonnegative by the Boolean implications and elementary-symmetric log-concavity:
\[
 E_3^2\ge E_2E_4,\quad E_2E_3\ge E_1E_4,\quad
 E_2^2\ge E_1E_3,\quad E_2^2\ge E_4,\quad E_1E_2\ge E_3.
\]
For the $tv_i$ coefficient, if $d=1$ then both $\delta_i$ and $\beta_i\alpha_j$ equal one, giving at least $E_2^2-E_4$. If $d=0,b=1$, their sum is at least one and $E_2^2\ge E_1E_3$ suffices; the remaining case is immediate. This also handles the symmetric coefficient.
\end{proof}

Let $y_i=x_iA_i$. By Lemma~\ref{pgr:s:fb:lem:pair} and \eqref{pgr:s:fb:eq:reciprocal},
\[
 \frac{\gamma_2\gamma_4}{p^2}\le\sum_{i<j}y_iy_j,
 \qquad \frac{\gamma_3}{p}=\sum_i y_i.
\]
Cauchy's inequality on four variables gives $(\sum_i y_i)^2\ge\tfrac83\sum_{i<j}y_iy_j$, proving the required last gap. More explicitly,
\[
 3\gamma_3^2-8\gamma_2\gamma_4
 =p^2\sum_{i<j}(y_i-y_j)^2
   +8p^2\sum_{i<j}x_ix_j(A_iA_j-ZB_{ij})\ge0.
\]
Multiplication by $p^2$ removes all reciprocal denominators; alternatively, continuity includes further zero tails. The fifth physical vertex has retained its head role throughout.

\subsection{An ordinary last-gap proof with at most four sinks}
For an arbitrary five-core, the following inequalities hold coefficientwise:
\begin{equation}\label{pgr:s:fb:eq:corelemmas}
 3fd\ge bg,\qquad 4fe_3+2d^2\ge be_4+4cg,
 \qquad 2de_3\ge2e_2g+ce_4.
\end{equation}
The appendix supplies complete ordinary proofs. Newton's inequalities in five tail variables give $e_3^2\ge2e_2e_4$. With at most four positive sink heads, pad zeros to four entries and use
\begin{equation}\label{pgr:s:fb:eq:finiteNewton}
\begin{gathered}
 E_1E_2\ge6E_3,\qquad E_2^2\ge\tfrac94E_1E_3,\qquad
 E_3^2\ge\tfrac83E_2E_4,\\
 E_1E_3\ge4E_4,\qquad E_2E_3\ge2E_1E_4.
\end{gathered}
\end{equation}
The last two constants can in fact be strengthened to $16$ and $6$. All displayed bounds follow from binomial-normalized Newton inequalities, including zeros~\cite{BH}.

Put $X=fE_1,Y=dE_2,Z=e_3E_3$. Using \eqref{pgr:s:fb:eq:corelemmas}--\eqref{pgr:s:fb:eq:finiteNewton}, the six terms of $\gamma_2\gamma_4$ are bounded in four groups:
\begin{align*}
 bgE_3&\le\tfrac12XY,\\
 be_4E_4+cgE_1E_3
 &\le(\tfrac14be_4+cg)E_1E_3
 \le(fe_3+\tfrac12d^2)E_1E_3
 \le XZ+\tfrac29Y^2,\\
 ce_4E_1E_4+e_2gE_2E_3
 &\le(\tfrac12ce_4+e_2g)E_2E_3\le YZ,\\
 e_2e_4E_2E_4&\le\tfrac{3}{16}Z^2.
\end{align*}
Since $\gamma_3=X+Y+Z$, subtraction gives the exact lower bound
\begin{align*}
 \gamma_3^2-\tfrac83\gamma_2\gamma_4
 &\ge X^2+\tfrac23XY+\tfrac{11}{27}Y^2+\tfrac12Z^2
       -\tfrac23XZ-\tfrac23YZ\\
 &=\left(X+\frac{Y-Z}{3}\right)^2
   +\frac8{27}\left(Y-\frac34Z\right)^2+\frac29Z^2\ge0.
\end{align*}
This proves the last comparison for the bounded-sink branch without any numerical or computer-assisted premise. Together with the preceding section, it covers all of $\gamma_5=0$.


\subsection{Actual degree and reproducibility}

\mnote{source~06 proves here, for every loopless directed relation with
nonnegative independent roles and actual degree $\rho\ge2$, the first
comparison $(\rho-1)\gamma_1^2\ge2\rho\gamma_2$ by Motzkin--Straus on the
graph of positive-weight arcs with mass shifting; this is
Theorem~\ref{pgr:thm:first} and is replaced by that pointer.}

At actual degree four this gives $3\gamma_1^2\ge8\gamma_2$. Along with \eqref{pgr:s:fb:eq:gaps}, these are precisely the three binomial-normalized inequalities. Positive-weight submatchings of a positive top support show there are no internal zeros, completing Theorem~\ref{pgr:s:fb:thm:main}.

Run \texttt{python verify.py} in the supplied package for a fresh standard-library exact replay. The driver checks integrity, unpacks the compressed certificate data to a fresh temporary directory, and runs the independent support, polynomial, and orbit checks. Neither a numerical solver, a prior research archive, nor network access is needed. Integrity checking by itself is not a proof replay. Separate exact audits corroborate the ordinary Boolean implications and square identities, but the ordinary arguments do not depend on those computations.

The theorem does not infer the degree-three inequality $\gamma_2^2\ge3\gamma_1\gamma_3$ from a padded quartic normalization. It therefore closes the actual-degree-four boundary only. Nonuniversal exterior neighborhoods and larger physical cores are outside its scope. No general weighted-preorder, real-rootedness, stability, or novelty claim is made.
\mnote{Added 2~October 2026, batch~77: the degree-three inequality
$\gamma_2^2\ge3\gamma_1\gamma_3$ on the face $\gamma_4=0$ is proved by
source~35 (Theorem~\ref{pgr:s:fq:thm:main}), without any padded
normalization; with it, five-vertex cores are rank-ULC at every actual
degree (Section~\ref{pgr:s:fq:sec}).}


\subsection{Appendix of source 06: ordinary coefficientwise core comparisons}\label{pgr:s:fb:app:core}
\mnote{printed once for sources 34 and~06, in source~06's stronger form;
see the note at the end of Section~\ref{pgr:s:fc:sec}.}

We prove the three comparisons in \eqref{pgr:s:fb:eq:corelemmas} without finite enumeration.

\subsubsection{The comparison involving three head roles}
For a repeated-head monomial $v_j^2v_k$, every feasible two-tail set $B$ in $b_{jk}$ makes $B\cup\{k\}$ feasible for head $j$. Thus $f_{jk}d_j\ge b_{jk}g_j$ coefficientwise.

For three distinct heads $j,k,l$, a fixed tail monomial in $bg$ has coefficient at most three, since choosing the head belonging to $g$ determines its four-tail factor and hence its two-tail factor. It suffices to show that each such monomial occurs in $fd$ at least once.

Fix a summand of $b_{jk}g_l$ and a witnessing two-tail matching onto $j,k$, and write $C=\{j,k,l,a,b\}$. If the two tails are $a,b$, choose an in-neighbor $r$ of $l$. If $r$ is one of $a,b$, pair $l$ with the original head whose source is the other tail. If $r$ is one of $j,k$, pair $l$ with the other original head. The resulting head pair has a feasible three-tail complement, and division of the original tail monomial by that complement gives a feasible three-tail set for the remaining head, using its original incoming arc.

Otherwise the two tails are $l,a$, after renaming $a,b$. Relabel $j,k$ so the matching is $l\to j,a\to k$. If $l$ has an in-neighbor in $\{j,k,a\}$, use $f_{jk}d_l$. Otherwise its nonempty neighborhood contains $b$; use $f_{kl}d_j$, with witnesses $a\to k,b\to l$ and $l\to j$. Each product is exactly the original tail monomial. Thus $3fd\ge bg$.

\subsubsection{Elementary facts on three tails}
For three activities $z_a,z_b,z_c$ write $t=\ee_1(z),s=\ee_2(z),p=\ee_3(z)$.
For $U\subseteq\{a,b,c\}$ let $L_U=\sum_{i\in U}z_i$ and
$B_U=\sum_{|S|=2,\ S\cap U\ne\varnothing}z_S$.
For nonempty $U$,
\begin{equation}\label{pgr:s:fb:eq:aux}
 sB_U\ge p(t+L_U),\qquad
 B_UB_V\ge2pL_{U\cap V},\qquad
 B_UB_V\ge b(U,V)s,
\end{equation}
coefficientwise, where $b(U,V)$ sums pairs that admit distinct representatives in $U,V$.
For singleton $U$, the first bound follows by expansion; for $|U|\ge2$, use $B_U=s$ and $s^2\ge2pt$. For the second bound, the two ordered products of distinct pairs containing $i\in U\cap V$ supply $2pz_i$. For the last bound, if both neighborhoods have size at least two use $b(U,V)\le s$. If just one is a singleton, the bound is immediate. For distinct singletons $\{a\},\{b\}$ it is
$z_a(z_b+z_c)z_b(z_a+z_c)=z_az_b(s+z_c^2)\ge z_az_bs$.
For equal singletons, $b(U,V)=0$.

\subsubsection{The comparison involving two head roles}
We prove $4fe_3+2d^2\ge be_4+4cg$. For the coefficient of $v_j^2$, it is enough that $d_j^2\ge2c_jg_j$. If $j$ has at least two in-neighbors among its four permitted tails $z$, then $d_j=\ee_3(z)$, $c_j\le\ee_2(z)$, $g_j=\ee_4(z)$, and
$\ee_3(z)^2\ge2\ee_2(z)\ee_4(z)$ coefficientwise. With just one in-neighbor of activity $a$, use $c_j=at,d_j=as,g_j=ap$ and $s^2\ge2tp$. Empty neighborhoods contribute zero.

For distinct heads $j,k$, put $x=u_j,y=u_k$ and use $t,s,p$ on the remaining set $A$ of three vertices. Let $P,Q\subseteq A$ be their incoming neighborhoods from $A$. If $k\to j$ put $U=A$, otherwise put $U=P$; if $j\to k$ put $V=A$, otherwise put $V=Q$. Write $\epsilon_P=1_{P\ne\varnothing}$, $\epsilon_Q=1_{Q\ne\varnothing}$, $\eta_P=1_{U\ne\varnothing}$, $\eta_Q=1_{V\ne\varnothing}$. Let $\epsilon$ indicate that $P,Q$ admit distinct representatives, and let $b=b(P,Q)$. Then
\[
\begin{array}{lll}
 c_j=B_P+yL_U,&d_j=\epsilon_Pp+yB_U,&g_j=\eta_Pyp,\\
 c_k=B_Q+xL_V,&d_k=\epsilon_Qp+xB_V,&g_k=\eta_Qxp,
\end{array}
\]
while $f_{jk}=\epsilon p$, $e_3=p+(x+y)s+xyt$, and $e_4=(x+y)p+xys$.
The desired coefficient is
\[
 R=4\epsilon pe_3+4d_jd_k-be_4-4(c_jg_k+c_kg_j).
\]
Its constant term is $p^2(4\epsilon+4\epsilon_P\epsilon_Q)\ge0$.
Its $x$ coefficient divided by $p$ is
$4\epsilon s+4\epsilon_PB_V-b-4\eta_QB_P$.
If $\epsilon=1$, use $B_P\le s$ and $b\le B_V$. If $\epsilon=0$, then $b=0$. The expression vanishes unless $\epsilon_P=\eta_Q=1$; in that case either $V=A$, or $P=Q$ is a singleton, so $B_V\ge B_P$. The $y$ coefficient is symmetric.

The $xy$ coefficient is
\[
 4\epsilon pt+4B_UB_V-bs-4p(\eta_QL_U+\eta_PL_V).
\]
It vanishes if $\eta_P\eta_Q=0$. Otherwise both indicators are one. If $\epsilon=0$ and one cross arc is present, assume $U=A$ and use $4sB_V\ge4p(t+L_V)$ from \eqref{pgr:s:fb:eq:aux}. If neither cross arc is present, $P=Q$ is a nonempty singleton, and $B_U^2\ge2pL_U$ suffices. Finally, if $\epsilon=1$, use $B_UB_V\ge bs$, $B_UB_V\ge2pL_{U\cap V}$, and $L_U+L_V-t\le L_{U\cap V}$ to obtain
\[
 4B_UB_V-bs+4p(t-L_U-L_V)
 \ge3B_UB_V-4pL_{U\cap V}\ge0.
\]
Thus every coefficient of $R$ is nonnegative.

\subsubsection{The comparison involving one head role}
Fix a head $j$ and put $x=u_j$. Write $z_r=\ee_r(z)$ for the other four tails. The constant and $x$ coefficients of
$2d_je_3-2e_2g_j-c_je_4$ are
\[
 Q_0=2d_jz_3-(2z_2+c_j)g_j,\qquad
 Q_1=2d_jz_2-2z_1g_j-c_jz_3.
\]
If the neighborhood has at least two vertices, $d_j=z_3,c_j\le z_2$, and the coefficientwise bounds $z_3^2\ge2z_2z_4$ and $z_2z_3\ge3z_1z_4$ prove both comparisons. With a singleton neighborhood of activity $a$, let $t,s,p$ refer to the other three tails. Direct expansion gives
\[
 Q_0=a^2(2s^2-3tp),\qquad
 Q_1=a^2(st-2p)+a(2s^2-3tp),
\]
which are coefficientwise nonnegative since $s^2\ge2tp$ and $st\ge3p$. The empty-neighborhood case is zero. Summing over heads proves $2de_3\ge2e_2g+ce_4$.


\section{Five-vertex cores at actual degree at most three (source 35)}\label{pgr:s:fq:sec}

Source~35 is \emph{The cubic boundary for five-vertex cores with universal
sinks} (5 pages as delivered, dated 1~October 2026; ``Research note'', no
author line). It is the one manuscript of this report that arrived in
ProveIt's batch~77: batch-77 manuscript~71, the archive
\nolinkurl{universal-sink-five-core-cubic-package} (six ZIP parts in one
archive), arrival commit \texttt{096ee7b87}, placement commit
\texttt{089d2b825}. It was written right after sources 34 and~06 and
continues their five-vertex line. Its files are shipped with the prefix
\texttt{35-five-core-cubic-}. Its $7{,}901$ certificate files, delivered in
six compressed shards of 95.1~MB, are not shipped: they are kept in the
arrival commit, and no shipped program can regenerate them (README,
``Reconstructing what is not shipped''). Its structural input
\nolinkurl{inputs/structural-9608.json} is shipped as the derived container
\nolinkurl{data/35-five-core-cubic-inputs-structural-9608.json.xz}, which
decompresses to the delivered file byte for byte.

\paragraph{Abstract of source 35, as delivered.}

We prove rank ultra log concavity at actual degree three for every arbitrary directed five-vertex core with a finite independent cloud of universal sinks and independent nonnegative tail and head activities. Combined with the established degree-four and degree-five results, this closes every actual degree from zero to five in this family. Each feasible ordered pair of disjoint endpoint supports is counted once. Ordinary reductions leave a two-sink problem on 9,608 core classes. Its exact hybrid proof uses established structural criteria, the earlier preorder theorem, and rational square certificates; 94 classes use a positive tail-sum multiplier. The complete finite replay is independent of the discovery code and numerical optimization. The cubic constant is sharp, even for preorder cores.

\mnote{Added 2~October 2026, batch~77. This section answers the statement
in the introduction to Part~III that actual degree three was open for
five-vertex cores. Several of its steps are already printed in this report
and are replaced by pointers: its Boolean coefficient formulas are
source~34's \eqref{pgr:s:fc:eq:coeff}; its first comparison is
Theorem~\ref{pgr:thm:first}; its argument that the one-tail criterion
applies to physical supports is the proof of
Theorem~\ref{pgr:r:ir:thm:oneside}. Its degree-four and degree-five inputs
are sources 06 and~34, which it cites and does not re-prove. Everything
else is printed below.}

\paragraph{Notation of this section.}
Source~35 works in the setting of Part~III and uses source~34's letters
$a,b,c,d,f,g$ for the Boolean core polynomials and $e_r$, $E_r$ for the
elementary symmetric functions of the core tail and sink head activities.
Three of its letters would clash with letters of other sources or with each
other. They are renamed here, and nothing else is.
\begin{center}\small
\begin{tabular}{p{50mm}p{24mm}p{66mm}}
\toprule
Source~35 & Here & Reason\\\midrule
$m$, the number of positive core tails & $n_u$ & source~14's $m$
  (Section~\ref{pgr:s:lg:sec}) counts positive \emph{sinks}\\
$q$, the number of positive sinks & $n_w$ & $q_\ell$ are the weights of
  a certificate\\
$p=\prod_{i\ne h}u_i$ (reduction, case~4) & $u_{C\setminus h}$ & source~34's
  notation $u_S=\prod_{i\in S}u_i$; $p$ is also the hub of a one-tail
  cover\\
$s=w_0+w_1$, $p=w_0w_1$ (two-sink lemma) & $E_1$, $E_2$ & these are the
  elementary symmetric functions of the two sink activities\\\bottomrule
\end{tabular}
\end{center}
The hub $p$ of a one-tail cover $(p,Q)$ keeps its letter, as in
Theorem~\ref{pgr:r:ir:thm:oneside}. The letter $h$ names a zero-tail core
vertex, as in source~06; $H$ is the core relation, and $P_H$, $M_H$, $R_H$
are indexed by it. The multiplier $M_H$ is not the side matroid $M$ of
Theorem~\ref{pgr:r:ir:thm:oneside}. A tempting misreading should be
avoided: the theorem below asserts the cubic constant three only on the
face $\gamma_4=0$. Across the whole five-vertex family that constant is
false; source~34's example in Section~\ref{pgr:s:fc:sec} has $\gamma_4>0$.

\subsection{Model, statement and proof dependencies}
Let $C=\{0,1,2,3,4\}$ carry an arbitrary loopless directed relation $H$. Let $W$ be a finite independent set disjoint from $C$. The full relation $D$ consists of $H$ and every arc $C\to W$, with no other arcs. Assign nonnegative tail activities $u_i$ and head activities $v_i$ independently on $C$, and head activities $w_x$ on $W$. Put
\[
 \Gamma_D(t)=\sum_{k=0}^5\gamma_k t^k,\qquad \gamma_0=1,
\]
where $\gamma_k$ is the sum of
$u_Sv_{T\cap C}w_{T\cap W}$ over ordered pairs $(S,T)$ of disjoint physical endpoint sets, $|S|=|T|=k$, that admit a directed perfect matching from $S$ to $T$. Every feasible pair is counted once, regardless of its number of witnesses. Write $\rho=\deg\Gamma_D$ for the actual surviving degree.

\begin{theorem}[Cubic boundary]\label{pgr:s:fq:thm:main}
On the entire nonnegative activity face $\gamma_4=0$,
\begin{equation}\label{pgr:s:fq:eq:cubic}
 \gamma_2^2\ge 3\gamma_1\gamma_3.
\end{equation}
Consequently the family is rank ultra log concave at every actual degree $\rho\le3$: the sequence $(\gamma_k/\binom\rho k)_{k=0}^{\rho}$ has no internal zeros and is log-concave. The coefficient $3$ in \eqref{pgr:s:fq:eq:cubic} is sharp, even for preorder cores.
\end{theorem}

No transitivity is assumed. Adjoining the diagonal has no effect on physically disjoint support pairs. Thus the result applies to universal-sink extensions of finite preorder cores. Reversal of every arc, accompanied by interchange of tail and head activities, gives the universal-source version.

\mnote{relation to the other sources. For a preorder core the theorem is a
special case of source~22's Theorem~\ref{pgr:r:ir:thm:main}: every
independently weighted preorder of actual degree at most three is
rank-ULC, and source~35 itself uses that theorem for 49 core classes
below. The new content is the non-transitive core. For four-vertex cores
the constant three holds on every face (source~10,
Theorem~\ref{pgr:s:ad:thm:main}); for five-vertex cores it holds on the face
$\gamma_4=0$ but not everywhere. Source~34's counterexample to the
constant three (Section~\ref{pgr:s:fc:sec}) has
$\gamma_4=4{,}830{,}450{,}000>0$ and actual degree five, so it lies off this
face and does not contradict the theorem.}

The ordinary reduction uses two established companion results, stated here to fix their scope:
\begin{itemize}
\item[(R3)] Every finite loopless directed relation with at most three active tail roles has actual-degree rank ultra log concavity under independent nonnegative role activities. This is the one-shore case of the previously proved role-matching-rank-three theorem; transitivity is unnecessary.
\item[(C4)] Every arbitrary directed four-vertex core with an independent universal-sink cloud satisfies $\gamma_2^2\ge3\gamma_1\gamma_3$ for all nonnegative independent role activities, and is rank ultra log concave at every actual degree.
\end{itemize}
The finite verification additionally uses the established physical-cover-three, two-tail/two-head role-cover, and one-tail stable-side criteria, together with the established weighted-preorder theorem at actual degree at most three. Their exact statements, approved versions, and hashes are recorded among the proof dependencies. These are mathematical inputs; the new replay checks every application hypothesis. It does not rerun every companion theorem's earlier certificate collection.

\mnote{in this report these inputs are: (R3), the one-shore case of
source~22's Corollary~\ref{pgr:r:ir:cor:rolerank}, that is,
Theorem~\ref{pgr:r:ir:thm:three} for three active tails; (C4), source~10's
Theorem~\ref{pgr:s:ad:thm:main}; the physical-cover-three criterion,
source~11's Theorem~\ref{pgr:r:pc:thm:main}; the two-tail/two-head role
cover, source~20's Theorem~\ref{pgr:r:rc:thm:rolecover}; the one-tail
stable-side criterion, Theorem~\ref{pgr:r:ir:thm:oneside} (source~22, a
case of source~15's Theorem~\ref{pgr:r:ps:thm:pendant}); and the weighted
preorder theorem, source~22's Theorem~\ref{pgr:r:ir:thm:main}. Source~35
names none of them. It calls them ``previously proved'' or ``established''
and identifies them by archive hashes in its dependency ledger
(\nolinkurl{35-five-core-cubic-DEPENDENCIES.md}). Its five companion
archives are, byte for byte, batch-73 archives 10, 34, 07, 06 and~22, that
is, sources 10, 34, 06 and~22 of this report. Some of the pinned dependency
notes still say ``proposed'' or describe parts that were open when they
were written. The ledger explains that later explicit approvals pin the
same theorem bytes, and that only the approved statements are used.}

\subsection{Boolean coefficients and the complete degree-drop reduction}
For disjoint $S,J\subseteq C$, let $F_{kr}$ sum $u_Sv_J$ over $|S|=k$, $|J|=r$, for which distinct selected tails can cover all selected core heads $J$. Feasibility is Boolean. Every such partial matching extends to any chosen $k-r$ sinks. With
\[
 a=F_{11},\quad b=F_{22},\quad c=F_{21},\quad d=F_{31},\quad f=F_{32},
 \qquad e_r=\ee_r(u),\quad E_r=\ee_r(w),
\]
the coefficients $\gamma_1,\dots,\gamma_4$ are given by source~34's
\eqref{pgr:s:fc:eq:coeff}, where
\[
 g=F_{41}=\sum_{j\in C}\eta_jv_ju_{C\setminus j},
 \qquad \eta_j=\mathbf1\{\text{some }i\ne j\text{ has }i\to j\text{ in }H\}.
\]
\mnote{source~35 states these formulas as its own equation, without
$\gamma_5$. They are source~34's (and source~06's
\eqref{pgr:s:fb:eq:coeff}) in the same letters, and are printed once, in
Section~\ref{pgr:s:fc:sec}. Source~35 writes $g$ with the indicator
$\eta_j$; this is source~34's $g=\sum_jv_jg_j$, with $g_j=u_{C\setminus j}$
for a nonempty in-neighbourhood of $j$ and $g_j=0$ otherwise. Source~35's
justification of the $\gamma_4$ formula is kept.}
For $\gamma_4$, four core tails allow at most one core head; if that head is $j$, the tail set is exactly $C\setminus\{j\}$, and one incoming arc suffices. This proves its displayed formula without counting matching witnesses.

An arc is \emph{effective} when the product of its role activities is positive. Deleting ineffective arcs preserves every positive-weight support. A zero tail activity alone does not permit deletion of the corresponding head role. Zero-weight sink heads can be removed, since sinks have no outgoing arcs.

Put $n_u=|\{i:u_i>0\}|$ and $n_w=|\{x:w_x>0\}|$. It suffices to consider $\gamma_3>0$ and $\gamma_4=0$; otherwise \eqref{pgr:s:fq:eq:cubic} is immediate. A positive support contains positive submatchings, so these conditions give actual degree three.
\begin{enumerate}
\item If $n_u\le3$, apply (R3) to the effective relation. Its role graph has at most three nonisolated tail vertices. No physical vertex is discarded from an active head role.
\item If $n_u\ge4$ and $n_w\ge4$, universal arcs give a positive support of size four, contradicting $\gamma_4=0$. Hence $n_w\le3$ in the remaining cases.
\item Suppose $n_u=5$ and $n_w=3$. Then $E_3>0$, $E_4=0$, and \eqref{pgr:s:fc:eq:coeff} gives $g=0$. Every product $u_{C\setminus j}$ is positive, so $\eta_jv_j=0$ for every $j$. Thus every effective internal arc is absent, and $\gamma_k=e_kE_k$.
For the three positive sink weights $x,y,z$,
\[
 E_2^2-3E_1E_3=\tfrac12\bigl((xy-xz)^2+(xy-yz)^2+(xz-yz)^2\bigr)\ge0.
\]
Also $e_2^2-e_1e_3$ is coefficientwise nonnegative: its nonzero monomial types $u_i^2u_j^2$, $u_i^2u_ju_k$, and $u_iu_ju_ku_l$ have coefficients $1,1,2$, respectively. Therefore
\[
 \gamma_2^2-3\gamma_1\gamma_3
 =(e_2^2-e_1e_3)E_2^2+e_1e_3(E_2^2-3E_1E_3)\ge0.
\]
\item Suppose $n_u=4$ and $n_w=3$. Let $h$ be the unique zero-tail core vertex, so that $u_{C\setminus h}>0$. The fourth coefficient is $u_{C\setminus h}v_h\eta_hE_3$. Its vanishing forces $v_h\eta_h=0$. Every possible in-neighbor of $h$ is one of the four positive tails. Thus $h$ has no effective incoming arc and, because $u_h=0$, no effective outgoing arc. Both roles are now unusable in every positive support, so deleting $h$ preserves $\Gamma_D$ and reduces to (C4).
\item It remains that $n_w\le2$. Pad with zero sink weights to use two explicit variables $w_0,w_1$. If $n_w=0$, only five physical vertices are usable, so the degree is at most two. The two-sink inequality below (Lemma~\ref{pgr:s:fq:lem:finite}) covers all these cases, including further zero activities.
\end{enumerate}
This exhausts the entire face $\gamma_4=0$. In particular, no quartic-normalized inequality is substituted for the required cubic constant $3$.

\mnote{case~4 is the step that source~34's scope paragraph says cannot be
taken in general: a fifth core vertex with zero tail activity may remain a
head. Here $\gamma_4=0$ forces its head role to be unusable as well, and
only then is it deleted. Source~06 keeps such a vertex as a head-only
vertex at actual degree four (Section~\ref{pgr:s:fb:sec}).}

\subsection{The exact two-sink verification}
With two sinks, there are at most seven usable physical vertices and hence degree at most three. With $E_1=w_0+w_1$ and $E_2=w_0w_1$,
\[
 \gamma_1=a+e_1E_1,\quad\gamma_2=b+cE_1+e_2E_2,\quad\gamma_3=fE_1+dE_2.
\]
The target polynomial is
\begin{equation}\label{pgr:s:fq:eq:target}
 P_H=(b+cE_1+e_2E_2)^2-3(a+e_1E_1)(fE_1+dE_2)
\end{equation}
in the twelve variables $x=(u_0,\ldots,u_4,v_0,\ldots,v_4,w_0,w_1)$, homogeneous of total degree eight.

\begin{lemma}[Finite verification]\label{pgr:s:fq:lem:finite}
For every loopless directed relation on five labeled core vertices, $P_H\ge0$ on the nonnegative orthant.
\end{lemma}

\paragraph{Coverage and selected proof branches.}
There are $2^{20}=1,048,576$ labeled loopless relations, represented by 9,608 disjoint orbits under simultaneous permutations of the five physical core labels. The verifier checks the full orbit partition explicitly. It then selects exactly one proof for each representative, in this order:
\begin{center}
\begin{tabular}{lr}
\toprule
Proof branch&Selected classes\\
\midrule
Two-tail/two-head role cover&43\\
Physical vertex cover of size at most three&21\\
One-tail cover with a certified HPP side&1,743\\
Established preorder theorem, additional classes&49\\
Direct rational square certificate&7,658\\
Tail-sum-multiplied rational square certificate&94\\
\bottomrule
\end{tabular}
\end{center}
The structural/preorder union has 1,856 classes. Of the 139 preorder core classes (after adding the diagonal), 90 belong to the preceding branches and 49 are additional. Their two-sink extensions are preorders of physical degree at most three.

\mnote{the last two counts are macros of the delivered file
\nolinkurl{article/certificate-summary.tex}, which is not shipped; their
values are typed in. The 9,608 class representatives are the core
catalogue of sources 34 and~06, in the same order: the delivered
\nolinkurl{inputs/cores.txt} is byte-identical to the shipped
\nolinkurl{data/34-five-core-discovery-cores.txt}. The certificates are
nevertheless new: their target is the two-sink polynomial
\eqref{pgr:s:fq:eq:target}, not source~34's moment polynomial or
source~06's, and none of their files equals a file of sources 34 or~06.}

For the one-tail branch, a stored cover $(p,Q)$ satisfies: every arc has tail $p$ or head in $Q$, possibly after reversal. Its side transversal matroid has centers $Q$, one element for every physical tail with its neighborhood in $Q$, and one private dummy at each center. A basis is counted once by matching feasibility. The checker retains every repeated multi-neighbor column; only genuine loops and singleton-neighborhood parallel elements are simplified. It checks the complete basis set, not just numerical invariants.

The side HPP inputs comprise 14 cases on at most seven elements, 201 eight-element cases of nonmiddle rank with all required HPP minors, and 1,528 cases matched to the published nine-element positive list. The small classifications and positive databank are those of Kummer and Sert \cite{KS,KSdata}. Each nine-element match includes an explicit element permutation and equality of all 126 rank-four basis indicators; a rank-five side is dualized first. The new computation checks membership in a proved-positive list, not a complete nine-element classification.

\paragraph{Why the one-tail criterion applies to physical supports.}
\mnote{source~35 proves here that the stable-side criterion applies to the
physical supports of an arbitrary directed relation. The proof has four
steps. The split seed identity is
$\mu_B=(\prod_ia_i)(\prod_{j\in Q}v_j)F_M(d_j=b_j/v_j,\ t_i=-u_i/a_i)$ for
role-split monomer variables $a_i,b_j$. An uncovered head $x$ is added by
$\mu\mapsto b_x\mu-c_x\partial_{a_p}\mu=b_x\mu(a_p-c_x/b_x)$ with
$c_x=u_pv_x\mathbf1_{p\to x}$. The two roles of each physical vertex are
merged by $(\partial_{a_i}+\partial_{b_i})\mu|_{a_i=b_i=z_i/2}$. Finally
the diagonal $\mu(t,\ldots,t)=t^n\Gamma_D(-t^{-2})$ transfers the roots.
These are the four steps of source~22's proof of
Theorem~\ref{pgr:r:ir:thm:oneside}, which is itself source~15's
Theorem~\ref{pgr:r:ps:thm:pendant} with a single hub, so the argument is
replaced by that pointer. Source~35's own remarks are kept. Its
coefficient $c_x$ includes the arc indicator. Multiaffinity in $a_p$
gives the displayed factorization. The merge sends $a_ib_i$ to $z_i$,
either of $a_i,b_i$ to $1$ and the constant term to $0$, so it removes
simultaneous use of both physical roles and preserves exactly the Boolean
support polynomial. The empty-support coefficient is one, which rules out
a zero output and controls zero-activity limits. The criterion is used in
its arbitrary-directed-relation scope, which the shipped note
\nolinkurl{35-five-core-cubic-audits-ONE_TAIL_HPP_SCOPE.md} audits.}

\paragraph{Exact algebraic certificates.}
The remaining representatives use one of two explicitly tagged identities:
\begin{equation}\label{pgr:s:fq:eq:certificate}
 M_H P_H=R_H+\sum_{\ell}q_\ell x^{\alpha_\ell}
       \left(\sum_j c_{\ell j}x^{\beta_{\ell j}}\right)^2,
 \qquad M_H\in\{1,U\},\quad U=\sum_{i=0}^4u_i.
\end{equation}
Here every $q_\ell$ is a positive rational, every exponent is a nonnegative integer, each $c_{\ell j}$ is rational, and every coefficient of $R_H$ is nonnegative. The total degree is eight for $M_H=1$ and nine for $M_H=U$. If $U>0$, the second identity proves $P_H\ge0$ by division. If $U=0$, all core tails are zero and every positive-rank support has weight zero, so $P_H=0$ directly. This proves the same target on every nonnegative face.

The independent checker reconstructs $\gamma_1,\gamma_2,\gamma_3$ by literal disjoint endpoint-pair enumeration and a separate grouped Hall test. It expands every square with integer and rational arithmetic, including contributions at monomials absent from the target, and verifies the entire remainder. It rejects unknown multipliers, wrong degrees, and duplicate representative IDs across the two certificate directories. Numerical Gram matrices or optimization may suggest squares during discovery; neither numerical output nor any solver is trusted or needed in the proof replay.

\mnote{the tail-sum multiplier $M_H=U$ is new in Part~III: the
certificates of sources 10, 34 and~06 are unmultiplied. (Source~22 uses
positive multipliers for four of its direct certificates,
Section~\ref{pgr:r:ir:sec:finite}.) The producer that
found the squares is not delivered, so the certificates cannot be
regenerated, only checked.}

\paragraph{Completion receipt.}
The complete independent replay verifies all 9,608 classes with no missing or duplicate proof IDs. It checks every one of the 7,901 retained certificate files: 7,807 direct certificates and 94 tail-sum certificates. The total is 1,368,926 positive rational square summands and 11,514,499 positive remainder terms. Of the direct files, 149 overlap earlier structural/preorder branches; they remain included and checked for provenance, while the selected ledger is disjoint. All literal/Hall polynomial comparisons, structural witnesses, published positive-list matches, and the full permutation-orbit partition pass. The executable verifier uses only Python's standard library and exact integer/rational arithmetic. The release records source hashes and theorem-application checks separately from replays of the companion results' older certificate collections.

\mnote{verification at intake (2~October 2026). The integrity check of
the delivered package passed on the merged six parts (74 files and six
data shards, 4.9~s). Fifteen certificates (ids 0--7, 1814--1817 and
8637--8639) were decoded from shards 01, 02 and~06, matched against the
delivered per-file hashes and checked by the delivered checker in partial
mode, together with the full orbit coverage of all $2^{20}$ labeled
cores: their records equal those of the delivered per-certificate
manifest; the sample includes the tail-sum certificate~1817. The
delivered extended and multiplier self-tests passed on this sample; the
multiplier self-test needs the tail-sum certificate~762, which was added
for it. The structural checker re-verified all 1,807 structural records and the 139
preorder records (49 new beyond the structural branches, 1,856 combined)
in 103~s, and reproduced the delivered structural records. The symbolic
sharpness check passed. The full replay of all 7,901 certificates was not
run: it would take about 1.5~hours with four workers and 6~GB of
temporary disk on the intake machine (the delivered receipt records
279~s on the delivery machine). The delivered approvals say that a clean
replay of the final six ZIP files was a separate release step, which they
do not record. The delivered receipt, which this paragraph reports, is
therefore the pre-release replay. Shipped receipts:
\nolinkurl{data/35-five-core-cubic-audits-full-hybrid-20261001T164052Z-receipt.json}
and the approvals beside it.}

The supporting dependency ledger gives further counts. The replay checks
192,528 basis indicators of the published list and 3,216 one-element
minors, including 411,648 rank-function identities. This is a
theorem-application and exact databank-membership check. It does not
rediscover the published HPP proofs or rerun the published Gram searches.

\subsection{Actual normalization and all-degree synthesis}

\mnote{source~35 proves here, for every finite loopless directed relation
under nonnegative independent role weights and actual degree $\rho\ge2$,
the first comparison $(\rho-1)\gamma_1^2\ge2\rho\gamma_2$, by the weighted
clique bound of Motzkin and Straus \cite{MS} on the graph of effective
arcs, with a direct mass-transfer proof of that bound. This is
Theorem~\ref{pgr:thm:first} with Lemma~\ref{pgr:lem:ms}, by the same route,
and is replaced by that pointer.}

At actual degree three, \eqref{pgr:eq:first} is $\gamma_1^2\ge3\gamma_2$, and the other normalized comparison is exactly \eqref{pgr:s:fq:eq:cubic}. At degree two the needed constant is $4$, again supplied by \eqref{pgr:eq:first}; degrees zero and one are vacuous. A positive top support has a positive matching witness, and taking subsets of its arcs gives positive supports of every smaller size. Hence there are no internal zero coefficients.

\begin{corollary}[All actual degrees]\label{pgr:s:fq:cor:all}
Every arbitrary loopless directed five-vertex core with an independent universal-sink cloud is rank ultra log concave under independent nonnegative tail and head activities, normalized by its own actual surviving degree $0\le\rho\le5$.
\end{corollary}
\begin{proof}
Degrees at most three are settled above. The actual-degree-four boundary theorem, source~06's Theorem~\ref{pgr:s:fb:thm:main}, gives the two interior constants $9/4$ and $8/3$; the actual-degree-five theorem, source~34's Theorem~\ref{pgr:s:fc:thm:main}, gives $\gamma_2^2\ge2\gamma_1\gamma_3$, $\gamma_3^2\ge2\gamma_2\gamma_4$, and the stronger last comparison $\gamma_4^2\ge(25/8)\gamma_3\gamma_5$. Their first comparisons are \eqref{pgr:eq:first}. Their statements are applied only at their asserted actual degrees. No lower-degree normalization is inferred from a padded higher-degree result. The companion archive identities and exact scope are listed in the dependency ledger.
\end{proof}
\mnote{source~35 writes ``the previously delivered actual-degree-four
boundary theorem'' and ``the previously delivered actual-degree-five
theorem''. Its ledger pins their archives, which are those of sources 06
and~34. With source~14's last gap (Section~\ref{pgr:s:lg:sec}), Part~III
now proves rank-ULC at every actual degree for four- and five-vertex cores
with universal-sink clouds.}

\subsection{Sharpness of the cubic constant}
Take core arcs $0\to3$ and $1\to2$, with the universal arcs to two sinks. For $\varepsilon>0$, set
\[
 u=(\varepsilon,\varepsilon,0,0,1),\qquad
 v=(0,0,\varepsilon^{-1},\varepsilon^{-1},0),\qquad
 w_0=w_1=\tfrac12.
\]
This is a preorder core after adding diagonals. Direct Boolean enumeration gives
\begin{align*}
 (\gamma_0,\gamma_1,\gamma_2,\gamma_3)
 &=\left(1,3+2\varepsilon,
 3+\tfrac52\varepsilon+\tfrac14\varepsilon^2,
 1+\tfrac12\varepsilon\right),\\
 \gamma_2^2-3\gamma_1\gamma_3
 &=\tfrac92\varepsilon+\tfrac{19}4\varepsilon^2
  +\tfrac54\varepsilon^3+\tfrac1{16}\varepsilon^4.
\end{align*}
The degree is three for every $\varepsilon>0$, and $\Gamma_D\to(1+t)^3$ coefficientwise. Hence $\gamma_2^2/(\gamma_1\gamma_3)\to3$, excluding every larger uniform constant. An independent symbolic replay is included.

\mnote{the symbolic replay is the shipped
\nolinkurl{code/35-five-core-cubic-replay-check_sharpness_portable.py};
it passed at intake, and its output equals the shipped
\nolinkurl{data/35-five-core-cubic-audits-portable-sharpness-receipt.json}.
The example was also recomputed at intake by direct enumeration, with
exact arithmetic, at $\varepsilon=10^{-3}$. The limit $(1+t)^3$ attains
equality in every degree-three Newton inequality, so no constant above
three can hold. No other source of Part~III states sharpness for its
cubic constant ($3$ for source~10, $2$ for source~34, $9/4$ for
source~06). The statement here is for five-vertex cores on the face
$\gamma_4=0$: the example uses all five core vertices, and the sharpness of
source~10's constant three for four-vertex cores is not addressed.}

\subsection{Scope and limits}
The delivered package states its scope limits. Source~35 does not claim
real-rootedness for the full family, a result for arbitrary nonuniversal
sink neighbourhoods, a general directed physical-degree-three
classification, a formal proof-assistant verification, or literature
priority. A padded degree-four or degree-five normalization is never used
at degree three. Its final audit adds that no numerical proof premise,
sampling substitute or producer-code import is used, and no stability
claim is made for every core. The package's command
\texttt{python3 -B verify.py \mbox{--workers 4} \mbox{--output fresh-replay.json}} is
the fresh proof replay; its integrity-only mode, partial logs, exploratory
numerical searches and the stored reference receipt are not. Existing
companion theorems are mathematical inputs, and their earlier certificate
archives are not rerun. The local audits check exact applicability and
data membership for the Kummer--Sert inputs. They do not rerun those
authors' Gram searches or the entire earlier dependency archives.

\mnote{an ordinary, certificate-free proof of
Lemma~\ref{pgr:s:fq:lem:finite} would remove the computer assistance from
Theorem~\ref{pgr:s:fq:thm:main}. Source~35 states no research questions.}


\part*{Appendix}
\addcontentsline{toc}{part}{Appendix}
\appendix
\section{Provenance and merge record}\label{pgr:sec:provenance}

\subsection{The archives}\label{pgr:sec:prov:archives}
All eighteen archives arrived in ProveIt commit \nolinkurl{f8c3a392a} on
1~October 2026 and were placed in commit \nolinkurl{8f5c53106} (batch~73,
cluster~C2); archives 11, 12 and~15 were handed over to this cluster by the
cluster that placed batch-73 manuscript~26 in \mrn. Only sources 16 and~25
cite a ProveIt commit: \nolinkurl{d5e863bba} (\prp{} at ``Write batch 70
(4/5)''); the labels they cite are unchanged since. All other 40-digit
hexadecimal strings in the packages are SHA-256 digests of package files.
Every package's checksum ledger verified at intake.
Source~35's archive \nolinkurl{universal-sink-five-core-cubic-package}
(six ZIP parts in one archive) arrived later, in commit
\nolinkurl{096ee7b87} on 2~October 2026 (batch~77, manuscript~71), and
was placed in commit \nolinkurl{089d2b825} (batch~77, cluster~U). It
cites no ProveIt commit. Its package integrity check passed at intake on
the merged six parts, and the five companion archives it pins by SHA-256
are batch-73 archives 10, 34, 07, 06 and~22, byte for byte.

\begin{center}\small
\begin{longtable}{p{9mm}p{40mm}p{52mm}p{13mm}p{20mm}}
\toprule
Src. & Archive(s) & Manuscript (pages as delivered) & Pin & Printed\\\midrule
\endhead
25 & \nolinkurl{preorder-gamma-degree-three-package} & Rank ultra log concavity for preorder gamma polynomials of degree at most three (14) & \nolinkurl{d5e863bba} & Section~\ref{pgr:u:d3:sec}\\
16 & archives 16--19, \nolinkurl{preorder-degree-four-part-}$k$\texttt{-of-4}, $k=4,3,2,1$ & Degree four preorder support polynomials (17); the base & \nolinkurl{d5e863bba} & Section~\ref{pgr:u:d4:sec}\\
23 & \nolinkurl{template26-structural-supplement} & A smaller proof of the last Newton gap for template 26 (6) & --- & Section~\ref{pgr:u:t26:sec}\\
20 & \nolinkurl{role-cover-package} & Real stable monomers under a two by two role cover (9) & --- & Section~\pgrfwd{pgr:r:rc:sec}{in Part~II}\\
12 & \nolinkurl{three-core-rayleigh-result} & Coefficientwise Rayleigh inequalities for a three vertex shore (6) & --- & Section~\pgrfwd{pgr:r:tc:sec}{in Part~II}\\
11 & \nolinkurl{physical-cover-three-result} & Rank ultra log concavity under a three vertex physical cover (5) & --- & Section~\pgrfwd{pgr:r:pc:sec}{in Part~II}\\
15 & \nolinkurl{bipartite-and-pendant-support-result} & Lorentzian and real stable directed support polynomials (6) & --- & Section~\pgrfwd{pgr:r:ps:sec}{in Part~II}\\
22 & \nolinkurl{independent-role-degree-three-package} & Independently weighted preorder support polynomials of actual degree at most three (21) & --- & Section~\pgrfwd{pgr:r:ir:sec}{in Part~II}\\
14 & \nolinkurl{universal-sink-last-gap-package} & Sharp final coefficient bounds for directed cores with universal sinks (4) & --- & Section~\pgrfwd{pgr:s:lg:sec}{in Part~III}\\
10 & \nolinkurl{universal-sink-all-degrees-package} & Rank ultra log concavity for four vertex cores with universal sinks (5) & --- & Section~\pgrfwd{pgr:s:ad:sec}{in Part~III}\\
34 & \nolinkurl{universal-sink-five-core-package} & Rank five ultra log concavity for five vertex cores with universal sinks (6) & --- & Section~\pgrfwd{pgr:s:fc:sec}{in Part~III}\\
06 & archives 07 and 06, \nolinkurl{universal-sink-five-core-boundary-package-part-1}, \texttt{-part-2} & The rank four boundary for five vertex cores with universal sinks (7) & --- & Section~\pgrfwd{pgr:s:fb:sec}{in Part~III}\\
35 & batch~77, \nolinkurl{universal-sink-five-core-cubic-package} (six parts in one archive) & The cubic boundary for five-vertex cores with universal sinks (5) & --- & Section~\ref{pgr:s:fq:sec}\\
31 & \nolinkurl{weighted-degree-three-package} & Vertex-weighted preorder support polynomials of actual degree at most three (17) & --- & superseded by 22\\
13 & \nolinkurl{universal-sink-degree-four-package} & Weighted degree four rank ultra log concavity for universal sink extensions (5) & --- & superseded by 10\\\bottomrule
\end{longtable}
\end{center}

\subsection{What each source contributed}\label{pgr:sec:prov:contrib}
\begin{itemize}[leftmargin=*]
\item \emph{25}: unit-activity rank-ULC through actual degree three; the
  bounded-core normal form at every degree; the minimum order fifteen of a
  nonreal-rooted degree-three gamma polynomial and its uniqueness; the
  ordinal-sum family. Its first inequality is \prp's.
\item \emph{16} (base): unit-activity rank-ULC at actual degree four; the
  exterior-only Lorentzian construction; the outward-star last gap; the
  kernel formula and coverage, the bounded branches and the integer-population
  certificate schema. Its Section~3 and Appendix~A are source~20's article
  (printed from source~20); its first inequality is \prp's.
\item \emph{23}: a second, smaller route for the last gap of template~26.
\item \emph{20}: real stability of the signed monomer polynomial under every
  two by two role cover, with independent activities; the free-product
  corollary; the gamma transfer. It generalizes \mrn, Theorem
  \texttt{mrn:w:sb:thm:main} (the bipartite case with disjoint cover sides)
  by a different proof.
\item \emph{12}: the coefficientwise three-core boundary inequality
  (computer-assisted), the full coefficientwise Rayleigh property,
  negative correlation, and a four-core obstruction.
\item \emph{11}: rank-ULC under every physical vertex cover of size at most
  three, with arbitrary internal core arcs; bipartite strictness and
  equality.
\item \emph{15}: the bipartite physical Lorentzian theorem, overlapping
  role covers with an independent core, the pendant-head stability
  criterion and its elementary seed.
\item \emph{22}: rank-ULC of every finite preorder of actual degree at most
  three with independent role activities; the articulation theorem; the
  zero-activity role filter; the half-plane-property side-matroid
  criterion and the final rank-32 Gram certificate. Its three-core and
  one-tail theorems are special cases of sources 11 and~15. It supersedes
  archive~31.
\item \emph{14}: the sharp last-coefficient bound $K(r,m)$ for every core
  size with universal sinks.
\item \emph{10}: rank-ULC at every actual degree for four-vertex cores with
  universal sinks, through the certified inequality
  $\gamma_2^2\ge3\gamma_1\gamma_3$. It supersedes archive~13.
\item \emph{34}: rank-ULC at actual degree five for five-vertex cores with
  universal sinks, and a counterexample to the constant three there.
\item \emph{06}: the actual-degree-four boundary of the five-vertex family,
  and the strengthened ordinary core comparisons.
\item \emph{35} (batch~77): the cubic boundary
  $\gamma_2^2\ge3\gamma_1\gamma_3$ on the face $\gamma_4=0$ for
  five-vertex cores, hence, with sources 06 and~34, rank-ULC at every
  actual degree; the degree-drop reduction; the two-sink lemma, whose
  9,608 classes are settled by structural theorems and source~22's
  preorder theorem (1,856), direct rational square certificates (7,658) and
  tail-sum-multiplied ones (94); sharpness of the constant three. Its
  Boolean coefficients are source~34's, its first comparison is \prp's,
  and its one-tail step is Theorem~\ref{pgr:r:ir:thm:oneside}. It
  re-delivers, byte for byte, statements and approvals of sources 10, 11,
  20, 22, 34 and~06.
\end{itemize}

\subsection{Where the merge had to choose}\label{pgr:sec:prov:choices}
\begin{enumerate}[leftmargin=*]
\item \emph{One report, three Parts.} All sources answer \prp's Research
  question \texttt{mr:q:gamma}; they would be additions to \prp, which
  already has thirteen Parts and whose batch-72 continuation was already
  split off into \mrn. They form one new report beside \mrn.
\item \emph{Base.} Source~16, the farthest-reaching unit-activity result.
  Its Section~3 and Appendix~A are printed once, from source~20, in
  Part~II; source~16's different operator proofs are kept there as second
  routes (Section~\ref{pgr:u:d4:sec:rolecover}).
\item \emph{Common definitions and the first inequality} are printed once
  (Section~\ref{pgr:sec:object}); the sources' copies are pointers.
\item \emph{Same theorem, two sources.} Source~25's theorem is the unit case
  of source~22's: both proofs are printed, in Parts I and~II. Source~23
  re-proves one entry of source~16's ledger: both routes are described,
  source~23's printed. Source~14's sharp bound contains the last
  inequalities of sources 10 and~34 (and of archive~13): it is printed
  once, and those proofs, which are the same route at $r=4,5$, are
  pointers. Sources 11 and~15 contain source~22's three-core and one-tail
  theorems: source~22's three-core proof, an ordinary argument through
  Wagner's rank-three theorem, is kept as a second route; its one-tail
  proof is the pendant proof of source~15 and is a pointer. The
  three-squares identity of sources 11, 12 and~22 is printed once (in
  source~12's section); the moment envelope of sources 10, 34 and~06 once
  (in source~10's); the core comparisons of sources 34 and~06 once, in
  source~06's stronger form; the nonreal cubic of sources 16, 22, 25 and
  archive~31 once (in source~25's).
\item \emph{Superseded versions.} Archive~31 (vertex activities) is an
  earlier version of source~22: same skeleton, eleven byte-identical
  proof notes; its theorem is printed as the case $u_i=v_i$ of source~22's.
  Its only own material, a ledger of vertex-only certificates for a finite
  special case, survives in the arrival commit. Archive~13 is an earlier
  version of source~10 with a weaker middle inequality,
  $4\gamma_2^2\ge9\gamma_1\gamma_3$, certified by binomial squares only,
  and an alternative first-gap argument valid in degree four; it survives
  in the arrival commit.
\item \emph{Source numbers} of multi-archive packages are their lowest
  archive numbers (16 and~06), as are their file prefixes.
\item \emph{Source~35} (batch~77) is a new last section of Part~III, not a
  new report: it continues the five-vertex line of sources 34 and~06 in
  their setting, with their coefficient formulas and their core
  catalogue. Its number is the next after this report's highest,~34. Its
  Boolean coefficients (source~34's), its first comparison
  (Theorem~\ref{pgr:thm:first}) and its argument for the one-tail
  criterion (the proof of Theorem~\ref{pgr:r:ir:thm:oneside}) are
  pointers. Its references to ``previously delivered'' or ``established''
  results are replaced by labels. Its letters $m$, $q$, $s$ and two uses
  of~$p$ are renamed (Section~\ref{pgr:s:fq:sec}). Its 27 delivered files
  that are byte-identical to batch-73 material already here (17 shipped
  files, five members of source~22's shipped audit archive, five printed
  manuscripts) are not shipped again.
\item \emph{Not shipped} (README): embedded exact copies of other archives,
  PDFs, delivery READMEs, checksum ledgers, rendered page images, build
  logs and scripts of the delivered PDFs, the files of source~16 larger
  than 2~MB, source~06's two bounded-sink certificate archives and their
  manifest, and the regenerable outputs listed in the README. Source~34's
  9,610 certificate files are shipped as one derived container whose
  members are byte-identical to the delivered files. Source~35's six
  certificate shards (95.1~MB), which no shipped program regenerates, are
  kept in its arrival commit; its structural input is shipped as a
  derived \texttt{.xz} container that decompresses to the delivered
  bytes; four regenerable outputs and ledgers are not shipped (README).
\end{enumerate}

\subsection{Shipped files}\label{pgr:sec:prov:files}
Each source's code, data, audits and proof notes are shipped as delivered,
renamed with its prefix (\nolinkurl{25-degree-three-}, \nolinkurl{16-degree-four-},
\nolinkurl{23-template26-}, \nolinkurl{20-role-cover-},
\nolinkurl{12-three-core-rayleigh-}, \nolinkurl{11-physical-cover-},
\nolinkurl{15-pendant-support-}, \nolinkurl{22-independent-role-},
\nolinkurl{14-sink-last-gap-}, \nolinkurl{10-sink-all-degrees-},
\nolinkurl{34-five-core-}, \nolinkurl{06-five-core-boundary-},
\nolinkurl{35-five-core-cubic-}): code in
\texttt{code/}, data and records in \texttt{data/}, audits and proof notes
at the top of the report directory. The report README lists them, maps the
delivered names, lists what is not shipped and how to retrieve it from the
arrival commit, and gives rerun instructions that work on copies. The text
of delivered files still uses delivered paths.

\begin{thebibliography}{99}\small
\bibitem{PRP} ProveIt project report, \emph{A Reflexive Root-Polytope Model for
Preorder $h$-Polynomials} (\prp), cited by sources 16 and~25 at the pinned
source revision \nolinkurl{d5e863bba6b40738af292b45005a0cf0495903e4}, Parts IV
and XII.
\url{https://github.com/VladimirReshetnikov/ProveIt/blob/d5e863bba6b40738af292b45005a0cf0495903e4/SetTheory/Cardinals/docs/reports/enumerative-combinatorics/preorder-root-polytopes/article.tex}.
Current version: \nolinkurl{SetTheory/Cardinals/docs/reports/enumerative-combinatorics/preorder-root-polytopes/}.
\bibitem{AC} C. A. Athanasiadis and F. Chapoton, \emph{Polytopes and posets
associated to preorders}, arXiv:2605.26916v1 (2026), Section 5 and
Proposition 9.2. \url{https://arxiv.org/abs/2605.26916}. (Source~16 lists
the authors as F.~Chapoton and C.~A.~Athanasiadis.)
\bibitem{AXY} C. A. Athanasiadis, Q. Xiao, and X. Yan, \emph{Lattice point
enumeration of some arbor polytopes}, arXiv:2603.11654v1 (2026), Remark 2.3.
\url{https://arxiv.org/abs/2603.11654}.
\bibitem{AJ} P. Alexandersson and A. Jal, \emph{Rook matroids and
log-concavity of $P$-Eulerian polynomials},
\href{https://arxiv.org/abs/2410.00127v3}{arXiv:2410.00127v3}, June 30, 2026.
\bibitem{BB} J. Borcea and P. Br\"and\'en, \emph{The Lee--Yang and
P\'olya--Schur programs. I. Linear operators preserving stability},
Inventiones Mathematicae \textbf{177} (2009), 541--569.
\url{https://arxiv.org/abs/0809.0401}.
\bibitem{BH} P. Br\"and\'en and J. Huh, \emph{Lorentzian polynomials},
Annals of Mathematics \textbf{192} (2020), 821--891.
\href{https://doi.org/10.4007/annals.2020.192.3.4}{doi:10.4007/annals.2020.192.3.4};
\href{https://arxiv.org/abs/1902.03719}{arXiv:1902.03719}.
\bibitem{BrandenHPP} P. Br\"and\'en, \emph{Polynomials with the half-plane
property and matroid theory}, Advances in Mathematics \textbf{216} (2007),
no.~1, 302--320. \url{https://arxiv.org/abs/math/0605678}.
\bibitem{CS} H. Crapo and W. Schmitt, \emph{The free product of matroids},
European Journal of Combinatorics \textbf{26} (2005), no.~7, 1060--1065.
\url{https://arxiv.org/abs/math/0409080}.
\bibitem{KS} M. Kummer and B. Sert, \emph{Matroids on Eight Elements with the
Half-plane Property and Related Concepts}, arXiv:2111.09610v4, Sections 5
and 6, especially Theorem 5.2. \url{https://arxiv.org/abs/2111.09610}.
\bibitem{KSdata} M. Kummer and B. Sert, supplementary matroid data, Zenodo
record 6108027 (CC~BY~4.0). Certified-positive nine-element list
\texttt{n9r4Hpp.txt}. \url{https://zenodo.org/records/6108027}.
\bibitem{MS} T. S. Motzkin and E. G. Straus, \emph{Maxima for graphs and a new
proof of a theorem of Tur\'an}, Canadian Journal of Mathematics \textbf{17}
(1965), 533--540, Theorem~1.
\href{https://doi.org/10.4153/CJM-1965-053-6}{doi:10.4153/CJM-1965-053-6}.
\bibitem{RSW} J. Ross, H. S\"uss and T. Wannerer, \emph{Dually Lorentzian
polynomials}, Monatshefte f\"ur Mathematik \textbf{208} (2025), 495--524.
\url{https://doi.org/10.1007/s00605-025-02134-6}.
\bibitem{RU} F. R\"ohrle and M. Ulirsch, \emph{Logarithmic concavity of
bimatroids}, Annals of Combinatorics \textbf{30} (2026), 501--523. Theorem A
and Example 2.4. \url{https://doi.org/10.1007/s00026-025-00780-z};
\url{https://arxiv.org/abs/2402.15317}.
\bibitem{WagnerR3} D. G. Wagner, \emph{Rank-three matroids are Rayleigh},
Electronic Journal of Combinatorics \textbf{12} (2005), Note N8, Theorem 1.1.
\href{https://doi.org/10.37236/1975}{doi:10.37236/1975};
\url{https://arxiv.org/abs/math/0403216}. (Source~12 cites the 2004 arXiv
version as a background reference only: its proof of the boundary theorem
uses the directed-support classification.)
\bibitem{WagnerSurvey} D. G. Wagner, \emph{Multivariate stable polynomials:
theory and applications}, Bulletin of the American Mathematical Society
\textbf{48} (2011), 53--84. \url{https://arxiv.org/abs/0911.3569};
\url{https://www.math.uwaterloo.ca/~dgwagner/ceb_wagner.pdf}.
\bibitem{WW} D. G. Wagner and Y. Wei, \emph{A criterion for the half-plane
property}, arXiv:0709.1269 (2007), Theorem 3.
\url{https://arxiv.org/abs/0709.1269}.
\bibitem{West} D. B. West, \emph{A short proof of the Berge--Tutte Formula and
the Gallai--Edmonds Structure Theorem}, author manuscript, Theorem 5.
\url{https://dwest.web.illinois.edu/pubs/galledm.pdf}.
\end{thebibliography}

\end{document}
